% 本文件由 LaTeX 排版 Agent 根据有效 translation.json 自动生成。
% author=John Chrysostom; work_id=2301; title=厄弗所书讲道集

\chapter{厄弗所书讲道集}

% structure: p0001 source=h1 target=omitted reason=page-title-already-rendered-by-parent

\Emph{纲要}\newline{}\Emph{讲道一}\newline{}\Emph{讲道二}\newline{}\Emph{讲道三}\newline{}\Emph{讲道四}\newline{}\Emph{讲道五}\newline{}\Emph{讲道六}\newline{}\Emph{讲道七}\newline{}\Emph{讲道八}\newline{}\Emph{讲道九}\newline{}\Emph{讲道十}\newline{}\Emph{讲道十一}\newline{}\Emph{讲道十二}\newline{}\Emph{讲道十三}\newline{}\Emph{讲道十四}\newline{}\Emph{讲道十五}\newline{}\Emph{讲道十六}\newline{}\Emph{讲道十七}\newline{}\Emph{讲道十八}\newline{}\Emph{讲道十九}\newline{}\Emph{讲道二十}\newline{}\Emph{讲道二十一}\newline{}\Emph{讲道二十二}\newline{}\Emph{讲道二十三}\newline{}\Emph{讲道二十四}\par

\SourceRecord{Translated by Gross Alexander.；Nicene and Post-Nicene Fathers, First Series；Vol. 13.；Edited by Philip Schaff.；Buffalo, NY: Christian Literature Publishing Co.,；1889.；<http://www.newadvent.org/fathers/2301.htm>.}

% structure: p0001 source=h1 target=section reason=page-title

\section{致厄弗所人书讲道}

厄弗所是亚细亚的首府。此城奉献给狄安娜，那里的人特别崇拜她，视之为伟大的女神。崇拜者的迷信确实如此之深，以至当她的神庙被焚毁时，他们连纵火者的名字都不肯透露。\par

蒙福的圣史若望一生大部分时间在那里度过：他被放逐时在那里，死在那里。保禄也曾将弟茂德留在那里，正如他在致弟茂德书中所说：\BibleQuote{我曾劝你留在厄弗所。}\BibleRef{弟前 1:3}\par

大多数哲学家，尤其是那些在亚细亚活跃的，也曾在厄弗所；据说连毕达哥拉斯本人也来自那里；也许是因为他真正的故乡\Place{萨摩斯}是\Place{爱奥尼亚}的一个岛。那里也是\Person{巴门尼德}、\Person{芝诺}和\Person{德谟克利特}的弟子们聚集之处，甚至到今天你还可以在那里看到许多哲学家。\par

我提及这些事实，并非仅仅为了陈述，而是要显示保禄在给这些厄弗所人写信时，必须付出极大的辛劳和努力。据说他确实将最深奥的思想托付给了他们，因为他们已经是受过良好教导的人；而这篇书信本身也充满了崇高的思想和教义。\par

他自罗马写信，并且如他自己告诉我们的，是在锁链中写的。\BibleQuote{也为祈求，使我能开口宣讲，大胆地传扬福音的奥秘，为这福音我作了带锁链的使者。}\BibleRef{弗 6:19}这封信充满了无比崇高和伟大的情感。那些他在别处几乎不曾提及的思想，在此却明确宣告；例如他说：\BibleQuote{为使天上的率领者与掌权者，藉着教会，得知天主各样的智慧。}\BibleRef{弗 3:10}又说：\BibleQuote{祂使我们与基督一同复活，使我们一同坐在天上。}\BibleRef{弗 2:6}又说：\BibleQuote{这奥秘在以前的世代没有告知人类，如今却藉着圣神启示给祂的圣宗徒和先知，就是外邦人同为承继者，同为恩许的分享者，在基督耶稣内。}\BibleRef{弗 3:5}\par

\SourceRecord{Translated by Gross Alexander.；Nicene and Post-Nicene Fathers, First Series；Vol. 13.；Edited by Philip Schaff.；Buffalo, NY: Christian Literature Publishing Co.,；1889.；<http://www.newadvent.org/fathers/230100.htm>.}

% structure: p0001 source=h1 target=section reason=page-title

\section{论厄弗所书第一讲}

第一章 1—2节 \BibleRef{弗 1:1}\par

\BibleQuote{\Emph{保禄，因天主的旨意，成为基督耶稣的宗徒，致书于在厄弗所的圣人以及在基督耶稣内的信友。} \Emph{愿恩宠与平安，由我们的天主圣父和主耶稣基督赐与你们。}}\par

请注意，他将\InnerQuote{通过}一词应用于圣父。但然后呢？我们要说祂是较低等的吗？当然不是。\par

\BibleQuote{致在厄弗所的圣人，}他说，\BibleQuote{以及在基督耶稣内的信友。}\par

请注意，他称有妻子、儿女和仆役的人为圣人。因为从书信结尾可以看出，他所用这个名称指的就是这些人，如他说：\BibleQuote{妻子，当服从自己的丈夫。} \BibleRef{弗 5:22} 又说：\BibleQuote{儿女，当孝敬父母。} \BibleRef{弗 6:1} 以及：\BibleQuote{仆役，当服从你们的主人。} \BibleRef{弗 6:5} 试想，现今的我们是何等懈怠，德行何等罕见，而那时有德之人是何等众多，甚至连俗人都可称为\Quote{圣人和信友}。\BibleQuote{愿恩宠与平安，由我们的天主圣父和主耶稣基督赐与你们。} \Quote{恩宠}是他的话；他称天主为\Quote{父}，因为这名号正是那恩宠赐予的明确标记。何以见得？请听他在别处所说：\BibleQuote{因为你们是儿子，天主就派遣了祂儿子的圣神，进入你们的心，呼喊：阿爸，父啊。} \BibleRef{迦 4:6}\par

\BibleQuote{并由主耶稣基督。}\par

因为为了我们人类，基督诞生，并显现于肉身。\par

\BibleRef{弗 1:3} \BibleQuote{愿我们的主耶稣基督的天主和父受赞美。}\par

请注意：祂是那降生成人的那位天主，纵然你不愿意，祂也是天主圣言的父。\par

\BibleRef{弗 1:3} \BibleQuote{祂在基督内，以天上各种属神的祝福，祝福了我们。}\par

他在这里暗指犹太人的祝福；因为那也是一种祝福，但不是属神的祝福。因为那祝福是怎样说的？\BibleQuote{愿上主祝福你，祂必祝福你身所生的。} \BibleRef{申 7:13} 又说：\BibleQuote{祂必祝福你出也蒙福，入也蒙福。} \BibleRef{申 28:4} 但这里却不是这样，而是怎样？\Quote{以各种属神的祝福。} 你还缺少什么呢？你已成了不朽的，你已得了自由，你成了儿子，你成了义人，你成了兄弟，你成了同承继者，你与基督一同为王，你与基督同受光荣；一切都已白白赐给你。\BibleQuote{祂岂不也把一切同祂一起赐给我们吗？} \BibleRef{罗 8:32} 你的初果受天使、革鲁宾、色辣芬的朝拜！你还缺少什么呢？\Quote{以各种属神的祝福。} 这里没有属血肉的。因此祂摒除了所有那些先前的祝福，当祂说：\BibleQuote{在世界上你们要受苦难，} \BibleRef{若 16:33} 为引导我们归向这些祝福。因为正如那些拥有属血肉事物的人不能听闻属神的事，同样，那些追求属神事物的人，若不先远离属血肉的事物，便不能达到它们。\par

再者，什么是\Quote{天上属神的祝福？} 他说，这不是在地上，如同犹太人的情形。\BibleQuote{你们将吃地上的美物。} \BibleRef{依 1:19} \BibleQuote{流奶流蜜的地方。} \BibleRef{出 3:8} \BibleQuote{上主要祝福你的地。} \BibleRef{申 7:13} 这里我们却没有这类事物，但我们有什么？\BibleQuote{谁爱我，必遵守我的话，我父也必爱他，我们要到他那里去，并在他那里作我们的住所。} \BibleRef{若 14:23} \BibleQuote{所以，凡听了我这些话而实行的，就好像一个聪明人，把自己的房屋建在盘石上：雨淋，水冲，风吹，袭击那座房屋，它却不坍塌，因为基础是建在盘石上。} \BibleRef{玛 7:24} 那盘石是什么？不就是那些高于一切变迁的属天事物吗？\BibleQuote{所以，凡在人前承认我的，我必在我天上的父前也承认他；但谁若在人前否认我，我在我天上的父前也必否认他。} \BibleRef{玛 10:32} 又说：\BibleQuote{心里洁净的人是有福的，因为他们要看见天主。} \BibleRef{玛 5:8} 又说：\BibleQuote{神贫的人是有福的，因为天国是他们的。} \BibleRef{玛 5:3} 又说：\BibleQuote{为义而受迫害的人是有福的，因为你们在天上的赏报是丰厚的。} \BibleRef{玛 5:11} 请看，祂处处提到天上，从未提到地上或地上的事物。又说：\BibleQuote{我们的家乡是在天上，我们也等待从那里来的救主，即主耶稣基督。} \BibleRef{斐 3:20} 又说：\BibleQuote{你们要思念天上的事，不要思念地上的事。} \BibleRef{哥 3:3}\par

\BibleQuote{在基督内。}\par

这就是说，这祝福不是藉梅瑟的手赐下的，而是藉基督耶稣；因此我们不仅在被赐予的福分性质上超越他们，也在中保上超越他们。正如他在\WorkTitle{致希伯来人书}中所说：\BibleQuote{梅瑟在祂的全家中尽忠，不过是臣仆，为给那以后要传报的事作证；但基督却是儿子，治理自己的家；我们就是祂的家。} \BibleRef{希 3:5}\par

\BibleRef{弗 1:4} \BibleQuote{就如祂在创世以前，在祂内拣选了我们，为使我们成为圣洁，没有瑕疵，在祂面前，在爱中。}\par

他的意思大致如此。藉着祂，天主祝福了我们；也藉着祂，天主拣选了我们。因此，祂将要在日后赐予我们这一切赏报。祂正是那位审判者，要说：\Quote{我父所祝福的，你们来吧！承受自创世以来为你们预备的国度。}（\BibleRef{玛 25:34}）又说：\Quote{我愿意我在那里，他们也同我在一起。}（\BibleRef{若 17:24}）这正是他在几乎全部书信中竭力证明的一点：我们的信仰并非新奇制度，而是从起初就已如此预表，并非出于任何旨意的改变，而实在是天主的救世计划与预定。这显示出他对我们极大的关怀。\par

\Quote{祂在我们内拣选了我们}是什么意思？他指的是藉着对基督的信德，基督在我们出生之前——甚至在世界奠基之前——就已为我们幸福地安排了这一切。\OriginalTerm{根基}{foundation}一词用得很美，仿佛祂向那从极高处抛下的世界指去。确实，天主的至高是浩瀚而不可言喻的，其差距不在于空间，而在于那不可相通的本性；受造物与创造主之间的距离是何等辽阔！异端者听到此语理当羞愧。\par

然而祂为何拣选了我们？\Quote{为使我们在祂面前，成为圣洁无瑕的。}当你听到\Quote{祂拣选了我们}时，不要以为单凭信德就足够了，他接着指出了生活与行为。他说，为此目的祂拣选了我们，且以此条件，\Quote{成为圣洁无瑕的。}从前祂也这样拣选了犹太人。条件是什么？\Quote{上主拣选了这个民族，脱离其他所有民族。}（\BibleRef{申 14:2}）若人们在自己选择时选取最好的，那么天主更会如此。事实上，他们的被选既是天主慈爱的标志，也是他们道德美善的标志。因为祂必定会拣选那些经得起考验的人。祂亲自使我们成圣，但我们也必须保持圣洁。圣洁的人是分享信德的人；无瑕疵的人是度无可指责的生活的人。然而祂所要求的不仅仅是圣洁和无瑕疵，而是我们要在\OriginalTerm{在祂面前}{before Him}显出如此。因为有些人看似圣洁无瑕，却只在人眼中如此——那些像粉饰的坟墓、披着羊皮的人。然而祂所要求的并非这等，而是如同先知所说的：\Quote{按我双手的清白。}（\BibleRef{咏 17:24}）什么清白？那在\Quote{祂眼前}的清白。祂要求的是天主之目可以注视的圣德。\par

在如此谈论了这些人的善工之后，他又再次回到祂的恩宠。他说：\Quote{由于\OriginalTerm{爱}{love}，祂\OriginalTerm{预定}{predestinated}了我们。}因为这并非出于我们的任何努力或善工，而是出于爱；然而也并非仅出于爱，也出于我们的德行。因为若仅出于爱，那么人人都当得救；而若仅出于我们的德行，那么祂的来临和整个救世计划便是多余。实际上，这既非仅出于祂的爱，也非仅出于我们的德行，而是出于二者。宗徒说：\Quote{祂拣选了我们。}拣选者深知他所拣选的是什么。\Quote{由于爱，}他又说，\Quote{祂预定了我们。}因为若没有爱，德行绝不能拯救任何人。请问：保禄能有什么益处？他若没有从起初蒙天主召唤，并因祂的爱而被吸引到祂面前，他如何能显示出他所显示的一切？此外，祂赐予我们如此大的特权，也是祂爱的结果，而非我们的德行。因为即使我们成为有德之人、相信并亲近祂，这仍然是那召唤者本身的工作，然而也属于我们。但当我们亲近祂时，祂竟赐予我们如此崇高的特权，将我们从仇敌的地位一下子带到了义子的地位——这实在是至高超之爱的作为。\par

第4至5节。他说：\Quote{由于爱，祂预定了我们，藉着耶稣基督获得义子的地位，而归於祂。}\par

你注意到没有一件事是无基督而行的？没有一件是无父而行的？一位预定，另一位引我们亲近。他说这些话是为了强调所成就的事，正如他在别处也说的：\Quote{不但如此，我们还要藉着我们的主耶稣基督，因天主而欢跃。}（\BibleRef{罗 5:11}）因为所赐的恩惠确实伟大，但藉着基督而赐予则更加伟大——因为祂没有派遣任何仆人（尽管是对仆人派遣），而是派遣了独生子自己。\par

第5节。他继续说：\Quote{按祂旨意的恩许。}（\BibleRef{弗 1:5}）\par

这就是说，因为祂恳切地愿意。这可以说是祂热切的愿望。因为\Quote{善意}一词到处都指先前的旨意，因为还有另一个旨意。例如，第一个旨意是罪人不应当灭亡；第二个旨意是，如果人变得邪恶，他们就应当灭亡。因为祂惩罚他们绝不是出于必然，而是因为祂愿意。你甚至可以在保禄的话中看到类似的东西，他说：\BibleQuote{我愿众人像我一样。}\BibleRef{格前7:7}又说：\BibleQuote{我愿年轻的寡妇嫁人，生育儿女。}\BibleRef{弟前5:14}那么，他用\Quote{善意}指的是第一个旨意、恳切的旨意、带着热切愿望的旨意，如同我们身上的情况；因为为了向较单纯的人说清楚，我不拒绝使用甚至有些通俗的说法；因为我们也常常这样，为了表达旨意的专注，说按照我们的决心行事。那么他要说的是：天主恳切地追求、热切地渴望我们的救恩。那么，为什么祂如此爱我们，祂从哪里来这样的情感？单单出于祂的良善。因为恩宠本身就是良善的果实。为此，他说，祂预定了我们成为子女的义子；这是祂的旨意，也是祂热切愿望的对象，好使祂恩宠的荣耀得以彰显。接着他说：\Quote{按照祂旨意的善意}，\par

第6节。\BibleQuote{为使祂恩宠的光荣受到赞美，这恩宠是祂在爱子内白白赐予我们的。}\BibleRef{弗1:6}\par

祂说，为使祂恩宠的光荣得以彰显，这恩宠是祂在爱子内白白赐予我们的。既然祂因此向我们施恩，为赞美祂恩宠的光荣，为彰显祂的恩宠，我们就当常在此恩宠中。\Quote{赞美祂的光荣}是什么意思？这是谁赞美祂？谁光荣祂？是我们，是天使，是总领天使，还是整个受造界？那算什么呢？天主性毫无缺乏。那么祂为什么要我们赞美和光荣祂呢？这是为了使我们对祂的爱在我们内更热烈地燃烧起来。祂不需要我们能奉献的任何东西：不是我们的服事，不是我们的赞美，也不是任何别的东西，唯独是我们的救恩；这是祂所做一切事的目的。那赞美和惊叹自己所领受恩宠的人，将因此更虔诚、更热切。\par

他说：\Quote{祂曾白白赐予我们}，祂没有说\Quote{祂曾恩赐我们}（\OriginalTerm{曾恩赐}{\GreekText{ἐχαρίσατο}}），而是说\Quote{祂曾施恩于我们}（\OriginalTerm{曾施恩}{\GreekText{ἐχαρίτωσεν}}）。这就是说，祂不仅解除了我们的罪，还使我们成为祂爱的对象。就好像有人把一个患麻风病的人，被疾病、衰老、贫穷和饥荒所摧残，然后一下子使他变成一个优雅的青年，美貌超乎众人，面颊闪耀光芒，眼眸胜过日光；然后把他安置在青春盛年，再给他穿上紫袍和冠冕以及一切王室的服饰。天主就是这样装扮和装饰我们的灵魂，使它美丽，成为祂喜悦和爱的对象。这样的灵魂是天使、总领天使和所有圣者所渴望瞻仰的。祂在我们身上倾注了这样的恩宠，使我们如此亲近祂。圣咏作者说：\BibleQuote{君王必爱慕你的美貌。}\BibleRef{咏44:11}想想我们以前说出过多少有害的话语，再看看现在我们说出多少优雅的话语。财富不再吸引我们，地上的事物也不再吸引我们，只有天上的事物，天上的事物。当一个孩子有外在的美貌，并且在他所有的话语中还有一种迷人的优雅时，我们不是称他为美丽的孩子吗？信友们就是这样。看看那些已经领受入门圣事的人说出怎样的话语！有什么比那嘴唇更美丽呢？它以纯洁的心和纯洁的唇，洋溢着喜乐的信赖，参与如此奥妙的\Quote{筵席}？有什么比我们用以弃绝魔鬼的服事、加入基督的服事的话语更美丽呢？既在圣洗池前的宣认，也在圣洗后的宣认？让我们反思，我们中那些玷污了我们的圣洗的人，哭泣，使我们能再次修复它。\par

第6节。他说：\BibleQuote{在爱子内，我们藉着祂的血获得救赎。}\BibleRef{弗1:6}\par

这如何可能呢？不仅是祂赐下祂的儿子，这一奇迹，而且更进一步，祂以这样的方式赐下，以至于爱子本人竟被杀害！\par

是的，还有更卓越的！祂为那些被恨的人赐下了爱子。看，祂为我们付了多高的代价。如果我们在恨祂、与祂为敌时，祂赐下了爱子，那么现在，我们藉着祂在恩宠中与祂和好，祂还会不做什么呢？\par

第7节。他说：\BibleQuote{我们过犯的赦免。}\BibleRef{弗1:7}\par

他再次从高处降下低处：先谈论子女的义子、圣化、无瑕可指，然后谈论受难；这不是降低他的论述，从更大的事物降到更小的，不，而是提升它，从较小的事物升到更大的。因为没有比这更伟大的了：这位儿子的血为我们而流。这比子女的义子和所有其他恩宠的恩赐都更伟大：祂连自己的儿子也没有保留。因为罪的赦免确实伟大，但更伟大的事是藉着主的血成就的。因为这是远超一切的，请看他又如何感叹：\par

第7、8节。\BibleQuote{按照祂恩宠的丰富，这恩宠祂为我们充盈满溢。}\par

上述恩赐是财富，但这一点更是如此。他说：\Quote{祂使之向我們豐厚。} 二者都是 \Quote{财富}，且 \Quote{它们已豐厚}，也就是说，以不可言喻的尺度傾注。我们无法用言语表达我们实际上经历了何等福分。因为这确实是财富，豐厚的财富，祂豐厚地赐予了，不是人的财富，而是天主的财富，以致无论如何都无法表达。为了向我们展示祂如何如此豐厚地赐予，他继续说：\par

第 8、9 节。\Quote{以各種智慧和明達，使我們知道祂旨意的奧秘。}\par

这就是说，使我们成为有智慧和明達的，在真正的智慧中，在真正的明達中。何等奇妙的友谊！因为祂将祂的秘密告诉我们；他说，祂旨意的奧秘，就好像说，祂使我们知道祂心中的事。因为这确实是充满一切智慧和明達的奧秘。因为你能提到什么与这智慧同等呢？它找到一种方法，将那些一文不值的人提升到财富和丰裕。有什么能与这巧妙的安排相比呢？那曾是仇敌的，那曾被憎恨的，顷刻之间被高举。不仅如此，还有，这件事发生在特定的时刻，这又是智慧的工作；而且它是藉着十字架完成的。这里需要长篇论述来指出这一切如何是智慧的工作，以及祂如何使我们成为有智慧的。因此他再次重复这句话：\par

\Quote{照祂在祂裡面所預定的美意。}\par

这就是说，这是祂所渴望的，祂为之劬劳，如同可以说，祂才能向我们启示这奧秘。什么奧秘？就是祂要使人类被高举在高处。这事已经成就。\par

第 10 节。\Quote{爲時期圓滿的分施，使天上和地上萬有，總歸於基督元首，就是歸於祂。}\BibleRef{弗 1:10}\par

他意思是说，天上的事物曾与地上的事物分离。它们不再有一个元首。就受造界的体系而言，确实有一位天主统管万有，但就一家之管理而言，在外邦人错误广泛蔓延时，情况并非如此，它们已与祂的服从分离。\par

他说：\Quote{爲時期圓滿的分施。}\par

他称之为时期的圆满。请注意他说话多么精确。他指出起源、目的、意愿、最初的意图出于圣父，而实现和执行由圣子完成，但他从未将\Quote{仆役}一词用于圣子。\par

他说：\Quote{祂在祂裡面揀選了我們，預定我們藉著耶穌基督得著兒子的名分，歸於祂自己；} 以及 \Quote{使祂恩寵的光榮得到讚頌，在祂裡面我們藉著祂的血獲得救贖——這是祂在祂裡面所預定的，爲時期圓滿的分施，使天上和地上萬有總歸於基督；} 他从未称祂为仆役。然而，如果 \Quote{在裡面} 和 \Quote{藉著} 这两个词意味着仅仅是一个仆役，那看看结果会怎样。就在书信的开头，他使用了 \Quote{藉著聖父的旨意} 这个说法。他的意思是，圣父愿意，圣子实行。但既不能因为圣父愿意，就排除圣子的愿意；也不能因为圣子实行，就剥夺圣父的行动。因为圣父和圣子一切都是共有的。因为祂说：\Quote{凡是我的都是你的，你的也是我的。} \BibleRef{若 17:10}\par

然而，时期的圆满就是祂的来临。因此，在祂藉着天使、先知和法律的职事做了一切之后，一无所成，几乎到了人类被造是徒然，被带入世界是徒然，反而走向毁灭的地步；当所有人都绝对灭亡，比在大洪水时更可怕时，祂设计了这一分施，即藉着恩宠，使人被造不是徒然，不是毫无目的。他称此为 \Quote{時期的圓滿} 和 \Quote{智慧}。为什么如此？因为那时，当他们正要灭亡时，他们被拯救了。\par

他说 \Quote{祂可能總歸。}\par

这个词 \Quote{總歸} 的意思是 \Quote{編織在一起}。然而，让我们努力接近其准确含义。在我们日常谈话中，这个词的意思是将在长篇论述中所说的事情总结为简洁的形式，对详细描述的事情的简明陈述。它也有这个含义。因为基督已将长期进行的分施汇集在祂自己里面，也就是说，祂将它们缩短了。因为 \Quote{祂在義中成全祂的話，並將其縮短} \BibleRef{罗 9:28}，祂既包含了从前的分施，又增添了别的分施。这就是 \Quote{總歸} 的含义。\par

这也有另一层含义；这含义是什么呢？祂设立了同一个元首统管万有，即按肉身而言的\Person{基督}，同样统管天使和人类。也就是说，祂赐给天使和人类同一个治理：对前者是降生成人的那一位，对后者是天主圣言。就如有人说到一所房子，部分朽坏、部分完好，祂重建了这房子，也就是说，祂使其更为坚固，立下更稳妥的根基。同样在这里，祂使万有归于同一个元首之下。因为若能如此，合一便能实现，紧密的联结便能实现：如果所有一切都能被置于同一个元首之下，并由此从上获得某种约束性的统一纽带。我们既蒙如此大福、如此高贵的特权、如此深厚的慈爱，就不可羞辱我们的恩主，不可使如此大的恩宠归于无效。让我们效法天使的生活、天使的德行、天使的行事为人；是的，我恳求和敦促你们，务使这一切不归于我们的审判和定罪，而归于我们获享那些善事；愿天主恩赐我们众人皆能达到，在\Person{耶稣基督}、我们的主内，愿光荣与德能归于圣父及圣神，等等。\par

\SourceRecord{Translated by Gross Alexander.；Nicene and Post-Nicene Fathers, First Series；Vol. 13.；Edited by Philip Schaff.；Buffalo, NY: Christian Literature Publishing Co.,；1889.；<http://www.newadvent.org/fathers/230101.htm>.}

% structure: p0001 source=h1 target=section reason=page-title

\section{厄弗所书讲道第二篇}

第一章 第11-14节 \BibleRef{弗 1:11}\par

\BibleQuote{在祂内，我们也成了基业，是照那按祂旨意所预定的，因为祂按自己旨意的计划运行万事。}\par

保禄竭尽全力，在一切场合彰显天主对我们不可言喻的慈爱。因为要完全做到是不可能的，请听他自己的话：\BibleQuote{啊！天主的智慧和知识的富藏是多么深奥；祂的判断是多么不可测量，祂的道路是多么难以追寻。}\BibleRef{罗 11:33}尽管如此，他仍尽其所能地彰显。那么，他所说的\BibleQuote{在祂内，我们也成了基业，被预定}是什么意思呢？上文他用了\BibleQuote{祂拣选了我们}，这里他说\BibleQuote{我们成了基业}。但既然阄是出于偶然，而非出于刻意的选择或德行（它与无知和意外紧密相连，常常越过有德之人而将无用之辈显露出来），请看他如何纠正这一点：他说\BibleQuote{是按那按祂旨意运行万事的预定}。也就是说，我们并非仅仅成了基业，正如我们并非仅仅被拣选（因为是天主拣选），同样我们并非仅仅被分配（因为是天主分配），而是\BibleQuote{按旨意}。这也是他在\WorkTitle{罗马书}中所说的：\BibleRef{罗 8:28}\BibleQuote{为按祂旨意蒙召的人}，以及\BibleQuote{祂所召的，也成义了；所成义的，也受光荣了}。他先用\BibleQuote{为按旨意蒙召的人}这个说法，同时想表明他们相对于其余人类的特权，于是也提到了阄的继承，但又没有剥夺他们的自由意志。因此，他特别强调的正是那个更属于幸运的点。因为这种阄的继承不靠德行，而是可以说靠偶然情况。这仿佛是在说：阄已投出，祂拣选了我们；但一切都是刻意的选择。被预定的人，就是说，祂选择了他们，便将他们分开。祂在我们出生之前就已视我们为被选中的。因为天主的预知是奇妙的，在万物开始之前就认识一切。\par

但请注意，他在所有场合都努力指出，这不是任何旨意改变的结果，而是这些事从一开始就这样塑造，使我们在这一方面丝毫不比犹太人差；因此他做一切事都带着这个目的。那么，基督自己怎么说？\BibleQuote{我被派遣，只是为了以色列家迷失的羊。}\BibleRef{玛 15:24}祂又对祂的门徒说：\BibleQuote{外邦人的路，你们不要走；撒玛黎雅人的城，你们不要进。}\BibleRef{玛 10:5}保禄自己也说：\BibleQuote{天主的圣言本来必须先传给你们，但你们既然拒绝它，并断定自己不配得永生，我们看，就要转向外邦人。}\BibleRef{宗 13:46}我说，这些说法是带着这个目的，以免有人以为这事只是偶然发生。\BibleQuote{是按那按祂旨意运行万事的旨意}，他说。也就是说，祂没有事后运作；祂从起初就塑造了一切，然后\BibleQuote{按祂旨意的计划}引领一切。所以，并非因为犹太人没有听从，祂才召叫外邦人，也不是出于必然，也不是出于他们引起的任何诱因。\par

\BibleQuote{为使我们这些早先在基督内怀着希望的人，都归于祂光荣的赞颂。在祂内，你们也听到了真理之言，就是你们救恩的福音。}\par

也就是说，藉着祂。请注意，他在所有场合都称基督为万物的创造者，绝不以附属者或仆人的头衔称呼祂。同样，在别处，在\WorkTitle{希伯来书}中，他说：\BibleQuote{天主在古时曾多次并以多种方式，藉着先知对我们的祖先说过话，在这末后的日子里，却藉着祂的儿子对我们说了话。}\BibleRef{希 1:1}也就是说，\BibleQuote{藉着}祂的儿子。\par

\BibleQuote{真理之言}，他说，不再是预像的真理之言，也不是影像的。\par

\BibleQuote{你们救恩的福音}。他恰当地称它为救恩的福音，在一个词中暗示与法律的对比，在另一个词中暗示与未来惩罚的对比。因为这信息不就是救恩的福音，它不再毁灭那些应被毁灭的人。\par

\BibleQuote{在祂内，你们既然信了，就受了所应许的圣神的印证，这印证是我们基业的前赎。}\BibleRef{弗 1:14}\par

在这里，\Quote{印证}一词是特殊预谋的指示。他不仅说我们被预定，也不仅说我们被分配，更说我们被印证。因为正如一个人要使那些落在他阄中的人显明，同样天主也分开他们以便相信，并为他们印证，以承受将来的基业。\par

你看，随着时间的推移，祂使他们成为惊奇的对象。当他们还在祂的预知中时，他们不为任何人所知，但一旦被印证，他们就成了明了的，虽然不象我们这样明了；因为除了少数人，他们将是明了的。以色列人也曾被印证，但那是藉着割损，像牲畜和无理性的受造物一样。我们也被印证，但却是作为儿子，\BibleQuote{藉着圣神}。\par

但\BibleQuote{所应许的圣神}是什么意思？无疑是指我们按照应许领受了那圣神。因为有两个应许：一个藉着先知，另一个藉着儿子。\par

先知们。——请听岳厄尔的话：\BibleQuote{我将我的圣神倾注在一切血肉之人身上，你们的儿子和女儿要说预言，你们的老人要做梦，你们的青年要见神视。}\BibleRef{岳 2:28} 再听基督的话：\BibleQuote{但你们将由圣神降临在你们身上而领受能力，并要在耶路撒冷、犹太全地、撒玛黎雅，直到地极，为我作证人。}\BibleRef{宗 1:8} 诚然，宗徒的意思是，祂理当作为天主而被相信；然而，他并不以此为基础来肯定他的主张，而是如同审理与人有关的事件那样加以查验，正如他在致\WorkTitle{希伯来书}\BibleRef{希 6:18}中所说的那样，其中说道，\BibleQuote{藉着两件不可更改的事，天主决不可能说谎，我们便得到强有力的鼓励。} 同样在这里，他也以已赐予的恩惠作为将来应许的确实记号。为此，他进一步称之为\Quote{\OriginalTerm{凭据}{earnest}}\BibleRef{格后 1:22}，因为凭据是整体的一部分。祂已经赎买了我们所最关心的，就是我们的救恩；同时赐予我们凭据。为何祂不立即将全部赐下？因为我们在自己这方面也没有完成全部工作。我们相信了。这是一个开始；祂在祂那一方面也赐下了凭据。当我们以行为表明我们的信德时，祂必将其余部分加添。不仅如此，祂还赐下了另一个\Quote{\OriginalTerm{担保}{pledge}}，即祂自己的血，并应许了另一个。如同两国交战时交换人质一样：同样，天主也赐下祂的圣子作为和平与庄严盟约的\OriginalTerm{担保}{pledge}，此外还有出自祂的圣神。因为那些确实分沾圣神的人，知道祂是我们嗣业的\OriginalTerm{凭据}{earnest}。这样的人就是\Person{保禄}，他如今在此已预先尝到那里的福乐。这就是为何他如此热切，渴望脱离下界之事，并在内心叹息。他将整个心神转移到那里，以不同的眼光看待一切。你没有分沾这现实，所以不能理解这描述。倘若我们都按所应当分沾圣神，我们就将看到天堂，以及那里的事物的秩序。\par

然而，这是什么的凭据？是\par

\OriginalTerm{第14节}{Ver.~14}：\BibleQuote{救赎天主的子民}\BibleRef{弗 1:14}\par

因为我们的绝对救赎将在那时发生。如今我们在世上生活，易受许多人间变故，并与不敬神的人同居。但我们的绝对救赎将于那时实现，那时将没有罪恶，没有人间苦难，我们不再与各种人混杂在一起。\par

目前，却只有一个\OriginalTerm{凭据}{earnest}，因为目前我们远离这些福乐。然而，我们的国籍并非在地上；即使现在，我们已超越下界之事的范围。是的，我们甚至现在就是寄居者。\par

\OriginalTerm{第14节}{Ver.~14}：\BibleQuote{归于祂光荣的赞颂}\BibleRef{弗 1:14}\par

他紧接着加上这句。为什么呢？因为这将使听见的人获得完全的确证。他意思是，如果天主这样做只是为我们，或许会有疑虑。但若为祂自己的缘故，为彰显祂的良善，他就引出一个理由，作为某种见证，说明这些事绝不可能另有他途。我们在各处都看到同样的说法应用于以色列子民的身上。\BibleQuote{求你为了你的名为我们这样做，}\BibleRef{咏 108:21} 天主自己也说：\BibleQuote{我为了我自己而行，}\BibleRef{依 48:11} 梅瑟也说：\Quote{即使不为别的，也求你为你圣名的光荣这样做。} 这使听见的人获得完全的确证；他们得知，凡祂所应许的，为祂自己的良善，祂必最确实地实现。\par

伦理。——然而，不要让这听感使我们过分安逸；因为虽然祂为自己而行，但仍然要求我们尽本分。如果祂说：\BibleQuote{尊敬我的，我必尊重他；藐视我的，他必被轻视。}\BibleRef{撒上 2:30} 让我们反思，祂也向我们有所要求。诚然，拯救仇人是祂光荣的赞颂，但那些成为朋友后并继续作为祂朋友的人也是如此。因此，如果他们再次回到从前的敌对状态，一切都将徒劳无益。没有另一个圣洗，也没有第二次和好，只有\BibleQuote{一种等待审判的可怕景象，以及将吞灭敌人的烈火。}\BibleRef{希 10:27} 如果我们打算同时永远与祂为敌，却又向祂求宽恕，我们将永不停息地与祂为敌，放纵私欲，日益堕落，并对升起的\OriginalTerm{正义的太阳}{Sun of Righteousness}视而不见。难道你没有看见那当开你眼睛的光线吗？使它们良善、健全、敏锐吧。祂已向你展示了真光；如果你躲避它，又逃回黑暗，你的借口是什么呢？会有怎样的宽恕呢？从那一刻起，再也没有了。因为这标志着不可言喻的敌意。的确，当你不认识天主时，即使与祂为敌，无论怎样，你还有某种借口。但当你尝到良善与蜂蜜后，如果你又抛弃它们，转回到你的呕吐物，你除了证明极度的仇恨和藐视之外，还在做什么？\Quote{不，}你也许会这样说，\Quote{但我是被本性强迫的。我确实爱基督，但本性强迫我。}如果你是在被动和强迫下，你将得到宽恕；但如果你因懒惰而屈服，则一刻也不能。\par

那么，来吧，让我们探讨这个问题：罪究竟是强制与约束的结果，还是懈怠与极大疏忽的结果？法律说：\Quote{不可杀人。}这里有什么强制？有什么暴力？一个人必须使用暴力才能强迫自己去杀人，因为我们中间有谁会自愿将剑刺入邻人的喉咙，让双手沾满鲜血？没有一个人。那么，你看到，恰恰相反，罪更恰当地说是暴力和强制的事。因为天主在我们本性中植入了一种魅力，将我们彼此相爱地连结在一起。\BibleQuote{每种动物（圣经说）都爱自己的同类，每个人也爱自己的邻人。}\BibleRef{德 13:15} 你看到，我们从本性中获得了趋向德性的种子，而恶的种子却违背本性？如果后者占优势，这不过是我们极其懈怠的证明。\par

再者，什么是奸淫？有什么样的必要性引领我们犯此罪？无疑地，有人会说：情欲的暴政。但是，告诉我，为什么这样？难道每个人不能有自己的妻子，从而制止这种暴政吗？诚然，他会说，但我对邻人的妻子产生了一种热情。这里的问题不再关乎必要性。热情不是必然的事，没有人是出于必然去爱，而是出于审慎的选择和自由意志。本性的放纵或许可能是必然之事，但爱这一个女人而非另一个女人并非必然。你的问题不在于本性的欲望，而在于虚荣、放荡和无度的纵欲。因为哪一个是合乎理性的：一个男人应该有一位明媒正娶的妻子，她是他孩子的母亲，还是一位不被承认的女人？难道你不知道，亲密产生依恋吗？因此，这并非本性的过错。不要责怪本性的欲望。本性的欲望是为婚姻而赋予的；它是为生育子女而赐予的，并非为奸淫和败坏。法律也知道对那些出于必然的罪过予以宽恕——更确切地说，当罪过出于必然时，没有一件事是罪过，但所有的罪过都源于放荡。天主并没有将人的本性造得非犯罪不可，因为如果是这样，就没有惩罚可言。我们自己也不追究出于必然和强制所行的事，更何况满有慈悲和仁爱的天主呢？\par

再者，什么是偷窃？这是必然之事吗？是的，有人会说，因为贫穷导致它。然而，贫穷迫使我们工作，而非偷窃。因此，贫穷实际上产生相反的效果。偷窃是懒惰的结果，而贫穷通常产生的不是懒惰，而是对劳动的热爱。所以，这罪是懈怠的结果，你可以从这一点学到。请问，哪一个更困难、更令人厌恶：夜间不睡四处游荡，破门而入，在黑暗中行走，手中有性命，时刻准备杀人，因恐惧而颤抖、如同死了一般；还是从事日常的工作，充分享受安全和保障？后者更容易；正是因为更容易，所以大多数人都从事它而非前者。你看到，德性符合本性，恶违背本性，正如疾病与健康一样。\par

再者，什么是谎言和伪誓？它们可能暗示什么必要呢？完全没有，也没有任何强制；这是我们自愿进行的事。我们不被信任，有人会说。确实，我们不被信任，因为我们选择了这样。因为如果我们愿意，我们本可以因自己的品德而比因自己的誓言更受信任。请问，为什么我们不信任某些人，即使他们发誓，而认为另一些人即使没有誓言也值得信任？你看到，在任何情况下都不需要誓言。我们说：\Quote{当某人说话时，即使没有誓言，我也相信他；但你，即使有誓言，我也不信。}因此，誓言是不必要的；事实上，它更是不信任而非信任的证据。因为一个人过分轻易发誓，并不让我们对他的谨慎有很好的看法。所以，最常使用誓言的人，在任何场合都没有必要使用誓言；而从来不用誓言的人，本身充分享有使用誓言的益处。有人说，誓言是为了产生信任的必要；但我们看到，那些避免发誓的人更易被信任。\par

再者，如果一个人是暴力的，这是必然之事吗？是的，他会说，因为愤怒冲昏了他，在他里面燃烧，不让灵魂得安宁。人啊，行暴力不是愤怒的结果，而是心胸狭隘的结果。如果是愤怒的结果，那么所有人每当愤怒时，就会不停地施暴。我们被赋予愤怒，不是让我们对邻人行暴力，而是为了纠正那些犯罪的人，为了激励我们，使我们不懒惰。愤怒被植入我们心中，仿佛一种刺，让我们向魔鬼咬牙切齿，猛烈攻击他，而不是彼此对阵。我们有武器，不是让我们彼此相争，而是让我们全副武装对抗敌人。你容易发怒吗？那就对你的罪发怒：惩罚你的灵魂，鞭策你的良心，作严厉的审判官，对你的罪作出无情的判决。这是利用愤怒的方法。为此，天主把它植入我们内心。\par

然而，再问：掠夺是出于必需吗？绝对不是。告诉我，贪得无厌有什么必需？有什么迫不得已？有人会说，是贫穷造成的，是害怕缺少生活必需品。但这正是你不应贪婪的原因。如此得来的财富毫无保障。你的所作所为，正如一个人被问为何把房基安在沙土上时，回答说：\Quote{是因为霜冻和雨水}——而这恰恰是他不应安在沙土上的原因。正是那些雨水、暴风和风最迅速地掀翻这样的地基。所以，你若想富有，就绝不可贪婪；若想将财富传给子孙，至少要获取正义的财富——如果有这样的财富的话。因为这样的财富存留稳固，而不义的财富则迅速消耗消灭。告诉我，你想富有，就夺取他人之物吗？这当然不是财富：财富在于拥有属于自己的东西。占有他人之物的人，永远不能成为富有者；照此逻辑，连那些接收他人委托货物的丝绸商贩，也会成为最富有、最有钱的人了。虽然货物暂时属于他们，但我们并不称他们为富有。为什么呢？因为他们占有的是别人的东西。即使那匹丝绸本身属于他们，但它的价值却不属于他们。不，即使金钱在他们手中，这也不是财富。如果委托的货物因很快要交还而不能使人更富有，那么靠劫掠得来的东西又怎能使人富有呢？然而，你若真的渴望富有（此事并非必需），你究竟想享受什么更大的好处呢？是更长的寿命吗？可是，这种性格的人通常很快短命。他们往往以夭折作为掠夺和贪婪的惩罚；不仅因失去收益的享受而受罚，而且带着些许收益和地狱离开人世。他们也常因放纵、劳苦和焦虑而患病死去。我真想明白，为什么人类如此热切追逐财富。天主为我们的本性设定了界限和范围，正是为了让我们无需在此之外追求财富。例如，祂命令我们用一件或两件衣服遮蔽身体，再多一些就没有必要了。成千上万件衣服，还被虫蛀，有什么益处呢？胃有它的限定，超出此限的任何东西都必然毁掉整个人。那么，你的牛群、羊群和切割肉食又有什么用呢？我们只需要一个屋顶遮蔽。那么，你广阔的地基和豪华建筑又有什么用呢？你剥夺穷人，好让秃鹫和寒鸦有地方居住吗？这些事情该当何等地狱的惩罚？许多人常在大厦上建造闪亮的柱子和昂贵的大理石，而这些地方他们从未见过。有什么计划是他们没有采用的呢？然而，他们自己得不到益处，别人也得不到。荒凉使他们无法离开那里，但他们仍不罢休。你看，这些事并非为了利益，而是在所有这些事例中，愚昧、荒谬和虚荣是动机。我恳求你避免这些，好使我们也能避免其他一切邪恶，并得到那些爱祂的人所应许的美善事物，在我们的主耶稣基督内，愿光荣、力量、尊崇归于祂及圣父和圣神，直到永远。阿们。\par

\SourceRecord{Translated by Gross Alexander.；Nicene and Post-Nicene Fathers, First Series；Vol. 13.；Edited by Philip Schaff.；Buffalo, NY: Christian Literature Publishing Co.,；1889.；<http://www.newadvent.org/fathers/230102.htm>.}

% structure: p0001 source=h1 target=section reason=page-title

\section{厄弗所书讲道辞 第三篇}

第一章 第15-20节 \BibleRef{弗 1:15}\par

\Quote{因这缘故，我既听见你们在主耶稣里的信德，并向众圣人所显明的爱德，就不住地感谢你们，在祈祷中常常提到你们；愿我们的主耶稣基督的天主，即光荣的父，赐给你们智慧和启示的神，使你们认识祂；照明你们的心眼，使你们知道祂召叫的望德是什么，祂在圣徒中嗣业荣耀的丰盛是什么，以及祂对我们信者所显的能力是何等超越的浩大，是按照祂大能的运行，就是祂在基督身上所运行的，使祂从死人中复活。}\par

宗徒的渴望无与伦比，蒙福的保禄的同情和深情无与伦比，他为整个城市和人民做了每一次祈祷，并把同样的话写给了所有人：\Quote{我感谢我的天主，在祈祷中常常提到你们。}想象一下，他心中有多少人，光是记住他们就够费力的了；他在祈祷中提到了多少人，为所有人感谢天主，好像他自己得到了最大的祝福。\par

\Quote{因此，}他说，即因为将要来的事，因为为那些正确相信并生活的人所存的美好的事物。因此，合适的做法是感谢天主，既为人类从祂手中所领受的一切，无论是先前还是以后；也应为相信者的信德感谢祂。\par

\Quote{听说，}他说，\Quote{你们在主耶稣里的信德，并你们向众圣人所显明的。}\par

他在一切场合都将信德和爱德（这光荣的一对）结合在一起；他不只提到那个国家的圣人，而是所有圣人。\par

\Quote{我不住地感谢你们，在祈祷中常常提到你们。}\par

你的祈祷是什么，你的恳求是什么？是\par

\Quote{愿我们的主耶稣基督的天主，光荣的父，赐给你们智慧和启示的神。}他要求他们理解两件事，正如他们有责任理解的那样：他们被召叫到什么样的祝福，以及他们如何从以前的状态中释放出来。然而，他本人说这些点是三个。那怎样是三个呢？为了我们能够理解关于未来的事情；因为从为我们存留的美物中，我们将认识祂无可言喻和超越的丰富，而从理解我们曾经是谁以及我们如何相信，我们将认识祂的能力和主权，在使那些长期远离祂的人回心转意上，\Quote{因为天主的懦弱比人更刚强。}\BibleRef{格前 1:25}因为祂用来使基督从死者中复活的那同一力量，也吸引我们到祂自己那里。那力量并不限于复活，而是远超过它。\par

\Quote{并使祂坐在自己的右边，在天上，远超一切统治、权柄、能力和宰制，以及一切可称的名号；并将万有服在祂的脚下，使祂作教会万有之首，教会是祂的身体，是那充满万有者所充满的。}\par

祂使我们分享的奥秘和秘密真是浩大。若不分享圣神并领受丰厚的恩宠，我们不可能理解这些。因此保禄祈祷。\Quote{光荣的父，}即那赐给我们伟大祝福的那一位，因为他常根据所谈论的主题称呼祂，例如：\Quote{仁慈的父和赐与一切安慰的天主。}\BibleRef{格后 1:3}又如，先知说：\Quote{上主是我的力量和我的能力。}\BibleRef{咏 17:1}\par

\Quote{光荣的父。}\par

他没有名字可以用来代表这些事情，在一切场合他们都称之为\Quote{光荣，}这事实上在我们中间，是各种宏伟的名称和称呼。注意，他说，光荣的父\BibleRef{宗 7:2}；但基督的天主。那怎么样呢？圣子低于光荣吗？不，没有人，甚至疯子，会这样说。\par

\Quote{赐给你们，}\par

也就是说，可以提升和展开你们的悟性，因为否则无法理解这些事情。\Quote{因为属血气的人不接受天主圣神的事，因为在他看来，它们是愚拙。}\BibleRef{格前 2:14}因此，需要有属灵的\Quote{智慧，}好使我们理解属灵的事，看见隐藏的事。那圣神\Quote{启示}一切。祂将陈述天主的奥秘。然而，天主的奥秘的知识，只有圣神领悟，祂也探查天主深奥的事。不是说\Quote{天使、总领天使或任何其他受造的力量可以给与，}即把一种属灵的恩赐赐予你们。如果这是启示，那么论证的发现因而就是徒劳的。因为那学习了天主并认识天主的人，不再有争论。他不会说：这是不可能的，那是可能的，另一件事怎么发生的？如果我们认识天主，如同我们应该认识祂；如果我们从我们应当认识祂的那一位，即从圣神本身，认识天主；那么我们将不再有任何争论。因此他说：\par

\Quote{使你们的心眼得开，认识祂。}\par

\Quote{那学习了天主是什么的人，不会对祂的应许有疑虑，也不会对已经成就的事不信。他祈求，在那时赐给他们\InnerQuote{智慧和启示的神。}然而，他也尽力用论证和已存的事实来证实它，从\InnerQuote{已经}存在的事实。因为，他即将提到一些已经成就的事，和一些尚未发生的事；他用已经成就的事作为未成就之事的保证：意思大概是这样：}\par

\Quote{使你们知道，}他说，\Quote{祂召叫的望德是什么。}\par

他的意思是，这尚未显明，但对信者并非如此。\par

\Quote{并且}，又\Quote{祂在圣人们中所得的基业，那荣耀的丰盛是什么。}\par

这也仍是隐藏的。\par

但什么是清楚的呢？那就是我们因祂的能力相信了祂使基督复活了。因为说服灵魂，远比复活尸体更是奇迹。我愿将此阐明。请听。基督对死者说：\Quote{拉匝禄，出来！} \BibleRef{若 11:43}，他立刻顺从了。伯多禄说：\Quote{塔彼达，起来！} \BibleRef{宗 9:40}，她也没有拒绝。在末日，祂要亲口发言，众人都要复活，且是那么迅速，以致\Quote{那些还活着的人，决不会在那些已经长眠的人之先} \BibleRef{得前 4:15}，一切都将发生，一切都将同时成就，\Quote{在一刹那，一眨眼之间} \BibleRef{格前 15:52}。但在信仰之事上，并非如此；那么是怎样呢？请再听祂所说：\Quote{我多少次愿意聚集你的子女，但你们却不愿意。} \BibleRef{玛 23:37}。你看出后者更为困难。因此，他以此构建了整个论证；因为按人的推断，影响人的选择远比改变自然更难。原因在于，祂要我们自愿成为善人。所以他有理由说：\par

\Quote{祂的能力向我们这信的人所显的浩大。}\par

是的，当先知、天使、总领天使都无能为力，当整个可见与不可见的受造物都失败时（可见的摆在我们面前，却无力引导我们，还有相当多不可见的也是如此），祂就安排了自己的来临，以向我们显示这是一件需要天主能力的事。\par

\Quote{荣耀的丰盛。}\par

也就是说，那不可言喻的荣耀；因为有什么语言足以表达圣人们那时将要分享的荣耀呢？没有。但确实需要恩宠，好使理智能够领悟它，哪怕只接纳一点微弱的光。有些事他们以前已经知道；现在他希望他们学习更多，更清楚地知道。你看到祂行了何等大事？祂使基督复活了。这是小事吗？再看。祂使祂坐在自己右边。还有什么语言能描述这个？那出自尘土、比鱼更沉默、被魔鬼戏弄的，祂在一瞬间高举到高处。这真是\Quote{祂能力的浩大}。看，祂把祂高举到何处。\par

\Quote{在天上。}\par

祂使祂远超一切受造的本性，远超一切统治者和掌权者。\par

\Quote{远超一切统治者}，他说。\par

因此，实在需要圣神，需要明智认识祂的智慧。实在需要启示。思想一下，人的本性距离天主有多么遥远。然而，祂却从这卑贱的地位高举祂到那样的高位。祂不是逐级上升，先一级，再一级，然后第三级。令人惊叹！祂没有简单地说\Quote{以上}，而是\Quote{远超}；因为天主高于那些崇高的权能。然后祂把祂——我们中的一位——从最低点带到至高主权，带到那无可超越的尊荣之上，高于\Quote{一切}统治权，他说，不是指一个而非另一个，而是高于所有：\par

\Quote{统治者、掌权者、有能者、主治者，以及一切有名可称的。}\par

天上无论有什么，祂都已成为万有之上。这是指那位从死者中复活的那一位说的，这是值得赞叹的；因为论及天主圣言，那是不可能的，因为万物在祂面前如同昆虫之于人。如果全人类都被视为唾沫，被视如天平上的灰尘，那么无形权势也如同昆虫。但对我们中的那一位来说，这确实是伟大而惊人的。祂从地下最深处高举了祂。如果万民如同一滴水，那么一个人在那滴水中是何等微小！然而，祂使祂超越万有，\Quote{不但是今世的，连来世的也一概超过了。}因此，有些权能的名字是我们无法理解、不得而知的。\par

\Quote{并使万有屈服于祂的脚下。}\par

不仅使祂高居其上而受尊敬，也不是通过比较，而是使祂坐在它们之上如同奴隶。令人惊叹！这些事确实可畏；因着住在祂内的天主圣言，每一个受造的权能都成了人的奴隶。因为一个人可以高于别人，而不使别人屈服，只是优先于他们。但这里并非如此。不，\Quote{祂使万有屈服于祂的脚下。}而且不仅仅是使之屈服，而是最彻底的屈服，其下无物。因此他加上\Quote{在祂的脚下。}\par

\Quote{并使祂为教会作万有之首。}\par

再次令人惊叹，祂把教会提升到了何处？仿佛用某种机械举起它，祂把它举到极高之处，并安置在那宝座上；因为头在哪里，身体也在哪里。头和身体之间没有间隔；若有间隔，就不再是身体，也不再是头了。\Quote{万有之上，}他说。\Quote{万有之上}是什么意思？祂不允许天使、总领天使或任何其他存在高于祂。但祂不仅以此方式荣耀了我们——高举了我们自身的那一位——还预备了整个族类共同跟随祂、依附祂、伴随祂的侍从。\par

\Quote{就是祂的身体。}\par

为使你听见"头"时，不仅想到至高权威，也想到结合，并看到祂不仅是至高统治者，也是身体的头。\par

\Quote{那充满万有者所充满的}，他说。\par

好像这还不足以表明密切的关系和联系，他又加上什么？\BibleQuote{基督的圆满就是教会}。而且正确，因为头的补充是身体，身体的补充是头。请注意\Person{保禄}观察何等伟大的安排，他如何不吝惜任何一个词，以便能表达天主的荣耀。\Quote{补充}，他说，即头仿佛被身体所充满，因为身体是由其各个部分组合而成的，他引入祂需要每一个个体，而不仅仅是所有共同和一起；因为除非我们众多，一人是手，另一人是脚，另一人是别的肢体，整个身体就不被充满。因此，是由众人使祂的身体充满。那时头被充满，那时身体得以完美，当我们全都结合在一起、联合在一起时。你岂不明白了\BibleQuote{祂继承光荣的丰厚？祂对信仰者所施展的极强大能？你们蒙召的盼望？}\par

道德。让我们尊敬我们的头，让我们思考我们是何等样的头的身体——一个万有都屈服于其下的头。按照这个表征，我们应当比天使更优越，是的，比总领天使更伟大，因为我们被尊荣超过他们全体。天主\BibleQuote{没有握住天使}，正如他在\BibleBook{}{希伯来书}中所写的，\BibleQuote{却握住了\Person{亚巴郎}的后裔}。\BibleRef{希 2:16} 祂没有握住任何宰制者或掌权者、统治权或其他权柄，却拿起了我们的本性，并使之坐在祂的右边。我为何说，使之坐下？祂使之成为祂的衣服，不仅如此，还将万有置于祂的脚下。你们以为有多少种死亡？多少灵魂？一万？是的，一万个一万次，但你们将说没有什么能与之相比。祂做了两件事，最大的事。祂既亲自降卑至最低的羞辱深处，又将人提升至崇高的荣耀高处。祂以自己的血拯救了我们。祂先前首先说了前者，即祂如何如此极大地谦卑自己。现在祂说那比那更强的事——一件大事，万事的冠冕。的确，即使我们被认为不配得任何东西，那也是足够的。或者，即使我们被认为配得这尊荣，即使没有圣子的被杀，那也是足够的。但这里是两者兼有，什么语言能力能超越并胜过它呢？复活本身并不伟大，当我思想这些事时。祂是指着祂说\BibleQuote{我们的主\Person{耶稣基督}的天主}，不是指着天主圣言。\par

让我们对我们关系的密切感到敬畏，让我们恐惧，免得有人从这身体上被切断，免得有人从它堕落，免得有人显得不配它。若有人将冠冕戴在我们头上，一顶金冠，我们岂不尽力使自己配得上那无生命的珠宝？但现在，不是冠冕在我们头上，而是，远为伟大者，基督成了我们的头，我们却毫不在意。然而，天使尊敬那头，总领天使和所有那些天上的权能也尊敬。而我们，是祂的身体，难道既不因此事也不因那事而敬畏吗？那么，我们的救恩希望何在？设想那王座，设想荣誉的过度。这，至少我们若选择，更能使我们震惊，是的，甚至比地狱本身更甚。因为，即使没有地狱，我们被授予如此尊荣却显出卑劣和不配，这岂不带来何等惩罚、何等报复？想想你们的头坐在谁附近，（这个单独考虑对任何目的都充分足够），祂在谁的右边，远超一切宰制者、掌权者和异能者。然而，这个头的身体竟被魔鬼践踏。不，愿天主阻止如此；因为如果是这样，这样的身体就不能再是祂的身体了。你们自己的头，你们较尊敬的仆人尊敬它，而你们却使你们的身体成为侮辱它者的玩物？那么，你们该受何等严厉的惩罚？若有人将皇帝的双脚用绳索和锁链捆绑，他岂不将受极刑？你们把整个身体暴露给凶猛的野兽，却不战栗吗？\par

然而，既然我们的讲论是关于主的身体，来，让我们把思想转向它，甚至那被钉十字架、被钉、被祭献的身体。如果你们是基督的身体，就背起十字架，因为祂背起了它：忍受唾弃，忍受掌掴，忍受钉子。那身体就是如此；那身体\BibleQuote{没有犯罪，口中也没有诡诈}。\BibleRef{伯前 2:22} 祂的手做了为需要者利益的一切事，祂的口没有说不适当的话。祂听见他们说\Quote{你附了魔}，祂什么也没回答。\par

此外，我们的讲论是关于这身体，我们中凡领受那身体并品尝那血的人，都在品尝那与那身体毫无不同、也不分离的东西。请思想，我们品尝那坐在高处、被天使朝拜、紧邻那不腐坏之权能的身体。哀哉！有多少得救的途径向我们敞开！祂使我们成为祂自己的身体，祂将自己的身体分赐给我们，然而这些事没有一件使我们远离邪恶。哦，黑暗，深渊之深，麻木！\BibleQuote{你们要思念}，他说，\BibleQuote{上面的事，那里有基督坐在天主的右边。}\BibleRef{哥 3:1} 而在这之后，有人仍贪恋金钱或放纵，有人被情欲所掳！\par

你们岂不见，在我们自己的身体上，任何多余无用的部分，就被割下、被切除吗？即使它曾属于身体，一旦残缺、坏死、腐朽、损害其余部分，便毫无用处。那么，我们切莫因为曾一度成为这身体的肢体，就过分自信。如果我们的身体，虽是自然之体，尚且遭受截肢，那么当道德原则败坏时，它岂不要遭受更可怕的灾祸吗？当身体不领受这自然的食物，当毛孔堵塞，它便坏死；当通道闭合，它便瘫痪。我们在灵性上也是如此：当我们塞住耳朵，灵魂便瘫痪；当我们不领受属神的话语，当我们因邪恶的性情而不是腐败的体液而受害时，这一切都会引发疾病、危险的疾病、消耗的疾病。那时便需要那火，需要那割裂。因为基督不能容忍我们带着这样的身体进入婚宴厅。若祂把穿着污秽衣服的人带走并赶出去，祂对那将污秽附着于身体的人，将不会做什么呢？祂将如何处置他呢？\par

我见许多人领受基督的圣体，轻率而随意，更多的是出于习惯和形式，而非出于思考和领悟。有人说：\Quote{当神圣的四旬期来临时，无论何人，都可领受圣体圣事；或当主显节之日到来时。}然而，使人适宜领受的，并非主显节或四旬期，而是灵魂的真诚与纯洁。有此，随时可领受；无此，决不可。\Quote{因为你们每次吃这饼，喝这杯，} \BibleRef{格前 11:26} 他说，\Quote{你们就是宣告主的死亡}，即\Quote{你们纪念为你们所完成的救恩，以及我所赐的恩惠。}请思量那些在旧约下领受祭献的人，他们实行何等大的节制？他们如何行事？他们做了什么？他们总是洁净自己。而你，当接近一个连天使都战栗的祭献时，你竟以季节的更替来衡量此事吗？你怎能出现在基督的审判座前，竟以污秽的手和口妄用祂的身体？你不会用不洁的嘴去亲吻国王，却以不洁的灵魂亲吻天国的君王吗？这是对祂的侮辱。告诉我，你岂愿以未洗的手来接近祭献？我想不会。你宁可不来，也不愿以污秽的手来。那么，你在这小事上如此谨慎，却以污秽的灵魂前来，竟敢触摸祂？手只暂时拿住它，而灵魂却完全吸收它。难道你没看见那些圣器被彻底清洗，如此光彩夺目吗？我们的灵魂应当比它们更洁净、更圣洁、更明亮。为何如此？因为那些器皿是为我们而造的。它们不领受其中的祂，也不认识祂。但我们却确实领受。那么，你不愿使用污秽的器皿，却以污秽的灵魂前来吗？注意这极大的不一致。其他时候你不来，即使你常常是洁净的；但在复活节，无论你犯了何等重罪，你却来了。哦！习惯和偏见的威力！每日的祭献是徒然的，我们每日站在祭台前也是徒然；没有人领受。我说这些话，不是要你们随便领受，而是要你们使自己配得领受。你不配得祭献，也不配得参与吗？若是如此，你也不配得祈祷。你听见传令员站着说：\Quote{凡在补赎中的人，都祈祷。}凡不领受的人，就在补赎中。若你是那些在补赎中的人之一，你就不该领受；因为那不领受的，就是那些在补赎中的人。那么，为什么祂说\Quote{不配祈祷的人，出去吧}，而你竟厚颜无耻地站在那里？但不对，你不属于那类，你属于那些配得领受的人，却对此漠不关心，视为无足轻重。\par

看，我恳求：一张君王的筵席摆在你面前，天使在那筵席上服务，君王亲自在那里，你却张着嘴呆立？你的衣服污秽了，你竟不以为意？还是洁净的？那就俯伏领受吧。祂每天进来察看宾客，与众人交谈。是的，此刻祂正对你的良心说话：\Quote{朋友，你为何站在这里，没有穿婚宴礼服？}祂没有说：你为什么坐下了？不，在他坐下之前，祂就宣告他不配进来。祂不说：\Quote{你为什么坐下吃？}而是\Quote{你为什么进来？}这正是此刻祂对我们每一个如此厚颜无耻站在这里的人所说的话。因为凡不领受圣体圣事的人，都是厚颜无耻地站在这里。正因如此，那些在罪中的人首先被赶出去；正如当主人出席他的筵席时，那些得罪了他的仆人不应在场，而被遣走一样：同样，当祭献被呈上，基督——主的羔羊——被祭献时，当你听见\Quote{让我们一同祈祷}的话，当你看见幔子拉开，那时你当想象天从上面降下，天使正在降临！\par

既然不适宜有未领受入门圣事者在场，同样也不适宜有已领受入门圣事却同时被玷污者在场。请告诉我：若有人受邀赴宴，洗了手，坐下，一切就绪，却拒绝用膳，难道这不是在侮辱邀请他的人吗？这样的人，还不如根本不来！现在你正是这样来到这里的。你与众同唱了赞美诗，宣称自己属于配得之人的行列——因为你没有与不配的人一同离开。为何留下，却不用膳呢？你会说：\Quote{我不配。}那么，你也不配在祈祷中有份，因为圣神不仅藉着祭品降临，也藉着这些圣歌环绕我们。难道我们没有看见我们的仆人先用海绵擦拭桌子、打扫房屋，然后摆上宴席吗？祈祷和宣报者的呼声所做的正是这事。我们仿佛用海绵擦净了教会，好让一切都在纯洁的教会中摆设好，使教会\BibleQuote{没有斑点，没有绉纹}\BibleRef{弗 5:27}。的确，我们的眼睛不配看这些景象，我们的耳朵也不配听！经上说：\BibleQuote{连走兽若触到这座山，也该用石头砸死。}\BibleRef{出 19:13} 那时他们连踏上这山都不配，但后来他们却走近了，看见了天主所站之处。将来你也能走近、看见；然而，当主临在时，你却离开。你在此的资格并不比慕道者高。因为从未到达奥迹与到达后却绊倒、轻视它们、使自己不配于此，二者截然不同。还可以讲更多、更令人敬畏的要点，但为了不加重你的理解负担，这些已经足够。连这些都不能使一个人恢复理智，那么更多的也无济于事。\par

为了避免我成为加深你们定罪的原因，我恳求你们：不要停止前来，而要努力使自己配得在场并接近。请告诉我：若有一位国王下令说：\Quote{凡做这事的人，可享用我的宴席}，你们难道不会竭尽全力争取被接纳吗？祂已召叫我们赴天国、伟大奇妙之王的宴席，我们却退缩犹豫，而不是急切奔跑前去？那我们的救恩指望又在何处？我们不能归咎于软弱，也不能归咎于本性。是懈怠——唯有懈怠——使我们不配。\par

以上是我自行讲述的。愿那刺透人心、赐予痛悔之圣神的那一位，刺透你们的心，将种子深植其中，好使你们因敬畏祂而怀胎，生出救恩之圣神，坦然无惧地亲近。因为经上说：\BibleQuote{你的儿女围绕你的桌子，好像橄榄栽子。}\BibleRef{咏 127:3} 哦，但愿没有老旧的、野性的、粗糙的！因为适于结出好果子——我指的是橄榄树的好果子——的幼树正是如此。它们茁壮生长，以致全都围绕在桌旁，一同聚集于此，并非徒然或偶然，而是怀着敬畏与虔敬。因为这样，你们将来在天上也将坦然无惧地瞻仰基督自己，并被算为配得那天国。愿我们众人赖我们主耶稣基督，蒙天主恩赐，都能获得，祂与圣父及圣神，是光荣、权能、尊威，自今而永远，及世之世。阿们。\par

\SourceRecord{Translated by Gross Alexander.；Nicene and Post-Nicene Fathers, First Series；Vol. 13.；Edited by Philip Schaff.；Buffalo, NY: Christian Literature Publishing Co.,；1889.；<http://www.newadvent.org/fathers/230103.htm>.}

% structure: p0001 source=h1 target=section reason=page-title

\section{厄弗所书讲道集 第四讲}

第二章 第1-3节 \BibleRef{弗 2:1}\par

\Quote{祂使你们活过来，你们在过犯和罪恶中死了，那时你们在世风习尚中行走，顺从这世界的首脑，就是空中权势的首领，即现今在悖逆之子身上发生作用的邪神。}\par

我们知道，有肉体的死亡，也有精神的死亡。参与第一种死亡并非罪恶，其中也无危险，因为它不带有任何罪责，它是出于本性，而非出于自主的选择。它起源于第一个人的过犯，此后在繁衍中传给了整个人性，并且无论如何很快会结束；而精神的死亡则是出于自主的选择，带有罪责，且永无止境。请注意，保禄已经指出这是一件极其重大的事，以至于使一个死去的灵魂复活远远大于使死人复活，所以他如今再次以其全部的真实伟大性来阐述它。\par

\Quote{并你们，}他说，\Quote{当你们在过犯和罪恶中死了的时候，那时你们在世风习尚中行走，顺从这世界的首脑，就是空中权势的首领，即现今在悖逆之子身上发生作用的邪神。}你们注意保禄的温和，他处处鼓励听众，不过分严厉。因为他说，你们已达到了邪恶的极致（死亡的含义便是如此），为了不使他们过分悲伤（因为当过去的恶行被揭露时，即使已被赦免且不再有危险，人仍会感到羞耻），他给他们提供了一个同谋，以免他们认为这完全是自己的作为，而且这个同谋是强有力的。那么是谁呢？魔鬼。他在致格林多人书中也做了类似的事，在说了\Quote{你们不要受骗：无论是淫乱的、拜偶像的} \BibleRef{格前 6:9}并列举了所有其他恶行，并最后加上\Quote{都不能承继天主的国}之后，他接着又说\Quote{你们中有些人也是如此}；他没有绝对地说\Quote{你们曾是}，而是\Quote{你们中有些人曾是}，也就是说，你们在某种程度上曾是。这里异端攻击我们。他们告诉我们，这些说法（\Quote{空中一切权势的首领}等）是针对天主的，他们放纵其不羁的舌头，把这些属于魔鬼的东西加于天主。那么我们如何使他们沉默呢？凭他们自己所用的话：因为，如果他是正义的，如同他们自己所承认的，而他却做了这些事，那就不再是正义者的行为，而是最不义、最败坏者的行为；败坏的天主不可能存在。\par

再者，为什么称魔鬼为这世界的\Quote{首脑}？因为几乎全人类都臣服于他，所有人都是自愿并自主地做他的奴隶。而对于基督，虽然祂应许了无数的祝福，却无人听从；而对于魔鬼，虽然祂什么也没有应许，反而送他们下地狱，所有人都屈服于他。因此，他的王国就在这个世界里，并且除了少数人外，他拥有比天主更多的、更顺从的臣民，这是由于我们的懈怠。\par

\Quote{顺从这权势，}他说，\Quote{就是空中的、邪神的权势。}\par

这里他再次说明，撒殚占据了天空之下的空间，而无形的力量是空中的邪神，受他操纵。因为他的王国是属于这个时代的，即会随现世而终止，听他在书信末尾所说：\Quote{我们的战斗不是对抗血肉，而是对抗率领者，对抗掌权者，对抗这黑暗世界的霸主} \BibleRef{弗 6:12}；这里，为了不让你听到\Quote{世界霸主}而说魔鬼不是受造的，他在别处 \BibleRef{迦 1:4} 称歪曲的时代为\Quote{邪恶的世界}，不是指受造物。因为在我看来，他在天空之下拥有统治权，即使在他过犯之后也没有失去统治权。\par

\Quote{现今在悖逆之子身上发生作用的，}他说。\par

你们注意到，他赢得我们不是靠暴力或强迫，而是靠劝说；\Quote{悖逆}或\Quote{不顺服}是他的用词，仿佛是说，他用诡计和劝说将他所有的追随者吸引到自己身边。他不仅通过告诉他们有一位同谋来鼓励他们，还通过将自己与他们并列，因为他说：\par

\Quote{我们从前也都在他们中间生活过。}\par

\Quote{都}，因为他不能说任何人被排除在外。\par

\Quote{放纵肉体的私欲，随从肉体和心意的欲望，生来就是义怒之子，和别人一样。}\par

也就是说，没有属灵的情感。然而，以免他诽谤肉体，或以免有人认为过犯不大，请看他如何谨慎处理此事：\par

\Quote{随从，}他说，\Quote{肉体和心意的欲望。}\par

也就是说，享乐的私欲。他说，我们激起了天主的愤怒，我们惹动了祂的义怒，我们就是义怒本身，别无其他。因为正如人的儿子因本性是人一样，我们生来也是义怒之子，和别人一样；也就是说，没有人是自由的，我们都做了值得义怒的事。\par

第4节 \Quote{然而天主因祂丰富的仁慈。} \BibleRef{弗 2:4}\par

不仅是仁慈，而且是丰富的仁慈；正如在另一处所说：\Quote{祢丰厚的仁慈} \BibleRef{咏 68:17}。又说：\Quote{求祢按祢丰厚的仁慈怜悯我} \BibleRef{咏 50:1}。\par

第4节 \Quote{因祂爱我们的大爱。} \BibleRef{弗 2:4}\par

祂为什么爱我们？因为这些事不值得爱，而值得最严厉的义怒和惩罚。因此，这是极大的仁慈。\par

第5节 \Quote{当我们因过犯而死的时候，祂使我们同基督一起活过来。} \BibleRef{弗 2:5}\par

基督再次被引入，这是一件深值我们信服的事，因为若初果活着，我们也同样活着。祂复活了祂，也复活了我们。你看见这一切都是论及降生成人的基督吗？你岂不见\Quote{祂对我们相信的人所施展的至大权能}\BibleRef{弗 1:19}？那些死了的，那些本是义怒之子的，祂使他们复活了。你岂不见\Quote{祂召叫的望德}吗？\par

第6节。\Quote{祂使我们与祂一同复活，并使我们与祂一同坐在天上。}\BibleRef{弗 2:6}\par

你岂不见祂嗣业的荣耀吗？\Quote{祂使我们一同复活}是明显的。但祂\Quote{在基督耶稣内，使我们与祂一同坐在天上}如何成立？这就像祂使我们一同复活一样真实。因为如今实际上还没有人复活，只不过因元首已复活，我们也复活了，正如历史上，若瑟的梦（\BibleRef{创 37:9}）中，当雅各伯下拜时，他的妻子也向若瑟下拜。同样，\Quote{祂也使我们与祂一同坐下}。因为既然元首坐着，身体也与祂同坐，因此他加上\Quote{在基督耶稣内}。或者，若这句话不是这个意思，那就是指藉着圣洗的水，祂\Quote{使我们与祂一同复活}。那么在这种情况下，祂如何使我们\Quote{与祂一同坐下}呢？因为，他说，\Quote{如果我们与祂一同受苦，也必与祂一同为王}\BibleRef{弟后 2:12}；如果我们与祂一同死，也必与祂一同活。要理解这些奥迹的深度，确实需要圣神和启示。然后，为使你对此事不存任何怀疑，请看他接着补充的话。\par

第7节。\Quote{为在将来的世代，显示祂在基督耶稣内，向我们所施的恩宠的至极丰富，和祂对我们的慈惠。}\BibleRef{弗 2:7}\par

他刚才讲论的是关于基督的事，这些事可能与我们无关（因为有人会说，基督复活与我们何干？），因此他说明这些事也确实关系到我们，因为祂与我们成为一体。只是他把我们这方面的事分开论述。他说：\Quote{我们这些因过犯而死了的，祂使我们与祂一同复活，并使我们与祂一同坐在天上。}所以，正如我所说的，不要不信，接受他从过去的事、从祂的元首地位，以及从祂愿意显示祂的慈惠所提供的证明。因为若不实现这些，祂将如何显示呢？而且祂要在将来的世代显示。什么？就是这些福泽既伟大，又比任何事物更为确实。因为如今所说的话，在不信者看来可能像是愚妄；但那时所有人都要明白。你还愿明白祂如何使我们与祂同坐吗？请听基督自己对门徒说的话：\Quote{你们也要坐在十二个宝座上，审判以色列十二支派}\BibleRef{玛 19:28}。又说：\Quote{但坐在我的右边和左边，不是我可以赐的，而是我父为谁预备了，就赐给谁}\BibleRef{玛 20:23}。所以这事已预备好了。他说的\Quote{在基督耶稣内，向我们施的慈惠}很有道理，因为坐在祂右边是超越一切荣耀的荣耀，是无可比拟的。他说，连我们也要坐在那里。这真是超越的富足，真是祂力量的极大施展，竟使我们与基督同坐。是的，即使你有一万个灵魂，难道不该为祂的缘故舍弃吗？是的，即使你须投入火焰，难道不该欣然忍受吗？祂自己也说：\Quote{我在哪里，我的仆人也将在那里}\BibleRef{若 12:26}。岂不更应如此吗？即使你们每天被肢解，难道不该为这些应许而欣然接受吗？想想，祂坐在哪里？在一切执政者和掌权者之上。和你同坐的是谁？是和祂。你是谁？一个已经死了的、生来就是义怒之子的人。你行了什么善事？没有。现在确实是高呼的时候：\Quote{啊，天主的智慧和知识的富厚，多么深奥！}\BibleRef{罗 11:33}\par

第8节。\Quote{因为你们得救是由于恩宠。}\BibleRef{弗 2:8}\par

为使所施恩惠的伟大不致过分抬高你，请看他是如何贬抑你：\Quote{你们得救是由于恩宠}，他说，\par

\Quote{藉着信德；}\par

然后，为使另一方面，我们的自由意志不受损害，他又加上了我们在这工程中的一部分，然而又取消了它，并加上，\par

\Quote{而且这并非出于你们自己。}\par

他意思是说，连信德也\Quote{并非出于你们自己}。因为祂若不来，若没有召叫我们，我们怎能相信呢？因为\Quote{他们若没有听到，怎能相信呢？}\BibleRef{罗 10:14}所以连信德本身的工作也不是我们自己的。\par

\Quote{是天主的恩赐}，他说，\Quote{不是出于功行。}\par

那么，你会说，信德足以拯救我们吗？不；但天主，他说，要求这一点，免得祂拯救我们时，我们是荒芜且毫无工作的。祂的意思是，信德拯救，但这是由于天主愿意，信德才拯救。因为，请问，信德如何在没有功行的情况下拯救？这本身就是天主的恩赐。\par

第9节。\Quote{为叫人没有可夸的。}\BibleRef{弗 2:9}\par

为在我们心中激起对这恩宠恩赐的恰当情感。\Quote{那么，}有人说，\Quote{难道祂自己阻碍了我们因功行而成义吗？}绝不然。但他说，没有人因功行而成义，为的是显示天主的恩宠和慈爱。祂并非因为我们有功行而拒绝我们，而是当我们无工可施时，祂以恩宠拯救了我们，使从此无人有可夸之处。然后，免于当你听到整个工程不因功行而因信德完成时，你变得懒散，请看他是如何继续：\par

第10节。\Quote{因为我们是祂的工艺品，在基督耶稣内受造，为行天主预备的善功，使我们在此中行走。}\BibleRef{弗 2:10}\par

请注意他所用的词句。他在这里暗指重生，这实际上是第二次创造。我们从未存在中被带到了存在。至于我们从前是什么，即旧人，我们已经死了。我们现在成为的，从前我们并非如此。因此，这项工作的确是一个创造，而且比第一次更为尊贵；因为从第一次创造，我们得以存在；但从这最后一次，我们更获得了善在。\par

\Quote{因为天主预备了善工，叫我们行在其中。}\par

不仅是要我们开始，更是要我们行在其中，因为我们需要一种持久的德行，延续到去世之日。假如我们要走一条通往皇城的道路，走过了大部分路程，却在快到终点时疲惫坐下，这对我们毫无益处。这就是我们蒙召的望德；他说：\Quote{为善工}。否则，对我们毫无益处。\par

道德教训。因此，他喜悦的不是我们只做一件善工，而是做一切善工；因为正如我们有五种感官，并应当适时使用它们，同样我们也应当运用各种德行。假如一个人节制却不慈悲，或慈悲却贪婪，或确实不拿别人的东西，却不施舍自己的，这一切都是枉然。因为单靠一种德行不足以使我们在基督的审判台前具有胆量；不，我们需要德行伟大、多样、全面且完整。请听基督对门徒所说的话：\Quote{你们要去使万民成为门徒，教导他们遵守我所吩咐你们的一切。}\BibleRef{玛 28:19} 又说：\Quote{无论谁废除这些最小诫命中的一条，在天国里要称为最小的。}\BibleRef{玛 5:19} 那是指复活的时候；不仅如此，他不能进入天国；因为主也惯称复活的时候为天国。\Quote{若废除一条，}祂说，\Quote{就称为最小的。} 因此我们需要一切德行。请注意，没有慈悲之工就不能进入天国；但即使只有这一项缺乏，我们也要被投入火中。因为祂说：\Quote{可咒骂的，离开我，到那给魔鬼和他的使者预备的永火里去！} 为什么呢？\Quote{因为我饿了，你们没有给我吃的；我渴了，你们没有给我喝的。}\BibleRef{玛 25:42} 你看，没有其他罪状，仅因这一项他们就灭亡了。也仅因这一项，那些童女也被排除在婚筵之外，尽管她们确实拥有清醒。正如宗徒所说：\Quote{并要追求圣化，没有它，没有人能见到主。}\BibleRef{希 12:14} 那么要思考，没有清醒就不可能见到主；但有了清醒也未必就能见主，因为常有别的事情阻挡。再者，如果我们凡事做得完全正确，却不为邻人服务，那么同样不能进入天国。这从何而知？从那比照塔冷通的仆人的比喻可知。因为在那个例子中，那人德行各方面都完整，一无所缺，但因为他在事业上懒惰，就被正当地驱逐出去。不，甚至只凭辱骂也可能坠入地狱。\Quote{因为凡}基督说，\Quote{对他的兄弟说\InnerQuote{你是愚人}的，就要受地狱之火的刑罚。}\BibleRef{玛 5:22} 如果一个人凡事都正确，却伤害他人，他也不能进入。\par

且莫将残忍归于天主，因祂将在此事上失败者排除于天国之外。因为在人间，若有人做任何违反律法的事，他便被逐出王前。若他违犯既定的法律之一，若他诬告他人，他便失去他的职位。若他犯奸淫，且被察觉，他便蒙受耻辱，纵然行了万千善行，他也算完了；若他杀人，且被定罪，这也足以毁灭他。今若人间的法律尚且如此严守，天主的法律岂不更当如此。有人却说：\Quote{但祂是良善的。}我们还要说这愚昧之言多久？我说愚昧，不是因为祂不善，而是因为我们一直以为祂的良善会为我们这些目的而效力，尽管我已再三用万千论据述说此事。请听圣经所说的：\BibleQuote{不要说，祂的仁慈广大，祂必因我众多罪恶而平息。}\BibleRef{德 5:6}祂并不禁止我们说：\BibleQuote{祂的仁慈广大。}这不是祂所吩咐的；相反，祂要我们不断说此言，为此保禄举出种种论据，但他的目的是下面的话。他的意思是，不要存此念想而仰慕天主的慈爱，为了犯罪，并说：\BibleQuote{祂的仁慈必因我众多罪恶而平息。}因为我也正是为此目的而讲论祂的良善如此之多，不是让我们仗恃它而任意妄为，因为那样这良善将有损于我们的救恩，而是让我们不因自罪而绝望，反而悔改。因为\BibleQuote{天主的慈爱领你悔改}\BibleRef{罗 2:4}，而非更大的邪恶。若你因祂的良善而堕落，你反而在人面前否认祂。我看见许多人如此诋毁天主的宽容；因此你若不善用它，你必受惩罚。天主是慈爱的天主吗？是的，但祂也是公义的审判者。祂是宽容罪恶的吗？诚然，但祂也要照各人的行为施行报应。祂是放过不义、涂抹过犯的吗？诚然，但祂也追究。那么，这些事怎能不矛盾？它们并不矛盾，若我们按时间区分它们。祂在此处消除不义，借着圣洗之泉和忏悔。祂在彼处以火与痛苦追究我们所行的。有人或许说：\Quote{既然如此，我被驱逐，丧失天国，无论我行恶行万千或仅一次，为何我不能行各种恶事呢？}这是忘恩仆人的论调；尽管如此，我们仍将继续解答这一点。绝不可为行善而作恶；因为我们都将同样失落天国，但在地狱中，我们将不会都受同样的刑罚，而是一人更严厉，另一人更轻缓。因为现今，若你与另一人\BibleQuote{藐视天主的慈爱}\BibleRef{罗 2:4}，一次在多次，另一次在少数，你们将同样丧失天国。但若你们并非同样藐视祂，而是一人在更大程度上，另一人在较小程度上，那么在地狱中，你们将感受到差别。\par

那么，有人或问，祂为何要威胁那些没有行怜悯事工的人，他们要去火里，且不仅是火里，而是那\BibleQuote{为魔鬼及其天使预备的}\BibleRef{玛 25:41}火里？这是为何？因为没有什么如此激起天主的义怒。祂将此置于一切可怖之事之前；因为若我们有义务爱仇敌，那么连爱自己的人都要回避的人，该受何等惩罚？他在这方面比外邦人更坏。因此，在此情形下，罪恶之大将使这样的人与魔鬼一同离去。经上记着说：不行哀矜的人有祸了；若这是在旧约之下，那么在新约之下更是如此。如果在那获取财富、享受财富、看管财富都被允许的地方，尚且如此规定以周济穷人，那么在这需要我们舍弃所有的时代，岂不更要如此？因为古人什么没有做过？他们献十分之一，再十分之一，为孤儿、寡妇、外方人；而有人曾惊讶地对我说起另一个人：\Quote{怎么，某人献了十分之一。}这句话蕴含何等耻辱！因为对犹太人不足为奇的事，在基督徒中竟成了奇事。若那时忽略十分之一尚有危险，试想现今危险该有多大。\par

又，醉酒者不得继承天国。然而多数人说什么呢？\Quote{哦，若我和他都处于同样情况，这倒是不小的安慰。}那么怎样？首先，你和他不会受同样的惩罚；即使不然，那也不是安慰。苦难中的同伴，当痛苦有一定比例时，确有安慰；但当痛苦超乎比例，并使我们措手不及时，就不再允许我们接受任何安慰了。告诉那个受酷刑、已陷入火焰中的人，说某人也同样受苦，他仍不会感到安慰。难道所有以色列人不是一起灭亡的吗？那给他们带来了什么安慰？不，岂非正是这事使他们痛苦？他们因此不断说：我们丧亡了，我们毁灭了，我们耗尽了。那么这里有什么安慰呢？我们以这样的希望安慰自己是徒然的。唯有一种安慰：避免落入那永不熄灭的火；但已落入其中的人不可能找到安慰，那里有咬牙切齿，有哀哭，有永远不死的虫和不灭的火。请告诉我，当你处于如此大的苦难和困苦时，你还能感到任何安慰吗？那时你还能保持自我吗？我恳求并劝告你们，不要以这样的论调虚妄地自欺并自我安慰；不，让我们实践那些能拯救我们的德行。我们的目标是与基督同坐，而你们还在为这些琐事虚耗光阴吗？为什么？即使没有别的罪，仅仅为了这些言论本身，我们岂不应该受极大的惩罚吗？因为我们如此麻木、可怜、懒惰，面对如此巨大的特权，竟说出这样的话来！哦！当你们那时想起行善的人，将多么哀叹！当你们看见奴隶和出身低微的人，在此世稍作劳苦，在那里却分享君王的宝座，这些对你们岂不是比折磨更可怕吗？因为即使在现在，当你看见任何人享有盛誉，即使你没有遭受任何痛苦，你也认为这比任何惩罚更糟，仅此就使你消耗殆尽，自怜自艾，痛哭流涕，认为这比死一万次还坏；那么那时你将遭受什么？为什么？即使根本没有地狱，单是想到天国，岂不足以毁灭和吞噬你吗？事情将会如此，我们自己的经验足以教导我们。因此，不要以这样的言论虚妄地恭维自己的灵魂；不，让我们警惕，关心自己的救恩，重视德行，激励自己实践善工，使我们堪当获得这无比的荣耀，在我们的主耶稣基督内——愿光荣、权能、尊崇归于祂，偕同圣父及圣神，自今而彼，及于万世。阿们。\par

\SourceRecord{Translated by Gross Alexander.；Nicene and Post-Nicene Fathers, First Series；Vol. 13.；Edited by Philip Schaff.；Buffalo, NY: Christian Literature Publishing Co.,；1889.；<http://www.newadvent.org/fathers/230104.htm>.}

% structure: p0001 source=h1 target=section reason=page-title

\section{关于\WorkTitle{厄弗所书}的第五篇讲道}

第二章 11-12节\par

\BibleQuote{所以你们应当记得，从前你们按肉体说是外邦人，被那称为\OriginalTerm{割礼}{Circumcision}的——即人手在肉身上所行的割礼——称为\OriginalTerm{未受割礼}{Uncircumcision}；那时你们没有基督，与以色列子民的团体隔绝，对预许的盟约是外人，在这世上没有希望，没有天主。}\par

有许多事显明天主的慈爱。第一，祂亲自拯救了我们，并亲自以这样的方式拯救了我们。第二，祂拯救了我们，正如我们原本的样子。第三，祂提拔了我们到我们现在的境地。这一切本身都包含了祂慈爱的最伟大证明，而这正是保禄现在在书信中论述的题目。他之前说过，当我们死于过犯，生为义怒之子时，祂拯救了我们；他现在进一步告诉我们，祂使我们与谁等同。因此，他说：\BibleQuote{所以，}他说，\BibleQuote{应当记得；}因为这是我们众人常有的情况：当我们从极其卑微的地位被提升到相应甚至更大的尊荣时，我们几乎不再记得先前的状况，沉浸在这新的荣耀中。因此他说：\BibleQuote{所以，你们应当记得。}——为什么\Quote{所以}？\BibleQuote{所以}？因为我们被创造为行善，这足以促使我们修养德行；\BibleQuote{记得}——因为那记忆足以使我们感恩于恩主——\BibleQuote{你从前是外邦人。}请注意他是如何降低犹太人的优越地位，而强调外邦人的劣势；的确，那并非劣势，但他是在根据各人的特性和生活方式进行论证。\par

\BibleQuote{那称为\OriginalTerm{未受割礼}{Uncircumcision}的。}\par

这样看来，犹太人的尊荣在于名称，他们的特权在于肉体。因为未受割礼算不得什么，受割礼也算不得什么。\par

\BibleQuote{那称为割礼的}他说，人手在肉身上所行的\OriginalTerm{割礼}{Circumcision}，那时你们与基督分离，与以色列子民的团体隔绝，对预许的盟约是外人，没有希望，在这世上没有天主。\par

他说，你们就是被犹太人这样称呼的。但是，当他即将表明恩赐在于与以色列子民共融时，为什么他要贬低以色列的特权？他并非贬低它。在本质上他提高了它，但只在那些他们没有共融的方面，他贬低它们。因为他在后面说：\BibleQuote{你们已成了圣徒的同乡和天主的家人。}请注意他离贬低它有多远。他说，这些事情是无关紧要的。他说，绝不要因为你们碰巧未受割礼，而如今在未受割礼中，就以为这里面有任何区别。不，真正的麻烦是：成为\BibleQuote{没有基督}，成为\BibleQuote{与以色列子民的团体隔绝}。而这里的割礼并不是\BibleQuote{这个团体}。此外，成为预许盟约的外人，没有未来的希望，在这世上没有天主，这些都是他们状况的一部分。他在谈论天上之事，他也谈论地上的事，因为犹太人对这些有很高的评价。同样，基督在安慰祂的门徒时，先说了：\BibleQuote{为义而受迫害的人是有福的，因为天国是他们的。}然后补充了较小的安慰：\BibleQuote{因为，}祂说，\BibleQuote{他们也迫害了在你们以前的先知。}\BibleRef{玛 5:10}因为这件事与那件大事相比小得多，但就亲近和相信而言，它伟大而充分，具有很大的力量。这就是在团体中的分享。他的话不是\BibleQuote{分离}，而是\BibleQuote{与团体隔绝}。他的话不是\BibleQuote{你们不感兴趣}，而是\BibleQuote{你们不但没有份，而且是外人}。这些表述极其有力，表明分离非常之大。因为以色列人自己也曾没有这个团体，但不是作为外人，而是作为对此漠不关心的人；他们从盟约中跌落了，但不是作为外人，而是作为不配的人。\par

但什么是\BibleQuote{预许的盟约}？祂说：\BibleQuote{我要将这片土地赐给你和你的后裔，}\BibleRef{创 17:8}以及其他祂所预许的一切。\par

他补充说：\BibleQuote{没有希望，}和\BibleQuote{没有天主。}虽然他们确实崇拜过神明，但那不是神明：\BibleQuote{因为偶像算不得什么。}\BibleRef{格前 10:19}\par

第13-15节。\BibleQuote{但如今，在基督耶稣内，你们从前远离的人，藉着基督的血，已成为亲近的了。因为祂是我们的和平，祂使双方合而为一，并拆毁了中间隔断的墙，在祂的肉身内废除了仇恨。}\BibleRef{弗 2:13}\par

那么，这可能是伟大的特权吗？我们会说，我们被接纳进入以色列子民的团体？你在说什么？\BibleQuote{祂将天上和地上的一切总归于自己，}而你现在却和我们谈以色列？是的，他会说。那些更高的特权我们必须凭信德领悟；这些则凭事物本身。\BibleQuote{但如今，}他说，\BibleQuote{在基督耶稣内，你们从前远离的人，已成为亲近的了，}这是指着团体而言。因为\BibleQuote{远离}和\BibleQuote{亲近}只是意志和选择的问题。\par

\BibleQuote{因为祂是我们的和平，祂使双方合而为一。}\par

这是什么意思：\Quote{两者为一？}祂不是说，祂把我们提升到他们那样高的出身；而是说，祂把我们和他们都提升到更高的地位。只不过对我们而言，恩赐更大，因为对犹太人曾有过应许，他们比我们更近；对我们没有应许，我们甚至更远离。因此他说：\Quote{使外邦人因祂的仁慈而光荣天主。}\BibleRef{罗 15:9} 诚然，祂把应许给了以色列人，但他们不配；对我们，祂没有应许，不，我们甚至是外人，与他们毫无共同之处；然而祂使我们成为一，不是将我们联结于他们，而是将我们和他们一同联结于一身。让我给你们一个比喻。假设有两尊雕像，一尊是银的，另一尊是铅的；然后两者都被熔化，两者都变成金子。看，这样祂使两者成为一。或者再换一种说法：假设两者之中，一个是奴隶，另一个是养子；两者都得罪了祂，一个如同被剥夺继承权的孩子，另一个如同逃亡者，一个从未认识父亲的人。然后两者都被立为继承人，都是嫡子。看，他们被提升到同一个尊位，两者成为一，一个来自更远的距离，另一个来自更近的距离，奴隶变得比他得罪之前更尊贵。\par

他接着说：\Quote{拆毁了隔断的中间墙。}\par

什么是\Quote{隔断的中间墙}？他解释说：\Quote{以祂的肉身废掉了冤仇，就是那规条中所包含的诫命的律法。}有些人确实认为他指的是犹太人与希腊人之间的墙，因为那墙不容许犹太人与希腊人交往。然而在我看来，似乎不是这个意思，而是他称\Quote{肉身中的冤仇}为中间墙，因为它是一个共同的屏障，同样将我们与天主隔开。正如先知说：\Quote{你们的罪恶使你们与天主隔绝。}\BibleRef{依 59:2} 因为祂对犹太人和外邦人所怀的冤仇，就像一道中间墙。这道墙在律法存在时，不仅没有被废除，反而被加强；\Quote{因为律法，}宗徒说，\Quote{激起忿怒。}\BibleRef{罗 4:15} 正如他在那处说\Quote{律法激起忿怒}时，并非将全部效果归于律法本身，而是因为我们违背了它；同样，此处他称它为中间墙，因为因着违命，它造成了冤仇。律法是一道篱笆，但这是为了安全而设，因此被称为\Quote{篱笆}，为了形成一道围墙。请再听先知的话：\Quote{我围着它挖了沟。}\BibleRef{依 5:2} 又说：\Quote{你拆毁了它的篱笆，使所有过路的人都摘取它。}\BibleRef{咏 79:12} 因此此处它意味着安全；同样，\Quote{我要撤去它的篱笆，它必被践踏。}\BibleRef{依 5:5} 又说：\Quote{祂为他们赐下律法作为护卫。}\BibleRef{依 8:20} 又说：\Quote{上主施行公义，向以色列人显明祂的道路。}\BibleRef{咏 103:6} 然而，它却成为中间墙，不再使他们在安稳中，反而将他们与天主隔绝。这隔断的中间墙就是由篱笆形成的。为了解释这是什么，他接着说：\Quote{以祂的肉身废掉了冤仇，就是律法的诫命。}\par

如何废掉的？因祂被杀害，在其中消灭了冤仇。不仅以此方式，而且因遵守律法。但如果我们从先前的过犯中得释放，却又被迫遵守它，那又将如何？那样就会重蹈覆辙，而祂却毁灭了律法本身。因为他说：\Quote{废掉了那规条中所包含的诫命的律法。}哦！令人惊奇的爱！祂赐给我们律法，要我们遵守；当我们不遵守，本应受罚时，祂甚至废除了律法本身。如同一个人，他把孩子交给老师，如果孩子变得不顺服，他甚至从老师那里释放孩子，把他带走。这是多么大的慈爱！所谓\Quote{以规条废掉}，\par

\Quote{以规条废掉？}\par

因为他将\Quote{诫命}与\Quote{规条}作了广泛的区分。那么，他或者是指\Quote{信德}，称它为\Quote{规条}（因为祂只凭信德拯救了我们），或者是指\Quote{训令}，如同基督所说的：\Quote{我却对你们说：你们不要发怒。}\BibleRef{玛 5:22} 也就是说：\Quote{你若相信天主使祂从死者中复活，你必得救。}\BibleRef{罗 10:6} 又说：\Quote{这话离你很近，就在你口里，在你心里。不要说：谁要升到天上去？或谁要下到深渊里去？}或是，谁\Quote{把祂从死者中带上来？}祂引入信德，取代了某种生活方式。因为祂不愿白白拯救我们，所以祂自己承受了惩罚，同时也要求人以教义为信德。\par

\Quote{为要把双方在自己内造成一个新人。}\par

你当注意：外邦人并未成为犹太人，而是两者都进入了另一种状况。祂不仅是为了使后者改变，更是为了将两者重造。他处处使用\Quote{造}这个字，而不说\Quote{改变}，为要指出所成就之事的权能，并指出即使这创造是不可见的，却同样真实，我们从此不应从这创造中退却，如同对待自然事物一样。\par

\Quote{为要在自己内把双方}\par

这就是说，由他自己。他没有将此使命交给别人，而是他自己，亲自将两者融合，产生了一个荣耀的个体，且比第一次创造更大；而那个个体，首先就是他本人。因为这就是\Quote{在他自己内}的意思。他亲自先给出了原型和榜样。他一手抓住犹太人，另一手抓住外邦人，自己站在中间，将他们调和，使他们之间的一切隔阂消失，并从高处用水与火重新塑造他们；不再用水与土，而是用水与火。他通过割礼成为犹太人，成为受诅咒者，成为没有律法的外邦人，并统御外邦人和犹太人。\par

\Quote{一个新人，}他说，\Quote{如此缔造和平。}\par

对于双方，和平指向天主，也指向彼此。因为只要他们仍然是犹太人和外邦人，就无法和好。如果他们各自不脱离自己特殊的状态，就无法达到另一个更高的状态。因为当犹太人成为信徒时，他才与外邦人联合。这就像人在一所房子里，楼下有两个房间，上面有一个宽大华丽的大厅：他们只有上到上面，才能彼此看见。\par

\Quote{缔造和平，}更特别是指向天主；因为上下文显示这一点，因为他说了什么？\par

第16节。\Quote{并藉着十字架使双方在一个身体内与天主和好。} \BibleRef{弗 2:16}\par

他说，不单是\Quote{和好}（\OriginalTerm{和好}{\GreekText{καταλλάξῃ}}），而是\Quote{彻底和好}（\OriginalTerm{彻底和好}{\GreekText{ἀποκαταλλάξῃ}}），表明在此之前人性很容易和好，例如在圣人们身上以及在律法时代之前。\par

\Quote{在一个身体内，}他说，且是他自己的身体，\Quote{与天主。} 这是如何实现的？他意指，藉着亲自承担该受的惩罚。\par

\Quote{藉着十字架，在其中杀死了仇恨。}\par

没有什么比这些话更果断、更富有表现力了。宗徒说，他的死亡已经\Quote{杀死了}仇恨。他\Quote{击伤}并\Quote{杀死}了它，不是将使命交给别人，也不仅凭他所行的事，也凭他所受的苦。他没有说\Quote{解散了}，他说\Quote{取消了}，但比一切都更强的是\Quote{杀死了}，使它永不再起来。那么，它怎么又起来呢？是由于我们极其的败坏。因为我们只要住在基督的身体内，只要我们联合，它就不会再起来，而是死了；更确切地说，那先前的仇恨根本不再起来。但如果我们生出另一种仇恨，那就不再是因为他，他已经毁掉并杀死了先前的仇恨。实在是你，产出了一个新的仇恨。\Quote{肉性的思念，}他说，\Quote{是与天主为敌；} \BibleRef{罗 8:6}如果我们丝毫不随从肉性，就不会产生新的仇恨，而\Quote{和平}将常存。\par

道德教训。试想，当天主运用了如此多的方法使我们与祂和好，并且已实现了和好，而我们竟再次陷入仇恨中，这是何等巨大的恶事！这种仇恨，等待它的不是新的圣洗，而是地狱本身；不是新的赦免，而是严苛的审判。血肉的思念是奢侈和怠惰，\Quote{血肉的思念}是贪得无厌和一切罪恶。为什么称之为血肉的思念呢？然而血肉若没有灵魂，什么也不能做。他说这话并不是贬低血肉，正如他说\Quote{\OriginalTerm{属血气的}{\GreekText{ψυχικά}}}\BibleRef{格前 2:14}并非用此表达来贬低灵魂；因为身体和灵魂本身，若没有接受那从高天而来的推动，都不能成就任何伟大或崇高的事。因此，他称灵魂单独完成的那些行为为\Quote{\OriginalTerm{属魂的}{\GreekText{ψυχικά}}}，而称身体单独完成的那些行为为\Quote{属血肉的}。不是因为它们是属魂的，而是因为它们没有接受来自天上的指引，因而灭亡。因此，眼睛是好的，但没有光就会犯无数错误；然而这是它们软弱的过错，而不是本性的过错。如果错误是本性的，我们就永远无法正确使用它们。因为凡是本性的，没有一样是恶的。那么他为什么称属血肉的情欲为罪呢？因为每当血肉自高自大，掌控了她的驭手时，就产生无数的祸害。血肉的美德在于服从灵魂；统治灵魂则是它的恶习。正如一匹马可能既好又敏捷，但没有骑手就显不出这点；同样，血肉也将在我们砍断它的跳跃时显示它的好处。但同样，如果骑手没有技巧，也显不出他。不仅如此，他本人还会造成比之前所说的更可怕的祸害。因此，无论如何我们必须有圣神在身边。圣神在身边将赋予骑手新的力量；这会使身体和灵魂都变得美丽。因为正如灵魂住在身体内时使它美丽，但当灵魂离开身体，使其失去自己本有的活力而离去时，就像画家混淆颜色一样，最令人厌恶的结果随之而来，身体的各个部分都趋于腐朽和分解；——同样，当圣神抛弃身体和灵魂时，随之而来的厌恶更加糟糕和严重。那么，不要因为身体低于灵魂就责骂它，因为我也不能容忍责骂灵魂，因为灵魂没有圣神就没有力量。如果一定要说点什么，那么灵魂比身体更应受责备；因为身体没有灵魂确实不能造成严重的伤害，而灵魂没有身体却能造成许多伤害。因为我们知道，即使当身体正在消耗，毫无放纵之时，灵魂却忙碌不停。正如那些巫术士、魔术师、嫉妒者、施咒者，尤其会使身体消耗。但除此之外，甚至连奢侈也不是身体需要的后果，而是灵魂疏忽的后果；因为食物，而不是宴饮，是身体需要的对象。因为如果我想装上坚固的嚼子，我就制止马；但身体却无法阻止灵魂行恶。那么他为什么称之为属血肉的思念呢？因为它完全属于血肉；因为当血肉掌控时，它就犯错，一旦它剥夺了自己的理性及灵魂的至上权。因此身体的美德在于服从灵魂，因为血肉本身既不善也不恶。因为身体本身能做什么呢？身体之所以是善的，是因为它的联系，是因为它的服从，但本身既不善也不恶，然而它有能力趋向二者，并具有同等的倾向。身体有自然的欲望，但不是淫乱或奸淫的欲望，而是快乐的欲望；身体有对食物的欲望，而不是对宴饮的欲望；对饮料的欲望，而不是对醉酒的欲望。为了证明醉酒不是身体的自然欲望，请注意，每当你超过度量，越过界限时，它连一刻也不能坚持。到此为止是属于身体的，但所有其余的过度，例如当她被带入肉欲时，当她变得麻木时，这些都是属于灵魂的。因为虽然身体是好的，但它仍然远低于灵魂，如同铅不如金贵重，然而金需要铅来焊接，同样灵魂也需要身体。或者如同一个高贵的孩子需要引导者，同样灵魂也需要身体。因为，正如我们谈论孩童之事，并非贬低孩童本身，而只是指在孩童时期所做之事；同样，我们现在谈论身体也是如此。\par

然而，只要我们愿意，就在我们能力之内，不再属于血肉，不，也不在地上，而在天上，在圣神内。因为我们在某处与否，与其说是由我们的位置决定，不如说是由我们的心境决定。至少，对于在某个地方的许多人，我们说他们不在那里，当我们说：\Quote{你刚才不在这里。以及你现在不在这里。}我为什么这样说呢？我们常说：\Quote{你不在（\OriginalTerm{在}{\GreekText{ἐν}}）你自己内，我不在（\OriginalTerm{在}{\GreekText{ἐν}}）我自己内，}然而有什么比一个人与自己相近更实质（更强烈的身体位置的实例）呢？然而尽管如此，我们说他不在自己内。那么，让我们在自己内，在天上，在圣神内。让我们安存在天主的平安与恩宠中，使我们脱离一切血肉之事，并能获得在我们的主耶稣基督内所应许的美善，愿光荣、权能、尊威归于祂，偕同圣父及圣神，自今而始，至于世世。阿们。\par

\SourceRecord{Translated by Gross Alexander.；Nicene and Post-Nicene Fathers, First Series；Vol. 13.；Edited by Philip Schaff.；Buffalo, NY: Christian Literature Publishing Co.,；1889.；<http://www.newadvent.org/fathers/230105.htm>.}

% structure: p0001 source=h1 target=section reason=page-title

\section{讲道集第六篇：论\BibleBook{厄弗所}{书}}

第二章 17-22节 \BibleRef{弗 2:17}\par

\BibleQuote{祂来了，向你们远处的人传了和平，也向近处的人传了和平；因为藉着祂，我们双方在一个圣神内，得以进到父面前。所以你们不再是外方人或旅客，而是与圣徒同国，是天主家里的人，已建在宗徒和先知的基础上，而基督耶稣自己就是这建筑物的角石。在祂内，整个建筑物彼此衔接，逐渐增长，在主内成为一座圣殿；在祂内，你们也一同被建筑，成为天主藉圣神居住的住所。}\par

宗徒说，祂没有藉别人的手差遣，也没有藉任何其他受造物宣布这信息，而是亲自以自己的身体成就了这一切。祂没有差遣天使或总领天使来执行这使命，因为修复如此众多且巨大的灾祸，并宣告所成就的事，不是任何其他受造物所能胜任的，而是需要祂亲自来临。主于是取了仆人的身份，甚至几乎取了仆役的身份，祂说：\BibleQuote{祂来了，向你们远处的人传了和平，也向近处的人传了和平。}他是对犹太人说的，他们与我们相比是近处的人。\BibleQuote{因为藉着祂，我们双方在一个圣神内，得以进到父面前。}\par

他说：\Quote{平安}，即那朝向天主的\Quote{平安}。祂已使我们与祂和好。因为主自己也说：\BibleQuote{我把平安留给你们，我将我的平安赐给你们。}\BibleRef{若 14:27} 又说：\BibleQuote{你们放心，我已经战胜了世界。}\BibleRef{若 16:33} 又说：\BibleQuote{你们因我的名无论求什么，我必要实行。}\BibleRef{若 14:14} 又说：\BibleQuote{因为父爱你们。}\BibleRef{若 16:27} 这些都是平安的明证。然而对外邦人如何呢？\Quote{因为藉着祂，我们双方在一个圣神内，得以进到父面前，}不是你们少而他们多，而是众人藉着同一恩宠。祂以自己的死亡平息了义怒，并藉着圣神使我们适于父的爱。请注意，这里的\Quote{在}意为\Quote{藉着}或\Quote{通过}。就是说，祂藉着自己和圣神将我们带到父面前。\BibleQuote{所以你们不再是外方人或旅客，而是与圣徒同国。}\par

你们明白吗？这并不仅仅是与犹太人，而是与那些圣洁而伟大的人物，如亚巴郎、梅瑟和厄里亚？正是与他们同属一城，我们被登录，我们宣告自己属于同一城。\BibleQuote{因为说这样话的人，}他说，\BibleQuote{是表明自己寻求一个家乡。}\BibleRef{希 11:14} 我们不再是远离圣徒的陌生人，也不是外方人。因为那不能获得天上福分的人，才是外方人。\BibleQuote{因为儿子，}基督说，\BibleQuote{永远存留。}\BibleRef{若 8:35}\par

\BibleQuote{并且是天主家里的人，}他接着说。\par

他们当初藉着许多劳苦和困难所得着的，如今因着天主的恩宠为你们成就了。请看你们蒙召的盼望。\par

\BibleQuote{建在宗徒和先知的基础上。}\par

请注意他如何融合一切：外邦人、犹太人、宗徒、先知和基督，并有时借用身体，有时借用建筑物来表明这合一。他说：\BibleQuote{建在宗徒和先知的基础上。}那就是，宗徒和先知是基础，并且他把宗徒放在前面，尽管就时间顺序而言他们是在后的，无疑是要表示和表达这一点：即两者同样都是基础，整个建筑是一个整体，有一个根源。请注意，外邦人也有圣祖作为基础。他说这话比说\Quote{接枝}更为有力。在那里他是把他们连接上去。然后他补充说，基督把所有的一切结合起来。因为角石既结合墙壁，又结合基础。他用了\BibleQuote{祂在祂自己内将两者造成一个新人}\BibleRef{弗 2:15}\par

\BibleQuote{在祂内，整个建筑物。}\par

请注意，他如何把一切接合起来，并一会表现祂从上面压住整个身体并使之凝聚，一会表现祂从下面支撑建筑物，成为根基或基础。然而他用了\BibleQuote{祂在祂自己内将两者造成一个新人}\BibleRef{弗 2:15}的表达，这清楚显示，基督藉着自己把两堵墙接合起来，并且是在祂内造成的。并且他说：\Quote{祂是一切受造物的首生者，}就是说，祂亲自支持万有。\par

\BibleQuote{在祂内，整个建筑物彼此衔接。}\par

无论你提到屋顶、墙壁还是任何其他部分，都是祂在支持整体。因此他在别处称祂为根基。\BibleQuote{因为除了已奠定的基础，即耶稣基督外，}他说，\BibleQuote{任何人不能再奠立别的。}\BibleRef{格前 3:11} \BibleQuote{在祂内，整个建筑物}他说，\BibleQuote{彼此衔接。}这里他显示其完美，并表明人除非以极大的精确度生活，否则无法在其中拥有位置。\BibleQuote{它逐渐增长，}他说，\BibleQuote{成为一座圣殿，在主内。}他又说：\BibleQuote{在祂内，你们也一同被建筑。}他连续地说：\BibleQuote{成为一座圣殿，成为天主在圣神内的居所。}那么这建筑物的目的是什么？就是天主居住在这圣殿内。因为你们每个人各自是一座圣殿，你们众人一起也是一座圣殿。祂住在你们内，如同住在基督的身体内，又如同住在属神的圣殿内。祂没有使用表示我们来就近天主的词（\OriginalTerm{来到}{\GreekText{πρόσοδος}}），而是使用了表示天主领我们来就祂的词（\OriginalTerm{被引领}{\GreekText{προσαγωγή}}），因为我们不是自己来的，而是被祂领来的。基督说：\BibleQuote{除非经过我，谁也不能到父那里去。}又说：\BibleQuote{我是道路、真理和生命。}\BibleRef{若 14:6}\par

他把他们与诸圣结合，又回到先前的比喻，绝不让他们与基督分离。无疑，这是一座将一直建造到主来临的建筑物。无疑，正因为如此，保禄才说：\BibleQuote{我像一位明智的工头，立定了根基。}\BibleRef{格前 3:10} 又说基督是根基。这一切是什么意思呢？你看到，这些比喻都涉及主题，我们不可按字面解释它们。宗徒是借类比说话，正如基督所做的那样，祂称圣父为农夫，\BibleRef{若 15:1} 而称自己为根。\BibleRef{默 22:16}\par

第三章第一节：\BibleQuote{为此，我保禄，为你们外邦人作基督耶稣囚徒的}\BibleRef{弗 3:1}\par

他提到了基督伟大而深情的关怀；现在他转向自己，这关怀虽微不足道，与基督的相比算不得什么，却足以使他们依附于他。为此，他说，我也被捆绑。因为如果我的主为你们被钉十字架，那我更该被捆绑。祂不仅自己被捆绑，还允许祂的仆人也受捆绑——\BibleQuote{为你们外邦人}。这极有分量；我们不仅不再厌恶你们，反而为你们的缘故被捆锁，他说，并且我分享这极大的恩宠。\par

第二节：\BibleQuote{你们曾听过天主恩宠的\OriginalTerm{分施}{dispensation}，就是为你们所赐给我的。}\BibleRef{弗 3:2}\par

他暗示了主在\Place{大马士革}对\Person{阿纳尼雅}关于他的预言，那时主说：\BibleQuote{你只管去，因为他是我所拣选的\OriginalTerm{器皿}{vessel}，要在外邦人和君王面前为我传扬我的名。}\BibleRef{宗 9:15}\par

所谓\Quote{恩宠的分施}，是指那启示给他的。这就是说，我不是从人学来的。\BibleRef{迦 1:12} 祂竟屈尊启示给我，一个微不足道的人，也是为了你们。因为祂亲自对我说：\BibleQuote{你走吧，我要打发你到远方外邦人那里去。}\BibleRef{宗 22:21} \Quote{你们曾听过}——原来那里是一个伟大的分施：召叫一个人，不受任何外来影响，直接从天上来，并对他说：\Quote{扫禄，扫禄，你为什么迫害我？} 并以那不可言喻的光使他失明！\Quote{如果你们曾听过，} 他说，\Quote{天主恩宠的分施，就是为你们所赐给我的。}\par

第三节：\BibleQuote{那奥秘借着启示显示给我，正如我以前在简短的字句中写过的。}\BibleRef{弗 3:3}\par

也许他曾通过某些人告诉他们，或者不久前在给他们写信。这里他指出一切都出于天主，我们毫无贡献。因为什么？我问，难道保禄本人，这位杰出的、精通法律的、曾在\Person{加玛里耳}足前按最完美的方式受教的人，不是因恩宠得救的吗？他也很有理由称此为奥秘，因为这是一个奥秘：瞬间将外邦人提升到比犹太人更高的地位。\Quote{正如我以前写过的，} 他说，\Quote{在简短的字句中，} 即简要地，\par

第四节：\BibleQuote{你们读时，便能明白。}\BibleRef{弗 3:4}\par

令人惊叹！所以他没有写全部，也没有写应当写的那么多。但这里是主题的性质妨碍了他。在其他地方，如\WorkTitle{希伯来书}\BibleRef{希 5:11}和\WorkTitle{格林多前书}\BibleRef{格前 3:2}的情况，是听众的能力不足。\Quote{你们读时，便能明白，} 他说，\Quote{我对基督奥秘的理解，} 即我如何知道、如何理解天主所说的事，或者，基督坐在天主的右边；并论及尊严，即天主\BibleQuote{没有这样对待任何民族。}\BibleRef{咏 147:20} 然后解释天主这样对待的是哪个民族，他接着说：\par

第五节：\BibleQuote{这奥秘在其他世代没有告诉世人，如今却在圣神内启示给祂的圣宗徒和先知们。}\BibleRef{弗 3:5}\par

那么，请告诉我，先知们不知道吗？那么基督怎么说梅瑟和先知们写了\Quote{这些关于我的事？} 又说\BibleQuote{如果你们相信梅瑟，也必相信我}\BibleRef{若 5:46}，又说\BibleQuote{你们查考圣经，以为其中有永生，正是这些为我作证}\BibleRef{若 5:39}。他的意思是说，或者这奥秘没有启示给所有人，因为他补充说：\Quote{这奥秘在其他世代没有告诉世人，如今却被启示了；} 或者，它没有通过事实和现实本身如此启示，\Quote{如今却在圣神内启示给祂的圣宗徒和先知们。} 请想一想。伯多禄若没有受圣神教导，绝不会到外邦人那里去。因为请听他所说的话：\BibleQuote{天主既赐给了他们圣神，如同赐给了我们一样}\BibleRef{宗 10:47}。是天主愿意藉圣神使他们领受恩宠。先知们曾发言，但他们没有这样完全认识；不仅这样，就连宗徒们在听到之后也未能完全领会。这远超出一切人的计算和通常的期待。\par

第六节：\BibleQuote{外邦人是同承继者、同为一体、同享应许的。}\BibleRef{弗 3:6}\par

这是什么意思？\Quote{同承继者、同享应许、同为一体？} 最后一点是大事，即他们应成为一体；这种与祂的极高亲密关系。因为他们确实知道要受召叫，但不知道如此伟大，他们尚不知道。因此他称之为奥秘。\Quote{同享应许。} 以色列人分享，外邦人也同为天主应许的分享者。\par

\Quote{在基督耶稣内，借着福音。}\par

也就是说，只因祂也被派遣到他们那里，并因他们的相信；因为经上没有简单地说他们是同为继承人，而是说\Quote{借着福音}。然而，这事其实并不那么伟大，实际上是一件小事，并且向我们揭示了另一件更大的事：不仅人不知道这事，而且连天使、总领天使或其他任何受造的能力也都不知晓。因为这确实是一个奥秘，并未被启示。他说：\Quote{你们能看出，}他说，\Quote{我的领悟。} 这或许暗指他在\WorkTitle{宗徒大事录}中对他们说过的，即他知道外邦人也被召叫。他说，这是他自己所知晓的，是\Quote{那奥秘的认知}，也就是他提及的，即\Quote{基督要在自己内使两者成为一个新人。} 因为他和伯多禄都是藉着启示被教导的，即他们不可轻视外邦人；他也在辩护中陈述了这一点。\par

第7节。\BibleQuote{我成了这福音的仆役，是按照天主恩宠的恩赐，这恩赐是依照祂大能的德能赐给我的。}\BibleRef{弗 3:7}\par

他曾说：\Quote{我是囚犯；} 但现在又再次说，一切都出于天主，正如他所说：\Quote{按照祂恩宠的恩赐；} 因为照恩赐的德能，才有这特恩的尊位。然而，单有恩赐还不够，除非它也将能力植入他内。\par

道德教训。这实在是一项大能的工作，具有强大的能力，是任何人的勤勉都无法相称的。因为他将三种资质带到了圣言的宣讲中：炽热而勇敢的热忱，准备好承受任何可能艰难的灵魂，以及知识兼智慧。因为他的冒险精神、无可指摘的生活，若没有领受圣神的能力，便毫无益处。且看这能力首先在他自己身上显现，或者更确切地说，听听他自己的话。\BibleQuote{免得我们的职务受人诋毁。}\BibleRef{格后 6:3} 又说：\BibleQuote{因为我们的劝勉，不是出于错谬，也不是出于不洁，也不是出于诡诈，也不是出于贪婪的掩饰。}\BibleRef{得前 2:3} 这就是你所见的他的无可指摘。又说：\BibleQuote{因为我们顾念光明的事，不但在主面前，也在人面前。}\BibleRef{格后 8:21} 此外，还有这些：\BibleQuote{我指着我在主基督耶稣里因你们所有的夸耀起誓，我是天天冒死。}\BibleRef{格前 15:31} 又说：\BibleQuote{谁能使我们与基督的爱隔绝呢？是患难吗？是困苦吗？是迫害吗？}\BibleRef{罗 8:35} 又说：\BibleQuote{在许多忍耐中，在患难中，在穷困中，在窘迫中，在鞭笞中，在监禁中，在不眠中。}\BibleRef{格后 6:4} 然后，是他的审慎与管理：\BibleQuote{对犹太人，我就成为犹太人；对法律以外的人，我就成为法律以外的人；对在法律下的人，我就成为在法律下的人。}\BibleRef{格前 9:20} 他也剃了头，\BibleRef{宗 21:24} 并做了许多类似的事。但这一切的冠冕在于圣神的能力。\BibleQuote{因为我不敢提说，}他说，\BibleQuote{除了基督借着我所做的事。}\BibleRef{罗 15:18} 又说：\BibleQuote{你们究竟在哪件事上不如别的教会呢？}\BibleRef{格后 12:13} 又说：\BibleQuote{因为我虽然在最伟大的宗徒中毫无所缺，虽然我什么都不是。}\BibleRef{格后 12:11} 没有这些，那工作是不可能的。\par

因此，使人成为信徒的并非他的奇迹；不，奇迹并没有成就这事，他也不是基于这些而提出自己的崇高要求，而是基于其他理由。因为一个人必须同时在品行上无可指摘，在待人接物上谨慎明智，不顾危险，且善于教导。他的成功大部分正是藉着这些资质取得的。有了这些，便无需奇迹。至少我们看到，他在无数类似情况下，远在使用奇迹之前，就已经成功了。但如今，我们不具备这些资质的任何一样，却想要命令一切。然而，如果其中一样与另一样分离，从此便无用处。一个人若是极其不顾危险，但生活却受人指责，那有什么益处？\BibleQuote{如果你眼中的光变成了黑暗，}基督说，\BibleQuote{那黑暗是多么大啊！}\BibleRef{玛 6:23} 同样，一个人若是生活无可指摘，但却懈怠懒散，又有什么益处？\BibleQuote{因为谁不背起自己的十字架跟随我，}祂说，\BibleQuote{不配做我的门徒；}\BibleRef{玛 10:38} 所以，\BibleQuote{善牧为羊舍命。}\BibleRef{若 10:11} 再说，既有这两样，若不同时谨慎明智，\BibleQuote{知道该如何回答每一个人？}\BibleRef{哥 4:6} 又有什么益处？即使奇迹不在我们能力之中，然而这两样资质仍在我们能力之内。尽管如此，虽然保禄自己贡献了这么多，他却将一切归于恩宠。这是感恩仆人的作为。若不是他被迫表白，我们根本不会听到他的善行。\par

那么，我们岂配提及保禄的名字吗？他，虽有恩宠相助，仍不满足，反而为工作献上了千万危险；而我们，既缺少那信心的源头，请问，我们指望什么去保全托付给我们的那些人，或是赢得那些尚未归栈的人呢？——我们这些人，是钻研自我放纵的，是到处寻找安逸的，是不能（更确切地说，是不愿意）忍受危险的一丝影子的，而且离他的智慧如同天地之远。因此，我们以下的人也远不及那些日子的人；因为那些日子的门徒比这些日子的导师更好，他们孤立在民众和暴君之中，四面都是敌人，却丝毫没有退缩或屈服。至少听听他对斐理伯人说的话：\BibleRef{斐 1:29} \BibleQuote{因为为了基督的缘故，赐给你们的恩典，不仅是信他，而且也是为他受苦。} 又对得撒洛尼人：\BibleRef{得前 2:14} \BibleQuote{因为弟兄们，你们成了犹太地天主教会众的效法者。} 又写信给希伯来人：\BibleRef{希 10:34} 他说：\BibleQuote{你们欣然忍受了财产的剥夺。} 对哥罗森人：\BibleRef{哥 3:3} 他作证说：\BibleQuote{因为你们已经死了，你们的生命与基督一同藏在天主内。} 实际上，他对这些厄弗所人也见证了许多危险和艰难。又写信给迦拉达人：\BibleRef{迦 3:4} 他说：\BibleQuote{你们受了那么多苦，难道徒然吗？如果真是徒然的话。} 你们也看见他们都在行善。因此，那时恩宠才有效地运行，因此他们生活在善工中。再听听他写给格林多人的话，他指责他们无数的过失；但他仍为他们作证，说：\BibleQuote{是的，它在你们身上产生了何等热忱，何等渴望！}\BibleRef{格前 7:11} 又，他在多少方面为他们作证？如今这些事，即使在导师中，也看不到了。全都消逝了。原因是：爱德冷淡了，罪人不受惩罚；（因为他写给弟茂德的话：\BibleRef{弟前 5:20} \BibleQuote{犯罪的人，你要在众人面前责备；}）原因是统治者处于病态；因为头若不健全，身体的其余部分怎能保持活力？但请看现今的混乱何等严重。那些生活有德、在任何情况下都能有自信的人，却占据了山顶，逃出了世界，把自己当作与仇敌和异己分离，而不是与他们所属的身体分离。\par

灾祸也满载着说不尽的祸害落到了教会身上。首席职务都成了可买卖的。因此无数邪恶发生，没有人纠正，没有人责备。不仅如此，混乱竟有了某种方法和一致性。有人做错了事，被控告了吗？他的努力不是证明自己无辜，而是尽可能找同谋。我们将会怎样？因为地狱是我们被威胁的部分。相信我，若不是天主在那里为我们积蓄了惩罚，你们每天都会看到比犹太人的灾难更深重的悲剧。那又如何？但请无人动怒，因为我不提名字；假设有人进入这座教堂，把你们此刻在这里的人（那些现在与我在一起的人）带出来，并审问他们；或者不如说不是现在，而是假设在复活节那天，某人拥有这样的精神，能彻底知道他们所做的事，仔细检查所有前来领圣体的人，以及在参与奥迹之后正在被洗涤（在圣洗中）的人；许多事会被发现，比犹太人的恐怖更令人震惊。他会发现那些行占卜、使用符咒、预兆和咒语的人，以及那些犯奸淫、通奸、酗酒、辱骂的人——贪吝的，我不愿再加，以免伤害任何在场者的感情。还有什么？假设任何人审查全世界所有的领圣体者，有什么过犯是他查不出的呢？如果他审查那些掌权者呢？难道他不会发现他们急切地追求利益？贩卖高位？嫉妒、恶毒、虚荣、贪吃、沦为金钱的奴隶？\par

那么，当这种不敬继续发生时，我们怎能不期待可怕的灾祸呢？为了确定那些犯下这样罪过的人会遭受何等严厉的惩罚，请思考古代的榜样。一个人，一个普通士兵，偷了圣物，全体都被击打。你肯定知道我所指的历史？我说的是加尔米的儿子阿干，那个偷窃被祝圣战利品的人。\BibleRef{苏 7:1} 还有先知发言的时候，正是他们的国家充满占卜者，如同培肋舍特人的时候。\BibleRef{依 2:6} 而现在，无数的邪恶已满盈，却无人恐惧。哦，从此让我们警醒。天主习惯于也惩罚义人与恶人；达尼尔和三个圣童就是如此，成千上万的其他例子也是如此，如今发生的战争也是如此。因为义人，无论他们有多少罪的负担，借此方式甚至可以卸下它；但恶人则不然。\par

鉴于这一切，我们要注意自己。你们没有看见这些战争吗？你们没有听到这些灾难吗？你们没有从这些事学到教训吗？列国和整座城市被吞没毁灭，无数倍的人又被奴役于野蛮人。\par

如果地狱仍不能使我们清醒，那么至少让这些事警醒我们。怎么，这些也仅仅是威胁吗？难道它们不是已经发生的事实吗？他们所受的惩罚是巨大的，但我们将受的惩罚更大，因为我们连他们的命运都不能使我们清醒。这篇讲论令人厌倦吗？我知道它本身确实如此，但如果我们留心听，它是有益的；因为它没有那种讨人喜欢的讲辞的特质——不仅如此，也永远不会再有，而是永远包含着那些能够使人谦卑、净化灵魂的话题。因为这些将为我们奠定以后来临的福分的基础，愿天主恩赐我们众人获得这些福分，在我们的主耶稣基督内，愿光荣与权能归于圣父，偕同圣神，从今世直到永远，世世无尽。阿们。\par

\SourceRecord{Translated by Gross Alexander.；Nicene and Post-Nicene Fathers, First Series；Vol. 13.；Edited by Philip Schaff.；Buffalo, NY: Christian Literature Publishing Co.,；1889.；<http://www.newadvent.org/fathers/230106.htm>.}

% structure: p0001 source=h1 target=section reason=page-title

\section{讲道七：论厄弗所书}

第三章　第八至十一节 \BibleRef{弗 3:8-11}\par

\BibleQuote{我原是众圣人中最小的，竟蒙受了这恩宠，得以向外邦人宣讲基督那深不可测的丰富，并使众人明白那奥秘的计划，这奥秘从永远以来就隐藏在创造万有的天主内：为使天上的率领者和掌权者，现在藉着教会，得知天主各样智慧，按照祂在基督耶稣内所定的永恒旨意。}\par

他们去医生那里的人，不应只是去那里就完了；还必须学习如何调治自己，并用药物。因此，我们这些来此的人，也不应只是这样做就完了，我们必须学习功课，就是保禄的超凡谦卑。什么？当他即将谈论天主恩宠的浩大时，听听他说什么：\Quote{我原是众圣人中最小的，竟蒙受了这恩宠。}这谦卑甚至在于哀哭他过去的罪，尽管已涂抹，并提及它们，持守己量，正如他称自己为\BibleQuote{亵渎者、迫害者、凌辱者；}\BibleRef{弟前1:13}然而没有比这更甚的：因为\Quote{从前}，他说，我就是这样；他又称自己为\BibleQuote{流产儿}\BibleRef{格前15:8}。但在如此众多伟大善行之后，并在那时，他竟如此自卑，称自己\Quote{比众圣人中最小的还小}，这实在是伟大无比的节制。\Quote{我原是众圣人中最小的；}他不是说\Quote{比宗徒们}。所以那句话比我们眼前的这句更弱。在那里他说：\BibleQuote{我不配称为宗徒}\BibleRef{格前15:9}。这里他说他甚至\Quote{比众圣人中最小的还小}；\Quote{我，}他说，\Quote{比众圣人中最小的还小，竟蒙受了这恩宠。}什么恩宠？\Quote{得以向外邦人宣讲基督那深不可测的丰富，并使众人明白那奥秘的计划，这奥秘从永远以来就隐藏在创造万有的天主内：为使天上的率领者和掌权者，现在藉着教会，得知天主各样智慧。}确实，对人并未启示；而你正在启明天使、总领天使、率领者和掌权者？是的，他说。因为\Quote{隐藏在}天主内，甚至\Quote{在创造万有的天主内}。而你敢说出这个？是的，他说。但天使如何得知此事？藉着教会。他又说，不仅\OriginalTerm{各样}{\GreekText{ποικίλος}}，而是\OriginalTerm{多种多样}{\GreekText{πολυποίκιλος}}智慧，即\Quote{增多而变化的}。那么这是什么？难道天使不知道？不，一点也不；如果率领者不知道，天使更不可能知道。那么难道总领天使也不知道？不，他们也不知道。但他们将如何知道？谁将启示？当我们被教导时，他们也藉着我们得知。因为听天使对若瑟所说：\BibleQuote{你要给祂起名叫耶稣，因为祂要把自己的百姓从罪恶中救出来}\BibleRef{玛1:21}。\par

保禄自己被派往外邦人，其他宗徒被派往受割礼的人。所以，他说，更奇妙、更令人惊奇的使命交托给了\Quote{比最小的还小的我}。这也是出于恩宠，使最微小的人受托最大的事；使他成为这些喜讯的宣讲者。因为成为更大喜讯的宣讲者，在这方面是伟大的。\par

\Quote{向外邦人宣讲基督那深不可测的丰富。}\par

如果祂的\Quote{丰富是深不可测的}，而且是在显现之后，那么祂的本质更是如此。如果它仍然是一个奥秘，那么在它被显明之前更是如此；他称之为奥秘，正是因为天使不知道它，也没有任何人知道。\par

\Quote{并使众人明白，}他说，\Quote{那奥秘的计划，这奥秘从永远以来就隐藏在创造万有的天主内。}\par

天使只知道这一点：\BibleQuote{上主的份子是祂的百姓}\BibleRef{申32:8}。又说：\BibleQuote{波斯国的君阻扰我}\BibleRef{达10:13}。所以，他们不知道这点并不奇怪；因为如果他们不知道从被掳归回的情况，就更不知道这些事了。因为这就是福音。\BibleQuote{祂要拯救，}它说，\BibleQuote{祂的百姓}\BibleRef{玛1:21}。没有提及外邦人。但关于外邦人的事，圣神启示了。天使知道他们被召，但不知道是归于与以色列相同的特权，甚至坐在天主的宝座上，这是谁曾预料？谁曾相信？\par

\Quote{这奥秘一直隐藏在}他说，\Quote{天主内。}\par

然而，他在\WorkTitle{致罗马人书}中更清楚地阐明了这\Quote{计划}。他接着说：\Quote{在天主内}，\Quote{祂藉着耶稣基督创造了万有。}他说\Quote{藉着耶稣基督}很好；因为祂藉着祂创造了万有的天主，也藉着祂启示这事；没有祂，祂什么也不造；因为经上说：\BibleQuote{没有祂，凡受造的，没有一样不是由祂而造}\BibleRef{若1:30}。\par

在提到\Quote{率领者}和\Quote{掌权者}时，他指的是天上的和地下的。\par

\Quote{按照永恒旨意。}意思是，现在已经实现，但不是现在才决定，而是从起初就已预先计划。\Quote{按照祂在基督耶稣内所定的永恒旨意。}也就是说，按照永恒的预知；预知将要发生的事，即未来的世代；因为祂知道将发生的事，就这样决定了。\Quote{按照世代的旨意}，可能是指祂藉着基督耶稣所造的世代，因为一切都是藉着基督造的。\par

第12节。他说：\BibleQuote{在他内，我们藉着对他所怀的信德，获得了放胆和怀着信赖的进路。}（见：\BibleRef{弗 3:12}）\par

他说：\Quote{获得进路}，不是作为囚犯，也不是作为请求赦免的人，更不是作为罪人；因为，他说，我们甚至拥有\InnerQuote{带着信赖的放胆}，也就是说，伴随着欢喜的信赖；这源于什么？\Quote{藉着对他所怀的信德。}\par

第13节。\BibleQuote{因此我求你们，不要因我为你们所受的患难而沮丧，这患难原是你们的光荣。}（见：\BibleRef{弗 3:13}）\par

为何是\Quote{为他们}？为何是\Quote{他们的光荣}？这是因为天主如此爱他们，甚至将自己的圣子赐给他们，并为他们的缘故使祂的仆人们受苦：正是为了使它们能获得如此多的祝福，保禄才被囚禁。这无疑出于天主对他们超越的爱：这也是天主论及先知时所说的话：\BibleQuote{我用我口中的话杀了他们。}（见：\BibleRef{欧 6:5}）但是当别人受苦时，他们为何沮丧呢？他的意思是，他们感到困扰、痛苦。这也是他在给得撒洛尼人写信时所说的：\BibleQuote{不要因这些患难而动摇。}（见：\BibleRef{得前 3:3}）因为我们不仅不应悲伤，反而应当喜乐。如果你们在预先警告中得到安慰，我们预先告诉你们，在此世我们必有患难。为何祈祷？因为主已经这样安排了。\par

第14-15节。\BibleQuote{因此，我在圣父面前屈膝，天上和地上的各家族都由他得名。}\par

他在这里显示了他为他们祈祷的精神。他不只是简单地说\Quote{我祈祷}，而是通过\Quote{屈膝}表明恳求是发自内心的。\par

这就是\Quote{各家族都由他得名}。\par

也就是说，他的意思是，不再按天使的数目来被计算，而是按那创造了天上和地上各支脉的天主，而非按犹太人的方式。\par

第16-17节。\BibleQuote{愿他按照他荣耀的丰富，藉着他的圣神，以大能坚固你们内在的人，使基督因着信德住在你们心中。}\par

请留意他以何等不满足的热切为他们呼求这些祝福，使他们不致被风浪摇动。但这如何实现呢？藉着\Quote{圣神在你们内在的人里，使基督因着信德住在你们心中}。这又如何可能呢？\par

第18-19节。\BibleQuote{为叫你们在爱中生根立基，能够与众圣徒一同领悟那阔、长、高、深是什么，并认识基督那超越知识的爱。}\par

因此，他现在的祈祷又与他开始时完全相同。因为他在开始时说了什么？\Quote{愿我们主耶稣基督的天主，荣耀的圣父，赐给你们智慧和启示的圣神，使你们认识他；光照你们心灵的眼目，使你们知道他的召叫有什么希望，他在圣徒中光荣的产业有什么丰富；以及他对我们相信的人所施展的无限伟大能力。}现在他又说了同样的话。\Quote{愿你们能与众圣徒一同领悟那阔、长、高、深是什么}，即完全认识那为我们所安排的神妙计划：\Quote{还有阔、长、高、深}——也就是说，天主那无限的爱，以及它如何延伸至各处。他借用有形物体的可见维度来描绘它，仿佛指向一个人。他包括了上下和侧面。他可以说，我确实这样说了，但教导你们这些事并非靠我的言语；那必须是圣神的工作。\Quote{靠着他的大能}，他说，你们必须\Quote{被坚固}，以面对等待你们的考验，并保持不动摇；因此，没有别的途径可以坚强，唯有依靠圣神，既因试探，也因属肉体的推理。\par

但基督如何住在心中呢？请听基督自己所说的话：\BibleQuote{我和我圣父要到祂那里去，并要在祂那里作我们的住所。}（见：\BibleRef{若 14:23}）祂住在那些忠信的心中，住在那些扎根于祂爱中、坚定不动摇的人心中。\par

他说：\Quote{愿你们}彻底\Quote{坚强}，因此需要极大的力量。\par

\Quote{愿你们充满天主的全部圆满。}\par

他的意思是这样的。虽然基督的爱超越一切人类认知的范围，但如果基督住在你们内，你们仍将认识它；是的，你们不仅从他那里得知这事，甚至要\Quote{充满天主的全部圆满}；他所谓\Quote{天主的圆满}，或者是指认识天主如何在圣父、圣子、圣神内受敬拜，或者是指激励他们如此尽一切努力，以充满天主所充满的全部美德。\par

第20节。\BibleQuote{愿那能按运行在我们身上的大能，充充足足地成就一切，超过我们所求所想的。}（见：\BibleRef{弗 3:20}）\par

天主\Quote{充充足足地成就一切，超过我们所求所想的}，这从宗徒自己所写的话中显而易见。他说：我固然祈祷，但祂自己，甚至不用我的任何祈祷，就要成就比我们所求更大的事，不只是\Quote{更大}，也不是\Quote{丰富地更大}，而是\Quote{超越丰富地}。这从\Quote{运行在我们身上的大能}可以看出来：因为我们从未求过这些事，也未曾期望过它们。\par

第21节。他总结说：\BibleQuote{愿荣耀归于祂，在教会内并在基督耶稣内，直到万代，永世无穷。阿们。}（见：\BibleRef{弗 3:21}）\par

他以祈祷和赞颂结束讲论是恰当的；因为那赐予我们如此浩大恩赐的那一位，理当受到荣耀和赞美，以致这甚至是我们对祂慈悲惊叹的一部分，为那藉着耶稣基督从天主手中赐予我们的事物而献上光荣。\par

\Quote{在教会内的荣耀。}他这样说很恰当，因为只有教会能够持续到永远。\par

似乎有必要说明所谓\Quote{家族}（\OriginalTerm{家族}{\GreekText{πατριαί}}）的含义。在地上，确实有\Quote{家族}，即由同一祖先繁衍的种族；但在天上，这如何可能，那里没有人由另一个人出生？那么，无疑，他所谓\Quote{家族}，是指天上生灵的集会与等级；正如我们在圣经中读到：\Quote{阿玛塔黎家族}：\BibleRef{撒上 10:21}；或者，也是指从祂那里，地上的父亲取得父亲之名。\par

然而，他并不是向天主祈求全部，而是向他们要求信德与爱德，且不只是爱德，而是\Quote{根基深厚且稳固}的爱德，使任何狂风都不能动摇它，任何事物都不能推翻它。他曾说，\Quote{磨难}是\Quote{光荣}；若我的磨难对你们如此，他会说，你们自己的磨难更将如此：因此，受苦并非天主遗弃人的标记，因为那位为我们行了如此伟大之事的天主，绝不会这样做。\par

再者，若要理解天主的爱，保禄需要祈祷，且需要圣神的居留，那么，谁仅凭理性推论能理解基督的本性呢？为何学习天主爱我们是一件困难的事？亲爱的，这极其困难。因为有些人甚至不知道这事，因此他们甚至说世上有无数灾祸；另一些人则不知道这爱的广度。确实，保禄并不是寻求知道它的广度，也不是为了衡量它；他怎能呢？他只是要理解这一点：它是超越的、伟大的。他说，他甚至能藉着赐予我们的知识显示出这事。\par

然而，比\Quote{获得能力而坚强}更高的是什么呢？是为了让基督居内。他说，我们祈求的是伟大的事，但天主甚至能超越这些，因此祂不仅爱我们，而且爱得如此深切。因此，亲爱的，让我们关心理解天主的爱。这确是一件大事；没有比这更对我们有益的，没有比这更深深触动我们的：这比地狱的恐惧本身更能说服我们的灵魂。那么，我们何时才能理解它呢？既从刚才提到的来源，也从每日发生的事。因为出于什么动机，这些事为我们成就了？出于祂的什么需要？完全没有。他反复把爱当作原因。但爱的最高程度是：人毫无先前的功劳而得益。\par

伦理。那么，让我们也效法祂；让我们善待我们的仇敌，善待恨我们的人，让我们亲近那些背弃我们的人。这使我们像天主。\Quote{因为若你们爱那爱你们的人，}基督说，\Quote{你们有什么赏报呢？}\Quote{连外邦人不也照样做吗？}\BibleRef{玛 5:46}。但爱的确凿证据是什么？是爱那恨你的人。我愿意给你们一个例子，（请原谅我），既然我在属神的人中找不到，我就从外面的人中引一个。你们没有看见那些恋爱中的人吗？他们的情妇对他们施加了多少侮辱，使用了多少伎俩，施加了多少惩罚：然而他们仍被束缚于她们，为她们燃烧，爱她们胜过自己的灵魂，整夜站在她们的门槛前。让我们从他们身上学习，当然不是爱那样的女人——我指的是妓女；不，而是这样爱我们的仇敌。因为告诉我，妓女对待她们的情人，难道不比世上所有的仇敌更傲慢吗？她们挥霍他们的财产，当面侮辱他们，给他们施加比对自己的仆役更卑贱的任务？然而他们仍然不放弃，虽然没有人有像情人在情妇中那样大的仇敌。是的，这个被爱的人轻视、辱骂，并且常常虐待他，而她越被爱，就越轻视他。有什么比这样的精神更残酷？然而他仍然爱她。\par

但或许我们也会在属神的人物中找到这样的爱，不是在我们这个时代（因为爱已\Quote{冷淡了}），\BibleRef{玛 24:12}，而是在那些伟大而荣耀的古圣中。梅瑟，有福的梅瑟，甚至超越了那些怀着人间热情去爱的人。怎样，以何种方式？首先，他放弃了宫廷、奢侈、随从和随之而来的荣耀，选择与以色列人在一起。然而，这不仅是没有人愿意做的事，甚至若被人发现，他也会因被视为与奴隶——不仅是奴隶，而且被视为可憎之人——有亲属关系而感到羞耻。但他不仅不以亲属为耻，而且全心捍卫他们，为他们冒险。\BibleRef{宗 7:24}如何？经上说，他看见一人欺负他们中的一个，就保护那受欺负的，并杀了那欺负人的。但这还不是为了仇敌。这行为本身固然伟大，但不如后来之事伟大。第二天，他看见同样的事发生，当他看见他保护的那人正伤害邻人时，他劝诫他停止作恶。但那人以极大的忘恩负义说：\Quote{谁立了你作首领和审判官管我们？}\BibleRef{宗 7:27}谁听到这些话不会发怒？如果先前的行为是出于激情和疯狂，那么他也会打击并杀死这人；因为那受他保护的人绝不会告发他。但经上说，因为他们是弟兄，他这样说话。当那人（希伯来人）受欺负时，他没有说这样的话：\Quote{谁立了你作首领和审判官管我们？}\Quote{你昨天为什么不这样说？}梅瑟会说，\Quote{你的不义，你的残忍，这些使我成为首领和审判官。}\par

但现在，请注意，事实上有些人甚至对天主自己说同样的话。当他们受欺负时，他们愿意天主是报复之神，并抱怨祂的容忍；但当他们自己作恶时，他们一刻也不容忍。\par

然而，有什么比这样的话更苦涩呢？尽管如此，在他被派遣到那忘恩负义、那不知感恩的民族那里之后，他去了，并未退缩。是的，即使在那些奇迹、在他亲手所行的奇事之后，他们屡次想用石头砸死他，他却从他们手中逃脱。他们也不断地抱怨，然而即便如此，他仍如此热爱他们，以至于在他们犯下那大罪时，他对天主说：\BibleQuote{如今，倘若祢愿意，就赦免他们的罪；若不然，求祢把我从祢所写的册上涂抹了吧。}\BibleRef{出 32:32}\Quote{我宁可与他们一同灭亡，}他说，\Quote{也不愿没有他们而得救。}看，这爱甚至到了疯狂的地步，真是无边的爱。\Person{梅瑟}，你说什么？你不顾天堂吗？\Quote{是的，}他说，\Quote{因为我爱那些亏负我的人。}你祈求被涂抹吗？\Quote{是的，}他说，\Quote{我有什么办法？因为这是爱。}这些事之后又怎样呢？听\WorkTitle{圣经}在别处说什么：\BibleQuote{因他们的缘故，\Person{梅瑟}遭了殃。}\BibleRef{咏 105:32}他们多少次放纵情欲？多少次拒绝他和他的兄弟？多少次想返回埃及？然而在这一切之后，他仍为他们燃烧，为他们痴迷，并准备为他们受苦。\par

因此，人应当如此爱仇敌：通过哀哭，通过不懈的忍耐，通过做一切事，通过显示一切恩惠，来致力于他们的得救。\par

再者，告诉我，\Person{保禄}做了什么？他岂不是甚至祈求代替他们受诅咒吗？\BibleRef{罗 9:3}但我们必须从主那里获得那伟大的榜样，因为祂自己也这样行，祂说：\BibleQuote{因为祂使太阳升起，光照恶人，也光照善人，}\BibleRef{玛 5:45}引用祂圣父的榜样；但我们却以基督自己为榜样。祂降生成人，来到他们中间，我说的是，祂为了他们成了奴仆，\BibleQuote{祂谦卑自下，空虚自己，取了奴仆的形体。}\BibleRef{斐 2:7}当祂来到他们那里时，祂自己没有\BibleQuote{转到外邦人的任何路上，}\BibleRef{玛 10:5}并同样吩咐祂的门徒，不仅如此，\BibleQuote{祂周游四方，医治各种病症，各种疾苦。}\BibleRef{玛 4:23}那么结果如何呢？其余的人都惊讶，赞叹说：\BibleQuote{那么，这人从哪里得到这一切呢？}\BibleRef{玛 13:56}但这些受惠者，祂的恩泽所及的人，却说：\BibleQuote{祂有魔鬼，}\BibleRef{若 10:20}\BibleQuote{亵渎天主，}\BibleRef{若 10:36}\BibleQuote{疯了，}是\BibleQuote{骗子。}\BibleRef{若 7:12}难道祂因此抛弃了他们吗？绝不，完全不是。当祂听到这些话时，祂甚至更显著地对他们施恩，径直走向那些即将钉死祂的人，为的是只拯救他们。而当祂被钉十字架后，祂的话是什么？\BibleQuote{圣父啊，宽恕他们吧，因为他们不知道自己做的是什么。}\BibleRef{路 23:34}在这之前受虐待，在这之后也受虐待，直到最后一口气，祂为他们做了一切，为他们祈祷。是的，即使在十字架之后，祂还为他们做了什么？祂不是派遣了宗徒吗？祂不是行了奇迹吗？祂不是震动了整个世界吗？\par

因此，我们应当这样爱仇敌，这样效法基督。\Person{保禄}就是这样做的。虽被石头打，受尽无数虐待，却仍为他们的益处做一切事。听他自己的话：\BibleQuote{我心中所切望和向天主的祈求，是为他们能得救。}\BibleRef{罗 10:1}又说：\Quote{我为他们作证，他们对天主有热诚。}又说：\BibleQuote{如果你这野橄榄枝被接上，何况这些本枝，岂不更会被接在自己的橄榄树上吗？}\BibleRef{罗 11:24}你想，这些表达所蕴含的情感是何等温柔，善意是何等浩大？无法表达，无法表达。\par

因此，我们应当这样爱仇敌。这就是爱天主，是祂命令了这事，是祂立了这法律。效法祂就是爱我们的仇敌。当想：你并非在造福你的仇敌，而是在造福你自己；你并非在爱他，而是在服从天主。因此，知道了这些事，让我们彼此坚固爱德，好使我们能完美地履行这一职责，并获得在基督耶稣我们的主内所应许的美善。愿光荣、权能、尊威归于祂，与圣父及圣神，自今至永远，世世无尽。阿们。\par

\SourceRecord{Translated by Gross Alexander.；Nicene and Post-Nicene Fathers, First Series；Vol. 13.；Edited by Philip Schaff.；Buffalo, NY: Christian Literature Publishing Co.,；1889.；<http://www.newadvent.org/fathers/230107.htm>.}

% structure: p0001 source=h1 target=section reason=page-title

\section{讲道第八篇：论厄弗所人书}

第四章，第一、二节\par

\Quote{所以我，\OriginalTerm{为主被囚的}{prisoner in the Lord}，恳求你们行事要与你们所蒙的召叫相称，凡事谦虚和温柔。}\par

教师的美德不在于追求来自下属的赞美和尊敬，而在于他们的救恩，并以这一目标去做一切事；因为若有人以别的为目标，他就不是教师，而是暴君。确实，天主立你们在他人之上，并非为使你们享受更大的尊崇和服事，而是为使你们自己的利益被忽略，他们每个人的利益得以建立。这是教师的职责。有福的保禄就是这样的人，他毫无虚荣，满足于做一个普通人，甚至是最微小的一个。因此他甚至自称为他们的仆人，并且常常以恳求的语气说话。看看他，即使现在写信，也没有一点专横、傲慢，而是全然的谦逊和克制。\par

他说：\Quote{所以我}，\Quote{为主被囚的，恳求你们行事要与你们所蒙的召叫相称。}请问，你恳求什么？是为了为你自己获得什么利益吗？不，他说，绝无此事；而是为了拯救他人。然而，那些恳求的人确实是为与自己重要的事而恳求。不错，他说，这对我来说也是重要的，正如他在其他书信中所说的：\BibleQuote{现在我们活着，如果你们在主内站稳，}\BibleRef{得前 3:8}因为他热切渴望他教导之人的救恩。\par

\Quote{我，为主被囚的。}伟大而崇高的尊严！比国王、执政官或任何其他人都更伟大。因此，他在写信给费肋孟时也使用了这个头衔：\BibleQuote{如保禄年迈，而且如今也为耶稣基督的囚犯。}\BibleRef{费 1:9}因为没有比为基督被捆绑更荣耀的了，正如捆住那些圣手上的锁链；为基督成为囚犯比作宗徒、教师、传福音者更为荣耀。凡爱基督的人，都会明白我所说的。凡为主所感动、热爱主的人，都明白这些锁链的能力。这样的人宁愿为基督成为囚犯，也不愿以天堂为居所。他所展示给他们的手比任何金子更荣耀，甚至比任何王冠更荣耀。是的，没有金冠冕束在头上能像为基督的铁链那样赐予荣耀。当时的监狱比王宫更荣耀，甚至比天堂本身更荣耀。为何我说比王宫？因为它包含了一位基督的囚犯。凡爱基督的人，都知道这个头衔的尊严，都知道这是什么德行，都知道天主赐予人类何等大的恩惠，甚至是为祂被捆绑。或许，为祂被捆绑比\BibleQuote{坐在祂的右边}\BibleRef{玛 20:21}更荣耀，比\BibleQuote{坐在十二个宝座上}\BibleRef{玛 19:28}更尊贵。\par

为何我说人间的荣耀？我羞于将世俗的财富和金饰与这些锁链相比。但暂且不提那些伟大而属天的荣耀，即便这事没有任何报酬，单是这牺牲本身便是巨大的报酬，这便是充分的补偿——为所爱者忍受这些苦难。那些爱人者，即使所爱不是天主，而是人，他们明白我所说的，因为他们更乐于为所爱者受苦，而非被所爱者尊荣。但为完全理解这些事，属于圣洁的团体，我是指宗徒们，并且只属于他们。因为请听有福的路加所说的话：\BibleQuote{他们离开公议会，心里欢喜，因被算作配为这名受辱。}\BibleRef{宗 5:11}在所有人看来，这似乎是愚蠢的事：受辱竟被视为配得，受辱竟令人欢喜。但对于那些明白基督之爱的人，这一切被视为最蒙福的。若有人给我选择，整个天堂或那条锁链，我宁愿选那条锁链。若有人问我，应当将我安置在与天使同在高处，还是与保禄同受捆绑，我宁愿选监狱。若有人要将我变成天上那些围绕宝座的掌权者之一，或是变成这样一个囚犯，我宁愿选囚犯。没有比为基督被捆绑更蒙福的了。但愿此刻我能在那地方（据说那条锁链仍在），以看见和钦佩那些人对基督的爱。但愿我能看见那使魔鬼恐惧战栗、天使敬重的锁链。没有比为基督受任何苦更尊贵的了。我认为保禄的幸福不是因为被\BibleQuote{提到乐园里}\BibleRef{格后 12:4}，而是因为被投入地牢；我认为他的幸福不是因为听见\BibleQuote{不可言说的话}，而是因为忍受了那些锁链。我认为他的幸福不是因为被\BibleQuote{提到第三层天上}\BibleRef{格后 12:2}，而是因为那些锁链。因为后者大于前者，听他本人如何知道这一点；他说并不是\Quote{我听见不可言说的话}恳求你们，而是什么？\Quote{我，为主被囚的，恳求你们。}我们也不奇怪他没有在所有书信中写这个头衔，因为他并非总在狱中，只在某些时候。\par

我认为为基督受苦比从基督手中接受尊荣更值得渴望。这是超越的尊荣，是超越万物的荣耀。如果祂自己为了我成为仆人，并\BibleQuote{空虚}\BibleRef{斐 2:7}了自己的荣耀，却认为祂自己被钉十字架为了我才是真正的荣耀，那么我不应该忍受什么？因为请听祂自己的话：\BibleQuote{父啊，求你荣耀我。}\BibleRef{若 17:1}你在说什么？你正与盗贼和掘墓者一同被带到十字架，你承受被诅咒的死亡；你将被吐唾沫和拳打；而你称此为荣耀？是的，祂说，因为我为我的挚爱者承受这些，我把它们完全视为荣耀。如果祂爱那可怜不幸的人，却称此为荣耀，即不坐在圣父的宝座上，也不在圣父的荣耀中，而是在羞辱中——如果这是祂的荣耀，并且祂将此事置于其他之前：那么我更应当视这些为荣耀。啊！那些有福的锁链！啊！那些被那锁链装饰的有福之手！保禄的手在\Place{吕斯特拉}抬起并扶起瘸子时，不如被那些锁链捆绑时那样有价值。如果我生活在那些时代，我会多么热切地拥抱它们，并将它们贴近我的眼珠。我绝不会停止亲吻那些被算为配得为我主捆绑的手。你惊奇于保禄，当毒蛇咬住他的手却没有伤害他？不要惊奇。它尊敬他的锁链。是的，整个大海也尊敬它；因为那时他也是在船难中获救时被捆绑的。如果有人此刻赐给我复活死人的能力，我不选择那能力，而选择这锁链。如果我没有教会的挂虑，身体强壮有力，我不会退缩于承担如此长的旅程，只为了目睹那些锁链，为了看见他被囚禁的监狱。他的奇迹的痕迹在全地众多，然而却不如他的伤疤那样珍贵。\BibleRef{迦 6:17}在圣经中，他行奇迹时并不如他受苦、被鞭打、被拖来拖去那样令我喜悦。以至于从他的身上人取走手帕或围裙。奇妙，真是奇妙，但这些还不如那些奇妙。\BibleQuote{他们打了他许多棍，便把他下在监里。}\BibleRef{宗 16:23}又有：在捆绑中，\BibleQuote{他们唱着赞美诗献给天主。}\BibleRef{宗 16:25}又有：\BibleQuote{他们用石头砸他，把他拖到城外，以为他死了。}\BibleRef{宗 14:19}你们愿意知道为了基督的缘故，捆绑在祂仆人身上的铁链是多么伟大吗？请听基督自己说的：\BibleQuote{你们是有福的。}\BibleRef{玛 5:11}为什么？当你们复活死人时？不。但为什么？当你们医治瞎子时？绝不是。但那是为什么？\BibleQuote{当人们辱骂你们，迫害你们，并且因为我而捏造各样坏话毁谤你们时。}\BibleRef{玛 5:11}如果被人毁谤使人如此有福，那么受恶待，那会成就什么不能成就的事？请听这位有福者自己在别处说的：\BibleQuote{从此以后，有公义的冠冕为我存留。}\BibleRef{弟后 4:8}但是，比这冠冕更荣耀的是锁链：他说，主会算我配得这锁链，我丝毫不关心那些事。对我而言，为基督受苦足以作为一切报偿。只要祂允许我说：\BibleQuote{我填补基督苦难所欠缺的。}\BibleRef{哥 1:24}，我便一无所求了。\par

伯多禄也配得了这锁链；因为我们读到，他被捆绑，交给士兵，且正在睡觉。\BibleRef{宗 12:6}然而他喜乐，没有失去理智，而且陷入沉睡，如果他极度焦虑就不可能这样。然而，他在两个士兵中间睡着；有一位天使来到他身边，拍他的肋旁，叫醒了他。那么，如果有人对我说：你愿意选择哪一个？你愿意做那拍打伯多禄的天使，还是被释放的伯多禄？我更愿意选择做伯多禄，因为为了他的缘故，连天使也来了。是的，我渴望享受那些锁链。你们说，既然他从大恶中被释放，他为何还祈祷？不要惊奇：他祈祷，因为他怕死；他怕死，因为他宁愿自己的生命仍是进一步受苦的题材。因为请听那有福的保禄自己说的。\BibleRef{斐 1:23}\BibleQuote{离开世界，与基督同在，是更好得多的；然而留在肉身内，为你们更是需要的。}\BibleRef{斐 1:23}这他甚至称为恩赐，在写信时这样说：\BibleQuote{因为你们蒙受了恩宠，\OriginalTerm{作为恩赐}{\GreekText{ἕ} \GreekText{χαρίσθη}}，不仅仅是为了相信祂，也是为了为祂受苦。}\BibleRef{斐 1:29}因此后者比前者更大：因为祂出于祂的自由恩宠赐予了它；诚然，这是一个恩赐，极其伟大，是的，比那些任何一件都大，比使太阳和月亮停止、比移动世界更大！这比拥有胜过魔鬼的能力或驱魔更大。魔鬼因被我们所运用的信德驱逐而忧伤，不如当他们看见我们为基督的缘故受苦、被囚禁时那么忧伤。因为这增加了我们的胆量。为基督的缘故被捆绑之所以是高贵的事，不是因为它为我们赢得了天国；而是因为它是为基督的缘故而做的。我祝福那些锁链不是因为它们引向天堂；而是因为它们是为上主而戴的。知道他被捆绑是为基督的缘故，这是何等夸耀！何等幸福，何等崇高的尊荣，何等辉煌的殊荣！我渴望永远沉浸在这些主题中。我渴望紧握这条锁链。我渴望，虽然实际上我没有能力，但仍在思想上，以类似他的性情将这锁链环绕我的灵魂。\par

\BibleQuote{牢狱的地基}，我们读到，\BibleQuote{都震动了}，当\Person{保禄}被囚时，\BibleQuote{众人的锁链也都解开了。}\BibleRef{宗 16:26}那么，你岂不见在锁链中有一个能消解锁链的本性吗？因为正如\Person{主}的死亡置死亡于死地，\Person{保禄}的锁链也释放了被锁的人，震动了囚牢，打开了门户。然而，这并非锁链的天然作用，而正好相反；锁链的天然作用是保住被锁者的安全，并非为他打开狱墙。不，一般的锁链并非如此，但为基督的缘故的锁链却是如此。\BibleQuote{狱警俯伏在\Person{保禄}和\Person{息拉}面前。}\BibleRef{宗 16:29}然而，这也不又是普通锁链的效果，使捆绑者俯伏在被捆绑者的脚下；不，相反地，是使后者置于前者的手下。而在这里，那自由的人伏在曾被锁的人脚下。捆绑者恳求他所捆绑的人，释放他脱离恐惧。告诉我，难道不是你捆绑了他？你不是把他投进内监？你不是把他的脚上了木狗？你为什么战栗？你为什么惊慌？你为什么哭泣？你为什么拔剑？他说：我从未捆绑过这样的人！我不知道基督的囚徒有如此强大的力量。你说的什么？他们领受了开启天国的权能，难道不能打开一座监狱？他们释放了被恶神捆绑的人，一块铁又怎能制服他们？你不认识这些人。所以你也蒙了宽恕。那囚徒是\Person{保禄}，众天使都尊敬他。他是\Person{保禄}，连他的手帕和围裙都能驱魔，赶走疾病。当然，来自魔鬼的锁链是金刚石般的，远比铁更不可解开；因为那锁链束缚灵魂，而铁只束缚身体。那么，那个释放了被缚灵魂的人，岂不能释放自己的身体？那个能挣断恶神的锁链的人，岂不能解开铁铐？那个甚至借着自己的衣服解开了那些囚徒、释放他们脱离魔鬼的魔咒的人，岂不能自己使自己自由？原来他先是自己受缚，然后释放了囚徒，好使你明白，基督的仆人在锁链中拥有的力量远大于那些自由的人。如果是一个自由人行这事，那就不如此奇妙了。所以锁链不是软弱的标志，反而是更大能力的标志，圣人的力量正是如此更辉煌地显示出来：当他在锁链中仍能胜过那些自由的人，当他在锁链中不仅使自己自由，也使被锁的人自由。墙壁有何用？把他投入内监有何益，而他竟连外门也打开了？又为何发生在夜间？且伴有地震？\par

啊，请容忍我片刻，容我不提宗徒的言语，却沉醉于宗徒的作为，以\Person{保禄}的锁链为宴席；容许我更长久地默想它。我抓住了那条锁链，没有人能使我与它分离。此刻我被情感捆绑，比他那时在木狗中更稳妥。这是无人能解开的束缚，因为它由基督的爱形成；这束缚连天使，不，连天国也没有能力解开。我们可以听见\Person{保禄}自己的话：\BibleRef{罗 8:38}\BibleQuote{无论是天使，是掌权者，是现在的事，是将来的事，是有能的，是高，是深，都不能使我们与天主在吾主耶稣基督内的爱相隔绝。}\par

那么，这事为何发生在半夜？又为何伴有地震？且听，惊叹天主眷顾的安排。众人的锁链都松开了，门也开了。而这事却是为了狱警的缘故，不是为了显耀，而是为了他的救恩：因为囚徒们不知道他们被松绑，从\Person{保禄}的呼喊可见；他说了什么？\BibleQuote{他大声呼喊说：不要伤害自己，因为我们都在这里。}\BibleRef{宗 16:28}但他们若看见门开了，自己得了自由，就绝不会都留在里面。这些人惯于穿墙、翻越屋顶和女儿墙，在锁链中也敢尝试各种手段；如今锁链松了，门也开了，狱警又睡着了，他们岂能忍受留在里面？不，睡眠的束缚对他们而言替代了铁的束缚。因此这事发生了，而奇迹对那将要蒙受救恩的狱警没有造成损害。此外，人在夜间被绑得最牢，而非白天；所以我们可以看见他们被绑得紧紧的，都在睡觉；但若这些事发生在白天，就会有大骚动和混乱。\par

再者，那建筑物为何震动？是为了唤醒狱警，让他目睹所发生的事，因为他独自堪当得救。我也求你，我祈祷，目睹基督恩宠的极大丰盛，因为在\Person{保禄}的锁链之中，也适宜提及天主的恩宠；不，其实这些锁链本身就是天主的恩赐和恩宠。的确有人抱怨：\Quote{为什么狱警得救了？}而他们本应赞美天主仁爱的那些情况，他们反而对此加以指摘。这并不足为奇。这样的人就像那些病弱之人，连滋养他们的食物也要挑剔，而他们本应珍视；他们断言蜜是苦的；又像那些视力模糊之人，被那本应照亮他们的事物反而弄瞎了眼睛。这并不是因为这些事物本身的性质所致，而是由于那些人软弱无力，无法善加利用。然而，我在说什么呢？当他们本该赞叹天主的仁爱，因祂接纳了一个陷入最绝望邪恶的人，并使他变得更好，他们反而指摘：\Quote{为什么？他怎么没有把这看作巫术和邪术，反而把他们关押得更紧，并大声喊叫呢？}许多因素共同阻止了这一点：首先，他听见他们向天主唱赞美诗。行邪术的绝不会唱这样的颂歌，因为据说他听见他们向天主唱赞美诗。其次，他们自己没有逃跑，甚至还阻止他自杀。如果他们是为自己着想，就不会仍留在里面，他们自己早就先逃走了。他们的仁慈也很大：他们阻止这个人自杀，即使他曾经捆绑了他们，这样做几乎是对他说：\Quote{的确，你曾牢牢地捆绑了我们，而且极其残忍，好使你自己能从最残忍的捆绑中获得释放。}因为每个人都因自己的罪被锁链桎梏；那些锁链是可诅咒的，而这些为基督的锁链却是可祝福的，值得许多热切的祈祷。因为这些锁链能够解开那些罪的锁链，祂通过感官所能觉察的事物向我们证明了这一点。你看见了那些被铁链捆绑的人获得释放了吗？你自己也必从其他磨人的捆绑中获得释放。这些锁链，即囚犯的锁链（我说的不是\Person{保禄}的），是那些罪的锁链的结果。那些被关在里面的人，是双重囚犯，而狱警自己也是一个囚犯。他们确实被铁和罪所捆绑，而狱警只被罪捆绑。\Person{保禄}释放了他们，以坚定狱警的信德，因为他所解开的锁链是可以看见的。基督自己也是如此；但顺序相反。在那次事例中，有双重的瘫痪。是什么？灵魂因罪而瘫痪，同时身体也瘫痪。那么主做了什么？祂说：\BibleRef{玛 9:3}\BibleQuote{孩子，放心吧，你的罪赦了。}祂先解开了真正瘫痪的锁链，然后进行另一项：因为当\BibleQuote{有几个经师心里说：这人亵渎了；耶稣知道他们的心意，说：你们为什么心里怀着恶念？什么更容易呢？是说：你的罪赦了，或是说：起来行走吧？但为叫你们知道，人子在地上有赦罪的权柄——于是对瘫子说：起来，拿起你的床，回家去吧！}行了看不见的奇迹后，祂用看得见的奇迹加以证实，用身体的治愈证实精神的治愈。祂为什么这样做？为应验经上所说的话，\BibleRef{路 19:22}\BibleQuote{恶仆，我要凭你的口定你的罪。}因为那些人说什么？\BibleQuote{除天主一个外，谁也不能赦罪。}当然，没有天使，也没有总领天使，或任何其他受造之物有能力。这是你们自己所承认的。那么又该说什么？如果显明我赦免了罪，那就充分证明我是天主。然而，祂并没有这样说，而是说什么？\BibleRef{玛 9:6}\BibleQuote{但为叫你们知道，人子在地上有赦罪的权柄——于是对瘫子说：起来，拿起你的床，回家去吧！}所以，当祂说，我行了更难的奇迹，你们就没有借口，没有余地反驳那个更容易的奇迹。因此，祂先行了那看不见的奇迹，因为有许多反对者；然后祂把他们从看不见的引导到看得见的本身。\par

那么，狱警的信德绝非轻率或仓促的。他看见了囚犯。他什么也没看见，什么也没听见不义之事；他看见没有一件是藉巫术成就的，因为他们正在向天主咏唱赞歌。他看见所做的一切都出自洋溢的仁慈，因为他们没有向他报复，尽管他们有能力这样做；因为他们有能力解救自己和囚犯并逃跑；如果不能解救囚犯，至少他们也能够解救自己；但他们没有这样做。他们就这样不仅藉奇迹，也藉行为博得了他的敬畏。因为保禄如何呼喊？\Quote{他大声喊叫说：不要伤害自己，因为我们都在这里。}你立刻看见他摆脱了虚荣和傲慢，以及他的同感。他没有说：\Quote{这些奇事是为我们行出来的}，而是好像他不过是囚犯之一，他说：\Quote{因为我们都在这里。}然而，即使他们之前没有自解捆绑，也没有藉奇迹如此行，他们本可以保持沉默，并释放所有被捆绑的人。因为倘若他们沉默不语，没有大声呼喊阻止他的手，他早已将剑刺穿自己的喉咙了。因此保禄也喊叫了，因为他被投入内监：仿佛他说了：\Quote{你这样做是伤害自己，因为你把那些能救你脱离危险的人推得那么远。}然而他们没有效法他所受到的对待；虽然如果他死了，所有人都会逃脱。你看，他们宁愿留在捆绑中，也不愿让他丧亡。由此他也在心中推理：\Quote{倘若他们是行巫术的，无疑他们就会释放其他人，并解除自己的捆绑：}（因为很可能许多这样的行巫术者也曾被监禁。）他更惊讶的是，他常负责看管行巫术者，却从未见过这样的事。行巫术者绝不会摇动地基，以致惊醒狱警，从而使自己的逃脱更加困难。\par

现在，让我们继续来看狱警的信德。\Quote{并且，}圣经说，\Quote{他要了灯，跳进去，战战兢兢地俯伏在保禄和息拉面前，领他们出来，说：诸位，我当做什么才能得救？}他拿着火和剑，呼叫说：\Quote{诸位，我当做什么才能得救？}\Quote{他们说：信靠主耶稣基督，你和你一家都必得救。} \BibleRef{宗 16:29} \Quote{这不是行巫术者的行为，}他会说，\Quote{来传授这样的道理。这里丝毫没有提到邪神。}你看他多么配得救恩：因为当他看见奇迹，从恐惧中解脱时，他没有忘记最要紧的事，反而在如此大的危险中，他关心那关乎他灵魂的救恩：他以应当来到导师面前的方式到他们面前：他俯伏在他们脚前。\Quote{并且他们讲了，}经文继续说，\Quote{主的话给他和他全家的人。当夜，就在那时候，他带他们去洗了伤；他和所有属于他的人立刻受了洗。} \BibleRef{宗 16:32} 请看这人的热诚！他没有迟延；他没有说：\Quote{等天亮吧，让我们看看，让我们思量一下；}而是以极大的热诚，他自己和他全家都受了洗。是的，不像如今大多数人，任凭仆役、妻子和儿女未受洗。我恳求你，要像狱警一样。我不是说在权力上，而是在心意上；因为权力若不能坚持，有何益处？那野蛮、不人道、习于无数恶行并以此为常的人，突然变得如此仁厚、如此体贴入微。\Quote{他洗了，}经上说，\Quote{他们的伤。}\par

另一方面，也要注意保禄的热诚。他被捆绑，被鞭打，就这样宣讲了福音。哦，那有福的锁链，那夜以何等大的产痛生产，产下了何等样的子女！是的，关于他们，他也可以说：\Quote{我在锁链中所生的。} \BibleRef{费 1:10} 你要注意，他如何夸耀，并愿意这样出生的子女因此而更加荣耀！你要注意，那些锁链的荣耀何等超绝，因为它们不仅给佩带者带来光彩，也给了当时被他所生的人。在保禄的锁链中所生的人有一些优势，我不是说在恩宠方面（因为恩宠是同一的），也不是在赦免方面（因为赦免对众人是同一的），而在于他们从一开始就被教导要因这样的事而喜乐和夸耀。\Quote{当夜，就在那时候，}经上说，\Quote{他带他们去洗了伤，并受了洗。}\par

现在请看果实。他立即以属世的事物回报了他们。\Quote{他领他们上了自己的家，摆上饭食，他和全家因信了天主而大为喜乐。}因为他还有什么不愿意做呢？如今，监狱的门开了，天堂也向他敞开了。他洗了导师的伤，在他面前摆上食物，并喜乐了。保禄的锁链进入了监狱，将那里的一切变为一座教会；它牵引着基督的身体，预备了属灵的筵席，为那使天使喜乐的生产而阵痛。那么，我先前说监狱比天堂更荣耀，难道没有理由吗？因为监狱成了天堂喜乐的缘由；是的，如果\Quote{在天堂上对一个悔改的罪人也有喜乐} \BibleRef{路 15:7}，如果\Quote{那里有两个或三个人因我的名聚在一起，我就在他们中间} \BibleRef{玛 18:20}，那么何况在保禄和息拉、狱警和他全家所在之处，以及他们那样恳切的信德呢！请注意他们信德的那种炽烈恳切。\par

但是，这监狱让我想起了另一座监狱。那是什么监狱呢？就是伯多禄被囚的那一座。不过，那里并没有发生类似的事情。没有。他被交给四班士兵看守，他没有唱歌，没有警醒，反而睡着了；另外，他也没有受过鞭打。然而，危险却更大，因为在我们前面的情况中，结局已经完成，囚犯保禄和息拉已经受了刑罚；但在他的处境中，刑罚尚未来到。因此，虽然没有鞭打来折磨他，但未来的预期却使他痛苦。请注意那里的奇迹。\BibleQuote{看，主的一位天使}，经上记载，\BibleQuote{站在他旁边，一道光照亮了牢房；他拍打伯多禄的肋旁，唤醒他说：\InnerQuote{快快起来。}他的锁链便从他手上脱落下来。\BibleRef{宗 12:7}} 为了使他不至于以为这事仅仅是光的作用，天使也击打了他。除了他自己，没有人看见那光，他认为那是一个异象。睡梦中的人对天主的仁慈如此迟钝。\BibleQuote{天使对他说：\InnerQuote{束上带子，穿上鞋。}他就这样做了。天使又对他说：\InnerQuote{披上你的外衣，跟我来。}他便出来跟着走，不知道天使所做的是真的，以为看见了异象。他们经过了第一道和第二道守卫，来到通往城里的铁门，那门自动为他们开了；他们出去，走过一条街，天使便立即离开了他。\BibleRef{宗 12:8}} 为什么这里没有发生像保禄和息拉那样的事呢？因为在他们的案件中，他们打算释放他们。因此，天主不愿他们以这种方式被释放。而在蒙福的伯多禄的案件中，他们打算将他带去处决。但有人可能会说，倘若他被带出来，被交到王的手中，然后从迫在眉睫的危险中被夺走，并且没有受到伤害，那岂不是更加奇妙吗？因为这样一来，士兵们也不至于丧命。关于此事，人们提出了一个重大的疑问。有人说，天主以他人的惩罚和毁灭来拯救祂的仆人吗？首先，那不是以他人的毁灭为代价；因为这不是出自上智的安排，而是出自法官的残忍。如何解释呢？天主的上智如此安排，不仅这些人不必灭亡，而且甚至连那法官也可以得救，就像狱警的情况一样。但是他并未善用这份恩赐。\Quote{到了天亮，}经上继续写道，\Quote{士兵们中间起了不小的骚动，不知道伯多禄到哪里去了。}后来怎么样呢？黑落德仔细追查此事，经上记载：\BibleQuote{他审问了卫兵，下令把他们处死。\BibleRef{宗 12:18}}其实，倘若他没有审问他们，处决他们或许还有一些借口。但事实上，他把他们带到面前，审问了他们，发现伯多禄曾被捆绑，监狱看守严密，看守者都在门前。没有墙壁被挖穿，没有门被打开，也没有任何其他欺诈的证据。他本该因此畏惧天主的大能，因为天主将伯多禄从危险中救出，并应崇拜那位能行如此大能者。可是相反，他却下令将那些人处死。那么在这种情况下，天主怎么会是原因呢？如果天主真的让墙被挖穿，从而救出伯多禄，那么这件事可能会被归咎于他们的疏忽。但如果天主如此上智地安排，使这事显明不是出于人的邪恶行为，而是出于天主的奇迹，那么黑落德为何这样做呢？因为如果伯多禄想逃跑，他会带着锁链直接逃走。如果他想逃跑，他在慌乱中绝不会如此周全，连鞋也会穿上，而会丢下它们。但事实上，天使对他说：\InnerQuote{穿上你的鞋}的目的是让他们知道，他做这事不是出于逃跑，而是从容不迫。因为被捆绑着，夹在两个士兵之间，他绝没有足够的时间解开锁链，特别是他也和保禄一样，囚在内牢。因此，看守者的惩罚是由于法官的不义。因为为什么犹太人没有这样做呢？因为现在我又想起了另一座监狱。第一座是罗马的那座，接下来是凯撒勒雅的那座，现在我们来到耶路撒冷的那座。当时，大司祭和法利塞人从他们派去带伯多禄出狱的人那里听说，\Quote{里面没有找到任何人}，但门都\Quote{关着}，而且\Quote{看守者站在门口}，为什么他们不仅没有处死看守者，反而\Quote{关于他们，非常困惑，不知这事将如何发展}。如果连那些心怀杀机的犹太人都不曾有此念头，那么你这个为了讨好犹太人而行事的人，就更不该有。因为这个不义的判决，惩罚很快临到了黑落德。\par

但如今若有人对此抱怨，那么也抱怨那些在大路上被杀的人，以及成千上万被不义处死的人，此外，还有基督诞生时被屠杀的婴孩；因为按你们的说法，基督也是他们死亡的原因。但这不是基督，而是黑落德之父的疯狂和暴政。你问：为何祂不将祂从黑落德手中夺走？诚然，祂本可以这样做，但这样做毫无益处。至少基督多少次从他们手中逃脱？然而这对那无情的人民有何益处？而在这里，所发生的事对信众甚至大有裨益。因为正如有记录在案，敌人自己也见证了事实，这见证是无可置疑的。因此，如同在那事例中，敌人的口无论如何都无法被堵住，只有那些承认事实的人才能做到，这里也是如此。因为这里的狱卒为何没有像黑落德那样做？不，黑落德所目睹的事，丝毫不比此人目睹的事逊色。就奇事而言，确信一个囚犯在门关闭时出来，与看见门被打开，同样奇妙。实际上，后者可能被看作是幻觉，而前者若被精确而详细地报告，则绝不可能是这样。所以，若此人像黑落德那样邪恶，他也会像黑落德杀害士兵那样杀害保禄；但他并非如此。\par

若有人问：\Quote{为何天主也允许婴孩被杀害？}恐怕我会陷入一篇比原本打算向你们讲的更长的讲论。\par

此时，让我们结束这篇讲论，衷心感谢保禄的锁链，因为它已成为我们如此多福泽的源泉，并劝勉你们，若为基督的缘故遭受任何事，不仅不要抱怨，反而要喜乐，如同宗徒们一样，是的，还要夸耀；正如保禄所说：\BibleQuote{所以我更喜欢夸耀我的软弱，} \BibleRef{格后 12:9} 因为为此他听到了那些话：\BibleQuote{我的恩宠为你足够了。} 保禄以锁链为荣；而你呢，以财富为傲？宗徒们因被配受鞭打而喜乐，而你寻求安逸和放纵？那么，你凭什么想与他们达到同样的结局，如果你在此世走的是与他们相反的道路？保禄说：\BibleQuote{现在，我为主被圣神所缚，往耶路撒冷去，不知道在那里要遇到什么事；但圣神在各城里向我证明说：有锁链和患难在等待我。} \BibleRef{宗 20:22} 那么，若锁链和患难等待你，你为何还要出发？正因如此，他说，我才能为基督的缘故被捆绑，为祂的缘故而死。\BibleQuote{因为我不但准备受捆绑，而且也准备为耶稣主的名字而死。} \BibleRef{宗 21:13}\par

伦理。没有比那灵魂更有福的了。他夸耀什么？在锁链、患难、桎梏、伤痕中；他说：\BibleQuote{我身上带着耶稣的烙印，} \BibleRef{迦 6:17} 仿佛它们是某种伟大的战利品。又说：\BibleQuote{为了以色列的希望，我才被这条锁链捆绑。} \BibleRef{宗 28:20} 又说：\BibleQuote{我为这事成了带锁链的使者。} \BibleRef{弗 6:20} 这是什么？你以囚犯的身份在世上行走，不羞愧，不惧怕吗？难道你不怕有人指责你的天主软弱？不怕有人因此不肯靠近你、加入羊栈吗？不，他说，我的锁链并非如此。它们甚至可以在王宫中闪耀光芒。他说：\BibleQuote{以致我的锁链在基督内，传遍了整个御营；并且大多数弟兄在主内，因我的锁链而确信不疑，更加放胆无惧地传讲天主的话。} \BibleRef{斐 1:13} 看，锁链中的力量比复活死人更强。他们看见我被捆绑，反而更加勇敢。因为哪里有锁链，那里必然也有大事。哪里有患难，那里确实也有救恩，确实也有安慰，确实也有伟大的成就。因为当魔鬼踢踹时，无疑他正被击中。当他捆绑天主的仆人时，圣道正是传播得最广。请看这无处不在的例子。保禄被囚；在监狱中他行了这些事，是的，他说，正是藉着我的锁链本身。他在罗马被囚，反而使更多人皈依信仰；因为他自己不仅放胆，许多人也因他而放胆。他在耶路撒冷被囚，在锁链中宣讲，使王惊奇，\BibleRef{宗 26:28} 使总督战栗，\BibleRef{宗 24:25} 因为害怕，据说他释放了他，而捆绑他的人竟不羞于从他被捆绑的人那里接受关于未来之事的训导。他在锁链中航海，修复了船只残骸，制伏了风暴。正是在他被锁链捆绑时，海怪扑向他，却从他手上脱落，没有伤他。他在罗马被捆绑，在锁链中宣讲，吸引成千上万的人归向他，用这锁链本身代替其他一切作为论据。\par

然而，如今我们并没有被捆绑的福分。但还有一种锁链，如果我们愿意佩戴它。那是什么？就是约束我们的手，不急于贪财。让我们用这锁链束缚自己。愿天主的敬畏成为我们铁链般的束缚。让我们释放那些被贫穷、患难所捆绑的人。打开监狱的门与释放被束缚的灵魂无可比拟。解开囚犯的锁链与\BibleQuote{使那受压制的得自由} \BibleRef{路 4:18} 无可比拟；后者远比前者伟大；因为前者没有赏报，而后者有千万倍的赏报。\par

\Person{保禄}的锁链已经证明很长，并使我们在其中停留了相当长的时间。是的，它确实很长，而且比任何金链都更美丽。这是一条锁链，借着一种无形的机械，仿佛一根金线放下来，将那些被它捆绑的人拉向天堂，并拉上诸天之上的天。奇妙的是：它虽在下面捆绑，却将它的俘虏向上拉；而就事物本身而言，这并非其本性。但当天主命令和安排时，就不要寻求事物的本性，也不要寻求自然的意义，而要寻求超乎本性和自然顺序的事物。\par

让我们学会不在苦难中沉沦，也不抱怨；请看这位有福的圣人。他曾受鞭打，且受重鞭，因为经上说：\Quote{当他们打了许多鞭在他们身上时。} 他也曾被捆绑，而且又是重重地捆绑，因为狱吏将他投入内监，并特别严密地看守。尽管他身处如此多的危险，在午夜，即使最清醒的人也沉睡时，另一种更牢固的束缚临到他们，他们却歌唱赞美主。有什么比这些灵魂更坚硬呢？他们想到三位圣童即使在火和火窑中也歌唱。\BibleRef{达 3:1} 或许他们就这样对自己推理：\Quote{我们至今还没有遭受过那样的事。}\par

但是我们的讲论做得很好，因为它这样将我们带到另一些锁链和另一个监狱。我该做什么呢？我本想保持沉默，但我不能。我发现了另一个监狱，比前一个更奇妙、更令人惊讶。但是，来吧，现在就振作起来，好像我刚开始讲论，用你们清新的头脑留意听。我本想中断讲论，但它不让我；因为正如一个人在饮宴中不能忍受中断他的饮宴，无论别人向他承诺什么；同样，我既已抓住这个为基督被捆绑之人的监狱的荣耀之杯，就不能停，不能缄默。因为如果\Person{保禄}在监狱中，在夜间，没有保持沉默，不，甚至在鞭打之下；那么我，大白天坐在这里，如此自在地说话，我能沉默吗？当那些被捆绑、在鞭打之下、在午夜都不能忍受沉默时？三位圣童没有沉默，不，不在火窑和火中，我们难道不因沉默而羞耻吗？那么，我们也来看看这个监狱。这里，他们也都被捆绑，但立刻，从一开始就明显，他们不是要被焚烧，只是进入一个监狱。因为你们为什么要捆绑那些即将被投入火焰的人呢？他们被捆绑，像\Person{保禄}一样，手和脚。他们被捆绑得和他一样猛烈。因为狱吏将他投入内监；王命令火窑被烧得极其炽热。现在让我们看结果。当\Person{保禄}和\Person{息拉}歌唱时，监狱被震动，门被打开。当三位圣童歌唱时，他们手脚的锁链都被解开了。监狱被打开，火窑的门也被打开：因为一阵带有露水的微风吹过其中。\par

但许多思绪围绕着我。我不知道该先说什么，后说什么。所以，我恳求，不要有人要求我按顺序，因为主题紧密相连。\par

与\Person{保禄}和\Person{息拉}一同被捆绑的人松了绑，然而他们却仍然睡觉。至于三个青年，却发生了别的事。把他们扔进去的人，自己倒被烧死了。然后，正如我乐意告诉你们的，王看见他们脱了绑，就俯伏在他们面前；他听见他们唱赞美歌，看见四个人行走，便呼唤他们。正如\Person{保禄}虽然有能力出来，却没有出来，直到把他关进去的人呼唤他，才把他带出来；同样，三个青年也没有出来，直到把他们扔进去的人命令他们出来。我们从中得到什么教训？不可过于急躁地招惹迫害，苦难中也不可过于急切地求解放，另一方面，当他们释放我们时，也不可继续留在苦难中。此外，狱警因为能进入圣人所在的地方，就俯伏在他们脚前。王只来到门口就俯伏。他不敢走进他为他们在火中预备的监牢。现在请注意他们的话。一个喊道：\Quote{诸位，我该做什么才能得救？}\BibleRef{宗 16:30}另一个，虽然不那么谦卑，却发出同样甜美的声音：\Quote{沙得辣客、默沙客、阿贝得乃哥，至高天主的仆人，出来，到这里来吧。}\BibleRef{达 3:26}多么伟大的尊严！\Quote{至高天主的仆人，出来，到这里来吧。}王啊，他们怎能出来呢？你把他们捆绑着扔进火里；他们在火中已经这么久了。如果他们是用金刚石做的，如果是金属块，在唱完那整首赞美诗时，岂不早已灭亡了？正是为此，他们得救了，因为他们赞美了天主。火尊敬了他们甘愿受苦的心，之后又尊敬了那美妙的歌声和他们的赞美诗。你用什么样的称号称呼他们？我先前说过：\Quote{至高天主的仆人。}是的，对天主的仆人，一切都有可能；因为如果有些人是人的仆人，尚且拥有权力、权威，并处理自己的事务，那么天主的仆人岂不更甚？他用他们最喜爱的名字称呼他们；他知道这样最能讨好他们。因为如果他们进入火中正是为了继续作天主的仆人，那么没有比这更令他们欣喜的声音了。即使他称他们为君王，称他们为世界的主宰，也比不上他称他们为\Quote{至高天主的仆人。}更能使他们喜乐。这有什么可惊奇的？\Person{保禄}在给那座强大的城市——她是世界的女主人，以自己崇高的尊严而自豪——写信时，把这称号视为与执政官、君王甚或世界帝国的尊严同等，不，远超过它们，无可比拟地更大，这称号：\Quote{保禄，耶稣基督的仆人。}\BibleRef{罗 1:1}\Quote{至高天主的仆人。}\Quote{是的，}他说，\Quote{如果他们表现出如此大的热心作仆人，毫无疑问，这就是我们要赢得他们的称呼。}\par

同样，也要注意这三个青年的虔诚：他们没有表现出愤慨、愤怒、抗议，而是走了出来。如果他们认为被扔进火窑是报复行为，他们就会对扔他们的人感到伤心；但事实并非如此；反而就像是从天堂出来一样，他们走了出来。先知论及太阳时说，\Quote{他好像新郎走出洞房，}\BibleRef{咏 18:5}这样说他们也不会错。但虽然太阳如此出来，他们却出来得更荣耀，因为太阳来是用自然光普照世界，而他们是用另一种方式，我是说，属灵的方式。因为他们，王立即颁布谕令，其中有这些话：\Quote{我以为最好宣布至高天主在我身上所行的神迹奇事。他的神迹何其伟大！他的奇事何其可畏！}\BibleRef{达 4:2}因此他们出来，发出更荣耀的光辉，确实在该地区照耀，但更重要的是，通过王的谕书，能够传遍世界，从而驱散各处弥漫的黑暗。\Quote{出来，}他说，\Quote{到这里来吧。}他没有下令扑灭火焰，而是由此特别尊重了他们，因为他相信他们不仅能在火中行走，甚至能在火仍燃烧时从中出来。\par

但是，如果你们觉得好，让我们再看一看狱卒的话：\Quote{诸位，我该做什么才能得救？}还有什么话比这更甜美？就连天使也因此欣喜跳跃。听到这样的话，连天主的独生子自己也成了仆人。起初相信的人曾对\Person{伯多禄}说了这样的话。\BibleRef{宗 2:37}\BibleQuote{我们该做什么？}他回答说：\BibleQuote{你们要悔改，并受洗。}听到犹太人说出这样的话，\Person{保禄}甘愿为他们下地狱，因他切望他们得救并服从。但请注意，他把这事完全交给他们，毫不浪费无用之功。然而，让我们看看下一件事。这里的君王没有说\Quote{我该做什么才能得救？}，但他的教训比任何言语都更清楚；因为他立刻成为宣讲者，不需要像狱卒那样受教。他宣扬天主，并承认祂的大能。\BibleQuote{你们的天主确实是万神之神、万君之主，因为祂派了祂的天使，解救了你们。}\BibleRef{达 2:47}结果如何？不是只有一个狱卒，而是许多人因君王的文字和眼见的事实而受教。因为君王不会说谎，这是每个人都清楚的，他决不愿为俘虏作这样的见证，也不愿推翻自己的作为；他决不愿招致如此疯狂的指控：所以除非真理充分显明，他绝不会在有如此多人见证的情况下写下这些话。\par

你们可看出锁链的力量何等伟大？在患难中歌唱的赞美何等有力？他们的心没有衰竭，没有沮丧，反而更加充满活力，他们的勇气也更加强大，这合情合理。\par

当我们思考这些事时，还有一个问题留给我们：为什么一方面在监狱里囚犯被释放，而另一方面在火窑中执行者却被烧死？因为那本应是君王的命运，因为捆绑他们的人或把他们扔进火窑的人，其罪并不大于下令做这事的人。那么为什么他们灭亡了？关于这一点没有必要仔细考察；因为他们是恶人。因此，这是天主的安排，为要显示火的力量，并使奇迹更加显著；因为如果火这样吞噬了外面的人，又如何显示里面的人毫发无伤呢？这是为使天主的大能得以彰显。也不要有人惊讶我把君王与狱卒并列，因为他做了同样的事；一个毫不比另一个更高贵，他们都得到了各自的报应。\par

但是，如同我说过的，义人在患难时，尤其是在捆绑中，更为充满活力；因为为基督的缘故受任何苦，是所有安慰中最甘甜的。\par

你们要我提醒你们另一个监狱吗？似乎有必要从这条锁链转向另一个监狱。你们想要哪一个？是\Person{耶肋米亚}的，\Person{若瑟}的，还是\Person{若翰}的？感谢\Person{保禄}的锁链，它为我们的讲论打开了多少监狱！你们愿意听\Person{若翰}的吗？他也曾为基督并为天主的法律而被捆绑。那么，他在监狱里无所事事吗？他不是从那里派门徒去问：\BibleQuote{祢是那要来的一位，还是我们该等待另一位？}\BibleRef{玛 11:2}看来即使在那里，他也在教导，因为他没有忽略自己的职责。再说，\Person{耶肋米亚}不是也曾预言关于巴比伦君王的事，并在监狱里完成他的工作吗？\Person{若瑟}呢？他不是在监狱里过了十三年吗？那又怎样？即使在那里，他也没有忘记他的美德。我还要提到一个人的捆绑，以此来结束我的讲论。我们的主自己曾被捆绑，祂解除了世界的罪恶。那双手被捆绑，那手行了无数善事。因为经上说：\BibleQuote{他们捆绑了祂，解去交给\Person{盖法}。}\BibleRef{玛 27:2}是的，祂被捆绑了，祂曾行了如此多的奇事。\par

默想这些事，让我们永不抱怨；无论我们是否在捆绑中，让我们喜乐；或我们不在捆绑中，让我们如同与祂一同被捆绑。看，捆绑是何等大的祝福！知道这一切，让我们为万事向天主献上感恩，藉着我们的主\Person{耶稣基督}，愿光荣、权能和尊威归于祂及父和圣神，从今世直到永远。阿们。\par

\SourceRecord{Translated by Gross Alexander.；Nicene and Post-Nicene Fathers, First Series；Vol. 13.；Edited by Philip Schaff.；Buffalo, NY: Christian Literature Publishing Co.,；1889.；<http://www.newadvent.org/fathers/230108.htm>.}

% structure: p0001 source=h1 target=section reason=page-title

\section{\WorkTitle{厄弗所书讲道集 第九篇}}

第四章 第1-3节 \BibleRef{弗 4:1}\par

\BibleQuote{我因此，在主内的囚犯，恳求你们，行事为人要与你们所蒙的召叫相称，完全以谦卑和温良，以坚忍，用爱德彼此担待；尽力以和平的联系保持圣神的合一。}\par

\Person{保禄}的锁链所显示的力量是伟大的，比奇迹更为荣耀。那么，看来并非徒然，也非无目的地，他在这里将其展示，而是作为最能打动他们的方法。他说什么呢？\BibleQuote{我因此，在主内的囚犯，恳求你们，行事为人要与你们所蒙的召叫相称。}这是怎样的呢？\BibleQuote{完全以谦卑和温良，以坚忍，用爱德彼此担待。}\par

成为囚犯本身并非光荣，而是为基督的缘故才光荣。因此他说：\BibleQuote{在主内}，即是为基督的缘故的囚犯。没有能与这相比的。但现在这锁链更加把我从主题中拖开，又把我拉回去，我无法抗拒，而是甘愿被拖曳——是的，更确切地说，全心乐意；但愿我永远有缘论说\Person{保禄}的锁链。\par

但如今不要打盹：因为我仍想解决另一个许多人所提出的问题，他们说：如果苦难是荣耀，\Person{保禄}自己在向\Person{阿格里帕}辩护时怎么说：\BibleQuote{我向天主祈求，不论少也好多也好，不但你，而且今天所有听我的人，都能像我一样，只是不要这些锁链。}\BibleRef{宗 26:29} 他这样说并不是认为这事应当避免；不；因为如果真是这样，他就不会以锁链、监禁和其他苦难为荣了；当他在别处写信时说：\BibleQuote{我更喜欢夸耀我的软弱。}\BibleRef{格后 12:9} 但实际如何？这正是证明他认为这些锁链有多么伟大；因为在写给格林多人时他说：\BibleQuote{我给你们喝奶，没有给肉，因为你们还不能吃；}\BibleRef{格前 3:2} 这里也是如此。他说话的对象不能承受那些锁链的美、可爱和福分。因此他加了：\BibleQuote{只是不要这些锁链。} 然而对希伯来人他并没有这样说，而是劝他们\BibleQuote{要与那些在锁链中的人一同被捆绑}\BibleRef{希 13:3} 因此他自己也因锁链而喜乐，并被捆绑，被带与囚犯一同进入内监。\Person{保禄}的锁链真是有力！这样的景象足以替代其他一切：看见\Person{保禄}被捆绑，从监牢里被带出来；看见他被捆绑，坐在里面，什么样的快乐能比得上这个？为了这样的景象，我有什么不愿付出呢？你们看见皇帝、执政官乘着战车，穿着金色服饰，他们的卫兵周身都是金子？他们的金戟、金盾、金衣、金马具？这景象比那个更令人愉快多少倍！我宁愿看见\Person{保禄}一次，与囚犯一同从监牢里出来，也比看见这些一万次，带着所有随从游行更愿意。当他这样被带出来时，你们想有多少天使在他前面引路？为了证明我所说的并非虚构，我将用一个古老的故事向你们说明。\par

\Person{厄里叟}先知（也许你们知道这人），当时（\BibleRef{列下 6:8}）\Place{阿兰}王与\Place{以色列}王交战，他坐在自己家中，揭露了\Place{阿兰}王在密室里与知情者所谋划的一切，使王的计谋失效，因为他预先告知秘密，不让\Place{以色列}王落入他所设的圈套。这严重困扰了王；他灰心丧气，陷入更大的困惑，不知如何找出那个泄露一切、密谋对抗他、破坏他计划的人。因此，正当他困惑并追查原因时，他的一名持戟侍卫告诉他，有一个先知名叫\Person{厄里叟}，住在\Place{撒玛黎雅}，他不让王的计划得逞，而是揭露一切。王以为他发现了全部真相。确实，没有人比他更悲惨地被误导了。他本应尊敬那人，敬畏他，惊叹他确实拥有如此大的能力，以至于坐在数里之外，就能知道王室内的一切，而无需任何人告诉他；但他没有这样做，反而被激怒，完全被情绪左右，他装备骑兵和士兵，派遣他们去把先知带到他面前。\par

现今\Person{厄里叟}有一个门徒，尚在预言的入门阶段，\BibleRef{列下 6:13}，还远远不配获得这类启示。国君的士兵来到那地方，想要捆绑那人，或者更确切地说，捆绑先知。——我又要坠入论锁链的话题了，因为这篇讲论完全与此交织。——门徒看见大队士兵，惊恐万分，战战兢兢地跑到师傅那里，告诉他这在他看来是灾祸的事，并告知他这不可避免的危险。先知笑他惧怕不当怕的事，命他放心。然而，门徒还不完美，不听他的话，仍因所见而惊惶，留在恐惧中。于是，先知做了什么？他说：\BibleQuote{主，请开启这青年的眼目，让他看见与我们同在的多过与他们同在的。}\BibleRef{列下 6:16} 他立刻看见先知当时所住的整座山，充满了如此众多的火马和火车。这些无非是天使的队伍。但若仅因这样一次机会，就有这样一支庞大的天使队伍陪伴\Person{厄里叟}，那么\Person{保禄}又有什么呢？这正是\Person{达味}先知告诉我们的。\BibleQuote{主的使者在敬畏他的人周围安营。}\BibleRef{咏 33:7} 又说：\BibleQuote{他们要用双手托着你，免得你的脚碰在石头上。}\BibleRef{咏 90:2} 我为什么提及天使呢？主亲自与他同在，当他前行时；因为不可能主显现给了\Person{亚巴郎}，却没有与\Person{保禄}同在。不，这是他的应许：\BibleQuote{我天天与你们同在，直到今世的终结。}\BibleRef{玛 28:20} 再者，当他向他显现时，他说：\BibleQuote{不要害怕，只管讲，因为我与你同在，没有人会下手害你。}\EditorialNote{96 9, 10} 又，他在梦中站在他旁边，说：\BibleQuote{放心吧，你怎样在耶路撒冷关于我作了见证，也必怎样在罗马为我作见证。}\BibleRef{宗 23:11}\par

圣人虽然在任何时候都是荣耀的景象，充满丰富的恩宠，但最是如此的是当他们为基督的缘故处于危险中、当他们是囚犯时；因为正如一位勇敢的士兵，在任何时候、出于自身，对观看者而言都是悦目的景象，但最是在他站立、列队在君王身旁时；同样，你们也想象\Person{保禄}，看见他在锁链中施教是何等伟大的事。\par

我是否可以顺便提一个此刻涌现在我心中的想法？真福殉道者\Person{巴比拉}曾被捆绑，而且他也是为了与\Person{若翰}相同的原因，因为他责备了一位国王的过犯。此人临终时吩咐，要将他的锁链与身体一起安放，并将身体捆绑着埋葬；直到今日，这些脚镣仍与他的骨灰混在一起，他对为基督的缘故佩戴过的锁链竟如此眷爱。正如先知论及\Person{若瑟}所说：\BibleQuote{他被锁在铁链中。}\BibleRef{咏 104:18} 就连妇女们以前也曾经历过这些锁链。\par

然而我们并不在锁链中，我也不是推荐这个，因为现在不是它们的时候。但你们，不要捆绑你们的手，而要捆绑你们的心和思想。还有别的锁链，那些不佩戴这一种的，将必须佩戴另一种。请听基督所说：\BibleQuote{捆住他的手和脚。}\BibleRef{玛 22:13} 但愿天主使我们免于经历那些锁链！愿祂使我们甚至饱尝这些锁链！\par

因这些缘故，他说：\BibleQuote{我这为主被囚的，劝你们，行事为人要与你们所蒙的召叫相称。} 但这召叫是什么？你们被召叫为祂的身体，经上说。你们以基督为头；虽然你们曾是\BibleQuote{仇敌}，犯了无数的过犯，然而\BibleQuote{祂已叫你们与祂一同复活，并叫你们与祂一同坐在天上。}\BibleRef{弗 2:6} 这是崇高的召叫，达到崇高的特权，不仅因为我们从先前状态被召叫，而且因为我们被召叫享有这样的特权，并借由这样的方法。\par

但如何可能\BibleQuote{相称地行事为人}？\BibleQuote{以各样的谦卑。} 这样的人行事为人相称。这是所有德行的基础。你若谦卑，并思想你是什么，以及你是如何得救的，你就会以此忆念作为导向一切德行的动机。你既不以锁链自傲，也不以我刚才提到的那些特权自傲，反而因知道一切都是出于恩宠，便谦卑自己。谦卑的人能同时成为慷慨和感恩的仆人。他说：\BibleQuote{你有什么不是领受的呢？}\BibleRef{格前 4:7} 再者，请听他的话：\BibleQuote{我比他们众人更劳苦；然而不是我，而是天主的恩宠与我同在。}\BibleRef{格前 15:10}\par

他说：\BibleQuote{以各样的谦卑}；不是口头的谦卑，也不是仅仅行为上的谦卑，而甚至是在人的举止和声调中：不是对一人谦卑，对另一人粗鲁；要谦卑地对待所有人，无论是朋友还是仇敌，是尊贵还是卑微。这就是谦卑。甚至在你的善行中也要谦卑；因为请听基督所说：\BibleQuote{神贫的人是有福的；}\BibleRef{玛 5:3} 祂将此放在首位。因此，宗徒自己也说：\BibleQuote{以各样的谦卑、温和与忍耐。} 因为一个人可能谦卑，却急躁易怒，如此一切都归于徒然；因为他常会被怒气控制，破坏一切。\par

他继续说：\BibleQuote{要以爱心互相宽容。}\par

如果一个人易怒或挑剔，怎么可能宽容？因此他告诉了我们方式：他说：\BibleQuote{以爱心。} 他可以说，如果你不宽容你的邻人，天主怎能宽容你呢？如果你不容忍你的同仆，主人又怎能容忍你呢？哪里有了爱，一切都能忍受。\par

他说：\Quote{殷勤致力}，又说：\Quote{以和平的联系，保持圣神的合一。}因此，用节制捆住你的双手。再说，那\Quote{联系}的美名。我们先前已经放下了它，它又自己回到我们这里。那是美好的联系，这也是美好的，而那另一个是这联系的果实。将你自己与你的兄弟绑在一起。凡在爱中彼此连结的，就轻松地承受一切。将你自己与他绑在一起，也把他与你绑在一起；你是两者的主人，因为无论我想要使谁成为我的朋友，我都能藉着和善而达成。\par

他说：\Quote{殷勤致力}；这不是一件容易做到的事，也不是每个人都能做到的。\par

他说：\Quote{殷勤致力}，接着又说：\Quote{保持圣神的合一。}什么是\Quote{圣神的合一？}在人身上，有一种神将各部分连结在一起，尽管肢体不同。这里也是如此；因为圣神被赐下，正是为了要使那些因种族和不同习俗而分离的人合一；因为老年与青年、富人与穷人、儿童与青年、女人与男人，以及每一个灵魂，都在某种程度上成为一体，而且比若有一个身体更完全。因为这种属灵的关系远高于那自然的关系，并且合一的完美更完整，因为灵魂的结合更完美，因为它既单纯又统一。那么，这种合一如何保持呢？\Quote{在和平的联系中。}在仇恨与不和之中，这是不可能的。因为他说：\Quote{你们中间有嫉妒和纷争，岂不是属乎肉体，照着世人的样子行吗？}\BibleRef{格前 3:3}因为火遇到干燥的木块，就把它们一起烧成一堆熊熊烈火，但遇到湿的木块，就完全不作用，也不能使它们合一；这里也是如此。凡是冰冷的本性都不能促成这种合一，而温暖的本性大多能。因此，至少产生了爱德的温暖；藉着\Quote{和平的联系}，他渴望将我们全都连结在一起。因为正如他所说，如果你要依附另一个人，除了使他依附于你自己之外，别无他法；如果你希望使这纽带加倍，他也必须反过来依附于你；同样，在此他希望我们彼此连结；不只是要我们和平相处，不只是要我们彼此相爱，而是要使所有人都成为同一个灵魂。这是光荣的纽带；藉此纽带，让我们将自己彼此连结，并连结于天主。这是一种不会擦伤、也不会束缚被绑之手的纽带，反而使它们自由，给它们充分的活动空间，并给它们比自由者更大的勇气。强者若与弱者捆绑，会扶持他，不让他灭亡；若再与懒惰者捆绑，反而会激发他、鼓励他。经上说：\Quote{兄弟互助，如坚城。}这链条没有任何距离可以中断，无论是天上、地上、死亡，还是任何其他事物，它比一切更强而有力。这链条即使只从一个灵魂发出，也能同时拥抱许多人；请听保禄所说：\Quote{你们在我们里面并不受限制，你们是在自己的情感中受限制；你们也要被扩大。}\BibleRef{格后 6:12}\par

那么，是什么破坏了这纽带？贪爱金钱、追求权力、贪图荣耀等，使它们松弛、分裂。那么我们如何确保它们不被切断？如果这些情绪被消除，并且那些毁灭爱德的事物没有中途来扰乱我们。因为请听基督所说：\BibleQuote{因为不法的事增多，许多人的爱就冷淡了。}\BibleRef{玛 24:12}没有什么比罪更与爱对立，我指的不只是对天主的爱，也包括对邻人的爱。但有人会说，甚至强盗之间也有和平吗？什么时候？告诉我。不是当他们怀着强盗的精神行事时；因为如果他们不在分赃的人之间遵守正义的规则，并给每个人应得的，你也会发现他们在战争和纷争中。所以即使在恶人中间也找不到和平；但在人过着正义和美德生活的地方，你到处都可以找到。但再说，有竞争关系的人之间会有和平吗？永远不会。那么你要我提谁呢？贪婪的人不可能与贪婪的人和平相处。所以，如果没有正直和善良的人（即使被他们冤枉）站在他们中间，他们的整个种族都会被撕碎。当两只饥饿的野兽，如果没有什么东西放在它们中间供它们吞噬，它们就会互相吞食。贪婪和邪恶的人也是如此。所以，不预先实行美德，就不可能有和平。如果我们愿意，让我们想象一座完全由贪婪之人组成的城市，给予他们平等的权利，没有人愿意被冤枉，但所有人都互相伤害。那座城市可能维持下去吗？不可能。再有，奸淫者之间有和平吗？不，你找不到任何两个人同心合意。\par

因此，回到正题：除此之外，再没有别的原因，只是因为\Quote{爱已冷却了}；爱已冷却的原因又在于\Quote{不法的事增多}。因为这导致自私，分裂和分离身躯，使它松弛并撕裂。但在有德行的地方，情况正相反。因为有德行的人也超越金钱；所以即使有一万这样的人贫困，他们也会平安相处；而贪婪的人，即使只有两个，也无法和平。所以，如果我们有德行，爱就不会消失，因为德行出自爱，爱出自德行。我将告诉你这是怎么回事。有德行的人不会把金钱看得比友谊更重要，也不会记仇，不会伤害邻人；他不傲慢，他大度忍受一切。爱就是由这些组成的。另外，爱的人也会忍受这些事，就这样他们相互产生。而且，爱出自德行，这一点可以从这里看出：因为我们的主说：\Quote{因为不法的事将增多，许多人的爱将冷却}，这明明告诉我们这一点。而德行出自爱，保禄告诉我们说：\Quote{爱邻人的，就完全了法律}。\BibleRef{罗 13:10} 所以，人必须是两者之一：要么非常温柔且被爱，要么非常有德行；因为拥有一者的人，必然拥有另一者；反之，不知道怎么爱的人，将因此犯许多恶行；而犯恶行的人，也不知道爱是什么。\par

劝言。让我们因此追求爱德；它是一种保护，不会允许我们受任何恶。让我们将自己绑在一起。让我们中间没有欺骗，没有虚假。因为在有友谊的地方，这些事都找不到。另一位智者也这样告诉我们：\Quote{你拔出刀来对付朋友，也不要失望，因为还可以回心转意。你张嘴对付朋友，不要害怕，因为还可以和好；但除非是责斥，或暴露秘密，或背叛的伤害，因为这些事使朋友离开}。\BibleRef{德 22:21} 因为他说：\Quote{暴露秘密}。现在，如果我们都是朋友，就没有必要保守秘密；因为没有人向自己保守秘密，无法对自己隐瞒任何事，所以也不会向朋友隐瞒。因此，在秘密不存在的地方，由此产生的分裂是不可能的。我们保守秘密，没有别的原因，只是因为我们不信任所有人。所以，爱的冷却产生了秘密。你有什么秘密呢？你想要伤害邻人吗？或者，你阻碍他共享某些利益，因此隐瞒它？但也许这些都不是。那么，是你害羞吗？如果是，那就是不信任的征兆。现在，如果有爱，就不会有\Quote{暴露秘密}，也不会有\Quote{责斥}。因为，请告诉我，谁会责斥自己的灵魂呢？即使这样的事发生，也是为了某种好处；因为我们知道，我们责斥孩子，是为了让他们有感觉。同样，基督也曾开始责斥那些城市，说：\Quote{唉，苛辣匝因！唉，贝特赛达！}\BibleRef{路 10:13} 为的是解救他们免受责斥。因为没有什么比这更能抓住心志，更能激励它，或在它松弛时紧张它。让我们从不要只为了责斥而相互责斥。因为什么呢？你会因为金钱的事责斥你的朋友吗？当然不会，如果你至少把你拥有的看作是公共的。那么，你会因为他的过失责斥吗？也不是，而是你会纠正他。或者，就像后文所说的：\Quote{为了背叛的伤害}；谁会杀害自己，或伤害自己呢？没有人。\par

让我们因此\Quote{追求爱德}；他不是简单地说\Quote{我们要爱}，而是说\Quote{让我们追求爱德}。\BibleRef{格前 14:1} 这需要大大的热情；她很快就消失在视线之外，她的飞逃极快；生活中有太多事伤害她。如果我们追赶她，她就不会超越我们并逃走，而是我们将快速夺回她。天主的爱是连接地与天的。天主的爱让人坐在王权的宝座上。天主的爱让天主显现在地上。天主的爱让主成为仆人。天主的爱让被爱者被交付为祂的敌人，圣子为恨祂的人，主为祂的仆人，天主为人，自由者为奴隶。它并未止于此，还召唤我们到更大的事。是的，它不仅解放我们脱离先前的恶，还应许赐予我们更大的福分。因此，让我们感谢天主，并追求每一种德行；其中首要的，让我们严格地实践爱德，以便我们能堪当承受所应许的福分；藉着我们的主耶稣基督的恩宠和仁慈，愿光荣、权能和尊威归于他，与圣父及圣神同在，自今以至永远。阿门。\par

\SourceRecord{Translated by Gross Alexander.；Nicene and Post-Nicene Fathers, First Series；Vol. 13.；Edited by Philip Schaff.；Buffalo, NY: Christian Literature Publishing Co.,；1889.；<http://www.newadvent.org/fathers/230109.htm>.}

% structure: p0001 source=h1 target=section reason=page-title

\section{厄弗所书讲道 第十篇}

\BibleRef{弗 4:4}\par

\BibleQuote{只有一个身体，一个圣神，正如你们蒙召，同有一个蒙召的望德。}\par

当受颂扬的保禄劝勉我们去做任何特别重要的事时，他是一位如此真正智慧且属神的人，便将他的劝勉建立在天上的事物上；这本身就是他从主那里学到的教训。因此他在别处也说：\BibleQuote{你们要在爱德中行走，正如基督也爱了我们。}\BibleRef{弗 5:2} 又说：\BibleQuote{你们该怀有基督耶稣所怀有的心情，祂虽具有天主的形体，却不以自己与天主同等为持守不舍的。}\BibleRef{斐 2:5} 他在这里也是这样做的，因为每当他摆在我们面前的榜样是伟大的，他就充满了热忱与情感。那么，当他激励我们合一时，他说了什么？\BibleQuote{只有一个身体，一个圣神，正如你们蒙召，同有一个蒙召的望德。}\par

\BibleQuote{只有一个主，一个信德，一个圣洗。}\BibleRef{弗 4:5}\par

那么这个身体是什么？是全世界所有的信徒，无论是过去、现在还是将来的。而且，在基督来临之前悦乐天主的那些人，也是\Quote{一个身体}。如何做到的呢？因为他们也认识基督。这从何显明？祂说：\BibleQuote{你们的父亲亚巴郎，}祂说，\BibleQuote{欢欣喜乐地要看见我的日子，他看见了，便欢喜快乐。}\BibleRef{若 8:56} 又说：\BibleQuote{如果你们相信梅瑟，}祂说，\BibleQuote{你们也必相信我，因为他曾论及我而写作。}\BibleRef{若 5:46} 先知们也不会写一位他们不知所言的对象，相反他们既认识祂，也朝拜祂。这样，他们也是\Quote{一个身体}。\par

身体与灵魂并非分离，否则它便不成其为身体。我们也常有这样的习惯，论及那些联合在一起、有极大一致性与连贯性的事物时，就说它们是同一个身体。同样，对于合一，我们也将那处于同一首领之下的视为一个身体。如果只有一个首领，那么就只有一个身体。身体由尊贵与不尊贵的肢体组成。只是较大的肢体不可对最小的肢体傲慢，后者也不可嫉妒前者。它们并非都贡献相同的一份，而是各自按照需要的比例。既然所有肢体都是为着必要和不同的目的而形成的，它们都有同等的尊荣。的确，有些是主要的肢体，其他则不那么主要：例如，头比身体其他所有部分更是主要的肢体，因为它容纳所有感官和灵魂的统治原则。没有头就不可能存活，而许多人失去了脚还能活很长时间。所以头不仅因其位置，还因其生命力和功能而优于脚。\par

我为什么说这些？在教会中有大量的人：有些人像头一样被高举，像眼中的眼睛一样审视天上的事物，他们远离地面，与地面毫无共通之处；而另一些人则占据脚的位置，行走在地上——是健康的脚，因为在地上行走对于脚并非罪过，但奔跑作恶却是罪。先知说：\BibleQuote{他们的脚奔跑作恶。}\BibleRef{依 59:7} 因此，他说，让这些头也不可对脚高傲，脚也不可嫉妒头。因为这样，各自的独特美丽就被破坏了，其功能的完美也受到了阻碍。这很自然，因为那对邻人设圈套的人首先要设圈套给自己。如果脚不愿将头带到任何必要之旅，它们的懒惰和懈怠也将同时伤害自己。或者，如果头不愿照顾脚，它自己将首先遭损。然而，这些肢体不会彼此相争；这不可能，因为自然已经如此安排。但对于人，怎能不使人攻击人呢？我们知道，没有人攻击天使，因为他们也不攻击总领天使。同样，无理性的受造物也不能在我们面前傲慢自大；但哪里本性同等尊贵，恩赐唯一，且谁也不比谁多，这又如何阻止呢？\par

然而，这些也正是你不该起来反对邻人的理由。因为如果万物都是共有的，没有谁比别人多什么，那么这疯狂从哪里来？我们享有同一本性，同样分享灵魂与肉身，呼吸同样的空气，食用同样的食物。这彼此对抗的叛乱从何而起？\newline{}诚然，凭自己的德行克服无形之能，足以导致傲慢；或者更确切地说，这并非傲慢，因为我理所当然地高昂着头，甚至极其高昂地对抗那恶灵。看，保禄如何对那恶灵高傲。当恶灵讲论关于他的伟大奇妙之事时，他使之沉默，连谄媚也不容忍。因为那个\Quote{有占卜之神的}使女喊着说：\BibleQuote{这些人是至高天主的仆人，他们向我们传报救恩的道路。}\BibleRef{宗 16:16}，他严厉斥责她，制止了她放肆的舌头。他又在别处写道：\BibleQuote{天主将很快把撒殚踏碎在你们脚下。}\BibleRef{罗 16:20}\newline{}本性的差异会产生任何影响吗？你难道没有察觉本性的差异毫无影响，唯有意志的差异才起作用吗？因此，由于他们的原则，他们远比一切更坏。\newline{}有人说，但我并不起来反对天使，因为我的本性与其相距甚远。然而，你实在不该起来反对人，也不该反对天使；因为天使在本质上与你不同，这对他既不是荣耀，对你也不是耻辱；而人与人之间在本性上毫无差异，只在原则上不同；甚至在人类中也有天使。所以，如果你不起来反对天使，就更不该反对人，反对那些在我们这本性中成为天使的人；因为若有人在德行上如同天使，那人比你优越的程度远胜于天使。为什么？因为那属于本性的，他人凭自己的意志实现了。再者，因为天使的家乡在距离上也远离你，居住在天堂；而这人同你生活在一起，激励你效仿。他实际上比天使离你更远。因为宗徒说：\BibleQuote{我们的国籍是在天堂。}\BibleRef{斐 3:20}为了让你明白这人的家乡更为遥远，请听他的头坐于何处——宝座上，他说，那王座上！那宝座离我们越远，他也越远。\par

你会说，但我看他享受尊荣，便生嫉妒。这正是颠倒是非、不仅给世界也给教会带来无数麻烦的根源。正如狂风突袭平静的港湾，使其比任何礁石或海峡都危险；同样，贪图光荣进入，倾覆并搅乱一切。\par

你常在大型房屋火灾中到场。你见过浓烟升天；若无人前来遏制灾祸，人人都只顾自己，火焰便自由蔓延，吞噬一切。有时整座城的人会围在四周，他们确实作为祸害的旁观者围观，却不协助救援。在那里，你可以看到他们所有人都站在那里，什么事也不做，只是每人伸出手，向新来的人指出此刻从窗口飞出的火舌，或倒塌的梁木，或整圈墙壁被推倒并猛烈地砸在地上。更有胆大冒失者，甚至敢于在燃烧时靠近房屋本身，不是为了伸手相助以制止灾祸，而是为了更充分地欣赏景象，能从近处仔细观察那些远处常被忽略的一切。\newline{}如果那房屋又大又华丽，对他们来说便是可怜的场景，令人落泪。确实，我们有可怜的场景可看：柱头化为灰尘，许多柱子本身被击碎，有的被火吞噬，有的甚至被立起它们的手推倒，以免助长火焰。曾经优雅地站立、承受天花板塑像的雕像，如今全暴露在外，屋顶被掀开，丑陋地裸露在露天中。为何还要描述里面存储的财富？金织物和银器皿？而那家主与主妇几乎不进入的地方，织物和香料的宝库，珍贵宝石的匣子——全成为一团烈火，如今在里面的是浴室工人、街道清洁工、逃跑的奴隶和所有人；里面一切都成为火与水的混合物、泥、尘土和半烧焦的梁木！\par

如今我为何要如此详细地描绘这幅景象？并非仅仅因为我想要向你们展示一所房屋的火灾，（这与我何干？）而是因为我想要尽可能生动地将教会的灾难呈现在你们眼前。因为确实，就像一场大火，或像从高处掷下的雷霆，它们落到了教会的屋顶上，却仍没有唤醒任何人；然而，当我们天父的房屋正在燃烧时，我们却仿佛沉睡，一种深沉而愚蠢的沉睡。然而，有谁不被这场火触及呢？教会的雕像中哪一个呢？因为教会无非是由我们人的灵魂建造的一座房屋。这座房屋并非处处具有同等的尊荣，而是构成它的石头，有些明亮发光，有些则较小较暗淡，但又优于其他。在那里，我们可能会看到许多人也是处于金子的位置，那装饰天花板的金子。另一些人，我们可能会看到，他们赋予了雕像所产生的美丽和优雅。许多人，我们可能会看到，像柱子一样矗立着。因为他也习惯于称人为\Quote{柱子}\BibleRef{迦 2:9}，不仅是因为他们的力量，也是因为他们的美丽，他们增添了极大的恩宠，并且他们的头上覆盖着金子。我们可以看到一大群人，通常构成宽阔的中部空间和整个圆周的延伸；因为整体身体占据了那些构成外墙的石头的地位。或者更恰当地说，我们必须转向更辉煌的图画。我所说的这座教会，并非由我们周围所见到的这些石头建成，而是由金、银和宝石建成，各处散落着丰富的金子。但是，哦，这令人痛哭！因为虚荣的无法律统治已经吞噬了这一切；那吞噬一切的火焰，还没有任何人能控制它。我们站着惊愕地凝视火焰，却再也不能扑灭这邪恶；或者，如果我们扑灭了它一小段时间，两三天后，从灰烬中吹起的火花又会推翻一切，消耗与先前一样多的东西，这里也是如此：因为这就是在这种火灾中通常会发生的事。至于原因，它已经吞噬了教会柱子的支撑物；那些支撑屋顶的人，以及从前使整个建筑保持在一起的人，它已将之包裹在火焰中。因此，也很容易蔓延到外墙的其余部分：因为在建筑物中也是如此，当火抓住木料时，它更有力地攻击石头；但当它推倒柱子并夷为平地时，则无需更多就能将一切余烬吞噬在火焰中。因为当上部的支柱和支撑物倒下时，那些部分本身也会很快随之倒下。此刻教会也是如此：火已抓住每一个部分。我们寻求来自人的荣誉，我们为光荣而燃烧，我们不听约伯的话，他说：\par

\Quote{如果像亚当（或像人一样）我遮盖了我的过犯\newline{}将我的不义藏在怀中，\newline{}因为我惧怕那大群的人。}\par

看哪，一种有德的精神？他说，我不以在全众面前承认我无心的罪为耻；若他毫不羞愧地承认，我们更当如此行。因为先知说：\Quote{提出你的理由，好让你成义。}\BibleRef{依 43:26}这邪恶的势力极大，一切都因此被倾覆消灭。我们离弃了主，成了荣誉的奴隶。我们不再能责备那些在我们管辖之下的人，因为我们自己也患上了与他们相同的狂热。我们被天主指派去医治别人，自己却需要医生。当医生本身也需要别人的医治之手时，还有什么康复的希望呢？\par

我说这些并非没有目的，也不是无谓的哀叹，而是为了让全体，连同我们的妇女和孩子，撒上灰烬，披上麻衣，长期守斋，恳求天主亲自向我们伸出祂的手，止住这危险。因为确实需要祂的手，那大能而奇妙的手。我们需要比尼尼微人更大的事。先知说：\Quote{还有三天，}\Quote{尼尼微就要倾覆。}\BibleRef{纳 3:4}一个可怕的消息，充满了巨大的威胁。怎么会不是呢？预期三天之内，城市成为他们的坟墓，所有人都要灭亡于同一个审判。因为若在同一所房屋里有两个孩子同时死去，这艰难便无法忍受；而对约伯来说，最不能忍受的似乎是屋顶塌落压在他所有孩子身上，使他们丧命；那么，看到不是一所房屋，也不是两个孩子，而是十二万人的民族被埋葬在废墟之下，又当如何呢？\par

你们知道这是何等可怕的灾难，因为近来这个警告就发生在我们身上，并非有先知发出声音——因为我们不配听到这样的声音——而是那警告从高处发出，比任何号角更为清晰响亮的呼喊。然而，正如我所说，先知说：\BibleQuote{三天之后，尼尼微必要毁灭。}这警告实在可怕，但现在我们甚至没有类似的情况；不，再也没有\Quote{三天}了，也没有\Place{尼尼微}将被毁灭，而是许多日子已经过去，普世教会已被毁灭，夷为平地，所有人都同样被邪恶淹没；不仅如此，那些身居高位者所受的压力甚至更大。因此，若我劝你们做比尼尼微人更大的事，不要惊奇；为什么呢？不仅如此，我现在不仅宣告斋戒，而且向你们提出那个也曾扶起那正在陷落的城市的补救方法。那是什么呢？先知说：\BibleQuote{天主看见了他们的作为，就是他们离开了邪恶的道路，天主就后悔了，不再降祂所说的灾祸给他们。} \BibleRef{纳 3:10}让我们这样做，无论是我们还是你们。让我们弃绝对财富的贪欲，对荣耀的贪求，恳求天主伸出祂的手，扶起我们跌倒的肢体。我们很应该如此，因为我们的恐惧与他们的对象不同；因为那时确实是只有石头和木材将要倒塌，身体将要灭亡；但现在这些都不是；不，而是灵魂将被交付于地狱的火。让我们恳求，让我们向祂认罪，让我们为过去的事感谢祂，让我们祈求未来的事，好使我们堪当被救脱离这凶残可怖的怪兽，并将我们的感恩献予仁爱的天主圣父——愿光荣、权能和尊威归于祂，归于圣子，偕同圣神，从现在，从今世，直到永远。阿们。\par

\SourceRecord{Translated by Gross Alexander.；Nicene and Post-Nicene Fathers, First Series；Vol. 13.；Edited by Philip Schaff.；Buffalo, NY: Christian Literature Publishing Co.,；1889.；<http://www.newadvent.org/fathers/230110.htm>.}

% structure: p0001 source=h1 target=section reason=page-title

\section{\WorkTitle{厄弗所书讲道集 第十一篇}}

\BibleRef{弗 4:4-7}\par

\BibleQuote{只有一个身体和一个圣神，正如你们蒙召，同有一个蒙召的希望；一主，一信，一洗；一个天主和万有之父，祂超越万有，贯穿万有，并在万有内。但恩宠赐给了我们每一个人，是按基督恩赐的尺度。}\par

\Person{保禄}所要求我们的爱不是普通的爱，而是那种将我们紧密结合、使我们彼此不可分离、并产生如同肢体与肢体之间那样完美合一的爱。正是这种爱结出伟大而荣耀的果实。因此他说，有 \BibleQuote{一个身体}；既因同情，也因不反对他人的益处，并分享他们的喜乐——他一下子用这个比喻表达了这一切。然后他优美地加上，\BibleQuote{一个圣神}，表明从这一个身体将产生一个圣神；或者说，有可能真有同一个身体，却没有同一个圣神，例如，若其中任何成员是异端者的朋友；或者，他通过这句话使他们因一致而羞愧，好像说：\Quote{你们领受了同一个圣神，被引到同一个泉源饮水，就不该在心灵上分裂}；或者此处\Quote{圣神}指他们的热忱。接着他加上：\BibleQuote{正如你们蒙召，同有一个蒙召的希望}，意思是，天主以同样的条件召叫了你们所有人。祂没有给一个人多于另一个人。祂无偿地赐给了所有人不朽、永生、不朽的荣耀、兄弟关系、嗣业。祂是众人的共同元首；\BibleQuote{祂使众人复活}，\BibleQuote{使他们与祂同坐}。\BibleRef{弗 2:6} 那么，你们在属灵世界里享有如此平等的特权，为什么还要骄傲自大？难道因为一个人富有，另一个人强壮？这是多么可笑！请告诉我，如果某天皇帝带来十个人，给他们全都穿上紫袍，让他们坐在御座上，赐予他们同样的尊荣，你们想，其中会有谁胆敢指责另一个人比自己更富有或更显赫？绝不会。我还没有说全；因为天上的差距并不像我们在地上相差那么大。有 \BibleQuote{一主，一信，一洗}。请看 \BibleQuote{你们蒙召的希望。一个天主和万有之父，祂超越万有，贯穿万有，并在万有内}。难道可能：一个呼号一位更大的神的名字，另一个呼号一位较小的神？一个因信德得救，另一个因行为？一个在洗礼中获得了赦免，另一个却没有？\BibleQuote{只有一个天主和万有之父，祂超越万有，贯穿万有，并在万有内}。\BibleQuote{祂超越万有}，意即主，在万有之上；\BibleQuote{贯穿万有}，意即提供、安排万有；\BibleQuote{并在你们万有内}，意即居住在你们众人内。现在，他们承认这是圣子的属性；因此，如果这是低劣性的论据，就绝不会用在圣父身上。\par

\BibleQuote{但恩宠赐给了我们每一个人}\BibleRef{弗 4:7}\par

那么，他说，这些不同的神恩从何而来？因为这个问题不断地使厄弗所人自己、格林多人以及许多其他人，要么陷入虚妄的傲慢，要么陷入沮丧或嫉妒。因此，他处处带着这个身体的比喻。所以现在他也提出了它，因为他要提到不同的恩赐。他确实在 \WorkTitle{格林多前书} 中更充分地进入了这个主题，因为那个病症在那里最为盛行；但在这里他只是暗示它。注意他所说的话：他没有说 \Quote{按各人的信德}，以免使那些没有很大成就的人陷入沮丧。但他说了什么呢？\BibleQuote{是按基督恩赐的尺度}\BibleRef{弗 4:7}。他说，所有主要和首要之点——圣洗、因信德得救、以天主为父、我们众人同领一个圣神——这些都是共同的。如果某人在任何神恩上有超越之处，不要因此而忧愁；因为他的劳苦也更重。那领了五个塔冷通的，被要求交还五个；而领了两个的，只带来两个，却仍得到与另一个同等的赏报。因此，宗徒在这里也根据同样的理由鼓励听者，表明恩赐的赐予不是为了抬高一人胜过另一人，而是为了教会的工作，正如他后面所说：\par

\BibleQuote{为成全圣徒，以从事服务工作，以建树基督的身体}\BibleRef{弗 4:12}\par

因此，连他本人也说：\BibleQuote{如果我不传福音，我就有祸了}。\BibleRef{格前 9:16} 例如：他领受了宗徒职分的恩宠，但正因这个缘故，\Quote{他有祸了}，因为他领受了它；而你们却没有这个危险。\par

\BibleQuote{按尺度}\BibleRef{弗 4:7}\par

\Quote{按尺度}是什么意思？意思是\Quote{不是按我们的功德}，因为那样的话，没有人会领受他所领受的；但我们都从恩赐中领受了。那么为什么一个人多而另一个人少？他会说，没有什么原因，只是事情本身无关紧要；因为每个人都为\Quote{建筑}作出贡献。他也借此表明，一个人领受多、另一个人领受少，并非出于他自身固有的功德，而是为了别人的缘故，正如天主自己衡量了它；因为他在别处也说：\BibleQuote{但如今天主却按祂的意愿，将肢体一个一个地安置在身体上了}。\BibleRef{格前 12:18} 他没有说明理由，以免使听者沮丧或灰心。\par

第8节。\Quote{因此经上说：\InnerQuote{祂升上高天时，掳掠了俘虏，将恩赐赐给人}}\BibleRef{弗 4:8}\par

好像他曾说：\Quote{你们为何骄傲？}一切都是由他而来的。先知在 \WorkTitle{圣咏} 中说：\BibleQuote{祢从人间领受了恩赐}\BibleRef{咏 67:18}，而宗徒却说：\BibleQuote{祂将恩赐赐给人}\BibleRef{弗 4:8}。二者是一回事。\par

第9、10节。\Quote{祂升上去了，若不是也曾降到地下最底层，又是什么呢？那降下的，正是那升到诸天之上，为要充满万有的那一位。}\par

当你听到这些话时，不要想这是从一处到另一处的单纯转移；因为保禄在《斐理伯书》\BibleRef{斐 2:5}所确立的论据，正是他在这里所坚持的。正如在那里，当他劝勉谦卑时，以基督为榜样，他在这里也是如此，说：\Quote{祂降到地下最底层。}若非如此，他所用的这句话\Quote{祂服从至死}\BibleRef{斐 2:8}就是多余的；然而从祂的上升，他暗示了祂的下降，而\Quote{地下最底层}按人的观念是指\Quote{死亡}，正如雅各伯也说：\Quote{你必将我的白发忧伤地带下阴间。}\BibleRef{创 32:48}又如圣咏所说：\Quote{免得我成为那些下到深渊的人}\BibleRef{咏 142:7}，即如死人一般。他为何在此详细谈论这地方？他说的俘掳是什么？是指魔鬼的俘掳；因为祂俘获了暴君，即魔鬼，以及死亡、诅咒和罪恶。看啊，祂的战利品和凯旋。\par

\Quote{祂升上去了，若不是也曾降到地下最底层，又是什么呢？}\par

这是针对撒摩沙特人保禄及其学派的抨击。\par

\Quote{那降下的，正是那升到诸天之上，为要充满万有的那一位。}\par

他说，祂降到地下最底层，此外再没有更低之处；祂升到万有之上，到那之外再没有更高之处。这是为了显示祂的神性能力和至高统治。因为万有早已被充满。\par

第11、12节。\Quote{祂所赐的，有做宗徒的，有做先知的，有做传福音的，有做牧者和教师；为成全圣徒，以从事服事的工作，以建立基督的身体。}\par

他在别处所说的\Quote{因此，天主极其举扬祂}\BibleRef{斐 2:9}，也正是他在这里所说的。\Quote{那降下的，正是那升上来的。}祂降到地下最底层并未伤害祂，也未妨碍祂成为远超诸天之上。所以一个人越谦卑，就越被高举。正如水的情况，人越向下压，就越向上涌；人越向后退以投掷标枪，就越命中目标。谦卑也是如此。然而，当我们论及天主的升高时，必须先设想下降；但论及人时，则完全不必。接着他进一步显示祂的眷顾和智慧，因为成就了这一切的那一位，拥有如此大能，且为我们的缘故拒绝下到那些底层，祂绝不会无目的地分施这些神恩。他在别处告诉我们这是圣神的工作，说：\Quote{圣神立了你们为监督，牧养天主的教会。}而这里他说是圣子；别处又说是天主。\Quote{祂赐给教会一些宗徒，一些先知。}但在致格林多人书中，他说：\Quote{我栽种，阿颇罗浇灌，然而使之生长的是天主。}又说：\Quote{栽种的和浇灌的原是一体，但各人要按自己的劳苦领取自己的报酬。}\BibleRef{格前 3:6}这里也是如此；因为即使你带来的很少，你已经领受了这么多。首先，他说\Quote{宗徒}，因为他们拥有一切神恩；其次\Quote{先知}，因为有些人虽不是宗徒，却是先知，如阿加波；第三\Quote{传福音的}，他们并不到处走动，而只是宣讲福音，如普黎史拉和阿桂拉；\Quote{牧者和教师}，即受委托照管整个团体的人。那么，牧者和教师是低一级的吗？是的，当然；那些固定在一处工作的人，如弟茂德和弟铎，比那些周游天下宣讲福音的人要低。然而，从这段经文无法得出从属和优先的次序，而是从另一封书信。他说\Quote{祂赐下}，你不可说反驳的话。或者他所谓的\Quote{传福音的}是指撰写福音的人。\par

\Quote{为成全圣徒，以从事服事的工作，以建立基督的身体。}\par

你们看到了职位的尊贵吗？每个人都建立，每个人都成全，每个人都服事。\par

第13节。他接着说：\Quote{直到我们众人达到信仰的统一和对天主之子的认识，达到成熟的人，达到基督圆满年龄的度量。}\BibleRef{弗 4:13}\par

他这里所说的\Quote{度量}是指完美的\Quote{认识}；因为正如一个人站得稳固，而孩童的心智被摇摆不定，信徒也是如此。\par

\Quote{达到信仰的统一。}\par

那就是，直到我们被证明拥有同一信德：因为这就是信德的合一，当我们都合而为一，当我们都同样承认那共同的纽带。在此之前，你必须努力达到这一目标。如果你为此领受了恩赐，好使你能建立别人，要当心不要因嫉妒别人而倾覆自己。天主尊荣了你，又指定你，使你能建立另一个人。是的，宗徒也正是为此忙碌；先知为此说预言、劝人，福音传道者为此传福音，牧者和教师也为此；所有人都从事一项共同的工作。因为不要向我讲论神恩的差异，而是说所有人都有一个工作。当我们都同样相信时，那时才会有合一；因为很显然，这就是他所称的\Quote{一个成全的人}。然而他在别处又称我们为\BibleQuote{婴孩}\BibleRef{格前 13:11}，即使我们已成年；但在那里，他是在与另一种比较对照，因为在那里，是与我们未来的知识相比较，他称我们为婴孩。因为他说了\BibleQuote{我们现在所知道的，只是局部的}\BibleRef{格前 13:9}，又加上了\Quote{模糊地}等话；而在这里他是针对另一件事，即变化无常，正如他在别处也说：\BibleQuote{但固态食物是为了成全的人}\BibleRef{希 5:14}。你看见在那里他也在什么意义上称他们为成全的吗？也要观察他在这里用什么意义称人为\Quote{成全}，通过随后添加的话，他说：\BibleQuote{使我们不再作小孩子}。他的意思是，我们要保守那点我们或许已经领受的小份额，以全部勤奋、坚定和稳固。\par

14节。\BibleQuote{使我们不再}。——\Quote{不再}这个词表明他们从前曾经处在这样的情形中，而且他更将自己视为需要改正的对象，从而改正自己。为此，他会说，有这么多的工人，好使建筑物不被摇动，不被\Quote{飘来飘去}，使石头得以牢固固定。因为这就是小孩子的特点：摇来摇去，被飘动，被摇动。\BibleQuote{使我们不再作小孩子，被飘来飘去，被各种教义之风吹动，借着人的诡计，以狡诈随从错谬的诡计。}\BibleRef{弗 4:14}他说：\Quote{被风吹动。}他用了这个比喻，以指出疑惑的灵魂处在何等大的危险中。他说：\Quote{借着人的诡计，以狡诈随从错谬的诡计。}\Quote{诡计}这个词指的是赌徒的伎俩。\Quote{狡猾的人}就是这样的，每当我们抓住更单纯的人时。因为他们也改变并转移一切。他在这里也暗指人类的生活。\par

15至16节。他说：\BibleQuote{惟用爱心说诚实话，凡事长进，连于他，即元首基督，从他，}（即从基督，）\BibleQuote{全身配合适当并连结在一起，借着每一关节的供应，按照各部分按尺度的运作，使身体增长，以致在爱中建立自己。}\par

他表达得十分晦涩，因为他想一下子说出一切。然而他的意思是这样的。如同从大脑降下的精神或生命力，将感觉能力通过神经传达给各肢体，不是简单地给所有肢体，而是按每一肢体的比例，给接受多的就多，给接受少的就少（因为精神是根源）；同样，基督也是如此。因为人的灵魂作为肢体依附于他，他的眷顾及按每一肢体尺度的适当比例所供应的神恩，促成它们的增长。但\Quote{借着供应的接触}是什么意思？就是说，借着感觉能力。因为从头部供应给肢体的这精神，\Quote{接触}每一肢体，从而推动它。好像有人会说：\Quote{身体按各肢体的比例接受供应，因而增长}；或换句话说，\Quote{肢体按各自尺度的比例接受供应，因而增长}；或者又如这样：精神从上丰沛地涌流，接触所有的肢体，并按每一肢体所能\Quote{接受它}的程度供应他们，\Quote{因而增长}。但是他为什么加上\Quote{在爱中}？因为没有别的方式，那圣神能降临。因为，假如一只手偶然从身体上砍下，来自大脑的精神会寻找那肢体，若找不到，不会从身体跳出来，飞向手，而是若找不到它在原位，就不接触它；同样，如果我们不在爱中结合，也会这样。凡此种种表达，他都用作促进谦卑。因为，他似乎在说，如果这个人或那个人比别人领受得多，又怎样呢？他领受了同一圣神，从同一元首发来，在所有人内同等地有效运作，同等地通传自己。\par

\BibleQuote{配合适当并连结在一起。}\par

这就是说，对它倍加关照；因为身体不可随意拼凑，而需极高的技巧与精细，因为一旦错位，便不复存在。因此，每肢体不仅须与身体结合，还须占据其适当位置；因为你若越出此界限，便不与身体结合，也领受不到圣神。难道你们没有看见，在意外发生的骨骼脱位中，当一根骨头偏离其正确位置而占据了另一根的位置时，它如何损伤整个身体，且常致人死亡？因此，有时会发现它不再值得保留。因为在许多情况下，许多人会将其切除，在其处留下空位；因为凡事过度皆为恶。元素亦然，若失去适当比例而过度，便损害整个系统。这就是\Quote{配合恰当，联络紧凑}的含义。那么，请思量各人留在自己适当位置，而不侵犯毫不属于己的他人位置，是何等重要。你把肢体组合在一起，祂从上供应。因为身体中既有如前所述接受器官，同样在圣神内也是如此，整个根源或源头来自上。例如，心脏接受气息，肝脏接受血液，脾脏接受胆汁，其他器官各自接受不同之物，但这一切都源自脑。天主也是如此行，极大地尊荣人，不愿远离人，祂确实使自己成为人依赖的源头，并使他们与自己同工；祂指派一些人任此职，另一些人任他职。例如，宗徒是整个身体中最关键的导管，从祂领受一切；祂使永生通过他们流到众人，如同通过血管和动脉，我是指通过他们的宣讲。先知预言将来之事，而唯有祂安排同样之事；你把肢体组合在一起，但祂以生命供应他们，\Quote{为成全圣徒，为服事的工作}。爱德建立人，使人彼此依附，并牢牢联结在一起。\par

道德教训。因此，我们若渴望领受从那元首而来的圣神的益处，就让我们彼此依附。因为从教会身体分离有两种方式：其一，当我们爱德冷淡时；其二，当我们胆敢做出与那身体不相称的事时；因为无论哪种方式，我们都会将自己从\Quote{基督的圆满}中割除。但若我们也被指派去建立他人，那么那些首先制造分裂的人，又岂能不受惩罚？没有什么比贪爱权力更能分裂教会。没有什么比分裂教会更激起天主的愤怒。是的，即使我们成就了千万件光荣之事，但若我们撕裂了教会的圆满，所受的惩罚将丝毫不亚于那些撕裂祂身体的人。因为那件事（基督受难）确实为世界的益处而发生，即便他们并无此意；而此事却毫无益处，且伤害极大。我这些话说给统治者，也说给被统治者。有一位圣人曾说了一句看似大胆的话，然而他还是说了出来。这话是什么？他说，连殉道的血也不能洗去这罪。请告诉我，你为何殉道？难道不是为了基督的光荣？那么，你为基督舍弃生命，又怎能摧残那基督为其舍弃生命的教会？请听保禄所说：\Quote{我不配称为宗徒\BibleRef{格前 15:9}，因为我曾迫害天主的教会并加以摧残。}\BibleRef{迦 1:13}这伤害并不小于来自敌人的伤害，不，更是大得多。因为那伤害反使她更显荣耀，而此事——当她被自己的儿女攻击时——却在敌人面前羞辱了她。因为在敌人看来，那些在她内出生、在她怀中养育、并完全学会她秘密的人，竟突然转变，做起敌人的工作，这是极大的虚伪标志。\par

我说这些话，是针对那些不加分辨地投身于分裂教会的人。因为一方面，若那些人持有与我们相反的教义，那么就更不应该与他们往来；另一方面，若他们持有相同的意见，则更不应与他们往来。为何如此？因为此时的病症源自对权力的贪欲。你们不知道\Person{科辣黑}、\Person{达堂}和\Person{阿彼兰}的下场吗？\BibleRef{户 16:1}我只说了他们吗？不也是那些与他们同伙的人吗？你们要说什么？难道要说：\Quote{他们的信德相同，他们和我们一样正统}？若是如此，他们为何不与我们在一起？\Quote{只有一个主，一个信德，一个圣洗。}若他们的主张正确，则我们的错误；若我们的正确，则他们的错误。他说：\Quote{小孩子，被各种风吹动，飘来飘去。}请告诉我，你们认为说他们是正统就足够了吗？那么圣职的祝圣就过去而废止了吗？若不严格遵守这一点，其他事又有什么益处呢？因为正如我们必须为信德争辩，也同样必须为这事争辩。因为若按古人的说法，任何人可以\Quote{满手}成为司铎，就让所有的人都来供职吧。这祭台白白被建立，教会的圆满白白存在，司铎的数目也白白有了。让我们把他们除掉并毁灭吧。你们会说：\Quote{绝对不可！}你们正在做这些事，却还说\Quote{绝对不可}？当这些事正在发生时，你们怎能说\Quote{绝对不可}？我说话并作证，不是顾自己的利益，而是为你们的救恩。若有人漠不关心，他自己必须留意；若这些事无人关心，我却关心它们。他说：\Quote{我栽种了，}\Person{阿颇罗}浇灌了，但天主使之生长。\BibleRef{格前 3:6}我们如何承受外邦人的嘲笑？若他们因我们的异端而责备我们，对这些事他们还会说什么呢？他们说：\Quote{若他们拥有同样的教义，同样的奥迹，为何一个教会的主管侵犯另一个？你们看，}他们说，\Quote{基督徒间的一切都充满了虚荣！他们中间有野心和伪善。除去他们，}他们说，\Quote{的人数，他们就什么也不是。割掉那疾病，腐化的群众。}你们要我告诉你们他们对我们的城市说了什么，他们如何指责我们的随和迁就？他们说，任何选择的人都可以找到追随者，而且永远不会缺少他们。哦，那是何等的嘲笑，这些事是何等的耻辱！然而嘲笑是一回事，耻辱是另一回事。若我们中有人被证明犯有最可耻的行为，并将面临刑罚，那么各处就有极大的惊慌和恐惧，人们说，惟恐他离开而投靠另一边。是的，让这样的人离开一万次，让他加入他们。我不只说那些犯罪的人，若有人无罪，且有意离开，也让他离开。我确实为此忧伤、痛哭、哀悼，心如刀割，仿佛自己被夺去一个肢体；然而我并不如此忧伤，以致因这种恐惧而被迫做任何错事。亲爱的，我们\Quote{不是主宰你们的信德}\BibleRef{格后 1:24}，也不是作为主上吩咐你们这些事。我们被任命为圣言的教师，不是为权力，也不是为绝对权威。我们居于顾问的地位来劝告你们。顾问说出自己的意见，不强逼听众，而让他在所说的事上完全自主选择；只有在没有表达出所出现的事时，他才该受责难。为此我们也说这些事，我们主张这些事，为使你们在那一天不能说：\Quote{没有人告诉我们，没有人命令我们，我们无知，我们以为这根本不是罪。}因此我主张并抗议，在教会中制造分裂的恶，并不亚于陷入异端。请告诉我，假设某国王的臣民，虽然没有投靠另一位国王，也没有将自己交给任何其他人，却拿起并抓住国王的王袍，从扣环上完全撕下，撕成许多碎片；他受的刑罚会比那些投靠他人的人更轻吗？此外，若他还掐住国王的喉咙杀死他，并将其身体肢解，他应受什么刑罚才能与他的罪相称？若对同是仆人的国王这样做，他所犯的罪已达任何刑罚都无法触及的程度；那么，杀害基督并将其肢解的人，该当什么地狱呢？是那受威胁的地狱吗？不，我想不是，而是另一个更可怕的地狱。\par

你们在场的妇女——因为这通常是妇女的缺点——请说话，将我所做的\OriginalTerm{比喻}{similitude}讲给那些不在场的人听；惊吓她们。如果有人想以此使我悲伤而报复我，就让他们知道这些是徒劳的。因为如果你想报复我，我给你一个方法，你可以不伤害自己而报复；或者更确切地说，没有伤害就不可能报复，但至少伤害较小。击打我，妇人，向我吐唾沫，当你在路上遇见我时，向我挥拳。你听到这个战栗吗？当我叫你击打我时，你战栗，而你撕裂你的主和主人却不战栗？你撕扯你的主和主人的肢体，却不战栗？教会是我们圣父的家。\BibleQuote{有一个身体，一个圣神。}但是你想要报复我吗？那么请停在我身上。为什么你要代替我报复基督？不，为什么你要用脚踢刺？事实上，报复从来不好也不对，但一人犯错而去攻击另一人则更糟。是我亏待了你吗？那么为什么使没有亏待你的人受苦？这真是极度的疯狂。我即将说的话并非讽刺，也非无的放矢，而是按我真实所想和所感。我希望那些和你们一起对我发怒、并因这发怒而伤害自己、离开去别处的人，每一个人都将拳头打在我脸上，剥我衣服、鞭打我，无论他的指控是公正或不公正，把怒气发泄在我身上，超过他们目前敢于做的那些事。如果这些事做了，也是无所谓的；一个微不足道、毫无价值的人受到这样的对待，没有什么。此外，我，受委屈和受伤害的人，可以呼求天主，祂或许会赦免你的罪。并非因为我如此大的信心；而是因为当受委屈的人为做错的人祈求时，他获得很大的信心。\Quote{若一人得罪另一人，}经上说，\Quote{则要为他祈祷}\BibleRef{撒上 2:25}；如果我无能，我可以找其他圣人，恳求他们，他们可以做到。但现在当天主被我们激怒时，我们还能恳求谁呢？\par

请注意这种一致性：因为那些属于这教会的人，有些人从不前来领圣体，或者一年只一次，且漫无目的，随随便便；另一些人确实更规律，但也过于粗心和无目的，且一边谈话，一边开无聊的玩笑；而那些看似认真的人，正是制造这祸害的人。然而，如果你们是为了这些事而认真，那么你们也最好加入那些漠不关心之人的行列；或者更好的是，既让他们不再漠不关心，也让你们不再像现在这样。我不是说你们在场的人，而是说那些正在离弃我们的人。这行为是奸淫。如果你们听不得别人这样说他们，同样也听不得我们这样说你们。要么是这一方，要么是另一方，必然犯了法。如果你们对我有这些怀疑，我准备辞去职务，并把它让给你们所选的任何人。只愿教会是合一的。但若我是合法被祝圣的，就请你们恳求那些非法登上主教座的人辞职。\par

我这样说，并非想命令你们，只是为了确保和保护你们。既然你们每人都已成年，将来必须为自己所做的事交账，我恳求你们不要把一切责任推给我们，而认为自己不负责任，免得你们徒然自欺，最终哀叹。我们确实要为你们的灵魂交账，但那是当我们失职的时候，当我们未能劝诫、未能警告、未能抗议的时候。但在说了这些话之后，让我也说\BibleQuote{我与众人的血是无干的}\BibleRef{宗 20:26}，并且\BibleQuote{天主必拯救我的灵魂}\BibleRef{则 3:19}。你们随便说，说出一个正当的理由为什么离开，我会回答你们。但你们不会说出来的。因此我恳求你们，从今以后努力高贵的抵抗并带回那些已退出的人，好让我们同心合意地向天主献上感恩；因为光荣永永远远属于祂。阿们。\par

\SourceRecord{Translated by Gross Alexander.；Nicene and Post-Nicene Fathers, First Series；Vol. 13.；Edited by Philip Schaff.；Buffalo, NY: Christian Literature Publishing Co.,；1889.；<http://www.newadvent.org/fathers/230111.htm>.}

% structure: p0001 source=h1 target=section reason=page-title

\section{厄弗所书讲道集 第十二篇}

\BibleRef{弗 4:17}\par

\Quote{因此我说，并在主内作证，你们不要再像外邦人也那样行，在\OriginalTerm{心的虚妄}{vanity of their mind}中，他们的理智昏暗了。}\par

教师的责任不仅是通过劝诫和教导，还要通过警醒他们并将他们交托给天主，来建造并恢复门徒们的灵魂。因为当由人作为同仆所说的话不足以点燃灵魂时，就必须将案件移交给天主。保禄也正是这样做的；在谈论了谦卑、合一以及我们不应彼此争斗之后，请听他所说的话。他说：\Quote{因此我说，并在主内作证，你们不要再像外邦人也那样行。}他没有说\InnerQuote{你们从此不要再像现在这样行}，因为那样说会打击过重。但他清楚地暗示了同样的事，只是用他人作为例子。同样，他在写给得撒洛尼人的书信中也做了同样的事，他说：\BibleQuote{不要像不认识天主的外邦人那样，在情欲的激情中行事} \BibleRef{得前 4:5}。他的意思是说，你们在教义上与他们不同，但那完全是天主的工作；我要求你们的是按照天主的生活和行为方式。这是你们自己的事。我请主为我所说的话作证，我没有退缩，而是告诉了你们应当如何行事。\par

他说：\Quote{在他的心的虚妄中。}\par

什么是心的虚妄？就是忙于虚妄之事。而现世生活中的一切不都是虚妄的吗？训道者说：\BibleQuote{虚妄的虚妄，一切都是虚妄} \BibleRef{训 1:2}。但有人会说：如果它们是虚妄和空虚，为什么被造？如果是天主的工程，怎么是虚妄？关于这些事有很大争议。但请听，亲爱的：他所说的虚妄不是天主的工程；绝不会！天不是虚妄的，地不是虚妄的——绝不会！——太阳、月亮、星辰，以及我们自己的身体也不是。不，所有这些都\BibleQuote{非常好} \BibleRef{创 1:31}。那么什么是虚妄？让我们听训道者自己说什么：\BibleQuote{我栽种了葡萄园，我得了歌唱的男女，我修造了水池，我拥有大群的牛羊，我也积攒了金银，我看出这些都是虚妄} \BibleRef{训 2:4}。又说：\BibleQuote{虚妄的虚妄，一切都是虚妄} \BibleRef{训 12:8}。也听先知说：\BibleQuote{他积聚财富，却不知谁来收取} \BibleRef{咏 38:6}。这就是\OriginalTerm{虚妄的虚妄}{vanity of vanities}：你华丽的建筑，你浩大丰盈的财富，在广场上拥来拥去的成群奴隶，你的炫耀和虚荣，你的高傲思想和你的铺张。因为这些都是虚妄的；它们不是出于天主的手，而是我们自己的创造。但为什么它们是虚妄的呢？因为它们没有有用的目的。财富用于奢侈时是虚妄的；但若\BibleQuote{分散给予穷人} \BibleRef{咏 111:9}，就不复是虚妄。但当你把它们花在奢侈上时，看它们的结果是什么——身体肥胖、胀气、喘息、腹饱、头重、肌肉松软、发烧、精力衰竭；因为一个人如果往\OriginalTerm{漏水的容器}{leaking vessel}里注水，就是徒劳，同样，过着奢侈放纵生活的人也是在往漏水的容器里注水。但还有，被称为\Quote{虚妄}的，是指那些本应装些东西却装不住的，人们称为\OriginalTerm{空虚的希望}{empty hopes}的东西。一般来说，凡是无目的、无用的，就被称为\Quote{虚妄}。那么让我们看看所有人间事物是否不都是这类。\BibleQuote{我们吃喝吧，因为明天我们就要死了} \BibleRef{格前 15:32}。那么，告诉我，结局是什么？腐坏。我们穿衣服吧。结果呢？没有。希腊人的生活就是这样。他们哲学思考，但却是徒劳。他们展示了一种艰苦的生活，但只是艰苦，不指向任何有益的目的，而是指向虚荣和众人的称赞。但众人的称赞是什么？是虚无，因为如果那些给予称赞的人自己都会消亡，那么称赞就更会消亡。一个给予别人荣誉的人，应当先给予自己荣誉；因为如果他自己不获得荣誉，又怎能给别人荣誉呢？然而我们现在甚至寻求从卑贱可鄙的人那里得到荣誉，这些人本身可耻，受人指责。那么这是什么荣誉？你们看，万事不都是\BibleQuote{虚妄的虚妄} \BibleRef{训 1:2}吗？因此，他说：\Quote{在他们的心的虚妄中。}\par

再者，他们的宗教难道不就是这样，是木头和石头吗？祂使太阳照耀，作为灯来照亮我们。谁会崇拜自己的灯呢？太阳供给我们光，但在太阳不能照到的地方，灯却可以。那么，为什么不崇拜你的灯呢？\Quote{不，}有人会说，\Quote{我崇拜火。}哦，多么可笑！如此荒谬，然而再看另一个荒谬。为什么要熄灭你所崇拜的对象？为什么要毁灭，为什么要消灭你的神明？为什么你不让你的房屋充满祂？因为如果火是神，就让祂吞食你的身体。不要把你的神放在锅或釜的底下。把祂带进你的内室，带进你的丝绸帷幔中。然而你不仅不带祂进去，而且如果祂偶然闯入，你会把祂从各处赶出去，召集所有人，好像有野兽闯入一样，你哭泣哀号，称你的神降临为大灾难。我有一位天主，我尽一切努力将祂怀抱在胸中，我认为真正的幸福不是祂造访我的住所，而是我能将祂吸引到我的心里。你也把火吸引到你的心里吧。这是愚妄和虚浮。火是有用的，是用来使用的，不是用来朝拜的；它是用来服侍和服务的，是我的奴隶，不是我的主人。它为我而造，我并非为它而造。如果你崇拜火，为什么你自己躺在卧榻上，却命令你的厨师站在你的神面前？你自己去学烹饪吧，如果你愿意，去当面包师，或铜匠，因为没有什么比这些技艺更高贵了，因为正是这些技艺是你的神所临在的。为什么你认为那技艺是耻辱，而你的神却是万有呢？为什么你把它交给你的奴隶，而自己不渴望它呢？火是好的，因为它是善的造物主的作品，但它不是天主。它是天主的作品，它没有被称作天主。你难道没有看见它的本性是多么难以驾驭吗？——当它抓住一座建筑时，它停止在任何地方吗？但如果它抓住连续的东西，它就毁灭一切；除非工匠或其他人的手扑灭它的狂暴，它不认识朋友也不认识敌人，而是同样对待一切。这就是你的神吗？你不羞愧吗？他实在说得好：\BibleQuote{在思念上虚妄。}\BibleRef{厄 4:17}\par

但他们说，太阳是神。告诉我，如何以及为何。是因为它散发出大量的光吗？但你难道没有看见它被云遮蔽，受制于自然的必然，被月食遮掩，被月亮遮蔽吗？云比太阳弱，但常常胜过它。而这正是天主智慧的作品。天主必须是全备的，但太阳需要许多东西，这不像神。因为它需要空气来照耀，而且是稀薄的空气，因为空气若太稠密，就不让光线通过。它还需要水和其他抑制力量，以防止它消耗一切。因为若不是泉水、湖泊、河流、海洋通过散发水汽产生一些湿气，没有什么能阻止普遍的火灾。那么，你们说，你们看它是一位神吗？何等愚妄，何等疯狂！一位神，你们说，因为它有能力造成伤害。不，相反，正因如此它不是神，因为它在造成伤害时不需要任何东西；而在行善时，它还需要许多其他东西。而行善是天主的本性；行恶却与天主的本性相悖。既然情况相反，它怎能是神呢？你们难道没有看见有毒的药物伤害人时不需要任何东西，而当它们要行善时却需要许多东西吗？它之所以如此，是为了你们的缘故，既是善的，也是无能的；善，好让你们承认它的主；无能，好让你们不说它是主。\Quote{但是，}他们说，\Quote{它滋养植物和种子。}那么，照此推理，难道粪土不也是神吗？因为连它也滋养植物。那么，照此推理，镰刀和农夫的手不也是神吗？请向我证明，太阳独自滋养作物，不需要土地、水或耕种的帮助；只要播下种子，它发出光芒，就会长出麦穗。但现在，如果这工作不是它独自完成的，而是也要靠雨水，那么水不也是神吗？但我暂不提这个。为什么土地不也是神？为什么粪土和锄头不也是神？那么，请问，我们要崇拜这一切吗？唉，多么无聊！事实上，没有太阳可能还能长出麦穗，但没有土地和水却不行；植物和其他一切也是如此。如果没有土地，这些东西根本不可能出现。如果有人像孩子和妇人那样，把一些土放在盆里，装满粪肥，放在屋檐下，植物也会从里面长出来，尽管可能很弱。所以，土地和粪肥的贡献更大，因此我们应该崇拜它们而不是太阳。它需要天空，需要空气，需要这些水来防止它造成伤害，如同缰绳来遏制它力量的凶猛，约束它不要把光线像烈马一样释放到世界上。现在告诉我，夜晚它在哪里？你们的神去了哪里？因为被限定和限制不是神的特征。这实际上只是物体的属性。但他们说，它里面有一种力量，它有运动。请问，这力量本身是神吗？那么它为什么自身不足，为什么不能约束火？因为我又回到同一个论点。但那力量是什么？是产生光，还是通过太阳发光，而本身却不具有这些性质？如果是这样，那么太阳比它优越。我们要绕多远才能解开这个迷宫？\par

再者，水是什么？他们不也说那是神吗？这确实是一场荒谬的争论。他们说，我们用于这么多用途的东西，难道不是神吗？同样，对于土地也是如此。他们实在\BibleQuote{在思念上行走于虚妄，理智被蒙蔽。}\BibleRef{厄 4:17-18}\par

但这些话如今他是就生活和行为而说的。希腊人都是淫乱者和奸淫者。当然。他们为自己塑造出这样的众神，自然就会行这一切事；他们若能逃过人的眼目，就没有人约束他们。因为复活的论证若在他们看来不过是神话，又有什么用处呢？是的，地狱的酷刑又有什么用处呢？——它们也不过是神话罢了。且留意这撒殚的念头。当人们告诉他们那些行淫乱的众神时，他们否认这些是神话，反而相信它们。然而，若有人对他们谈论惩罚，他们会说：\Quote{这些}，他们说，\Quote{是诗人们，他们把一切都变成寓言，要使人的幸福状态处处被颠覆。}\par

但哲学家们，据说，发现了某种真正伟大的、远胜于这些的东西。怎么？他们引入了命运，告诉我们没有天主的眷顾，没有任何事物受到照顾，万物都由原子组成？或者，另一些人又说天主是一个身体？或者，是谁，告诉我？是那些要把人的灵魂变为狗的灵魂，并且要使人相信一个人曾经是狗、狮子、鱼的人吗？你们要胡说到什么时候才肯罢休，\Quote{心智昏暗}？因为他们行事说话如同真的在黑暗中，无论教义方面还是生活行为方面都是如此；因为在黑暗中的人看不见眼前的东西，常常看到一根绳子就以为是活蛇；或者，被灌木丛钩住，就以为是人或邪灵抓住了他，于是大为惊恐，大大慌乱。他们所怕的就是这类东西。经上说：\Quote{他们在那里大为恐惧，其实并无恐惧} \BibleRef{咏 52:5}；但他们本该害怕的东西，他们却不害怕。正如在乳母怀中的孩子粗心地把手伸进火里，也大胆地伸进蜡烛里，却害怕一个穿着麻衣的人；同样，这些希腊人就像永远是小孩子（他们自己中有人也说过，希腊人永远是小孩子），他们害怕那些不是罪的东西，例如身体的污秽、葬礼的玷污、床铺、守日之类；而那些真正是罪的东西，非本性的情欲、奸淫、淫乱，他们却毫不在意。不，你可以看到一个人洗去死人玷污，却从不洗去死行；又看到他在追求财富上大发热心，却以为一切都会被一声鸡叫破坏。\Quote{他们的心智就这样昏暗了。} 他们的灵魂充满了各种恐惧。例如：一个人会说：\Quote{那人}，\Quote{是我出门时第一个遇见的人}；当然，必有千万灾祸随之而来。又一次，\Quote{那个该死的仆人递给我鞋子时，先递了左脚的那只}，——可怕的不幸与灾祸！\Quote{我自己出门时先迈左脚}，这也是不祥之兆。这些就是家中发生的不幸。然后，我出门时，右眼皮跳动，这是流泪的确实征兆。再者，妇女们，当芦苇撞击支架发出响声，或者她们自己被梭子划伤时，也把这也当作征兆。又当她们用力投梭，梭子撞击顶部芦苇，撞到支架发出响声，她们又把这当作征兆，以及无数其他荒谬可笑的事。因此，如果驴叫，或公鸡啼鸣，或人打喷嚏，或任何其他事情发生，他们就像被万重锁链捆绑的人，或者像我在说的，像被困在黑暗中的人，对一切都猜疑，比世上所有奴隶都更奴性十足。\par

但愿我们不如此。而应藐视这一切，如同生活在光明中、拥有天上户籍、与世界毫无共同之处的人，只将一件事视为可怕，那就是罪，和得罪天主。若没有这事，就藐视其余一切，以及那引入这些事的魔鬼。让我们为这些事感谢天主。让我们殷勤，不仅自己永不被这奴役所困，而且若有任何我们所爱的人被囚，就折断他的锁链，将他从这最苦最可鄙的囚禁中释放，使他自由无羁地奔向上天，提起他低垂的翅膀，并教导他为生命和道理而成为智者。让我们为一切事感谢天主。让我们恳求祂不以我们所领受的恩赐为不配的，同时我们自己也要努力尽本分，不仅用言语，也用行动来教导。如此我们才能获得祂无数的祝福，愿天主使我们大家都能被算为堪当，在我们的主耶稣基督内，偕同圣父和圣神，一切光荣、德能与尊威，自今、从今世，直到永远。阿们。\par

\SourceRecord{Translated by Gross Alexander.；Nicene and Post-Nicene Fathers, First Series；Vol. 13.；Edited by Philip Schaff.；Buffalo, NY: Christian Literature Publishing Co.,；1889.；<http://www.newadvent.org/fathers/230112.htm>.}

% structure: p0001 source=h1 target=section reason=page-title

\section{厄弗所书讲道集 第十三篇}

\BibleRef{弗 4:17-19}\par

\BibleQuote{因此，我这样说，并在主内作证：你们不要再像外邦人那样，憑虚妄的心思行事；他们的理智昏暗了，因着心中的无知和心硬，与天主的生命隔绝了。他们既已麻木，便放纵情慾，贪婪地行各种不洁之事。}\par

这些话不仅是对厄弗所人说的，如今也是对你说的；而且，这不是出于我，而是出于保禄；或更好说，既不是出于我，也不是出于保禄，而是出于圣神的恩宠。因此，我们该当感觉，彷彿那恩宠亲自在说话。现在请听它说什么。\BibleQuote{因此，我这样说，并在主内作证：你们不要再像外邦人那样，憑虚妄的心思行事；他们的理智昏暗了，因着心中的无知和心硬，与天主的生命隔绝了。}既然如此是无知、是心硬，为何要责备呢？一个人无知，本应不是因此而受恶待或受责备，而是应受教导他所不知道的事。但请看，他如何立刻断绝了他们的一切借口。\BibleQuote{他们既已麻木，}他说，\BibleQuote{便放纵情慾，贪婪地行各种不洁之事；但你们却不是这样学了基督。}在此他向我们显示，他们心硬的原因是他们的生活方式，而他们的生活方式是他们自己的懒惰和麻木的结果。\par

\BibleQuote{他们既已麻木，}他说，\BibleQuote{便放纵自己。}\par

每当你们听见\BibleQuote{天主就任凭他们陷于可弃绝的心思}\BibleRef{罗 1:28}时，当记住这句话：\Quote{他们放纵自己。}既然他们放纵自己，天主又如何任凭他们呢？如果天主任凭他们，他们又怎能放纵自己呢？你看这表面上的矛盾。因此，\Quote{任凭}一词的意思是他允许他们被放弃。你看见吗，不洁的生活也是类似教义的基础？主说：\BibleQuote{凡作恶的，便恨光，不来接近光。}\BibleRef{若 3:20}因为一个放荡的人，一个比在泥沼中打滚的猪更沉迷于滥交的人，一个贪财的人，一个对节制毫无渴望的人，怎能进入这样的生活呢？他说，他们以这事为他们的\Quote{工作}，因此有了他们的\Quote{心硬}\BibleRef{弗 4:19}，因此有了\Quote{理智的昏暗}。有一种情况是，即使光亮照耀，眼睛也会在黑暗中；眼睛变弱，要么是由于体液流入，要么是因泪液过多。同样，在这里也如此：当这世上的事务的强劲洪流吞没了理解力的感知能力时，它就陷入黑暗。就像如果我们被置于深水中，由于大量的水像屏障一样覆盖在我们上面，我们就无法看见太阳；同样，在理解力的眼中也会发生心的盲目，即麻木，每当没有恐惧激动灵魂时。经上说：\BibleQuote{在他的眼中没有对天主的敬畏}\BibleRef{咏 35:1}；又说：\BibleQuote{愚妄人心里说：没有天主。}\BibleRef{咏 13:1}现在，盲目除了来自麻木之外，没有其他原因；这堵塞了通道；因为每当体液凝结并聚集在一处时，肢体就变得死寂麻木；无论你烧它、割它，或以任何方式对待它，它都感觉不到。那些人也是如此，一旦他们放纵于情慾，尽管你像火或刀一样向他们应用这话，却什么也触动不了他们，什么也达不到他们；他们的肢体完全死了。除非你能除去麻木，以便触及健康的肢体，否则你所做的一切都是徒劳。\par

\BibleQuote{贪婪地，}他说。\par

在此，他完全剥夺了他们的借口；因为他们有能力，只要他们愿意，不\Quote{贪婪}、不\Quote{放纵}、不贪食，却仍能享受他们的欲望。他们能够适度地享有财富，甚至享乐和奢侈；但一旦他们毫无节制地放纵，就毁灭了一切。\par

\BibleQuote{行各种不洁之事，}他说。\par

你看，他如何通过提及\Quote{行不洁之事}来断绝他们的一切借口。他的意思是，他们不是偶然犯罪，而是刻意从事这些可怕的行为，并把这事当作学问。\Quote{各种不洁之事}：不洁包括所有的奸淫、淫乱、不自然的慾望、嫉妒、各种放荡和纵慾。\par

第20、21节。他接着说：\BibleQuote{但你们却不是这样学了基督，如果你们真听过他，并在他内受过教导，正如真理是在耶稣内。}\par

\Quote{如果你们真听过他}这句话不是出于怀疑，而是强烈肯定；正如他在别处也说：\BibleQuote{如果这是公义的，天主就以患难报复那加患难给你们的人。}\BibleRef{得后 1:6}这就是说，你们不是为了这些目的而\Quote{学了基督}。\par

第22节：\BibleQuote{就该脱去你们从前生活方式的旧人。}\BibleRef{弗 4:22}\par

这确实就是学基督——正确生活；因为生活邪恶的人不认识天主，也不被天主认识；听他在别处说：\BibleQuote{他们自称认识天主，但在行为上却否认他。}\BibleRef{铎 1:16}\par

\BibleQuote{正如真理是在耶稣内：就该脱去你们从前生活方式的旧人。}\par

这就是说，你们不是在这些条件下订立盟约的。在我们中间的不是虚妄，而是真理。正如教义是真实的，生活也是如此。罪是虚妄和谎言；而正直的生活是真理。因为节制确实是真理，因为它有伟大的结局；而放荡却归于虚无。\par

他说：\Quote{这败坏，} \Quote{随从欺骗的私欲。} 正如他的私欲变得败坏，他自己也是如此。那么，他的私欲如何变得败坏呢？死亡使一切消散；因为听先知如何说：\BibleQuote{就在那一天，他的思念都归于无有。} \BibleRef{咏 145:4} 不仅死亡如此，还有许多其他事物；例如，美貌，在疾病或衰老来临之际，便消退、衰亡并遭受败坏。身体的活力同样被这些所摧毁；即使在老年，奢侈本身也不再提供同样的快乐，从\Person{巴尔齐来}的例子可以清楚看出；你们无疑知道这段历史。或者，在另一种意义上，私欲败坏并摧毁那旧人；正如羊毛被产生它的同样手段所摧毁，旧人也是如此。因为贪图荣耀摧毁他，享乐也常常摧毁他，\Quote{私欲}将完全\Quote{欺骗}他。因为这并非真正的快乐，而是苦涩和欺骗，全是伪装和外观。事物的表面确实光鲜，但事物本身却充满痛苦、极端的悲惨、可憎和彻底的贫乏。摘下假面具，揭开真面目，你就会看到骗局，因为这是骗局：当存在的不显现，不存在的东西却被展示出来。就这样，欺诈得以施行。\par

宗徒为我们描绘了四个人。我将对此加以解释。在此处他提到了两个，如此说：\Quote{脱去旧人，在你们心神的灵里更新，并穿上新人。} 而在\WorkTitle{罗马书}中，又提到另外两个，如他所说：\Quote{但我看出在我的肢体中另有一条法律，与我心神的法律交战，并把我掳去，隶属于那在我肢体中罪的律法。} \BibleRef{罗 7:23} 后者与前者有关联：\Quote{新人}与\Quote{内在的人}，\Quote{旧人}与\Quote{外在的人}。然而，这四个中的三个是受败坏支配的。或者更准确地说，有三个：新人、旧人，以及这个人——在其实体和本性上的人。\par

第23节。他说：\Quote{并且你们要在你们心神的灵里更新。} \BibleRef{弗 4:23}\par

他说：\Quote{你们要更新。} \Quote{在你们心神的灵里。}\par

第24节。\Quote{并穿上新人。} \BibleRef{弗 4:24}\par

你们看见吗？主体是同一个，但衣服是两件：脱去的和穿上的。他继续说：\Quote{这新人，是按照天主在真理的正义和圣洁中所创造的。} 那么，为什么他将德行称为\Quote{人}？又为什么将恶习称为\Quote{人}？因为一个人不能没有行动而显明；所以这些，如同本性一样，显明一个人是善还是恶。正如脱衣穿衣是容易的，因此我们也可以看到德行和恶习也是如此。年轻人是强壮的；因此，我们也要为行善而变得强壮。年轻人没有皱纹，因此我们也不应有。年轻人不摇动，也不易患病，因此我们也不应如此。\par

请注意，他如何将德行的实现，即从无中产生，称为\Quote{创造}。但如何？那之前的创造不也是按照天主吗？不，绝不是，而是按照魔鬼。他是罪的唯一创造者。\par

这是怎么回事？因为人从此被创造，不是由水，也不是由地，而是\Quote{在真理的正义和圣洁中}。这如何解释？他的意思是，他立即创造了他，使他成为儿子：因为这是从圣洗圣事开始的。这就是现实，\Quote{在真理的正义和圣洁中}。从前在犹太人中也有一种正义，也有一种圣洁。然而，那正义并不在真理中，而是在预像中。因为身体洁净是洁净的预像，而非洁净的真理；是正义的预像，而非正义的真理。他说：\Quote{在正义和圣洁中}，这些是\Quote{属于真理的}。\par

并且这个表述是针对虚谎的；因为有许多人，在局外人看来似乎是正义的，却是虚谎的。现在，正义是指普遍的德行。因为请听基督如何说：\Quote{除非你们的正义超过经师和法利塞人的正义，你们决进不了天国。} \BibleRef{玛 5:20} 再者，一个人被称为正义，是指无人向他控诉；因为在法庭上，我们也说那人是正义的，如果他不正义地受对待，却没有不义地报复。因此，如果我们在可怕的审判台前，彼此能够显为正义，我们或许会得到一些仁慈。其实在天主面前，无论我们有什么可呈上，都不可能显为正义。因为祂在正义上处处得胜，正如先知也所说：\Quote{祢在审判中得胜。} 但如果我们不违背彼此之间的正义，那么我们就是正义的。如果我们能够表明我们受到了不义的对待，那么我们就是正义的。\par

他如何对已经穿上衣服的人说\Quote{穿上}呢？他现在是在谈论那源于生活和善行的衣服。之前，那衣服来自于圣洗圣事，而现在则来自日常生活和善工；不再是\Quote{随从欺骗的情欲}，而是\Quote{随从天主}。但\Quote{圣洁}这个词是什么意思呢？它是纯洁的，是应得的；因此我们也用这个词来表示对逝者的最后义务，犹如说：\Quote{我不再欠他们什么，我没有什么要负责的了。}因此，我们惯常说\Quote{我已尽了一切义务}等等，意思是\Quote{我不再欠什么了。}\par

道德。那么我们的责任就是永不脱下正义之衣，先知也称之为\BibleQuote{救恩之衣}\BibleRef{依 61:10}，这样我们才能像天主。因为天主确实穿上了正义。让我们穿上这件衣服。现在，\Quote{穿上}这个词除了表明我们绝不可脱下它之外，别无他意。请听先知的话，他说：\BibleQuote{他披上诅咒如同自己的外衣，诅咒就进入他的内脏。}\BibleRef{咏 108:18}又说：\BibleQuote{你以光为衣，如同披上外衣。}\BibleRef{咏 103:2}此外，我们也常谈论某人\Quote{穿上}了某人。因此，不是一天、两天或三天，而是他要我们永远以德行为衣，永不脱下这件衣服。因为一个人被剥去衣服时并不像被剥去德行时那样丑陋。在前者，他的同伴目睹他的裸体；在后者，他的主和天使看见他。如果你曾看见有人赤裸裸地走过广场，告诉我，你难道不感到痛苦吗？那么，当你被剥去了这件衣服，赤身露体地四处游荡时，我们该说什么呢？你难道没有看见那些我们称为流浪者的乞丐，他们四处漂泊，我们甚至可怜他们？然而他们仍没有借口。当他们在赌博中失去衣服时，我们不原谅他们；那么，如果我们失去这件衣服，天主岂能原谅我们？因为每当魔鬼看见一个人被剥去了德行，他立刻就会改变并毁坏他的面容，击伤他，并驱使他陷入极大的困境。\par

让我们脱去我们的财富，使我们不被剥夺义德。财富的外衣损坏了这件衣裳。它是荆棘的袍子。荆棘的本性如此；它们缠绕我们越紧，我们就越赤身露体。邪淫剥夺了这件衣裳；因为它是火，这火要烧毁这件衣裳。财富是蛀虫；蛀虫吞吃一切，连丝衣也不放过，财富也是如此。因此，让我们脱下这一切，好使我们能成为正义的，使我们能\Quote{穿上新人}。让我们不保留任何旧的、任何外在的、任何\Quote{腐败的}。德行并非艰难，并非难以达到。你们不见那些在山中的人吗？他们舍弃了房屋、妻子、儿女和一切高位，将自己与世隔绝，身穿麻衣，铺灰于下；他们脖子上挂着锁链，把自己关在狭小的斗室里。他们还不止于此，还以禁食和不断的饥饿折磨自己。如果我现在命令你们也这样做，你们岂不都退避吗？你们岂不说：这是不能忍受的吗？但是不，我并没说必须做这类事——我固然希望如此，但我并不立下法律。那么如何呢？享受你们的沐浴，照顾你们的身体，自由地进入世界，持家，有仆役侍候，随意享用你们的饮食！但在每件事上都要驱除过度，因为过度导致罪恶，同一件事，无论是什么，一旦过度，就变成罪恶；所以过度不是别的，正是罪恶。因为请看，当愤怒被激起超过适当限度时，它就冲出口出恶言，做出各种伤害；对美色、财富、荣耀或任何其他事物的无节制贪恋也是如此。请别对我说，我所说的那些人确实是坚强的；因为有许多比你更软弱、更富有、更奢侈的人，已经承担了那种艰苦严峻的生活。我何必提男人呢？还有未满二十岁的少女，她们一生都待在闺房里，过着娇柔安逸的生活，在充满香膏和香料的闺房中，躺在柔软的卧榻上，她们本身柔嫩，又因过度放纵而更加柔弱，整日里除了打扮自己、佩戴珠宝、享受各种奢华之外别无他事，从不自己服侍自己，而是有许多侍女站在旁边，穿着比皮肤还柔软的轻软衣服，上等细麻和精致的织物，不断沉浸在玫瑰和类似甜美的香气中——是的，正是这些人，在瞬间被基督的火焰抓住，脱去了所有的懒惰，甚至脱去了她们的本性，忘记了她们的娇嫩和青春，如同许多高贵的摔跤手，脱去了那柔软的衣服，冲入战场的中心。也许我所说的看来难以置信，但仍然是真实的。那么，这些极其娇嫩的少女，正如我亲耳所闻，已经将自己锻炼到如此严格的程度，以至于她们用最粗糙的马毛布包裹自己赤裸的身体，赤着娇嫩的脚不穿鞋，躺在树叶铺成的床上；不仅如此，她们彻夜守望，不在意香料或任何其他旧日的享乐，甚至让她们曾经精心打理的头发不加特别护理，只是简单朴素地束起来，以免凌乱。她们唯一的餐食在晚上，不是蔬菜，也不是面包，而是面粉、豆子、豆类、橄榄和无花果。她们不停地纺线，比在家里的女仆还要辛苦。还有呢？她们会担当照顾生病的妇女，抬她们的床，洗她们的脚。甚至许多人也做饭。基督火焰的力量如此之大；她们的热心如此超越她们的本性。\par

然而，我并不要求你们这样的事，因为你们甘心被女人超过。但至少，如果有些工作并不太费力，至少也做这些：约束粗鲁的手和不节制的眼。告诉我，有什么如此艰难？有什么如此困难？行公义正直的事，不亏负任何人，无论你们是穷人还是富人，是店主还是雇工；因为不义甚至可能延伸到穷人。或者你们没有看见，有多少争吵由此而生，把一切弄得天翻地覆？自由地结婚，生儿育女。保禄也曾如此吩咐这样的人，也曾给他们写信。我所说的那场争斗是否太大，那磐石是否太高，其顶端是否近于天堂，而你们无法达到那样的高度？那么至少抓住较小的事，追求那些较低的事。你们没有勇气舍弃自己的财富？那么至少不要夺取别人的东西，也不要亏负他们。你们不能禁食？那么至少不要放纵自己。你们不能睡在树叶铺的床上？那么至少不要为自己准备镶银的床榻；要使用不是为了炫耀，而是为了安息的床榻和被褥，而不是象牙床。使自己谦卑。为什么要用过重的货物装满你们的船？如果装备轻省，你们就一无所惧，没有嫉妒，没有盗贼，没有伏击。因为你们的钱财其实并不比忧虑更多。你们在财产上不如在焦虑和危险上丰富，\BibleQuote{这些事引人陷入许多诱惑和欲望。}\BibleRef{弟前 6:9} 这些事是渴望获得大财富的人所忍受的。我不是说，要你们服侍病人；但至少，吩咐你们的仆人去作。那么你们看，这岂非不是劳苦的工作？不，因为当柔弱的少女胜过我们如此遥远的距离时，这怎么可能呢？我恳求你们，让我们羞愧吧；因为在世俗的事上，我们确实没有在哪一点上输给她们，无论是在战争还是在竞赛中；但在属灵的争战中，她们却占了我们的上风，率先夺取了奖赏，像许多鹰一样飞得更高；而我们，像寒鸦一般，永远生活在蒸汽和烟雾中；因为把心思放在膳食和厨师上，确实是寒鸦和贪食的狗的事。请听关于古代妇女的事：她们是伟大的人物，伟大的女人，令人钦佩；如撒辣、黎贝加、辣黑耳、德波辣和亚纳；在基督的时代也有这样的人。然而她们在任何情况下都没有超过男人，而是居于第二等。但现在恰恰相反，女人超过并胜过我们。多么可鄙！这是何等的羞辱！我们占据头的位置，却被身体超过。我们被指定统治她们；不但要统治，还要在善德上统治；因为那统治的人，尤其应当在这方面统治——以美德超越；但如果他反而被超过，他就不再是统治者。你们明白基督来临的权能有多大吗？祂如何解除了诅咒？因为现在女人中守贞的比从前更多，那些贞女更端庄，寡妇也更多。没有女人会轻易说出一个不雅的字眼。那么，告诉我，为什么你们要说污秽的话？不要对我说，她们是出于沮丧或绝望而守贞。\par

女性喜爱装饰，这是她们的缺点。然而连这一点，你们作丈夫的也胜过她们，因为你们甚至以她们为你们自己的装饰而自豪。我不认为妻子如此炫耀她的珠宝，像丈夫炫耀他妻子的珠宝那样。他为自己金色的腰带而骄傲，不如为妻子佩戴金饰而骄傲。所以，连这一点也是你们造成的，是你们点燃了火花，燃起了火焰。但更重要的是，这在女人身上不如在男人身上罪大。你们被指定去管理她；在每一方面你们都要占有优越的地位。那么，在这一方面也向她表明，你们通过自己的衣着，对她的这种奢华不感兴趣。女人修饰自己比男人更合宜。如果你们自己不能逃脱诱惑，她怎能逃脱呢？她们也有虚荣心，但这与男人相同。她们在某种程度上容易动怒，这又与男人相同。至于她们所擅长的那些事，就不再与男人相同了：我指的是她们的圣德、热忱、虔诚和对基督的爱。那么，有人会说，为什么保禄把她们排除在教师的职位之外？这又一次证明她们与男人的差距有多大，也证明当时的妇女是伟大的。因为，请告诉我，当保禄或伯多禄或那些古代圣人教导的时候，一个女人介入这职务，是否合适呢？而我们却沉沦到这样的地步，以致要问为什么女人不能作教师。我们确实变得和她们一样软弱。我说这些话，并不是想抬高她们，而是要使我们羞愧，责罚和告诫我们，好使我们重新获得属于我们的权威，不是因为我们身材高大，而是因为我们的远见、对她们的保护和我们的美德。因为这样，身体也将在它合宜的秩序中，当有最好的头来管辖时。愿天主恩赐，使所有妻子和丈夫都能按祂的圣意生活，好使我们在那可怕的日子，都配得享受我们恩主的慈爱，并获得在主耶稣基督内所应许的美物。愿光荣、权能和尊崇归于祂，偕同圣父及圣神，从今世直到永远。阿们。\par

\SourceRecord{Translated by Gross Alexander.；Nicene and Post-Nicene Fathers, First Series；Vol. 13.；Edited by Philip Schaff.；Buffalo, NY: Christian Literature Publishing Co.,；1889.；<http://www.newadvent.org/fathers/230113.htm>.}

% structure: p0001 source=h1 target=section reason=page-title

\section{厄弗所人书 第十四篇讲道}

\BibleRef{弗 4:25-27}\par

\BibleQuote{因此，你们应弃绝谎言，每人应与邻人说真话，因为我们彼此互为肢体。你们生气，却不要犯罪；不可让太阳在你们的愤怒中落下；也不要给魔鬼留余地。}\par

在一般地谈论了\Quote{旧人}之后，他接着又详细地描绘了他；因为这种教导当我们通过具体细节来学习时更容易领会。他说了什么？\BibleQuote{因此，弃绝谎言。}是什么样的谎言？他是指偶像吗？当然不是；它们固然也是谎言。但他现在不是在谈论它们，因为这些人与它们无关；而是谈论人与人之间的事，即指欺骗和虚假的事。\BibleQuote{你们各人要说真话，}他说，\BibleQuote{与邻人}；接着更触动良心的还是：\BibleQuote{因为我们彼此互为肢体。}不可让人欺骗他的邻人。正如圣咏作者在几处所说：\BibleQuote{他们用谄媚的嘴唇和双关的心说话。}\BibleRef{咏 11:2}因为没有什么，是的，没有什么比欺骗和诡诈更能产生敌意了。\par

请看，他如何随处用身体的比喻来使他们羞愧。他说：不要让眼对脚说谎，也不要让脚对眼说谎。例如，如果有一个深坑，坑口横放着芦苇，上面盖着土，外表给眼看的是坚固地面的假象，眼不会借用脚去探试，看它是否下陷、下面是空心的，还是坚实的、扛得住吗？脚会说谎，不如实报告实情吗？再者，如果眼看见了蛇或野兽，它会向脚说谎吗？岂不会立刻通知它，脚得到通知就不会继续前行？再者，当眼和脚都无法分辨，全凭嗅觉，例如一种药是否有毒；嗅觉会对口说谎吗？为何不？因为它也会毁灭自己。它如实地报告它看起来的样子。再者，舌头会对胃说谎吗？当东西是苦的，它吐掉；如果是甜的，它吞咽下去。请看服务与交换；请看从真相中产生的眷顾，可以说，是发自内心的自发如此。同样，我们也应当这样；我们不要说谎，因为我们是\BibleQuote{彼此互为肢体}。这是友谊的确切标记；相反则是敌意的标记。那么，如果一个人背叛你，你会问？听真理。如果他背叛，他就不是肢体；然而他说：\Quote{不要对肢体说谎。}\par

\BibleQuote{你们生气，却不要犯罪。}\par

请看他的智慧。他既说话阻止我们犯罪，若我们不听，他仍不放弃我们；因为他的父性慈悲没有遗弃我们。正如医生给病人开药方，若病人不遵从，他仍不轻视他，继续用劝说加添建议，再继续治疗；保禄也是如此。因为若有人行事不同，只求名誉，若被忽视就烦躁；而那位凡事以病人的康复为目标的，只着眼如何恢复病人，扶起他。这就是保禄所做的。他曾说：\BibleQuote{不可说谎。}但如果说谎引发了愤怒，他又继续治疗这个。他说了什么？\BibleQuote{你们生气，却不要犯罪。}完全不动怒固然好。但若有人落入激愤中，仍不可落入如此严重的程度。\BibleQuote{不可让太阳，}他说，\BibleQuote{在你们的愤怒中落下。}你希望能尽情发怒吗？一小时、两小时、三小时对你已经足够；不要让太阳离去，留下你们彼此为敌。太阳是因天主的良善而升起；不要让它离去，照耀在不配的人身上。因为主以祂的大仁慈派遣了它，祂自己已宽恕了你的罪，而你却不宽恕你的邻人，看，这是何等的恶！此外还有另一件事。有福的保禄害怕夜晚，唯恐黑夜独自抓住那受委屈的人，仍燃烧着愤怒，再次燃起火焰。因为白天有许多事驱散愤怒，你可以放纵它；但傍晚一到，你就该和好，趁这恶尚新鲜时扑灭它；因为若黑夜赶上了它，明天将无法扑灭在黑夜中累积的更严重的恶。即使你切断了大部分，却仍不能全部切断，它会从剩余部分再次为下一夜供应燃料，使火焰更猛烈。正如太阳若在白天的炎热无法软化和驱散在夜间凝结成云的那部分空气，它就会给暴风雨提供材料，黑夜赶上余下的部分，再用新的水汽喂养它；愤怒也是这样。\par

\BibleQuote{也不要给魔鬼留余地。}\par

因此，彼此相争就是\Quote{给魔鬼留地步}；因为，我们本应紧密团结，一同对抗他，却松懈了对他的敌意，反而发出转向彼此攻击的信号；因为魔鬼的\Emph{地位}从未像在我们的敌意中那样稳固。由此产生的祸害不可胜数。正如建筑中的石头，只要紧密贴合不留缝隙，就坚固不动；但如果有一根针的缝隙，或比头发还窄的裂缝，就会毁坏整个建筑；魔鬼也是如此。只要我们紧密团结在一起，他就无法施展任何诡计；但当他使我们稍一松懈，他就如洪水般冲进来。在每种情况下，他只需一个开始，而这是最难做到的；但一旦做到，他就四面八方为自己打开道路。从此，他开启耳朵听毁谤，说谎的人更受信任；他们有敌意做辩护者，而不是真理做公正的审判者。正如友谊之处，即使是真实的恶事也显得虚假；同样，敌意之处，即使是虚假的也显得真实。那里有另一种心态，另一种审判台，不公正地听审，而带有极大的偏见和偏袒。正如天平中，若铅被放入秤盘，就会压垮整个秤；这里也是如此，只不过敌意的重量远重于任何铅。因此，我恳求你们，尽一切可能在日落前熄灭你们的敌意。因为如果你在第一天未能制伏它，那么第二天，甚至常常一整年，你都会拖延它，而敌意从此会自我增长，无需任何帮助。因为它使我们怀疑别人说的话本意并非如此，还有手势以及一切，它激怒我们，激怒我们，使我们比疯子更失常，无法容忍说出一个名字或听到它，而是用谩骂和辱骂说尽一切话。那么我们如何平息这种激情？我们如何熄灭这火焰？通过反省我们自己的罪，以及我们应向天主交账的多少；通过反省我们报仇的对象不是敌人，而是我们自己；通过反省我们使魔鬼喜悦，加强了我们的敌人，那真正的敌人，并且为他而错待了我们自己的肢体。你想报复并怀有敌意吗？那就怀有敌意吧，但要对魔鬼，而不是对你的肢体。为此，天主用愤怒武装我们，不是要我们把剑刺向自己的身体，而是要把整个剑刃浸入魔鬼的胸膛。在那里把剑一直埋到柄，是的，如果你愿意，连柄一起埋进去，再也不要拔出来，反而再加上一把又一把。当我们对自己属灵家庭的人仁慈、彼此和平时，这确实会发生。让金钱消亡，让光荣和名誉消亡；我自己的肢体比这一切都更宝贵。让我们这样对自己说；让我们不要为了获得财富、为了获得光荣而对我们的本性施暴。\par

第28节。\Quote{偷窃的，}他说，\Quote{不要再偷窃。}\BibleRef{弗 4:28}\par

你们看到旧人的肢体是什么吗？谎言、报复、偷窃。为什么他不说：\Quote{偷窃的}该受惩罚、受折磨、受拷打；而是\Quote{不要再偷窃}？\Quote{反而要劳力，亲手做有益的事，好能有所分施给有需要的人。}\par

那些自称为纯洁的人在哪里？他们满身污秽，却还敢给自己取这样的名号？因为这是可能的，非常可能的，摆脱羞辱，不仅靠停止犯罪，还要行善事。你们觉悟我们该如何脱离罪吗？\Quote{他们偷窃了。}这是罪。\Quote{他们不再偷窃。}这并没有除去罪。但如何除去呢？如果他们劳力，并爱德地分施给别人，这样就能除去罪。他不仅希望我们工作，而且要\Quote{工作}到\Quote{劳力}的地步，好使我们能\Quote{分施}给别人。因为窃贼确实也工作，但做的是恶事。\par

第29节。\Quote{不要让任何腐败的话从你们口中出来。}\BibleRef{弗 4:29}\par

什么是\Quote{腐败的话}？就是别处也被称为\Quote{闲话、毁谤、污秽的谈话、戏笑、愚妄的谈论}。你们看他如何斩断愤怒的根源？谎言、偷窃、不合时宜的谈话。然而，\Quote{不要再偷窃}这句话，他与其说是原谅他们，不如说是安抚受害者，并规劝他们满足，如果永不再遭受同样的事。他也很好地给了关于谈话的劝告；因为我们不仅因行为，也要因言语受罚。\par

他继续说：\Quote{只要是有益的，按需造就，好能恩宠听者。}\par

这就是说，只说你邻居有益的事，一句多余的话也别说。因为天主给你口舌正是为此，让你感谢祂，让你建立你的邻居。所以如果你毁坏了那建筑，还不如沉默，永不说话。因为正是工人的手，如果不筑墙而学会拆墙，就理应被砍掉。因为圣咏作者也说：\Quote{主要剪除一切谄媚的嘴唇。}\BibleRef{咏 11:3}口——这是万恶之源；或者更确切地说，不是口，而是那些邪恶使用它的人。从此产生侮辱、辱骂、亵渎、情欲的诱惑、凶杀、奸淫、偷窃，一切都源于此。你会说，凶杀如何产生？因为从侮辱走向愤怒，从愤怒走向殴打，从殴打走向凶杀。那么奸淫又如何？一个人会说：\Quote{某个女人爱你，她说了些好话。}这立刻瓦解了你的坚定，于是你的情欲就在你内燃起。\par

因此保禄说：\Quote{那善的}。既然话语如此浩大，他理所当然地泛泛而言，吩咐我们使用这类言辞，并给我们一个交流的范本。那么这是什么？通过说\Quote{为了建树}，他要么是指这个，即那听见你的人可能感激你：例如，一个弟兄犯了奸淫；不要张扬这个过犯，也不要幸灾乐祸；你这样做对听者毫无益处；相反，很可能你会伤害他，因为你给了他刺激。然而，你劝告他该做什么，就是对他施以极大的恩惠。你训导他如何保持缄默，教导他不可辱骂任何人，你就给他上了最好的一课，你就对他施以最大的恩惠。与他谈论痛悔、虔诚、施舍；这一切都会软化他的灵魂，为这一切他都会承认你的恩情。而激起他的笑声，或通过污秽的交谈，你反而会煽动他。称赞邪恶，你就要倾覆并毁灭他。\par

或者他这样意思是：\Quote{好使他们，听者，充满恩宠}。因为正如香膏使享用的人得到恩宠，同样好言语也是如此。因此，有一个人曾说：\Quote{你的名字如同倒出的香膏}。\BibleRef{歌 1:3}它使他们散发那股馨香。你看，他继续推荐的是什么，现在他也在说，命令各人按各自的能力建树邻人。那么，你给出这样劝告给他人，岂不更该给自己！\par

第30节。他又说：\Quote{不要叫天主的圣神担忧}。\BibleRef{弗 4:30}\par

这是一件更可怕、更惊人的事，正如他在致得撒洛尼人的书信中也这样说；因为在那里他也使用了这样的一种表述。\Quote{那拒绝的，拒绝的不是人，而是天主。}\BibleRef{得前 4:8}这里也是如此。如果你说了一句辱骂的话，如果你打击你的弟兄，你不是打击他，你是\Quote{使圣神担忧}。然后还进一步增添了所赐的恩惠，以加重责备。\par

他说：\Quote{不要叫圣神担忧，你们原是受了他的印记，以待救赎的日子。}\par

他印记我们为一群王室的羊群；他，将我们从一切先前的事物中分别出来；他，不让我们与那些遭受天主义怒的人同列——而你竟使他担忧？看，他的话在那里多么惊人：\Quote{因为那拒绝的，}他说，\Quote{拒绝的不是人，而是天主：}而在这里多么刺人：\Quote{不要叫圣神担忧，}他说，\Quote{你们原是受了他的印记。}\par

道德。让这个印记常在你的口上，永远不要消灭这个印记。属灵的口从不说出这类事。不要说：\Quote{这不算什么，如果我说了一句不雅的话，如果我侮辱了某人。}正因为这似乎不算什么，所以是极大的恶。因为那些似乎不算什么的事，就这样容易被轻视；而那些被轻视的事，就会不断增长；那些不断增长的事，就成为无可救药。\par

你有一个属灵的口。想想你刚一出生就说了什么话——什么话与你的口相称。你呼求天主为\Quote{圣父}，而你却立刻辱骂你的弟兄？想想看，你为何呼求天主为\Quote{圣父}？是出于本性吗？不，你绝不能这样说。是出于你的良善吗？不，也不是这样。那么是出于什么？是出于纯粹的慈爱，出于温柔，出于祂的大慈大悲。因此，每当你呼求天主为\Quote{圣父}时，不仅要想到，通过辱骂你正做与你这崇高出身不相称的事，而且也要想到，正是出于慈爱你才有了这崇高出身。那么，在从纯粹的慈爱中领受之后，不要通过对你弟兄的残忍而玷污它。你呼求天主为\Quote{圣父}，却辱骂？不，这不是圣子的事。这些离祂很远。圣子的事是宽恕祂的仇敌，为那些钉死祂的人祈祷，为那些恨祂的人流祂的血。这些是配得上圣子的事：使祂的仇敌——忘恩负义者、不诚实者、鲁莽者、背信弃义者——使他们成为弟兄和继承人：而不是像对待奴隶那样侮辱已成为弟兄的人。\par

想想你的口说出过什么样的话——这些话配得上什么样的餐桌。想想你的口触碰过什么，品尝过什么，领受了什么样的食物！你竟以为辱骂你的弟兄不是什么严重的事吗？那你又怎么称他为弟兄呢？如果他不是弟兄，你如何说\Quote{我们的天父}？因为\Quote{我们}一词指许多人。想想你在举行奥迹时与谁站在一起！与革鲁宾、与色辣芬！色辣芬并不辱骂：不，他们的口只履行这一项职责，就是唱赞美颂，光荣天主。那么，若你用口辱骂，你又怎能与他们一同说\Quote{圣、圣、圣}？请告诉我。假设有一只御用器皿，总是盛满御膳，专为此目的而设，然后有一个仆役拿来盛粪。那器皿装满粪后，他还会再胆敢把它放回其他专用器皿中吗？当然不会。辱骂就是这样，毁谤就是这样。\Quote{我们的天父！}但这是什么？仅此而已吗？还要听接下来的话：\Quote{你是在天上的。}你一说\Quote{我们的天父，你是在天上的}，这话就提升你，给你的心神加上翅膀，指示你有一位在天上的父。那么，不要做任何地上的事，不要讲任何地上的话。祂将你列于那高天的大军之中，将你编入天上的歌咏团。你为何把自己拖下来？你站在君王的宝座旁，却辱骂？你不怕君王认为这是冒犯吗？如果在我们中间，一个仆役殴打或攻击同仆，即使他做得公正，我们也立即斥责他，并认为这是冒犯；而你呢，与革鲁宾一同站在君王宝座旁，却辱骂你的弟兄？你没有看见这些圣器吗？它们不是只用于一个目的吗？有谁敢把它们用于其他用途？然而你比这些器皿更圣洁，是的，远更圣洁。那么，你为何玷污、为何污染自己？你站在天上，却辱骂？你与天使同籍，却辱骂？你配得上主的亲吻，却辱骂？天主以如此多、如此大的事物恩赐你的口：天使的赞美诗、食物——不是天使的，不，而是超越天使的——祂自己的亲吻，祂自己的拥抱，而你却辱骂？哦，不，我恳求你。由此产生的祸害极大；愿基督徒的灵魂远离它。我说话时不能说服你，不能使你羞愧吗？那么我现在有责任警告你。请听基督的话：\Quote{凡向弟兄说\InnerQuote{愚笨的人}的，应受地狱的火。}\BibleRef{玛 5:22}如果最轻微的言语都导致地狱，那么说狂妄话的人该当受什么？让我们训练我们的口保持静默。由此而来的益处甚大，恶语带来的祸害甚大。我们不可在此耗费我们的财富。让我们给口装上门户和门闩。若有不快的话从我们口中溜出，让我们痛责自己。让我们恳求天主，让我们恳求我们所辱骂的人。不要认为这样做有失身份。我们伤害的是我们自己，不是他。让我们用祈祷和与所辱骂之人的和好作为补救。若我们在言语上要如此谨慎，更应在行为上为自己制定律法。是的，如果我们有朋友，无论他们是谁，若他们向任何人说恶言或辱骂，就要要求他们并令其赔偿。我们必须学会：这种行为本身就是罪；因为我们若学会这个，就会很快远离它。\par

愿平安的天主保守你的心思和你的口舌，用可靠的围篱——即祂的敬畏——围护你，藉着我们的主耶稣基督，愿光荣归于圣父及圣神，直到永远。阿们。\par

\SourceRecord{Translated by Gross Alexander.；Nicene and Post-Nicene Fathers, First Series；Vol. 13.；Edited by Philip Schaff.；Buffalo, NY: Christian Literature Publishing Co.,；1889.；<http://www.newadvent.org/fathers/230114.htm>.}

% structure: p0001 source=h1 target=section reason=page-title

\section{厄弗所书讲道词 第十五篇}

\BibleRef{弗 4:31}\par

\BibleQuote{一切毒辣、怨恨、忿怒、争吵、毁谤以及一切邪恶，都要从你们中除掉。}\par

正如蜜蜂决不会在一个不洁净的器皿中安家——这就是为什么擅长此事的人用香料、香油和甜味喷洒场地；他们还在蜂箱（蜜蜂孵化后即将定居的柳条篮）上喷洒香酒和所有其他甜物，以免有任何臭味惹恼它们，使它们再次飞走——同样，圣神也是如此。我们的灵魂好比一种器皿或篮子，能够接收灵性恩赐的蜂群；但如果其中有胆汁和\BibleQuote{毒辣和怨恨}，蜂群就会飞走。因此，这位有福而有智慧的园丁彻底洁净我们的器皿，既不吝惜刀子，也不吝惜任何铁器，并邀请我们到这群灵性蜜蜂面前；当他聚集时，他用祈祷、劳苦和一切其余来洁净我们。请注意他怎样洁净我们的心。他已逐出谎言，他已逐出愤怒。现在，他又指出那邪恶如何可以更彻底地被根除；如果我们不是，他说，\BibleQuote{毒辣}的灵。因为正如我们的胆汁常发生的情况：如果只有一点点，那么即使容器破裂，也只会引起一点点紊乱；但如果这种物质的强度和酸性过度，那么先前容纳它的器皿就不再能容纳它，就像被灼热的火烧穿一样，它不再能将它控制在指定范围内，而是被其极度的尖锐撕裂，让它溢出并伤害整个身体。它就像一只被带进城市的极其凶猛可怕的野兽；只要它被关在为其制作的笼子里，无论它多么狂暴，多么咆哮，都不能伤害任何人；但如果它被愤怒压倒，冲破中间的栅栏，能够跳出来，它就会使城市充满各种混乱和骚动，使所有人奔逃。胆汁的性质也是如此。只要它保持在适当的界限内，就不会对我们造成大的危害；但一旦包裹它的薄膜破裂，没有任何东西阻止它立即散布到整个系统，那么，我说，在那一刻，尽管量非常小，但由于其质量的过度强度，它以其特有的毒性污染我们本性的所有其他成分。例如，它发现血液在位置和性质上都与它相近，使血液中的热变得更加酸涩，事实上，它附近的一切也是如此；它从正常的温度溢出界限，将一切变成胆汁，并立即攻击身体的其他部分；这样，它将其有毒品质注入一切，使人说不出话，导致他断气，驱逐生命。现在我为什么如此详细地陈述这一切？为的是，从这身体上的毒辣，我们可以理解那灵魂上的毒辣的不可容忍的邪恶，以及它如何首先摧毁那产生它的灵魂本身，使一切变得毒辣，从而我们可以避免经历它。因为如同前者燃烧整个体质，后者也燃烧思想，并将它的俘虏带到地狱的深渊。因此，为了通过仔细考察这些事来逃避这邪恶，勒住这怪兽，或更确切地说，彻底根除它，让我们听从\Person{保禄}所说的：\BibleQuote{让一切毒辣都要}（不是消灭，而是）\BibleQuote{除掉}从你们中。我何必费心去约束它？当我有能力将它从我的灵魂中驱逐、移除并放逐时，还有什么必要看守一个怪兽？让我们听从\Person{保禄}，他说：\BibleQuote{让一切毒辣从你们中被除掉。}但是，啊，那控制我们的乖谬！虽然我们应当尽一切努力做到这一点，却有些人如此愚昧，竟然因这邪恶而自我祝贺，并以此为傲，以此为荣，还被人嫉妒。\Quote{某某人，}他们说，\Quote{是一个毒辣的人，他是蝎子、蛇、毒蛇。}他们视他为可怕的人。但是，好人哪，你为什么害怕那毒辣的人？\Quote{我害怕，}你说，\Quote{他伤害我，他毁灭我；我无法抵抗他的恶意，我害怕他把我这个单纯的人、无法预见他任何阴谋的人，扔进他的陷阱，把我们缠在他为了欺骗我们而设的网罗里。}我不禁微笑。为什么呢？因为这些是小孩子的论点，他们害怕那些不该怕的东西。真的，没有什么比一个毒辣和怀恶意的人更值得我们鄙视，更值得我们嘲笑。因为没有什么比毒辣更无力。它使人变得愚昧无知。\par

难道你们不见恶意是盲目的吗？难道你们从未听过，为邻人挖坑的，是为自己挖坑吗？也许有人说：我们岂不该惧怕一个充满骚动的灵魂吗？如果我们真的该像惧怕恶灵、愚人和疯子那样惧怕心怀怨恨的人（因为他们确实随意行事），我自己也承认；但如果我们该像惧怕善于处理事务的人那样惧怕他们，那绝不可能。因为对于妥善处理事务，没有比明智更必需的了；而没有什么比邪恶、恶意和虚伪更能阻碍明智。看看那些患黄疸的人，他们多么难看，所有的光彩都枯萎了。他们多么虚弱、瘦小，一无所用。同样，这类灵魂也是如此。邪恶是什么，不就是灵魂的黄疸病吗？那么邪恶没有一点力量，确实没有。你们愿意我再用一个实例，通过摆出奸诈者和真诚者的画像，让我的话更明白吗？阿贝沙隆是个奸诈者，他\BibleQuote{偷得了所有人的心}。\BibleRef{撒下 15:6}请看他的奸诈多么大。记载说：\Quote{他走来走去，}记载说，\Quote{并且说：\InnerQuote{你们没有判断吗？}}想要争取每个人归向他。但达味是真诚的。那么怎样？看他们二人的结局，看前者多么充满疯狂！因为他只顾伤害父亲，在其他事上都瞎了眼。但达味却不如此。因为\BibleQuote{行正直路的，步步安稳} \BibleRef{箴 10:9}；而且合情合理：他不以过分精明的方式处理任何事情，不谋恶。那么，让我们听从真福保禄，让我们怜悯，是的，让我们为心怀怨恨的人哭泣，让我们用尽一切方法，竭尽全力从他们的灵魂中根除这一恶习。因为当我们体内有胆汁时（尽管那确实是有用的元素，因为没有胆汁人不可能存在，我指的是作为他本性元素的胆汁），我说，我们竭尽所能摆脱它，虽然我们从它获益甚多，却不为摆脱灵魂中的那种东西做任何事，也不付出任何努力，尽管它在任何情况下都无益，反而有极大的害处，这岂非荒谬？他说：\BibleQuote{你们中若有人自以为有智慧，该变为愚妄，好成为有智慧的} \BibleRef{格前 3:18}。再听路加说什么：\BibleQuote{他们怀着欢乐和诚实的心进食，赞美天主，获得全民的爱戴} \BibleRef{宗 2:46}。怎么，我们岂不见现今纯朴真诚的人也受到众人的尊敬吗？没有人嫉妒这样一个人顺遂时，没有人践踏他逆境时，反而众人与他同乐顺遂，与他同忧不幸。而每当心怀怨恨的人兴旺时，众人无不哀叹，如同发生了坏事；但他若不幸，众人无不欢喜。那么，让我们怜悯他们，因为他们在全世界都有共同的敌人。雅各伯是个真诚的人，却战胜了奸诈的厄撒乌。\BibleQuote{因为智慧不进入存心不良的灵魂里} \BibleRef{智 1:4}。\BibleQuote{一切怨恨都该从你们中间除掉}。不要让一丝残余存留，因为若被搅动，它会像阴燃的火炭一样，将内心的一切转为熊熊烈火。那么，让我们清楚理解这怨恨是什么。例如，虚伪的人、狡猾的人、伺机作恶的人、多疑的人。从这样的人那里，常会产生\Quote{忿怒}和\Quote{愤怒}；因为这样的灵魂不可能安宁，而\Quote{愤怒}和\Quote{忿怒}的根源正是\Quote{怨恨}。这种性格的人既愠怒，从不放松自己的灵魂；他总是郁郁寡欢，总是阴沉。因为正如我所说，他们自己首先收获自己恶行的果实。\par

\Quote{以及喧嚷，}他补充说。\par

怎么，连喧嚷也要除掉吗？是的，因为温和的人必须具有这样的性格，因为喧嚷承载愤怒，如同马承载骑手；绊倒马，你就会摔下骑手。\par

教训。这一点尤须女人注意，她们在任何场合都高声喊叫和吵闹。只有一件事情上高声喊叫是有用的，就是在宣讲和教导时。但在任何其他情况，绝对没有，甚至在祈祷中也没有。如果你要学一个实际的功课，永远不要高声喊叫，那么你就永远不会生气。看，这就是保持温和的方法；因为正如不喊叫的人不可能发怒，同样，喊叫的人不可能不发怒。因为你莫要对我说一个人是无法平息、报复心强、有天生的怨恨和天生的暴躁。我们现在说的是这种情欲的突然发作。\par

为此，养成灵魂绝不提高声音、从不呼喊的习惯，对达到这一目的不无小补。你断绝呼喊，就能剪去愤怒的翅膀，你就能抑制内心的第一次冲动。因为正如人不可能不举起双手就摔跤一样，同样，人也不可能不提高声音就陷入争吵。你捆住拳击手的双手，然后命令他击打——他做不到。同样，愤怒也会被解除武装。但呼喊会激起愤怒，即使它并不存在。因此，女性尤其容易陷入其中。妇女每当对女仆发怒时，便以喊声充满整个房屋。而且常常，如果房子恰好建在狭窄的街道旁，所有路人都能听到主母责骂，女仆哭泣哀号。还有什么比那哭号的声音更可耻的呢？世间发生了什么事？周围所有的女人立即窥视进来，其中一个说：\Quote{某人正在打自己的女仆。}还有什么比这更无耻的呢？\Quote{那么，难道根本不应该打吗？}不，我不是这样说（因为打是必要的），但不能频繁，不能过度，也不能为自己的任何过错，正如我常说的，也不能为她服务上的任何小过失，而只能在她伤害自己的灵魂时。如果你为此类过错惩罚她，人人都会赞许，无人责备你；但若为你自己的任何理由，众人都将谴责你的残忍与严苛。而比这一切更卑劣的是，有些人如此凶猛粗野，以致鞭打她们到瘀伤数日不消。她们竟剥去少女的衣服，为此召唤自己的丈夫，且常常将她们绑在床榻上。哀哉！那时，告诉我，难道没有一个关于地狱的念头临到你吗？什么？你剥去女仆的衣服，将她暴露在你丈夫面前？你难道不羞于他因此责备你吗？然后你又激怒他，威胁要将她锁起来，先用万千辱骂的名称指责那可怜可悲的人，称她为\Quote{\Place{帖撒利}女巫、逃奴和娼妓}？\par

因为她的愤怒不容她甚至顾惜自己的口舌，她只专注于一个目标：如何对另一个人报仇雪恨，即便自己蒙羞。然后，这些事后，她当然会像任何暴君一样端坐，召唤儿女，召来愚昧的丈夫，把他当作刽子手对待。这些事应该发生在基督徒家中吗？\Quote{是的}，你们说，\Quote{但奴隶是一个麻烦、鲁莽、无耻、不可救药的种族。}不错，我自己也知道，但有其他方法维持秩序：用惊吓、威胁、言语——这些既能更有效地触动她，又能使你不致蒙羞。你，一个自由妇女，口出秽言，难道你不是更羞辱自己吗？然后，如果她要去浴场，背上便有瘀伤，裸露时她带着你残忍的印记。\Quote{但是}，你们说，\Quote{整个奴隶族类若受宽容，便无法忍受。}不错，我自己也知道。但正如我所说，用别的方法纠正她们，不单用鞭子和恐吓，也用奉承和善行。如果她是信徒，她就是你的姊妹。要想到你是她的主母，她服侍你。若她放纵无度，就断绝醉酒的机会；叫你的丈夫来，告诫她。难道你不觉得一个女人被打是多么可耻吗？那些为男子设立无数刑罚的人——火刑柱和刑架——也几乎从不吊死一个女人，而是将男人的怒气限制在打她耳光；他们对女性如此尊重，以至于即便绝对必要时，若她恰巧怀孕，也常常不吊死她。因为男人打女人是耻辱；男人打尚且如此，同性打就更甚了。此外，正是这些事使妇女在丈夫眼中成为可憎。\Quote{那么}，你们可能说，\Quote{如果她行为娼妓呢？}就让她嫁给一个丈夫；断绝通奸的机会，不让她过于养尊处优。\Quote{那么，如果她偷窃呢？}看管她，监视她。——\Quote{过分！}你们会说，\Quote{什么，难道我要作她的看守？多么荒谬！}请问，为什么你不该作她的看守？她不是与你拥有同样种类的灵魂吗？她不是蒙受天主同样的恩赐吗？她不与你分享同一张桌子吗？她不是与你同享同样的高贵出身吗？\Quote{但是那么}，你们会说，\Quote{如果她是个诽谤者、或饶舌者、或醉鬼呢？}可是，有多少自由妇女也是如此啊？现今，对于妇女的一切缺点，天主已命男人担负：只是，祂说，不要让妇女成为娼妓，除此之外的一切缺点都要容忍。是的，即使她是醉鬼、或诽谤者、或饶舌者、或恶眼者、或浪费者、挥霍你的财产，你仍要她作你生活的伴侣。训练并约束她。你必须如此。正是为此你是头。因此，管理她，尽你的本分。是的，即使她不可救药，是的，即使她偷窃，也要保管你的财物，不要过分惩罚她。如果她饶舌，就使她沉默。这是最高的智慧。\par

但如今，有些人竟堕落到如此无耻的地步，以至露出头来，并拉着女仆的头发。——你们为何都脸红呢？我并非对你们全体说话，而是对那些陷入如此野蛮行为的人。保禄说：\BibleQuote{女人不可不蒙头。}\BibleRef{格前11:5}那么，你竟完全除去她的头饰吗？你看出你是如何侮辱你自己了吗？如果她确实以头赤裸地出现在你面前，你会称之为侮辱。那么，当你自己赤裸其头时，你却说不令人震惊吗？然后你会说：\Quote{如果她不改正呢？}那么就用杖和鞭子责打她。然而，你自己也有多少缺失，而你自己却不改正！我说这些话不是为了她们，而是为了你们这些自由妇女，叫你们不做任何如此不相称的事，不做任何使你们蒙羞的事，叫你们不伤害自己。如果你在家中对待女仆时学会这功课，不苛刻而是温柔宽厚，那么你在对待丈夫时更会如此。因为她虽有权柄，却不做这类事，那么在有约束的情况下就更不会做了。所以，对女仆所用的纪律，对你们获得丈夫的善意将大有裨益。因为祂说：\BibleQuote{你们用什么量器量给人，也必用什么量器量给你们。}\BibleRef{玛7:2}要在你们口上套上辔头。如果你受训练能勇敢地忍受女仆的回嘴，你就不会因同等之人的无礼而烦恼，且在超越烦恼中，你将达到最高的哲学。然而有些人甚至加上誓言，但没有什么比一个如此愤怒的女人更令人震惊的了。但你会说，如果她打扮得花枝招展呢？那么，禁止这个；我同意；但先以你自己为开始来遏制它，不是靠恐惧而是靠榜样。在一切事上你自己做一个完全的模范。\par

他说：\Quote{一切诽谤，}他说，\Quote{要从你们中除掉。}观察恶行的进程。苦毒产生忿怒，忿怒产生怒气，怒气产生喧嚷，喧嚷产生诽谤，即辱骂；接着从恶言发展到殴打，从殴打发展到伤害，从伤害发展到死亡。然而保禄不愿提及这些，只说：\Quote{这，}他说，\Quote{要从你们中除掉，连同一切恶意。}什么是\Quote{连同一切恶意}？它以此终结。因为有些人像那些暗中咬人的狗，它们对走近的人完全不吠，也不发怒，却摇尾乞怜，露出温和的样子；但当我们不防备时，它们便会咬住我们。这些人比那些公开为敌的更危险。既然也有一些人像狗，既不喊叫，也不发怒，受冒犯时也不威胁我们，却在暗中策划，图谋无数恶事，不在言语上报复而在行动上报复；他暗指这些人。他说：让那些事从你们中被除掉，\Quote{连同一切恶意。}不要节省你的言语，随后在行动上报复。我管束我的舌头、减少喧嚣的目的，是为了防止它点燃更猛烈的火焰。但如果你没有喧嚣却做着同样的事，你心里藏着火和燃烧的炭，你的沉默又有什么益处呢？难道你不知道那些内部燃烧、外表不显的火灾是最具破坏性的吗？那些从不暴露到表面的伤口是最致命的；那些烧灼内脏的热病是最严重的吗？同样，这种吞噬灵魂的愤怒是最危险的。但他说，连这个也要从你们中除掉，\Quote{连同一切恶意，}包括各种各类，大的小的。那么让我们听从他，让我们驱除一切\Quote{苦毒和一切恶意，}免得我们\Quote{使圣神忧郁。}让我们消灭一切苦毒；让我们将它连根拔起。一个苦毒的灵魂绝不能产生任何好、任何有益的东西，只有不幸，只有眼泪，只有哀哭和悲叹。你们难道没有看见那些吼叫或呼叫的野兽，我们如何躲避它们？例如狮子、熊？但对羊却不是这样；因为羊没有吼叫，而是柔和温顺的声音。同样，在乐器中，那些响亮刺耳的最令人不悦，如鼓和号；而那些不是这样却柔和悦耳的，则令人愉快，如笛、竖琴和箫。那么让我们预备我们的灵魂，使它永不高声喊叫，这样我们也能克制我们的愤怒。当我们剪除了这个，我们自己将首先享受宁静，我们将驶入那和平的港湾，愿天主恩赐我们众人皆能抵达，在我们的主耶稣基督内，愿光荣、权能、尊崇归于父，偕同圣神，从今世直到永远。阿们。\par

\SourceRecord{Translated by Gross Alexander.；Nicene and Post-Nicene Fathers, First Series；Vol. 13.；Edited by Philip Schaff.；Buffalo, NY: Christian Literature Publishing Co.,；1889.；<http://www.newadvent.org/fathers/230115.htm>.}

% structure: p0001 source=h1 target=section reason=page-title

\section{\WorkTitle{厄弗所书讲道第十六篇}}

\BibleRef{厄 4:31-32}\par

\BibleQuote{一切毒辣、怨恨、忿怒、争吵、毁谤以及一切邪恶，都要从你们中除掉；彼此要善良，要仁慈，互相宽恕，如同天主在基督内宽恕了你们一样。}\par

如果我们想达到天国，单单弃绝邪恶是不够的，还必须大量行善。确实，要免受地狱之苦，我们必须弃绝邪恶；但要达到天国，我们必须紧守德行。难道你们不知道，就连外邦人的法庭上，当审查人的行为、全城聚集时，也是如此吗？不但如此，外邦人中有个古老习俗：用金冠冕加冕——不是给那没有对国家作恶的人，因为那本身不过足以使他免于惩罚——而是给那提供了伟大公共功绩的人。只有这样，一个人才会被提升到这种荣誉。但是，我特别需要说的话，不知怎的，几乎从我口中溜走了。因此，对我刚才所说的话稍作纠正后，我收回这一划分的第一部分。\par

因为当我说弃绝邪恶足以防止我们掉进地狱时，正在说话间，一个可怕的判决悄悄临到我，这判决不仅对那些胆敢作恶的人施加惩罚，而且惩罚那些错过任何行善机会的人。那么，这是什么判决呢？当那一天，那可怕的日子来到，既定时间已到，审判者坐在审判座上，把绵羊放在右边，山羊放在左边；他对绵羊说：\BibleQuote{我父所祝福的，你们来吧！承受自创世以来给你们预备了的国度吧！因为我饿了，你们给了我吃的。}\BibleRef{玛 25:34}到目前为止很好，因为这样的慈悲应得这奖赏。然而，那些没有将自己所有的分给有需要的人，他们不仅应失去福乐，还要被投入地狱之火，这有何公正理由呢？当然，这也有其公平理由，与前者一样：因为我们从中得知，行善的人将享受天上的福乐，而那些虽未被指控有恶行，却忽略了行善的人，将和行恶的人一起被投入地狱之火。除非有人可以说，不行善本身就是恶的一部分，因为它源于懒惰，而懒惰是恶习的一部分，或者更确切地说，不是一部分，而是它的一个根源和有害的根。因为懒惰是万恶之师。因此，我们不要愚蠢地问这样的问题：那既没有作恶也没有行善的人会占据什么位置？因为不行善本身就是作恶。告诉我，如果你有一个仆人，他不偷窃，不辱骂你，不违抗你，并且不酗酒，不作任何恶习，但却总是闲坐，不尽仆人应尽的对主人的任何职责，你会不会惩罚他？你会不会把他放在刑架上？告诉我。然而，他确实没有作恶。不，但这就是作恶。但是，如果你愿意，让我们将此应用到生活中的其他情况。再比如一个农夫。他不损坏我们的财产，不图谋反对我们，他不是小偷，他只是双手被绑在身后，坐在家里，既不播种，也不犁一垄地，不把牛套在轭下，不照看葡萄树，事实上不从事耕作所需的任何其他劳动。现在，我说，我们难道不应惩罚这样的人吗？然而他没有对任何人做错事；我们无法指控他。不，但正是这一点他做错了。他做错是因为他没有为自己的那份公共福利做出贡献。再告诉我，如果每一个工匠或技工都不作恶——甚至对同行也不作恶——只是懒惰，我们的生活难道不会因此完全结束并灭亡吗？你希望我进一步以身体为例来展开论述吗？那么，让手既不打头，也不割掉舌头，不挖出眼睛，不做这类任何恶事，只是闲置着，不向全身提供应有的服务，难道它不应该被砍掉，而不是让人带着它闲置，成为整个身体的累赘吗？同样，如果嘴巴既不咬手，也不咬胸，但却未能履行所有适当的职责，难道它不是应该被封住更好吗？因此，如果在仆人、技工和整个身体的情况下，不仅作恶，而且不行善都是极大的不义，那么对于基督的身体来说，更是如此。\par

道德。因此，真福保禄在引领我们远离罪恶时，也引领我们走向德性。因为，请告诉我，如果好种子没有播下，仅仅割除所有荆棘又有什么益处呢？因为我们的劳作若半途而废，终将回到同样的祸害中。因此，保禄怀着对我们深切而慈爱的关切，并未让他的劝诫停留在根除和摧毁恶情上，而是立刻敦促我们展现出培植善情。因为他说了：\BibleQuote{一切毒辣、怨恨、忿怒、吵闹、毁谤，连同一切恶毒，都应当从你们中除掉}之后，又加上：\BibleQuote{彼此要以仁慈相待，要心存怜悯，互相宽恕。}因为这些都是习惯和性情。我们抛弃前者并不足以使我们稳定地操练后者，还需要新的动力，以及不亚于我们避免恶情时所付出的努力，才能获得善情。正如在身体方面，黑人如果摆脱了这种肤色，并不会立即变白。更确切地说，让我们不要用物理的论据来论证，而是从关乎道德选择的例子来取材。一个人若不是我们的敌人，并不一定是我们的朋友；还有中间状态，既非敌对也非友谊，也许大多数人对我们就是如此。一个人没有哭，并不一定就在笑，而是有中间状态。同样，这里也是如此。没有\Quote{毒辣}的人并不一定就\Quote{仁慈}，没有\Quote{忿怒}的人并不一定就\Quote{心存怜悯}；而需要明确的努力，才能获得这种卓越的品德。现在，请看真福保禄如何按照最优秀的农艺规则，彻底清理和耕作那位农夫托付给他的土地。他除去了坏种子；现在他劝我们保留好植物。他说：\Quote{要仁慈}，因为如果荆棘拔除后田地闲置，它又会生出无用的杂草。因此，有必要通过播下好种子和好植物来预先占据那空闲和休耕的状态。他除去\Quote{忿怒}，放入\Quote{仁慈}；他除去\Quote{毒辣}，放入\Quote{心存怜悯}；他铲除\Quote{恶毒}和\Quote{毁谤}，栽上\Quote{宽恕}来代替。因为\InnerQuote{互相宽恕}这句话的意思是：他意指，你们要存心彼此宽恕。这种宽恕比在金钱事务上表现的宽恕更为伟大。因为那宽恕借款者金钱债务的人，固然做了崇高可嘉的事，但这份仁慈仅限于身体，尽管他本人因那属灵且涉及灵魂的益处而得到充份的报答；然而那宽恕过犯的人，将同样造福自己的灵魂和接受宽恕之人的灵魂。因为如此行事，他不仅使自己，也使他人更富爱德。因为报仇并不如宽恕那样深刻地触动那些得罪我们之人的灵魂，从而羞辱他们，使他们惭愧。因为若选择另一条路，我们对自己和他们都无益处，反而通过寻求报复（如同犹太人的首领那样）并点燃他们内心的怒火而伤害双方；但若我们以温柔回报不义，我们将化解他的一切愤怒，并在他的心中设立一个法庭，作出有利于我们的判决，且比我们自己所能做的更严厉地谴责他。因为他会自责，定自己的罪，并寻找各种借口，以更充份的尺度回报所赐予他的那份忍耐，知道如果他以同等尺度回报，他就处于劣势，因为不是他先开始，而是从我们那里接受了榜样。因此，他会努力超越尺度，以回报的过度来掩盖他自己在回报行动中居于第二位的不利；而另一方若曾是首先受苦者，他所遭受的因时间而累积的不利，则将以过度的仁慈来弥补。因为人若心地正直，受到他们曾经伤害之人的善待，比受到伤害更受感动。因为一个人受到善待却不知回报，是可耻的罪恶，是受人责备和轻蔑的；而一个人受到虐待却不图报复，则得到所有人的赞扬、喝彩和好评。因此，这种品行比任何其他事情都更能触动他们。\par

所以，你若有意报复，就这样报复：以善报恶，使对方成为你的债务人，并赢得光荣的胜利。你受了恶吗？行善；这样你就向你的敌人报了仇。因为如果你图谋报复，众人会同样指责你和他。但如果你忍耐，情况就不同了：他们会称赞你、钦佩你，而他却会受到责备。对一个敌人来说，还有什么比看到他的敌人受到众人称赞和钦佩更痛苦的惩罚呢？对一个敌人来说，还有什么比在敌人面前受到众人责备更痛苦的呢？如果你报复他，你自己可能也会受到谴责，并且只有你一个人报复；但如果你宽恕他，所有人都将成为你的报复者。这比他可能遭受的任何恶事都更严厉：他的敌人竟有这么多人替他报仇。你若开口，他们就会沉默；但你若沉默，他们就会用成千上万只舌头打击他，使你获得更大的报复。如果你责备他，很多人反而会指责你（因为他们会说你的话是出于激情）；但当那些没有受过他伤害的人因你过分的忍耐而同情你时，就像他们自己受了委屈一样，这种报复就毫无嫌疑了。你会说：\\Quote{但那又将如何}，你会说，\\Quote{倘若没有人报复？}但人们不可能如此铁石心肠，看到这样的智慧而不钦佩。即使他们当时不报复他，过后当他们心情合适时，也会这样做，并且他们会继续嘲笑他、辱骂他。即使没有其他人钦佩你，那个人自己一定会钦佩你，尽管他可能不承认。因为我们对正义的判断，即使我们陷于极深的邪恶，仍然是公正不偏的。为什么你们以为我们的主基督说：\\BibleQuote{谁打你的右脸，连另一面也转给他}（\\BibleRef{玛 5:39}）？岂不是因为一个人越忍耐，他对自己和他人所施的恩惠就越显著吗？因此祂命令我们\\Quote{连另一面也转给他}，以满足发怒者的欲望。因为谁是这样怪物，能不立刻感到羞愧？据说就连狗也有这种感觉：如果它们吠叫并攻击一个人，而那人躺下不动，就能平息它们的愤怒。如果它们敬畏愿意受它们伤害的人，那么人类，既然更理性，更会如此。\par

然而，我们不应忽略不久之前我回忆起并引以为证的一点。那是什么呢？我们谈到犹太人以及他们中的首领如何因寻求报复而被指责。然而律法允许他们这样做：\\BibleQuote{以眼还眼，以牙还牙}（\\BibleRef{肋 24:20}）。确实如此，但并非要人们互相挖眼，而是要通过害怕遭受报复来遏制攻击的胆量，从而使他们既不对他人作恶，自己也不受他人之恶。因此，经上说\\BibleQuote{以眼还眼}，是为了束缚攻击者的双手，而不是让你放手攻击他；不仅要保护你的眼睛不受伤害，也要保护他的眼睛完好无损。\par

至于我所探究的——如果报复被允许，为什么那些实践它的人却被定罪？这究竟是什么意思？祂在这里说的是报复心；因为在一时冲动之下，祂允许受害者行动，如我所说，以遏制攻击者；但祂不再允许怀恨在心，因为此时的行为已不再是出于激情或沸腾的愤怒，而是出于预谋的恶意。现在，天主宽恕那些可能因受侮辱而冲动并急于报复的人。因此祂说：\\BibleQuote{以眼还眼}；然而又说：\\Quote{报复者的道路通向死亡}。如果在那允许以眼还眼的地方，报复者仍要受如此严重的惩罚，那么对于被命令甚至要忍受虐待的人，更要如何呢？因此，让我们不要报复，而要熄灭我们的怒火，使我们堪当来自天主的慈爱（因为基督说：\\BibleQuote{你们用什么量器量，也要用什么量器量给你们；你们用什么审判审判人，你们也必受什么审判}（\\BibleRef{玛 7:2}）），并使我们既逃脱今生的网罗，在即将来临的日子，也因我们的主耶稣基督的恩宠和慈爱，在祂面前获得宽恕。愿光荣、权能、尊崇归于祂，与圣父及圣神，从今世直到永远。阿们。\par

\SourceRecord{Translated by Gross Alexander.；Nicene and Post-Nicene Fathers, First Series；Vol. 13.；Edited by Philip Schaff.；Buffalo, NY: Christian Literature Publishing Co.,；1889.；<http://www.newadvent.org/fathers/230116.htm>.}

% structure: p0001 source=h1 target=section reason=page-title

\section{厄弗所书讲道辞 第十七篇}

厄弗所书 4:32 与 5:1-2\par

\BibleQuote{你们要彼此仁慈，有同情心，互相宽恕，如同天主在基督内宽恕了你们一样。因此，你们要效法天主，如同蒙爱的儿女；并在爱德中生活，如同基督也爱了你们，并为我们交出了自己，作为献与天主的祭品和牺牲，发出馨香之气。}\par

过去的事比将来的事有更大的说服力，显得更奇妙和更令人信服。因此保禄将他的劝诫建立在我们已经领受的恩惠上，因为这些事因基督的缘故更有力量。因为说：\BibleQuote{你们宽恕，你们也必蒙宽恕}\BibleRef{玛 6:14}，和\BibleQuote{你们若不宽恕，你们也决不蒙宽恕}\BibleRef{玛 6:15}——这是向有悟性的人、相信未来之事的人说的，很有分量；但保禄不仅以这些论据，也以已经为我们成就的事来呼吁良心。在第一种方式中，我们可以逃避惩罚，而在后一种方式中，我们可以分享一些正面的好处。你效法了基督。单是这一点就足以推荐美德，因为它就是\Quote{效法天主}。这是比另一个更高的原则：\BibleQuote{因为祂使太阳上升，光照恶人，也光照善人；降雨给义人，也给不义的人}\BibleRef{玛 5:45}。因为他不仅说我们是\Quote{效法天主}，而且是在我们领受这些恩惠的事上这样做。他要我们彼此怀有父亲般的慈心。因为此处的心是指慈爱和怜悯。既然我们作为人，难免会伤害人或受伤害，他就采取次佳的办法，想出了一个补救——即我们应彼此宽恕。然而这是无法相比的。因为如果你此刻宽恕了某人，他也会反过来宽恕你；而你对天主既没有给予，也没有宽恕什么。而且你宽恕的是同仆，而天主宽恕的是一个仆人、一个敌人和一个恨祂的人。\par

\BibleQuote{如同天主，}他说，\BibleQuote{也在基督内宽恕了你们。}\par

而且这含有一个深远的暗示。他的意思是：天主并非不冒任何风险或无需代价就宽恕了我们，而是牺牲了祂的圣子；因为祂为了宽恕你，牺牲了圣子；而你却常常即使看到宽恕既没有风险也无须代价，仍不给予宽恕。\par

\BibleQuote{因此，你们要效法天主，如同蒙爱的儿女；并在爱德中生活，如同基督也爱了你们，并为我们交出了自己，作为献与天主的祭品和牺牲，发出馨香之气。}\par

为了不让你以为这是出于必要，请听祂如何说：祂\BibleQuote{交出了自己}。正如你的主爱了你，你也要爱你的朋友。不，你也不能如此爱，但仍要尽力而为。哦，还有什么比这声音更有福呢！告诉我王权或其他你愿意的，都无法相比。宽恕别人，你就\Quote{效法天主}，你就被造成像天主一样。宽恕过犯比宽恕金钱债务更是我们的责任；因为如果你宽恕债务，你并没有\Quote{效法天主}；而你若宽恕过犯，你就是\Quote{效法天主}。然而，你怎么能说\Quote{我贫穷，不能宽恕它}，即债务，当你不能宽恕你能宽恕的东西，即过犯时？你难道认为在这件事上有任何损失吗？难道这不是财富，不是丰盈，不是丰富的储备吗？\par

看，还有另一个更高贵的激励：他说，\BibleQuote{如同蒙爱的儿女}。你还有另一个充分的理由效法祂，不仅因为你从祂手中领受了如此大的恩惠，而且因为你被称为祂的儿女。既然不是所有的儿女都效法他们的父亲，而是蒙爱的儿女才如此，因此他说，\BibleQuote{如同蒙爱的儿女}。\par

\BibleRef{弗 5:2} \BibleQuote{在爱德中生活。}\par

看，这是万事的根基！凡有爱德之处，就没有\Quote{愤怒、怒气、喧嚷、毁谤}，一切都消失了。因此他把要点放在最后。你如何成为儿女？因为你被宽恕了。基于你获得如此巨大特权的同一理由，也要宽恕你的邻人。告诉我，如果你在监狱中，要担当上万过犯，有人把你带进王宫；或者不谈这个，假设你发高烧，濒临死亡，有人用药使你康复，你会不珍视他超过一切，甚至珍视那药的名字吗？如果我们这样看重使我们受益的时机和场所，如同自己的灵魂，那么更应看重事物本身。因此，你要成为爱德的热爱者；因为你藉爱德而得救，藉爱德而成为儿子。如果你有能力拯救他人，难道你不会使用同样的良药，并向所有人提出这个建议：\Quote{宽恕，好使你得到宽恕}？\par

\BibleQuote{正如基督也，}他加上说，\BibleQuote{爱了你们。}\par

你只宽恕朋友，祂宽恕敌人。因此，来自我们主的恩惠要大得多。因为在我们身上，\Quote{如同}是如何保持的呢？显然，它将通过我们善待仇敌来实现。\par

\BibleQuote{祂为我们交出了自己，作为献与天主的祭品和牺牲，发出馨香之气。}\par

你看见了吗，为仇敌受苦是\Quote{馨香的芬芳}和\Quote{悦纳的祭献}吗？如果你死去，你将确实成为祭献。这就是\Quote{效法天主}。\par

\BibleRef{弗 5:3} \BibleQuote{至于邪淫，和一切不洁或贪吝之事，在你们中间连提都不提，正如圣徒相称。}\par

他谈论了苦毒的愤恨、怒气；现在他转向较小的恶：因为贪念是较小的恶，请听梅瑟如何在律法中也如此说：首先，\BibleQuote{不可杀人}\BibleRef{出 20:13}，这是怒气的工作；然后，\BibleQuote{不可奸淫}\BibleRef{出 20:14}，这是贪念的工作。因为正如\Quote{苦毒}、\Quote{喧嚷}、\Quote{一切的恶毒}、\Quote{毁谤}等，是易怒之人的工作；同样，\Quote{邪淫、不洁、贪婪}是纵欲者的工作；因为贪财和感官享乐源自同一情欲。但正如在前者中，他除去了\Quote{喧嚷}作为\Quote{愤怒}的载体；同样，现在他除去了\Quote{秽语}和\Quote{戏谑}作为贪念的载体；因为他继续说：\par

第4节：\BibleQuote{污秽、愚妄的言谈、或戏谑，这都不相宜；却当感恩。}\BibleRef{弗 5:4}\par

不可有戏谑之言，不可有猥亵之事，不论言语或行动，这样你便能扑灭火焰——\Quote{连提都不可提}，他说，\Quote{在你们中间}，也就是说，让它们在任何地方都不得出现。他写信给格林多人时也这样说：\BibleQuote{据说你们中间有邪淫的事}\BibleRef{格前 5:1}；也就是说，你们都要纯洁。因为言语是行为的途径。然后，为了不显得他是一位令人反感、刻板、并扼杀活泼性情的人，他继续补充理由，说：\Quote{这都不相宜}，这与我们无关——\Quote{却当感恩}。说一句戏谑之言有什么益处呢？你只是引人发笑而已。试问，鞋匠会忙碌于那些不属于或\Emph{相宜}于他行业的事务吗？或者他会购买那种工具吗？不，绝不会。因为不需要的东西，与我们无关。\par

伦理教训。不可有一句闲言，因为从闲言我们也会堕入秽语。现今不是轻浮嬉乐的时候，而是哀哭、患难、悲叹的时候：你竟要扮演小丑吗？哪一个摔跤手上场时忽略了与对手的搏斗，而口出戏言？魔鬼就在近旁，\BibleQuote{牠咆哮游行}\BibleRef{伯前 5:8}要捉拿你；牠调动一切，扭转一切来对付你的生命，妄图逼你离开退隐之处；牠咬牙切齿，怒吼咆哮；牠向你的救恩喷吐火焰；而你竟安坐那里，口出戏言，说\Quote{愚妄的话}，讲\Quote{不相宜的事}？那么，你将何等出色地胜过牠！我们是在嬉戏啊，亲爱的。你愿意知道圣人们的生活吗？请听保禄所说：\BibleQuote{有三年之久，我昼夜不停地流泪劝诫你们每一个人。}\BibleRef{宗 20:31}既然他为米利都和厄弗所的人如此热切，不是用愉悦的言词，而是用眼泪引入劝诫，那么对其他的人该说什么呢？再听他对格林多人所说的：\BibleQuote{我怀着极大的忧愁和痛苦的心，流了许多眼泪给你们写信。}\BibleRef{格后 2:4}又说：\BibleQuote{有谁软弱，我不软弱呢？有谁跌倒，我不焦急呢？}\BibleRef{格后 11:29}另在他处他也说，天天渴望，可谓渴求离开世界：\BibleQuote{我们在这帐棚里叹息}\BibleRef{格后 5:4}；而你竟又笑又玩？现在是战争时期，而你却在拨弄舞者的乐器？看看战场上人们的脸：阴暗紧锁的容貌，可畏而令人肃然起敬的额头。注视那严厉的眼睛，那急切、跳动、悸动的心，那集中、颤抖、极其焦虑的精神。在彼此对阵的营垒中，一切都是良好秩序，一切纪律严明，一切寂静无声。他们不说，我不是说无礼的话，就是连一个声音也不发出。如果那些有可见的仇敌、且绝不因言语而受伤的人，尚且保持如此巨大的沉默，那么你，你的战争、你战争的主要部分在于言语，你却让这一部分赤裸暴露吗？或者你不知道，正是在言语上我们最易中圈套？你自娱自乐，口出戏言，引人发笑，以为这是小事？有多少假誓、多少伤害、多少秽语，是从戏言中产生的！\Quote{不，}你会说，\Quote{戏谑并非如此。}然而，请听他如何排除一切戏谑。现在是战争与战斗的时候，是警醒与守护的时候，是武装与列阵的时候。嬉笑的时间在这里没有位置；因为那是属于世界的。听基督所说：\BibleQuote{世界要喜乐，但你们要忧愁。}\BibleRef{若 16:20}基督因你的灾祸而被钉十字架，你还笑吗？他因你的灾厄、那临到你身上的风暴而挨打、受尽极大痛苦，而你竟寻欢作乐？那么，你岂不是更激怒他吗？\par

但因此事在一些人看来似乎是无关紧要的，并且难以防范，让我们稍微讨论这一点，以向你们显示它是多么巨大的邪恶。因为这确实是魔鬼的作为，使我们忽视无关紧要的事。首先，即使它是无关紧要的，即使在这种情况下，当人知道最大的邪恶由它产生并增长，且它常常以淫乱告终时，也不应当忽视它。然而，它甚至不是无关紧要的，从此可见。那么让我们看看它从何产生。或者，让我们看看圣人应当是怎样的人：——温柔、温良、忧愁、悲伤、痛悔。那么，喜欢戏谑的人不是圣人。不，即使他是一个希腊人，这样的人也会被蔑视。这些事只允许那些在舞台上的人。哪里有污秽，哪里就有戏谑；哪里有不合时宜的笑声，哪里就有戏谑。请听先知所说：\BibleQuote{当以恐惧事奉主，且以战兢喜乐。} \BibleRef{咏 2:11} 戏谑使灵魂柔软而懒惰。它过分激动灵魂，常常充满暴力的行为，并制造战争。但还有什么？总之，你们不是来到人中间吗？那么\BibleQuote{要除去孩子气的事。} \BibleRef{格前 13:11} 为什么，你们不会允许自己的仆人在市场上说不敬的话；而你们，说自己是天主的仆人的人，却在公共广场上说俏皮话？如果那\Quote{清醒}的灵魂不被偷走，那还好；但一个松弛放荡的灵魂，谁能不把它带走？它将是自己的凶手，不需要魔鬼的诡计或攻击。\par

此外，为了理解这一点，也请看这名称本身。它意指多变的人、各种性情的人、不稳定的人、易弯的人、可以成为任何东西的人。但这与那些事奉盘石的人相去甚远。这样的性格很快转变和改变；因为他必须模仿动作和言语、笑声和步态，以及一切，是的，这样的人不得不编造笑话：因为他也需要这个。但这远离基督徒，去扮演小丑。再者，扮演小丑的人必然招致他随意嘲弄的对象们显著的仇恨，无论他们是在场，还是不在场而听说了这事。\par

如果这事是光荣的，为什么它留给江湖艺人？什么，你使自己成为江湖艺人，还不感到羞耻？为什么你们不允许你们的淑女这样做？岂不是你们把它视为不庄重的标志，而不是明智的标志？在好戏谑的灵魂中居住着巨大的邪恶，巨大的毁坏和荒凉。它的统一被打破，建筑朽坏，恐惧被放逐，敬畏消失。你有一条舌头，不是用来嘲笑别人，而是用来感谢天主。看看你们那些所谓的逗乐者，那些小丑。这些就是你们的戏谑者。我恳求你们，从你们的灵魂中驱逐这无恩宠的才能。这是寄生虫、江湖艺人、舞女、娼妓的事；远离一个慷慨的灵魂，远离一个高贵的灵魂，是的，甚至远离奴隶。如果有人失去了尊重，如果有人是卑贱的，那人也就是戏谑者。对许多人来说，这事甚至显得是一种美德，这实在令我们悲伤。正如情欲一点一点地驱使人走向淫乱，戏谑的癖好也一样。它似乎有一种优雅，但再没有比这更无优雅的了。因为听圣经所说：\BibleQuote{在雷声之前有闪电，在羞惭的人之前有恩宠。} 现在没有比戏谑者更无耻的了；所以他的口充满的不是恩宠，而是痛苦。让我们从我们的餐桌上驱逐这习俗。然而有些人甚至把它教给穷人！可怖啊！他们使受苦的人扮演小丑。为什么，这个祸害还能在哪里找到？它已经被带进教会本身。它已经抓住了圣经本身。我需要说什么来证明这邪恶的严重吗？我确实感到羞愧，但我仍然要说；因为我渴望显示这恶行已经发展到何等程度，使我不显得是在说笑，或是在对你们谈论一些小事；使甚至这样我也能引导你们脱离这错觉。没有人以为我在编造，但我要告诉你们我确实听到的事。有一个人曾与一个以知识自傲的人在一起（现在我知道我会惹人发笑，但我仍然要说）；当盘子放在他面前时，他说：\Quote{孩子们，拿去吃，免得你们的肚子发怒！} 还有，另一些人说：\Quote{祸哉，玛门，以及没有你们的人；} 戏谑引入了许多这样的恶行；如他们所说：\Quote{现在没有诞生。} 我说这些是为了显示这卑劣脾气的严重；因为这些都是一个缺乏所有敬畏的灵魂的表达。这些事难道不足以招致雷霆吗？人们还能找到许多由这些人说出的其他类似的话。\par

因此，我恳求你们，让我们彻底摒弃这一恶习，口出合宜之言。不要让圣洁的嘴唇说出可耻卑微之人的话语。\BibleQuote{因为正义与不义有什么相通？光明与黑暗有什么联系？} \BibleRef{格后 6:14} 我们若远离一切污秽之事，而得获所应许的福分，那便是有福的；决不可在我们身后拖着这样一串污秽，并以此玷污心灵的纯洁。因为喜作诙谐者，不久便会成为谩骂者；而谩骂者又会为自己招致无数其他的灾祸。当我们驯服了灵魂的这两种官能——愤怒与欲望（参\WorkTitle{斐德若篇}，第二十五、三十四章），使它们如同驯服的马匹置于理性的轭下，然后让理智如御车者驾驭它们，使我们能\BibleQuote{获得天主从高天召唤的奖赏}\BibleRef{斐 3:14}；愿天主恩赐我们众人皆能达到，藉着我们的主耶稣基督，与祂及圣神，一同归于圣父，光荣、威能与尊威，自今而永，及于万世。阿们。\par

\SourceRecord{Translated by Gross Alexander.；Nicene and Post-Nicene Fathers, First Series；Vol. 13.；Edited by Philip Schaff.；Buffalo, NY: Christian Literature Publishing Co.,；1889.；<http://www.newadvent.org/fathers/230117.htm>.}

% structure: p0001 source=h1 target=section reason=page-title

\section{厄弗所书讲道第十八篇}

\BibleRef{弗 5:5-6}\par

\BibleQuote{因为你们应该清楚知道：凡是淫乱的，不洁的，和贪婪的——即崇拜偶像的——在基督和天主的国内，都得不到产业。不要让人以空言欺骗你们，因为就是为了这些事，天主的义愤才降到那些悖逆之子身上。}\BibleRef{弗 5:5-6}\par

很可能在我们的先祖时代，也有一些人\BibleQuote{使人民的手发软}\BibleRef{耶 38:4}，并实行厄则克耳所提及的事——更确切地说，是行了那些假先知的事，他们\BibleQuote{为了几把麦子，在人民中间亵渎天主}\BibleRef{则 13:19}；这件事，依我看，即使在今天也有人在做。例如，当我们说凡称自己兄弟为傻子的，必被投入地狱之火时，有些人就说：\Quote{什么？称自己兄弟为傻子的就要被投入地狱之火？不可能！}他们说。同样，当我们说\BibleQuote{贪婪的人是崇拜偶像者}时，他们又在这事上打折扣，说这种说法是夸张的。他们就是用这种方式低估和曲解所有的诫命。蒙福的保禄在写\WorkTitle{厄弗所书}时，正是针对这些事说了这样的话：\BibleQuote{因为你们知道，凡是淫乱的、不洁的、贪婪的人——即崇拜偶像者——在基督和天主的国内都得不到产业}，并加上：\BibleQuote{不要让人以空言欺骗你们。}\BibleRef{弗 5:5-6} 所谓\Quote{空言}，就是那些暂时令人愉悦，但事实上毫无根据的话；因为整个情况是一场欺骗。\par

\BibleQuote{因为就是为了这些事，天主的义愤才降到那些悖逆之子身上。}\BibleRef{弗 5:6}\par

他所说的\Quote{因为淫乱}，意思是，\Quote{因为贪婪}，\Quote{因为不洁}，或者既因为这些事，也因为\Quote{诡诈}，因为有欺骗者。\Quote{悖逆之子}：他这样称呼那些完全悖逆的人，即那些不服从祂的人。\par

第7—8节：\BibleQuote{所以，你们不要与他们同流合污。因为你们从前曾是黑暗，但现在在主内却是光明。}\BibleRef{弗 5:7-8}\par

请注意他多么明智地催促他们前进：首先，从基督的思想出发，要你们彼此相爱，不伤害任何人；另一方面，从惩罚和地狱之火的思想出发。\BibleQuote{你们从前曾是黑暗，但现在在主内却是光明。} 这也是他在\WorkTitle{罗马书}中所说的：\BibleQuote{那时你们在现今所羞耻的事上，有什么果实呢？}\BibleRef{罗 6:21}，并提醒他们从前的邪恶。就是说，想想你们从前是什么样，现在成了什么样，不要回到从前的邪恶中，也不要\BibleQuote{轻慢天主的恩宠}\BibleRef{希 10:29}。\par

\BibleQuote{你们从前曾是黑暗，但现在在主内却是光明！}\BibleRef{弗 5:8}\par

他说，这不是由于你们自己的美德，而是通过天主的恩宠，这福分临到了你们。也就是说，你们也曾一度该受同样的惩罚，但现在不再如此了。因此，要\BibleQuote{行走}如同\BibleQuote{光明之子}。然而，\BibleQuote{光明之子}的意思，他后来会解释。\par

第9—10节：\BibleQuote{因为光明的果实在于各种良善、正义和真理，证实什么是主所喜悦的。}\BibleRef{弗 5:9-10}\par

他说：\Quote{在各种良善}，这是针对恼怒和怨恨的；\Quote{和正义}，这是针对贪婪的；\Quote{和真理}，这是针对虚假的享乐：不是前者，他说，即我所提到的那些，而是它们的对立面。\Quote{在一切}上；即圣神的效果应当在每件事上都显明出来。\Quote{证实什么是主所喜悦的}；这样，那些事就是幼稚和不完美心智的标志。\par

第11—13节：\BibleQuote{不要参与黑暗无果的作为，反而要谴责它们。因为他们暗中做的事，连提起都是可耻的。但一切事情，当受到谴责时，就被光明显明。}\BibleRef{弗 5:11-13}\par

他曾说：\BibleQuote{你们是光明。} 光明通过揭露黑暗中所发生的事来谴责。所以，他说，如果你们有德行、显明，恶人就不能隐藏。因为正如点着蜡烛，一切都显露出来，小偷无法进入；同样，如果你们的光照耀，恶人被发现，就会遭捕获。因此，揭露他们是我们的责任。那么，主如何说：\BibleQuote{不要判断，免得你们受判断}？\BibleRef{玛 7:1} 保禄没有说\Quote{判断}，他说的是\Quote{谴责}，即纠正。而\BibleQuote{不要判断，免得你们受判断}这句话，他是针对很小的过失说的。实际上，他还加上了：\BibleQuote{为什么你看见你兄弟眼中的木屑，却不理会自己眼中的大梁呢？} 但保禄所说的是这样：如同一个伤口，只要它深藏不露，外表被掩盖，在皮下蔓延，就不会被注意；同样，罪过只要被隐藏，仿佛在黑暗中，就会在充分安全感下被大胆犯下；但一旦\Quote{被显明}，就成为\Quote{光明}；当然不是罪过本身（这怎么可能？），而是罪人。因为当他被带到光明中，当他受到劝诫，当他悔改，当他获得宽恕，你难道没有清除他所有的黑暗吗？你难道没有治愈他的伤口吗？你难道没有使他的不结果实转为结果实吗？要么这是他的意思，要么就是我上面所说的：你的生命\Quote{被显明，就是光明}。因为没有人会隐藏无可指责的生命；而那些隐藏的事，是被覆盖的黑暗所隐藏的。\par

第14节：\BibleQuote{因此他说：你这睡眠的，醒来吧！从死者中起来，基督必光照你。}\BibleRef{弗 5:14}\par

他藉\Quote{睡覺者}和\Quote{死者}指的是那在罪中的人；因為他如死者一樣散發惡臭，又如睡覺者一樣不活動，而且像他一樣看不見什麼，只是做夢，產生幻覺和錯覺。有些人讀作\Quote{你將觸摸基督}，但另一些人讀作\Quote{基督將光照你}；而後者更為可能。離開罪，你就能夠看見基督。\BibleQuote{因為凡作惡的，恨光，不來就光。}（\BibleRef{若 3:20}）所以不作惡的人就來就光。\par

他這話不僅是針對不信者說的，因為許多信友，與不信者無異，甚至有些人更甚，固執於邪惡。因此，也必須向他們呼喊：\BibleQuote{睡覺的人，醒來吧，從死者中起來，基督必要光照你。}向這些人也適宜說這話：\BibleQuote{天主不是死人的天主，而是活人的天主。}（\BibleRef{瑪 22:32}）如果祂不是死人的天主，那麼就讓我們活着。\par

現在有些人說\Quote{貪婪的人是偶像崇拜者}這句話是誇張之詞。然而，這陳述並非誇張，而是真理。如何以及以何種方式呢？因為貪婪的人背離天主，正如偶像崇拜者一樣。為了避免你認為這只是一個空泛的主張，基督有話說：\BibleQuote{你們不能事奉天主而又事奉金錢。}（\BibleRef{瑪 6:24}）那麼如果不可能同時事奉天主和金錢，那些事奉金錢的人就把自己排除在天主的服務之外；而他們否認了祂的主權，事奉沒有生命的金子，很明顯他們是偶像崇拜者。\Quote{但我從未塑造過偶像，}有人會說，\Quote{也沒有設立祭壇，沒有宰殺羊，沒有奠酒；不，我進了教堂，向天主的獨生子舉起雙手；我領受奧蹟，參與祈禱，以及基督徒應盡的一切。那麼我怎麼會是偶像崇拜者呢？}他會說。是的，這正是最令人驚奇的事：當你經歷過，並\BibleQuote{嘗過天主的仁慈，且看過主是何等美善}（\BibleRef{詠 33:8}）之後，你竟會離棄那仁慈者，而轉向一個殘酷的暴君，並且假裝事奉祂，而實際上卻使自己順服於貪婪那艱難而惱人的軛下。你尚未告訴我你自己所盡的本分，只提及了主人所賜的恩惠。因為請告訴我，我懇求你，我們如何判斷一個士兵？是當他盡忠職守，保衛君王，由君王供養，被稱為君王的親兵時，還是當他只顧自己的私利時？假裝與君王同在，關心他的利益，同時卻在推進敵人的事業，我們認為這比脫離君王的服務而投靠敵人更為惡劣。現在你正在侮辱天主，正如一個偶像崇拜者所行的，不僅僅是用你自己的口，而是用被你傷害的成千上萬人的口。但你會說：\Quote{他不是偶像崇拜者。}然而，每當他們說：\Quote{哦！那個基督徒，那個貪婪的傢伙！}時，那麼他不僅用自己的行為進行侮辱，還常常迫使那些被他傷害的人也說這樣的話；如果他們不說，這就要歸功於他們的敬畏之心。\par

我們難道看不見事實如此嗎？什麼是偶像崇拜者？難道他不是也常常崇拜情慾，而不能控制情慾嗎？我的意思是，例如，當我們說外邦人崇拜偶像時，他會說：\Quote{不，那是維納斯，或是瑪爾斯。}如果我們說：這位維納斯是誰？他們之中較穩重的人會說：是快樂。或者瑪爾斯是什麼？是憤怒。同樣，你也崇拜金錢。如果我們說：金錢是誰？是貪婪，而這正是你所崇拜的。\Quote{我並不崇拜它，}你會說。為什麼不？因為你不俯伏叩拜？不，恰恰相反，你更是在行為和實踐上崇拜它；因為這是更高層次的崇拜。為了讓你明白這一點，看看天主的情況：誰更真正崇拜祂？是那些僅僅在祈禱中站立的人，還是那些實行祂旨意的人？很明顯，是後者。金錢的崇拜者也是如此：實行金錢旨意的人，才是真正崇拜金錢的人。然而，那些崇拜情慾的人常常能不受情慾的束縛。人們可能會看到一個崇拜瑪爾斯的人常常控制自己的憤怒。但對你來說卻並非如此；你使自己成為情慾的奴隸。\par

是的，但你不宰杀羊？不，你宰杀人，合理的灵魂，有些因饥荒，有些因亵渎。没有什么比这样的祭献更疯狂的了。谁曾见过灵魂被祭献？贪婪的祭台是多么可咒！当你从这里经过这座偶像的祭台时，你会看到它散发着公牛和山羊的血腥；但当你经过贪婪的祭台时，你会看到它散发着人类鲜血的骇人气息。在今世站在它面前，你看到的将不是鸟翅的燃烧，不是蒸汽，不是烟气升腾，而是人的身体在消亡。因为有些人投身悬崖，有些人系上绞索，有些人用匕首刺穿喉咙。你看到了残忍而不人道的祭献吗？你想看到比这些更骇人的吗？那么我将不再向你展示人的身体，而是在另一个世界中被屠杀的人的灵魂。是的，因为灵魂也可能遭受灵魂特有的屠杀；正如有身体的死亡，也有灵魂的死亡。先知说：\BibleQuote{那犯罪灵魂，它必死亡。} \BibleRef{则 18:4}然而，灵魂的死亡并不像身体的死亡；它远更骇人。因为这种身体的死亡，使灵魂与身体彼此分离，使灵魂脱离许多焦虑和劳苦，并将身体送到一个显明的居所；然后当身体在时间中分解和朽坏后，它又将在不朽中重新聚集，并收回它自己的灵魂。我们看到身体的死亡就是这样。但灵魂的死亡是可怕而恐怖的。因为这种死亡，当分离发生时，并不像身体那样让它过去，而是再次将它束缚于一个不朽的身体，并把它交付给不灭之火。这就是灵魂的死亡。因此，既然有灵魂的死亡，也就有灵魂的屠杀。身体的屠杀是什么？就是变成尸体，被剥夺来自灵魂的能量。灵魂的屠杀是什么？就是它也变成一具尸体。灵魂如何变成尸体？因为正如身体在灵魂离开它并使其失去自身生命能量时变成尸体，同样，灵魂在圣神离开它并使其失去属神能量时也变成尸体。\par

这些多半是在贪婪的祭台上所进行的屠杀。他们不满足，不止于人的血；不，贪婪的祭台并不满足，除非它同时也祭献灵魂本身，除非它同时接纳祭献者和被祭献者的灵魂。因为那祭献的人必须先被祭献，然后才祭献；死人祭献那仍然活着的人。因为当他出言亵渎，当他辱骂，当他恼怒时，这些岂不正是灵魂的无数不治之伤吗？\par

你看到这说法并非夸张。你愿意再听另一个论证，教你贪婪如何是偶像崇拜，并且比偶像崇拜更骇人吗？拜偶像者敬拜天主的受造物（经上说：\BibleQuote{他们敬拜并事奉受造物，而非造物主。} \BibleRef{罗 1:25}），但你在敬拜你自己的受造物。因为天主并未创造贪婪，而是你不知餍足的欲望发明了它。且看这疯狂与愚昧。那些拜偶像的人，也尊敬他们所拜的偶像；若有人以不敬或嘲笑的方式谈论它们，他们就会起来辩护；而你，却仿佛醉了一般，在敬拜一个对象，它非但免于指责，反而充满不虔敬。所以，你甚至比他们更邪恶。你永远无法以此为借口说它不是恶。如果他们已处于极端无可推诿，那么你更远甚于此，你永远在指责贪婪，辱骂那些投身其中的人，却仍然事奉并服从它。\par

若蒙允许，我们将查考偶像崇拜起源于何处。一位智者\BibleRef{智 14:16}告诉我们，有一位富人因儿子过早去世而悲伤，且无法从哀恸中得到安慰，便以这种方式慰藉自己的情感：他为死者造了一尊无生命的像，不断凝视，似乎通过这像仍保有已逝之人；而某些谄媚者，\Quote{他们的天主是自己的肚腹}\BibleRef{斐 3:19}，为向那富人表示尊敬而恭敬那像，便将这习俗演变成了偶像崇拜。因此，偶像崇拜源于灵魂的软弱、无意义的习俗和放纵。但贪吝却非如此：它确实出于灵魂的软弱，然而是一种更严重的软弱。并非有人失去了儿子，也非寻求悲伤的慰藉，也非被谄媚者引诱。那么它是如何产生的呢？我将告诉你们。加音由于贪吝而欺骗了天主；本应献给天主的，他留给了自己；本应留给自己的，他却献给了天主；于是这邪恶从天主那里开始了。因为如果我们属于天主，那么我们的财产初果更应属于天主。再者，人对女人的强烈情欲也源于贪吝。\BibleQuote{他们看见人的女儿}\BibleRef{创 6:2}，便纵身沉溺于情欲。从此又发展到贪财；因为希望拥有比邻舍更多世上财物的愿望，除了源于\Quote{因爱冷淡}之外，别无他源。希望占有多于自己份内的份额，除了源于鲁莽、厌世和对他人的傲慢之外，别无他源。请看大地，其幅员何其广阔？穹苍和天的宽广远超出我们的需要？正是为了终止你的贪吝，天主才如此宽阔地延展了造化的界限。而你竟仍如此贪求无厌？你听到贪吝就是偶像崇拜，竟不因此而战栗吗？你想继承大地吗？那么你在天上就没有基业。你急切要将产业留给别人，好使你自己丧失产业吗？告诉我，如果有人给你能力拥有一切，你难道不愿意吗？如今你若有心，这就在你权下。然而有些人说，他们为将遗产转给别人而悲伤，宁愿自己将它消耗掉，也不愿看别人成为它的主人。我也不开脱你们的这种软弱；因为这正是灵魂软弱的特征。不过，至少要做到这样：在你的遗嘱中，立基督为你的继承人。你确实应当在有生之年这样做，因为这才显示出正确的态度。尽管如此，无论如何，即使只是出于必需，也要稍微慷慨一些。基督确实以此目的命令我们施舍穷人：为了使我们在有生之年变得明智，诱导我们轻视金钱，教导我们藐视世俗的事物。吝惜将财物赠给这人那人，当人死后不再拥有它时，如你所想，这并不是轻视金钱。那时你不再给出你自己的东西，而是出于绝对的必要：感谢死亡，而不是感谢你。这不是爱的行为，这是你的损失。不过，即使如此，也要这样做；至少那时你要放弃你的情欲。\par

道德教训。请思考你犯了多少抢劫的行为，多少贪吝的行为。要四倍偿还。这样向天主申诉你的案情。然而，有些人已陷入如此疯狂和盲目，甚至那时仍不明白自己的责任；反而在所有情况下继续行事，仿佛特意加重天主对自己的审判。这就是为什么我们蒙福的宗徒写道并说：\Quote{作为光明之子而行。}如今贪吝之人比任何人都生活在黑暗中，并在周围一切事物上散布极大的黑暗。\par

他又说：\Quote{不要与黑暗无益的功业有分，反倒要责斥它们；因为他们在暗中所行的事，连提起来也是可耻的；但一切事一经责斥，就由光明显露出来。}我恳求你们，所有不愿意无故被恨、而愿意被爱的人，请听。\Quote{何必被人恨呢？}有人说。一个人抢劫，你不责斥他，却怕他的仇恨？但这并非无故被恨。但你公正地谴责他，却仍害怕仇恨吗？谴责你的兄弟，为你对基督所负的爱德，为你对兄弟所负的爱德而招致敌意。在他走向毁灭之坑的路上拦住他。因为请他入席、以客气的话语、问候和款待对待他，这些并非友谊的显著证据。不，我所提的那些才是我们必须赐予朋友的恩惠，为拯救他们的灵魂免于天主的义怒。当我们看见他们躺在邪恶的炉中，让我们扶起他们。但有人说：\Quote{没有用，他不可救药了。}然而，你要尽你的本分，这样你就在天主面前为自己辩护了。不要掩藏你的才能。你之所以有口才，有口舌，正是为了纠正你的邻人。只有不能说话、没有理性的动物才不关心邻人，不顾及他人。当你称天主为\Quote{父}，称你的邻人为\Quote{兄弟}，看见他犯下无数恶行时，你难道宁愿要他的好感而不顾他的福祉吗？不，我恳求你们，不要这样。再没有比从不忽略我们兄弟的罪过更真实的友谊证据了。你看见他们彼此为仇吗？使他们和好。你看见他们犯贪吝之罪吗？制止他们。你看见他们受冤屈吗？为他们辩护。那主要的益处不是加给他们，而是加给你自己。我们成为朋友正是为此，为能彼此有益。一个人听朋友说话，与听任何其他偶然之人说话，其心态是不同的。对偶然之人，他或许会怀疑，同样对待老师也是如此，但对朋友则不然。\par

\Quote{因為，}他說，\BibleQuote{那些他們在暗中所做的事，連提起也是可恥的；但一切事，一經責備，就由光顯明出來。}他這裡想說的是什麼？是這個意思。就是說，在這世界上，有些罪是暗中所犯的，也有些是公開的；但在另一個世界卻不是這樣。現在沒有一個人不對自己的一些罪有所自覺。這就是為什麼他說：\BibleQuote{但一切事，一經責備，就由光顯明出來。}那麼，這又是指拜偶像的事嗎？不是；這裡的論述是關乎我們的生活和我們的罪。\BibleQuote{因為凡顯明出來的，}他說，\BibleQuote{就是光。}\par

因此，我懇求你們，永不厭惡責備，也不因受責備而不悅。因為事情在黑暗中進行時，確實進行得更有把握；但當它有很多見證人目睹所發生的事時，它就被揭露了。所以，我們務要盡一切努力驅散我們弟兄們心中的死氣，驅散黑暗，並吸引那\BibleQuote{\OriginalTerm{義德的太陽}{Sun of righteousness}}到我們這裡來。因為如果有許多發光的光體，美德的道路對他們自己來說就會是容易的，而那些在黑暗中的也會更容易被發現，因為光被舉起，驅散了黑暗。反之，如果情況相反，就怕罪惡與黑暗的濃霧壓倒光，驅散它的透明度，那些發光的光體自身反而會熄滅。讓我們因此樂於彼此有益，好使我們全體一致，將讚頌與光榮歸於仁慈的天主，藉著獨生聖子的恩寵與仁慈，願光榮、能力、尊崇歸於祂，並與聖父及聖神，從現在直到永遠。阿們。\par

\SourceRecord{Translated by Gross Alexander.；Nicene and Post-Nicene Fathers, First Series；Vol. 13.；Edited by Philip Schaff.；Buffalo, NY: Christian Literature Publishing Co.,；1889.；<http://www.newadvent.org/fathers/230118.htm>.}

% structure: p0001 source=h1 target=section reason=page-title

\section{厄弗所书释义 第十九篇}

厄弗所书 5:15, 16, 17\par

\BibleQuote{所以，你们应谨慎行事，不要像愚昧人，却要像智慧人；要赎回光阴，因为日子是邪恶的。因此，不要作糊涂人，但要明白什么是主的旨意。}\par

他仍在清除苦毒的根，仍在割断愤怒的根基。他说什么？\BibleQuote{你们应谨慎行事。} \BibleQuote{他们是羊在狼群中，} 并且他嘱咐他们也要 \BibleQuote{像鸽子。} 因为他说：\BibleQuote{你们要像鸽子那样温驯。} \BibleRef{玛 10:16} 既然他们既处在狼群中，又奉命不可自卫，反而要忍受恶事，他们就需要这劝告。固然前一个足以使他们更强壮；但现在加上这两个，试想这多么增强。这里也留意他如何谨慎地保护他们，说：\BibleQuote{注意你们的行事。} 整个城市都与他们为敌；是的，这场战争也蔓延到家中。他们分裂了：父亲反对儿子，儿子反对父亲，母亲反对女儿，女儿反对母亲。那么怎么办？这些分裂从何而来？他们听见基督说：\BibleQuote{谁爱父亲或母亲超过我，不配是我的门徒。} \BibleRef{玛 10:37} 因此，免得他们以为他无缘无故引起战争和争斗（因为他们若报复，很可能产生许多愤怒），为防止这事，他说：\BibleQuote{谨慎行事。} 这就是说：\Quote{除了福音的宣讲，不要在其它任何事上给你们招来的仇恨留把柄。} 让这成为仇恨的唯一理由。不要让任何人再有别的指责；但当这无损于宣讲、不妨碍虔敬时，要表现出一切恭敬和服从。因为\WorkTitle{圣经}上说：\BibleQuote{凡人所当得的，就给他们：当得粮的，给他纳粮；当得税的，给他上税。} \BibleRef{罗 13:7} 当世上其余的人看见我们忍耐时，他们便会羞愧。\par

\BibleQuote{不要像愚昧人，却要像智慧人，赎回光阴。}\par

他给出这劝告，并非要你们变得狡猾和多变。他的意思是这样：这时间不属于你们。目前你们是客旅、寄居的、外方人和陌生人；不要寻求尊荣，不要寻求光荣，不要寻求权势，也不要寻求报复；忍受一切，这样 \BibleQuote{赎回光阴}；放弃许多东西，凡是他们可能要求的。试想，有一个人有一所华丽的房子，有人闯入要杀他，他付出一大笔钱从而救了自己。那时我们会说他赎回了自己。同样，你们拥有一所大房子，即保存着真正的信德。他们会来夺走一切。他们要求什么，就给他们什么，只保全主要的东西——即信德。\par

\BibleQuote{因为日子是邪恶的。}\par

日子的邪恶是什么？日子的邪恶应当属于日子。身体的邪恶是什么？疾病。灵魂的邪恶是什么？邪恶。水的邪恶是什么？苦味。每一事物的邪恶，是与该事物受邪恶影响的本性相关。如果日子有邪恶，它应当属于日子，属于时辰，属于日光。同样，基督说：\BibleQuote{一天的邪恶一天当就够了。} \BibleRef{玛 6:34} 从这话我们将理解另一处。那么他称 \Quote{日子邪恶} 是什么意思？\Quote{时间} 邪恶是什么意思？这不是事物的本质，不是受造物本身，而是在其中发生的事。正如我们习惯说：\Quote{我过了一个不愉快和糟糕的一天。} 但如何会不愉快呢？除非是由于其中发生的情况。其中发生的事，善的来自天主，恶的来自坏人。因此，时间中发生的恶是人所创造的，所以时间被称为邪恶。这样我们也称时间为邪恶。\par

第 17, 18 节。他接着说：\BibleQuote{因此，不要作糊涂人，但要明白什么是主的旨意；不要醉酒，酒能使人放荡。}\par

因为在这方面的无节制使人易怒、粗暴、性急、暴躁和野蛮。酒被赐予我们是为了喜悦，而不是为了醉酒。如今，一个人不醉酒竟显得没有男子气概和可鄙。那么还有什么救恩的希望？什么？告诉我，不醉酒是可鄙的，而醉酒本该是天下最可鄙的事？因为连一个普通人远离醉酒也是完全正确的；更何况是一个战士，一个生活在刀剑、流血和杀戮中的人；更不用说战士，他的脾气因其他原因——权力、权威、常处于计谋和战斗中——而变得更加尖锐。你想知道酒在何处有益吗？请听\WorkTitle{圣经}说：\BibleQuote{当将烈酒给将亡的人喝，将酒给心里苦毒的人喝。} \BibleRef{箴 31:6} 这是公道的，因为酒能减轻苛刻和忧郁，驱散眉头的阴云。\BibleQuote{酒能悦人心，} \BibleRef{咏 103:15} \BibleBook{}{圣咏}作者说。那么酒如何导致醉酒呢？因为同一事物不可能产生相反的效果。显然醉酒并非来自酒，而是来自无节制。酒赐予我们只是为了身体的健康；但这一目的也因过度使用而受阻。此外，听我们蒙福的宗徒给弟茂德写的：\BibleQuote{为你的胃病和屡次患病，稍微用点酒。} \BibleRef{弟前 5:23}\par

这就是天主以中庸比例塑造我们身体的原因，并使之满足于少量，由此立即教导我们：祂使我们适应于另一种生命。祂原本愿意从起初就将那生命赐予我们；但因我们自甘不配，祂就延迟了；在延迟期间，祂甚至也不容许我们过度放纵；因为一小杯酒和一个饼就足以满足人的饥饿。祂使人——所有受造动物的主人——形成所需食物比它们少，身体也比它们小；这无非向我们宣告：我们正快速走向另一种生命。\Quote{不要醉酒，}他说，\Quote{于酒，其中含有放荡}；因为这酒不拯救，反而毁灭；不仅毁灭身体，也毁灭灵魂。\par

第18、19、20、21节。\BibleQuote{但要以圣神充满；以圣咏、圣诗、属神歌曲彼此对谈，以你们的心向主歌唱、奏乐；常常为一切事物，因我们主耶稣基督之名，感谢天主父；以敬畏基督之心彼此服从。}\BibleRef{弗 5:18}\par

他说，你希望快乐吗？你希望利用这一天吗？我给你属神的饮料；因为醉酒甚至切断我们舌头的清晰声音；它使我们口齿不清、结结巴巴，扭曲眼睛，以及整个身体。学习唱圣咏，你将看到这工作的愉悦。因为唱圣咏的人充满圣神，如同唱撒殚歌曲的人充满不洁之神。\par

所谓\Quote{以你们的心向主}是什么意思？意思是，以紧密的注意和理解。因为那些不密切注意的人，只是唱歌，发出话语，而他们的心却游荡别处。\par

\Quote{常常，}他说，\Quote{为一切事物，因我们主耶稣基督之名，感谢天主父；以敬畏基督之心彼此服从。}\par

\BibleQuote{要让你们的请求藉著感恩被天主知晓}\BibleRef{斐 4:6}；因为没有什么比一个人感恩更令天主喜悦的了。但我们将最能够感恩天主，通过将我们的灵魂从前述的事中抽离，并通过用他告诉我们的方法彻底洁净它们。\par

他说：\Quote{但要充满，}\Quote{以圣神。}\par

那么，这位圣神就在我们内吗？是的，的确在我们内。因为当我们从灵魂中驱散了谎言、苦毒、淫乱、不洁、贪婪，当我们成为良善、心肠柔和、彼此宽恕，当没有戏谑，当我们使自己配得之时，还有什么能阻止圣神来临并光照我们呢？祂不仅会来到我们这里，还要充满我们的心；当我们内部点燃了如此大的光时，那么德行的道路将不再是难以达到的，而是变得容易和简单。\par

\Quote{常常感恩，}他说，\Quote{为一切事物。}\par

那么，我们是否要为所遭遇的一切感恩？是的，即使是疾病，即使是贫困。因为旧约中有一位智者曾这样劝告说：\BibleQuote{凡临到你的，你都要欣然接受；当你被贬抑时，要忍耐} \BibleRef{德 2:4}；在新约中更应如此。是的，即使你不明白这话，也要感恩。因为这就是感恩。但如果你在舒适和富足、成功和顺遂时感恩，那没有什么伟大，没有什么奇妙。所要求的是，当人在磨难、痛苦、沮丧时感恩。不要说其他话，只要说：\Quote{主，我感谢祢。} 我为何要提今世的磨难呢？我们甚至应当为地狱本身、为来世的折磨和惩罚而感谢天主。因为这对那些留意的人来说，确实是有益的，当对地狱的恐惧像缰绳一样放在我们心上。因此，让我们不仅为我们所见的祝福感恩，也为我们所未见的、为那些我们违背自己意愿而接受的祝福感恩。因为祂赐给我们许多恩惠，是我们不想要的，也不知道的。如果你们不信我，我立刻向你们讲明。请想想，我恳求你们，那些不虔敬的、不信的外邦人不是把一切归于太阳和他们的偶像吗？那么怎样呢？难道祂不也赐福给他们吗？使他们有生命、健康、子女等，不都是祂眷顾的工作吗？还有那些被称为马西翁派和摩尼派的人，他们岂不也亵渎祂？那么怎样呢？难道祂不每日赐福给他们吗？如果祂赐福给那些不认识祂的人，何况我们呢？因为天主特有的工作，若非对全人类行善，借着惩罚和享受，那又是什么呢？因此，我们不可只在顺遂时感恩，因为这没有什么伟大。魔鬼也深知这一点，因此他说：\BibleQuote{约伯岂是无缘无故敬畏天主？祢岂不是四面围护他和他的家，并他所有的一切？祢若伸手打击他所有的一切，他必当面诅咒祢！} \BibleRef{约 1:10} 然而，那被诅咒者并未得逞；愿天主也使他从我们身上不得逞；但每当我们处于贫困、疾病或灾难中时，让我们加倍感恩；我所说的感恩，不是指言语，也不是指口舌，而是指行为和工作，指心思和心灵。让我们全心感谢祂。因为祂爱我们胜过我们的父母；善恶之间的差别有多大，天主的爱与父亲的爱之间的差别就有多大。这不是我的话，而是爱我们的基督自己的话。请听祂自己说的话：\BibleQuote{你们中间有哪一个人，儿子向他求饼，反而给他石头？你们纵然不善，尚且知道把好东西给你们的儿女，何况你们的天父，岂不更把好东西赐给求祂的人？} \BibleRef{玛 7:9} 此外，请听祂在别处说的话：\BibleQuote{妇女岂能忘掉她吃奶的婴孩，不怜恤她亲生的儿子？她们或许能忘掉，但我不会忘掉你，这是上主说的。} \BibleRef{依 49:15} 因为祂若不爱我们，为何创造我们？祂有需要吗？我们为祂提供什么服事和职务？祂需要我们能提供什么？听先知说：\BibleQuote{我曾对上主说：祢是我的主，我没有祢以外的幸福。} \BibleRef{咏 15:2}\par

然而，忘恩负义、麻木不仁的人说，这才配得上天主的良善，即人人平等。告诉我，忘恩负义的凡人，你否认哪些事物是出于天主的良善，你所指的平等又是什么？你说：\Quote{某人自幼跛足；另一人疯狂，被附身；另一人到了极老年纪，却终生贫困；另一人受最痛苦的疾病：这些是眷顾的工作吗？一人耳聋，另一人哑巴，另一人贫穷，而另一人不虔敬，是的，极不虔敬，充满万千恶习，却享受财富，拥有妾侍和食客，拥有一座华宅，过着闲散的生活。} 他们列举了许多此类例子，并长篇累牍地抱怨天主的眷顾。\par

那麼，我們要對他們說什麼呢？倘若他們是希臘人，告訴我們宇宙由某一位管理，我們就要以同樣的話回覆他們：\Quote{那麼，事物都沒有天主眷顧嗎？你們如何敬拜神祇、崇拜精靈和英雄呢？因為若有眷顧，必有某一位統管萬有。}但若有人，無論是基督徒或異教徒，對此不耐煩而動搖，我們要對他們說什麼呢？請問，如此多美好事物，告訴我，能自然產生嗎？每日的光明？萬物中存在的美麗秩序和預先顧慮？星辰的迷宮般舞蹈？晝夜均等的運行？植物、動物和人類中本性的層級？告訴我，誰安排這些？若無一位統管者，萬物自行結合，那麼誰使這穹蒼運轉——如此美麗、如此廣大，我指的是天空——並將它安置在大地上，甚至還在諸水之上？誰賜予豐收的季節？誰在種子和植物中植入如此巨大的力量？因為偶然之事必然混亂；而有秩序則意味著設計。告訴我，我們周圍偶然之事，哪一件不是充滿巨大的混亂、巨大的喧囂與困惑？我不僅指偶然之事，也指那些暗示某種行動者、但行動者無技藝之事。例如，有木材、石頭，還有石灰；讓一個不諳建築的人拿起它們開始建造，努力幹活；他豈不是會毀掉一切？又如，一艘船沒有舵手，卻裝有一艘船應有的一切，卻沒有造船匠；我不是說它裝備不全或未完成，而是說即使裝備齊全，它也無法航行。那麼，廣闊的大地立在諸水之上，告訴我，豈能如此堅固、如此長久地站立，而沒有某種力量維繫它？這些觀點有任何道理嗎？想像這樣的念頭豈不是極端荒謬？若大地支撐著蒼天，看，又有一個重擔；若蒼天也承載於水之上，則又生出另一個問題。或者不如說，不是另一個問題，因為這是天主眷顧的工作。因為承載於水上的物體不應是凸形的，而應是凹形的。為什麼？因為任何凹形物體的整個身體都浸沒在水中，如同船隻；而凸形物體的身體完全在水之上，只有邊緣接觸水面；因此它需要一個堅硬、能承受的物體來支撐所施加的重擔。那麼，大氣支撐著蒼天嗎？大氣遠比水更柔軟、更易讓步，無法承受任何東西，連最輕微的東西都不能，更不必說如此巨大的體積了。總之，若我們願意追隨天主眷顧的論證，無論普遍地還是詳細地，時間本身也會不夠用。因為我現在要問那位提出上述問題的人：這些事是出於天主眷顧，還是缺乏眷顧？若他說不是出於眷顧，我便再問：那麼它們如何產生？但他永遠無法給出任何解釋。難道你不知道嗎？\par

那麼，你們更應當在涉及人的事上，不要質疑，不要過分好奇。為什麼呢？因為人比這一切更高貴，這一切都是為人而造的，人不是為它們而造的。如果你連祂的天主眷顧中可見的技藝和設計都不明白，又如何能明白祂親自作為主體的理由呢？告訴我，請問，天主為何將人造得如此渺小，遠低於天的高度，以致人甚至對自己頭頂之上的事物產生懷疑？為何北方和南方的氣候無法居住？告訴我，為何冬天黑夜更長、夏天更短？為何冷熱的程度如此？為何身體會死亡？我還要問你千萬個問題，若你願意，我將永不停止發問。在每一個問題上，你必定無法回答。而這是一切事物中最具天主眷顧的：事物之理對我們隱藏。因為若非如此，人豈不要想像自己是萬物的原因嗎？\par

\Quote{但這樣一個人，}你會說，\Quote{是貧窮的，貧窮是一種惡。又怎樣呢，生病是怎麼回事，殘廢是怎麼回事？}啊，人啊，這些都不是什麼。唯獨一件事是惡，那就是犯罪；這才是我們該深入探究的事。然而我們卻忽略了探究真正罪惡的原因，而忙於其他事。為什麼我們沒有一個人審視自己為何犯了罪？犯罪——是在我能力之內，還是不在我能力之內？我何必繞來繞去找許多理由？我將在自己內找尋。那麼，我有時克服了忿怒嗎？我是否因羞愧或畏懼人而克制了怒氣？一旦我發現這事做到了，我就會發現犯罪是在我能力之內。沒有人審視這些事，沒有人為此忙碌。反而如約伯所說：\Quote{人卻以一種全然不同的方式在水面上漂浮言語。}因為這與你有何關係，若某人瞎了，或某人貧窮？天主沒有命令你關注這些，而是關注你自己在做什麼。因為，一方面，若你懷疑有某種能力在統管世界，你是所有人中最愚昧的；但若你確信此事，為何還要懷疑我們該取悅天主呢？\par

\Quote{時常感謝，}他說，\Quote{為一切感謝天主。}\par

你去医生那里，当发现某人有了伤口时，你会看见他使用刀和烧灼器。但在你的事上，我连这样也说不上来；你去木匠那里吧。然而你不考察他的理由，尽管你不明白那里所行的一切事，而且许多事在你看来会成为难题；例如，当他刨空木头时，当他改变其外形时。不，我还要带你去看一种更易理解的技艺，比如画家的技艺，在那里你会头晕目眩。告诉我，他所作的岂不像是随意的吗？他的线条以及线条的转折和弯曲是什么意思呢？但当他涂上颜色时，艺术的美就会显现出来。然而，即便如此，你仍无法获得任何准确的理解。但我为何要提木匠和画家这些同仆呢？告诉我，蜜蜂是如何筑巢的，然后你再来谈论天主吧。掌握蚂蚁、蜘蛛和燕子的手工，然后你再来谈论天主吧。告诉我这些。但你不能，绝不可能。人哪，你还不停止你虚妄的探询吗？因为它们实在是虚妄的。你还不停止为许多事徒劳忙碌吗？没有什么比这无知更智慧的了，因为那些声称自己一无所知的人是最智慧的，而那些在这些问题上过分劳神的人是最愚昧的。所以，声称知识并非处处是智慧的标志，有时也是愚昧的标志。告诉我，假设有两个人，其中一个声称要拉出他的线条，测量天地之间的距离，另一个则嘲笑他，宣称他不明白；请问，我们应该嘲笑哪一个，是那个说他知道的，还是那个不知道的？显然，是说他知道的那个人。因此，无知者比声称知道的人更智慧。又怎样呢？如果有人声称要告诉我们海中有多少杯水，而另一个人承认自己无知，那么这里的无知岂不又比知识更智慧吗？确实，远为更甚。为什么会这样呢？因为那知识本身就是极度的无知。因为那个说自己无知的人，确实知道一些事。知道什么呢？就是这件事是人所不能理解的。是的，这并非微不足道的知识。而那个说自己知道的人，却偏偏不知道他声称知道的事，并因此完全可笑。\par

教训。哀哉！有多少事物教导我们约束这不合时宜的鲁莽和无聊的好奇心；然而我们仍不约束，反而对别人的生活感到好奇；例如，为何一人是瘸子，另一人贫穷。如此推理，我们会陷入另一种无穷无尽的琐碎，例如，为何某人是女人？为何不都是男人？为何会有驴？为何有牛？为何有狗？为何有狼？为何有石头？为何有木头？如此，论证将漫无止境。这实在是为什么天主为我们的知识划定了界限，并将其深植于本性的原因。现在，请注意这忙碌的好奇心过度到了何种程度。因为虽然我们仰望如此高远如从地到天，却丝毫不受影响；然而一旦我们登上高塔的顶端，稍微俯身向下看时，一种眩晕和头晕立刻就抓住我们。现在，告诉我这原因。不，你永远找不到原因。为什么眼睛拥有比其他感官更大的能力，并被更远的物体所吸引？通过与听觉的比较就可以看出。因为没有人能喊得那么大声，以至于充满眼睛所及的范围，也没有人能听到那么远。为什么不是所有的肢体都享有同等的尊荣？为什么不是所有的肢体都领受同一功用和同一位置？\Person{保禄}也探究过这些问题；更确切地说，他没有探究，因为他是明智的；但当他偶然提到这个话题时，他说：\BibleQuote{祂按自己的旨意，把每个肢体都安置在身体上了。}\BibleRef{格前 12:18}祂将一切归于祂的旨意。所以，让我们只需要\BibleQuote{凡事感谢。}\BibleQuote{因此，}他说，\BibleQuote{凡事感谢。}这是一个好仆人、明智的、有见识的仆人该做的；相反则是多言者、懒散者和好事者。我们难道没有在仆人中间看到，那些无价值、无用的仆人，既是多言者，又是琐碎者，并且窥探主人想要隐藏的事；而明智和良善的仆人只关心一件事，就是如何完成他们的服务。多言者无所作为；正如多作者从不说不合时宜的话。因此，\Person{保禄}在论及寡妇时写道：\BibleQuote{她们不仅懒惰，而且好说闲话。}\BibleRef{弟前 5:13}现在告诉我，最大的差别是什么：我们的年龄与儿童的年龄之间，还是天主与人之间？我们与蚊虫相比，还是天主与我们相比？显然是天主与我们之间。那么你为什么在这些问题上如此忙碌呢？\BibleQuote{凡事感谢。}\Quote{可是，}你说，\Quote{如果一个外邦人问这个问题，我该如何回答他？他想从我这里知道是否有天主眷顾，因为他本人否认有任何如此行使预见的存有。}那么，转过身去，你自己问他同样的问题。因此他会否认有眷顾。然而，确实有\OriginalTerm{天主眷顾}{Providence}，从你所说的话中可以明显看出；但它是不可理解的，这从那些我们无法找到理由的事上也显而易见。因为如果连由人安排的事，我们常常不理解安排的方式，而且实际上其中许多在我们看来似乎矛盾，但我们仍然默许，那么在天主的事上岂不更加如此？然而，在天主那里，没有什么是矛盾的，在信徒看来也不矛盾。所以，让我们\BibleQuote{凡事感谢，}让我们为一切事荣耀祂。\par

他说：\BibleQuote{你们要彼此服从，怀着对基督的敬畏}。因为如果你为了统治者的缘故、为了金钱的缘故或出于尊敬而服从，那么为了对基督的敬畏，就更应当服从了。让服务与服从互相交换。因为这样就不会有奴役性的服务。不要让一个人以自由人的身份坐下，而另一个人以奴隶的身份坐下；相反，主人和奴隶都彼此做仆人更好——这样作奴隶远比其他任何方式的自由更好；这一点从以下可以清楚看出。设想一个有一百个奴隶的人，他丝毫不为他们服务；再设想另一种情形，有一百个朋友，彼此互相伺候。哪一种生活更幸福？哪一种更有乐趣、更享受？在一种情况下，没有怒气、没有挑衅、没有忿怒，也没有任何类似的事；在另一种情况下，全是恐惧和忧虑。在一种情况下，一切都是被迫的；在另一种情况下，是出于自由的选择。在一种情况下，他们彼此服务是因为被迫如此；在另一种情况下，是彼此满足。天主愿意如此；为这缘故，祂洗了门徒的脚。不仅如此，如果你仔细考察这事，主人方面确实也有服务的回报。如果骄傲不让那服务的回报显现，那又怎样呢？然而，如果奴隶一方面提供身体的服侍，而你维持那身体，供应食物、衣服和鞋，这就是服务的交换：因为除非你也提供你的服务，否则他也不会提供他的，而是会自由，并且如果没有支持，就没有法律能强迫他这样做。如果对仆人来说是这样，那么对自由人也如此，又有什么荒谬呢？他说：\BibleQuote{你们要服从，怀着对基督的敬畏}。那么，当我们还将得到赏报时，这义务何等重大。但他不愿意向你服从？然而你还是要服从自己，不仅是让步，而是服从。要对所有人怀有这种情感，仿佛所有人都是你的主人。因为这样你很快就会使所有人成为你的奴隶，以最卑贱的奴役方式臣服于你。因为当你没有从他们那里接受任何东西，反而把你自己的给他们时，你会更确定地使他们成为你的。这就是\BibleQuote{你们要彼此服从，怀着对基督的敬畏}，好使我们能克制所有情欲，成为天主的仆人，并保持彼此亏欠的爱。那时，我们也将配得上从天主而来的慈爱，借着祂独生子的恩宠和怜悯，愿荣耀、权能、尊贵归于圣父及圣子及圣神，从现在直到永远。阿们。\par

\SourceRecord{Translated by Gross Alexander.；Nicene and Post-Nicene Fathers, First Series；Vol. 13.；Edited by Philip Schaff.；Buffalo, NY: Christian Literature Publishing Co.,；1889.；<http://www.newadvent.org/fathers/230119.htm>.}

% structure: p0001 source=h1 target=section reason=page-title

\section{厄弗所书讲道集 第20篇}

\BibleRef{弗 5:22-24}\par

\BibleQuote{妻子要服从自己的丈夫，如同服从主一样；因为丈夫是妻子的头，如同基督是教会的头——祂自己也是身体的救主。但教会怎样服从基督，妻子也该怎样在一切事上服从丈夫。}\par

一位智者列出许多福分时，也将此列为福分：\BibleQuote{与丈夫相合的妻子。}\BibleRef{德 25:1} 在其他地方，他又将此列为福分：\BibleQuote{妻子应与丈夫和睦同居。}\BibleRef{德 40:23} 的确，从一开始，天主似乎就特别为这一结合做了预备；并且论述这两者如同一体时，祂这样说：\BibleQuote{祂造了他们，有男有女}\BibleRef{创 1:27}；又说：\BibleQuote{并不分男女性别。}\BibleRef{迦 3:28} 因为人与人之间没有比夫妻更加亲密的关系，如果他们按应该的方式结合。因此，一位有福的人也这样说：当他想要表达超乎寻常的爱，为他所爱的、同心合意的人哀悼时，他没有提父亲、母亲、子女、兄弟或朋友，而是说了什么？他说：\BibleQuote{你的爱对我奇妙非凡，超越妇女的爱。}\BibleRef{撒下 1:26} 因为实在说来，这种爱比任何专制都更专制：因为别的爱也许坚强，但这激情不仅坚强，而且永不消退。因为我们本性中深深植根着一种爱，它不知不觉地将我们的身体结合在一起。因此，从起初女人就出自男人，之后又从男人和女人生出男人和女人。你看到这紧密的纽带和联系吗？天主如何不容许另一种不同的本性从外进入？请注意祂做了多少眷顾的安排。祂允许男人娶自己的姊妹；或者说不是姊妹，而是女儿；不，也不是女儿，而是比女儿更亲密的事物，就是他自己的血肉。这样，祂从同一个开端形成了一切，将一切聚集在一起，如同建筑物中的石头。因为祂既没有从外面造出女人，这是为了使男人不把她当作外人；也没有把婚姻限于她一人，免得她因自我封闭、一切以自我为中心而与其他部分隔绝。正如植物，最好的那些具有单一的茎干，并生发出许多枝条；（如果一切都局限在根上，那就毫无用处；而如果有很多根，树就不再值得赞美；）同样，这里也是如此。祂从一个人，即亚当，使全人类衍生，用最强烈的必要性阻止他们彼此分裂或分离；之后，使这关系更加严格，祂不再允许娶姊妹或女儿为妻，免得我们另一方面将爱局限于一点，从而以另一种方式彼此隔绝。因此基督说：\BibleQuote{那起初创造人的，造了他们有男有女。}\BibleRef{玛 19:4}\par

因为由此产生大恶和大善，无论对家庭还是对国家。因为没有什么比夫妻之爱更能焊接我们的生命。为此许多人会放下武器，为此他们愿意牺牲生命。保禄绝不会无缘无故、毫无目的地在这一点上花这么多力气，就像他在这里说的：\BibleQuote{妻子要服从自己的丈夫，如同服从主一样。} 为什么呢？因为当他们和谐时，孩子被良好抚养，家仆秩序井然，邻居、朋友和亲戚都感受到馨香。但如果相反，一切都将颠倒混乱。正如军队的将军彼此和平时，一切井然有序；反之，如果他们彼此不和，一切就会颠倒；这里也是如此。因此，他说：\BibleQuote{妻子要服从自己的丈夫，如同服从主一样。}\par

然而多么奇怪！因为别处怎么说：\BibleQuote{若有人不向妻子和丈夫告别，就不能跟从我。}\BibleRef{路 14:26} 如果他们有义务\BibleQuote{如同服从主一样}服从，为什么主说为了主的缘故他们必须离开他们？然而这确实是他们的义务，是当行的。但\Quote{如同}一词不一定普遍表示完全平等。他或者意味着：\Quote{\InnerQuote{如同}知道你们是主的仆人}（顺便说，这是他在别处说的，即使不为丈夫的缘故，也须主要为主的缘故）；或者意味着：\Quote{当你服从丈夫时，要如同服侍主。} 因为如果抗拒这些外在的权柄（我指的是政府的权柄），就是\BibleQuote{抗拒天主的安排}\BibleRef{罗 13:2}，那么不服从丈夫的妻子更是如此。这本是天主从起初的旨意。\par

让我们以此为基本原则：丈夫占据\Quote{头}的位置，妻子占据\Quote{身体}的位置。\par

第23、24节。然后，他继续论证，说：\BibleQuote{丈夫是妻子的头，如同基督是教会的头——祂自己也是身体的救主。但教会怎样服从基督，妻子也该怎样在一切事上服从丈夫。}\par

然后，在说了\Quote{丈夫是妻子的头，如同基督也是教会的头}之后，他又加上\Quote{并且祂是身体的救主}。实际上，头就是身体的救恩健康。他早已为丈夫和妻子预先立定了他们爱情的基础和供应，给每人分派了适当的位置：给丈夫的是权威和远见，给妻子的是顺服。那么，正如\Quote{教会}——即丈夫和妻子——\Quote{服从基督，同样，你们的妻子也要服从自己的丈夫，如同服从天主}。\par

第25节。\BibleQuote{丈夫们，爱你们的妻子，如同基督也爱了教会。}\BibleRef{弗 5:25}\par

你已经听到顺服有多么重大；你称赞并钦佩\Person{保禄}，宛如一位可敬而属神的人，将我们的整个生命凝聚在一起。你做得好。但现在也听听祂对你的要求；因为祂再次使用了同一个例证。\par

\Quote{丈夫们，}祂说，\Quote{爱你们的妻子，如同基督也爱了教会。}\par

你已经看到了服从的尺度，也来听听爱的尺度。你要你的妻子服从你，如同教会服从基督吗？那么你就当为妻子取同样的眷顾，如同基督为教会所取的。是的，即使你需要为她舍命，即使需要粉身碎骨一万次，即使需要忍受并承受任何苦难——也不要拒绝。即使你忍受了这一切，你也仍然、甚至那时也没有做过任何像基督所做的事。因为你确实是为一个已与你结合的人而做；但祂却是为那背弃祂、恨祂的人而做。因此，正如祂以不倦的柔情使那背弃祂、恨祂、藐视祂、轻慢祂的人俯伏在祂脚下，既非借着威吓，亦非借着暴力，也不是借着恐惧或任何类似的手段；同样，你也要如此待你的妻子。是的，即使你看到她轻视你、藐视你、嗤笑你，你也能靠着对她的极大关怀、柔情和善意，使她俯伏在你脚下。因为再没有比这些纽带更有力的支配力量，尤其对于丈夫和妻子。对于仆人，或许能用恐惧来束缚；但即使是对他也不能，因为他很快就会挣脱离去。然而，对于生命的伴侣、子女的母亲、一切喜乐的根基，绝不应以恐惧和威吓来捆绑，而应以爱和温柔。因为妻子在丈夫面前战栗，那是怎样的结合？丈夫自己若与妻子如同与奴隶同住，而不是与自由妇女同住，他又能享受怎样的快乐？是的，即使你因她而受苦，也不要责备她；因为基督也没有这样做。\par

第26节。\BibleQuote{并且祂交付了自己，}祂说，\BibleQuote{为她，为要圣化并洁净她。}\BibleRef{弗 5:26}\par

那么，她是不洁的！那么，她有瑕疵，那么，她是不好看的，那么，她是无价值的！无论你娶了怎样的妻子，你绝不会娶到像基督娶教会时那样的新娘，也不会有一位像教会远离基督那样远离你的新娘。然而，尽管这样，祂并不憎恶她，也不因她那超乎寻常的畸形而厌恶她。你愿意听她的畸形被描述吗？听\Person{保禄}说什么：\BibleQuote{因为你们从前是黑暗，}\BibleRef{弗 5:8}你看到她肤色的黑暗了吗？有什么比黑暗更黑？但再看她的狂妄：\Quote{活在恶意和嫉妒中，}祂说。\BibleRef{铎 3:3}再看她的不洁：\Quote{悖逆、愚蠢。}但我说什么呢？她既是愚蠢，又是口出恶言；然而，尽管她有这么多瑕疵，祂却为她、在她畸形时交付了自己，如同为一位青春焕发的人，如同为一位深爱者，如同为一位绝色美人。正是为此，\Person{保禄}惊叹地说：\Quote{为义人死，尚不常见；}\BibleRef{罗 5:7}又说：\BibleQuote{当我们还是罪人的时候，基督为我们死了。}\BibleRef{罗 5:8}祂虽如此，却接纳了她，使她美丽，洗净了她，甚至不惜为此交付自己。\par

第26、27节。祂接着说：\BibleQuote{为要圣化她，洗净她，借着水的洗涤和圣言，好把教会呈献给光荣的自己，没有瑕疵、皱纹或任何类似的缺陷，而是圣洁无瑕的。}\BibleRef{弗 5:26-27}\par

\Quote{洗涤或沐浴}祂洗去她的不洁。\Quote{藉着圣言}，他说。什么圣言？\BibleQuote{因父、及子、及圣神之名}\BibleRef{玛 28:19}。祂不仅装饰了她，而且使她成为\BibleQuote{荣耀的，没有斑点、皱纹，或任何这类毛病}。那么，我们也当追求这种美，我们就能创造它。不要从你妻子手中索取她无法拥有的事物。你看见教会从她的主手中领受了一切吗？藉着祂成为荣耀，藉着祂成为纯洁，藉着祂成为无瑕？不要因你的妻子有缺陷就背弃她。请听圣经说：\BibleQuote{飞禽中蜜蜂是小的，但它的产品却是甘甜之首}\BibleRef{德 11:3}。她是天主的创造。你所责难的并非她，而是造她的天主；那女人能做什么呢？不要因她的美貌而赞美她。基于美貌的赞美、憎恨和爱情属于未受节制的灵魂。你当寻求灵魂的美。效法教会的新郎。外在的美貌充满自负和极大的放肆，使男人陷入嫉妒，且常使你猜疑可怕之事。但它有什么快乐呢？或许第一个月、第二个月，至多一年；此后便不再有；因熟悉而产生的仰慕逐渐消逝。同时，由美貌产生的恶——骄傲、愚昧、轻蔑——依然存在。而在并非如此的人身上，就没有这类事。但基于正当理由开始的爱情，依旧炽热，因为其对象是灵魂的美，而非身体的美。请告诉我，有什么比天堂更好？有什么比星辰更好？任你提到什么身体，也没有这样美丽的。任你提到什么眼睛，也没有这样闪耀的。这些受造时，连天使也惊奇地注视，如今我们也惊奇地注视，却不如最初那样。熟悉就是如此；事物不再同样触动我们。更何况一位妻子呢！如果疾病又袭来，一切便立刻消逝。让我们在妻子身上寻求温情、谦逊、温柔；这些是美的特征。但我们不要寻求容貌的可爱，也不要因这些她无能为力的事而责备她；不，倒是根本不要责备（那是粗鲁），也不要急躁或愠怒。你们岂不见，有多少人与美貌妻子生活后可怜地结束一生，而又有多少人与不大美貌的人生活，却极其愉快地活到耄耋之年？让我们擦去内在的\Quote{斑点}，抚平内在的\Quote{皱纹}，除掉灵魂上的\Quote{瑕疵}。这就是天主所要求的美。让我们使她在天主眼中美丽，而非在我们眼中。不要追求财富，也不要追求外在的高贵出身，而要追求灵魂内的真高贵。不要容忍任何人藉妻子致富；因为这样的财富是卑劣可耻的；不，决不要有人从这一来源寻求致富。\BibleQuote{因为那些想要发财的人，陷于诱惑和罗网，以及许多愚昧而有害的私欲，并陷入毁灭和丧亡}\BibleRef{弟前 6:9}。因此，不要在你妻子身上寻求大量的财富，其他一切就都会顺利。请告诉我，谁会忽略最重要的事，而去关注不那么重要的事呢？然而，唉！这正是我们在每种情况下的感受。是的，如果我们有一个儿子，我们关心的不是如何使他成为有德之人，而是如何为他找到一位富有的妻子；不是如何使他举止得体，而是如何使他钱财充沛：如果我们经营一项生意，我们询问的不是它如何能不犯罪，而是如何能带给我们最大的利润。一切都变成了金钱；于是，一切都败坏、毁灭，因为那种情欲占据了我们的心。\par

第二十八节。\BibleQuote{照样，丈夫也应当爱自己的妻子，}他说，\BibleQuote{如同爱自己的身体}\BibleRef{弗 5:28}。\par

这又是什么意思？他来到了何等更大的类比和更强的例证；不仅这样，而且还到了一个何等更近、更清楚和更新的义务。因为那另一个没有很大的强制力，因为祂是基督，是天主，并献出了自己。如今他在不同的基础上进行论证，说：\Quote{丈夫也当如此}；因为这事不是恩惠，而是债务。然后，\Quote{如同他们自己的身体}。为什么？\par

第二十九节。\BibleQuote{从来没有人恨恶自己的肉身，反而培养抚育它}\BibleRef{弗 5:29}。\par

就是说，以极大的关怀照顾它。她如何是他的肉？请听；\BibleQuote{这才是我骨中的骨，}亚当说，\BibleQuote{肉中的肉}\BibleRef{创 2:23}。因为她是由取自我们的质料造成的。不仅如此，而且，\BibleQuote{他们要成为，}天主说，\BibleQuote{一体}\BibleRef{创 2:24}。\par

\Quote{正如基督对教会一样。}这里他回到先前的例子。\par

第三十节。\BibleQuote{因为我们是他身体、他的肉和他的骨的肢体}\BibleRef{弗 5:30}。\par

第三十一节。\BibleQuote{为此，人要离开父亲和母亲，依附自己的妻子，两人成为一体}\BibleRef{弗 5:31}。\par

请看又一个义务的第三个理由；因为他表明，一个人离开生养他、他由之而出生的人，与妻子结合；然后一体就是父亲、母亲和孩子，来自二者混合的实体。的确，通过他们种子的混合，孩子得以产生，因此三人成为一体。我们与基督的关系也是如此；我们藉着分享成为一体，而且我们比孩子更加是。为什么和如何如此？因为从一开始就是这样。\par

不要向我说事情是如此这般。你难道看不见我们自己的肉身本身就有许多缺陷吗？例如，一人瘸腿，另一人脚歪，第三人手枯萎，第四人某个肢体软弱；然而尽管如此，他并不为此忧伤，也不将其切断，反而常常偏爱它胜过其他肢体。这是很自然的；因为它是他自身的一部分。一个人对自己怀有多大的爱，他就希望我们以同样的爱来对待妻子。不是因为我们是同一本性；不，这义务的基础远大于此；那就是并非两个身体，而是一个；他是头，她是身体。那么他为何在别处说：\Quote{而基督的头是天主}\BibleRef{格前 11:3}？我也要说，正如我们是一个身体，基督与圣父也是一体。这样，圣父也被发现是我们的头。他设下两个比喻：自然身体的比喻和基督身体的比喻。因此他又接着说：\newline{}\par

第32节。\Quote{这是伟大的奥迹：但我是指基督和教会说的。}\BibleRef{弗 5:32}\newline{}\par

为什么他称此为伟大的奥迹？这是伟大而奇妙的事，蒙福的梅瑟，或更确切地说，是天主所暗示的。但现在，他说，我是指基督说的：祂离开了父，降来，来到新娘那里，成为一灵。\Quote{因为那与主结合的，便是与祂成为一灵。}\BibleRef{格前 6:17}他说得对，\Quote{这是伟大的奥迹。}然后，仿佛他在说：\Quote{但尽管如此，这寓意并不破坏爱，}他补充道：\newline{}\par

第33节。\Quote{但你们各人也要各自爱自己的妻子，如同爱自己一样；妻子也当敬重她的丈夫。}\BibleRef{弗 5:33}\newline{}\par

因为确实，实际上这是一个奥迹，是的，一个伟大的奥迹：一个人要离开那给予他生命者、生养他者、抚育他者，离开那为他忍受产痛和悲伤者，离开那些曾施予他如此多、如此大恩惠的人，离开那些他素来亲密交往的人，而与一个从未见过、与他毫无共同之处的人结合，并尊崇她胜过其他人。这确实是一个奥迹。然而父母在发生这事时并不悲伤，反倒在不发生时悲伤；他们欣喜于自己的财富为此而花费和挥霍。——真是一个伟大的奥迹！其中含有隐藏的智慧。梅瑟最早以此喻示了它；现在保禄也如此宣告，他说：\Quote{关于基督和教会。}\newline{}\par

然而，这话不仅是为丈夫说的，也是为妻子说的：\Quote{他当珍爱她如同自己的肉身，如同基督珍爱教会一样，}以及\Quote{妻子当敬重她的丈夫。}他不再只列举爱的义务，而是什么？\Quote{她当敬重她的丈夫。}妻子是第二权威；因此她不可要求平等，因为她是在头之下；丈夫也不可因她顺服而轻视她，因为她是身体；头若轻视身体，自身也必灭亡。而是在他那一方以爱作为她顺服的平衡。例如，让手、脚和所有其他肢体都服务头，但头要为身体提供所需，因为头自身拥有一切感官。没有什么比这结合更好的了。\newline{}\par

然而有人会说：有害怕的地方，怎能还有爱呢？我说，爱在此处会卓越地存在。因为她害怕和尊敬，也就爱；而那爱他的，也尊敬他如同头，并爱他如同肢体，因为头本身也是整个身体的一个肢体。因此他将一方置于顺服，另一方置于权威，为能有和平；因为哪里有平等权威，哪里绝不能有和平；无论是民主制的家庭，还是人人作主的情形；统治权必须是一元的。这在身体方面是普遍适用的，因为当人们属神时，就有平安。那时有\Quote{五千个灵魂}，没有一人说\Quote{他所拥有的任何东西是自己的}\BibleRef{宗 4:32}，而是彼此顺服；这是智慧和敬畏天主的标志。然而，他解释了爱的原则；却没有解释害怕的原则。请注意，他如何详述爱的原则，引用关于基督和关于自身肉体的论点，以及\Quote{为此，人要离开父母}\BibleRef{弗 5:31}这话。至于从害怕引出的论点，他就不再加详述。为什么呢？因为他宁愿这原则，即爱的原则，占上风；因为哪里有此，其他一切自然随之而来；但哪里有彼，却未必。因为爱妻子的丈夫，即使她不十分顺服，也会忍受一切。所以，同心合意是困难和不易实现的，除非人们被那建立在最高权威之上的爱所联结；无论如何，害怕未必能达到此效果。因此，他更多地讲论这强而有力的纽带。妻子看似因被命令要害怕而受损，实则获益，因为主要义务即爱是加在丈夫身上的。\Quote{但是，}有人会说，\Quote{如果妻子不敬重我呢？}没关系，你仍要爱，尽你的本分。因为即使别人应尽的义务不随之而来，我们当然仍应尽自己的本分。这是我所说的一个例子。他说：\Quote{要在敬畏基督中彼此顺服。}那么，若别人不顺服呢？你仍要遵守天主的法律。同样，在此也是如此。即使妻子不被爱，仍要敬重，以免自己有过失；丈夫即使妻子不敬重他，仍要爱她，以免自己有任何欠缺。因为各人都领受了自己的本分。\newline{}\par

这就是按基督结合的婚姻，是属神的婚姻，属神的诞生，不是由血气，不是由肉欲，也不是由肉体的意愿。基督的诞生就是如此，不是由血气，也不是由肉欲。\Person{依撒格}的诞生也是如此。请听圣经如何说：\BibleQuote{\Person{撒辣}已停止，不再有月经。}\BibleRef{创 18:11}是的，这是婚姻，不是情欲，也不是肉体的，而是完全属神的，灵魂与天主结合，是言语无法形容的结合，唯独祂知道。因此他说：\BibleQuote{那与主结合的，便是与祂成为一神。}\BibleRef{格前 6:17}请注意他多么热切地努力将肉身与肉身、灵魂与灵魂结合。而异端者在哪里？如果婚姻当受谴责，他绝不会称基督和教会为新娘和新郎；绝不会用劝勉的方式引出这些话：\Quote{人要离开父亲和母亲}，并又加上，那是\Quote{指着基督和教会说的}。因为诗篇作者论到教会也说：\BibleQuote{女子，请听，请看，请侧耳；忘掉你的民族和你的父家！君王将恋慕你的美色。}\BibleRef{咏 45:10}因此基督也说：\BibleQuote{我出自父，来到了世上。}\BibleRef{若 16:28}但我说祂离开了父，你们不要想这是人间的变迁，如同\Quote{出来}一词一样，并不是说祂真地出来了，而是指祂的降生成人，同样，\Quote{祂离开了父}也是如此。\par

现在，为什么他没有论到妻子说，她将依附于她的丈夫？我说，这是为什么？因为他是在谈论爱，是在对丈夫说话。因为他确实对妻子论到敬畏，说：\BibleQuote{丈夫是妻子的头}\BibleRef{弗 5:23}，又说：\BibleQuote{基督是教会的头。}而对他，他谈论爱，并将爱的领域交托给他，向他宣告与爱相关的事，这样把他与她绑定并紧密相连。因为一个男人为了妻子离开父亲，然后又离开这同一位妻子并抛弃她，他还能得到什么宽容呢？\par

难道你没有看出天主愿意她享受多大的荣耀，因为他使你离开你的父亲，将你与她联结？那么，有人会说，如果我们尽了本分，而她却未效法？\BibleQuote{倘若那不信的人离去，就由他离去吧；在这种情形下，弟兄或姊妹都不受束缚。}\BibleRef{格前 7:15}\par

然而，当你听到\Quote{敬畏}时，要求那种合乎自由女子的敬畏，而不是像对奴隶一样苛求。因为她是你自己的身体；如果你这样做，你就是因羞辱自己的身体而羞辱自己。这个\Quote{敬畏}是什么样的？就是不顶撞、不反叛、不争首位。让敬畏停留在这界限内就够了。但如果你如所命令的那样去爱，你就会使敬畏变得更甚。或者更确切地说，你将不再是通过敬畏来做这事，而是爱本身将发挥作用。女性天生有某种软弱，需要许多支持，许多迁就。\par

但那些在第二次婚姻中结合的人会说什么？我绝不是指责他们，绝无此意；因为宗徒本人也允许他们，尽管的确是出于迁就。\par

供应她一切所需。为她做一切事，忍受一切辛苦。这是你不可推卸的责任。\par

在这里，他认为不宜像他在许多情况下所做的那样，引用外邦人的例子来引入他的劝告。基督的例子，如此伟大有力，就已足够；尤其是在关于服从的论点上。他说：\Quote{人要离开}\Quote{他的父亲和母亲}。看，这例子来自外邦。但他没有说\Quote{要同住}，而是说\Quote{要依附}，以此显示结合的紧密和炽热的爱。不仅如此，他还进一步通过他的补充，以这样的方式解释服从，以致二人不再像是两个人。他没有说\Quote{一灵}，没有说\Quote{一魂}（因为那是显而易见的，任何人都可能做到），而是说成为\Quote{一体}。她是第二权威，确实拥有权威，并有相当大的同等的尊严；但同时也拥有某种优越性。家中的福祉主要在于此。因为他引用了先前的例子，即基督的榜样，表明我们不仅应当爱，还应当治理；他说：\Quote{使她成为}\Quote{圣洁无瑕的}。但\Quote{体}一词与爱有关，\Quote{依附}一词同样与爱有关。因为如果你使她成为\Quote{圣洁无瑕的}，其他一切都会随之而来。寻求天主的事，那些人的事就会随之而来。治理你的妻子，整个家就会和谐。请听保禄所说：\BibleQuote{她们若愿意学什么，可以在家问自己的丈夫。}\BibleRef{格前 14:35}如果我们这样管理自己的家，我们也将适合管理教会。因为家确实是一个小教会。这样，通过成为好丈夫和好妻子，我们就有可能超越所有的人。\par

试想\Person{亚巴郎}、\Person{撒辣}、\Person{依撒格}以及他家中生养的三百一十八人。\BibleRef{创 14:14}全家何等和谐地结合，如何充满虔诚，遵行了宗徒的训令。她也\Quote{敬重丈夫}；因为听听她自己的话：\Quote{直到现在，我还没有发生这样的事，而且我的丈夫也老了。}\BibleRef{创 18:12}而他又是如此爱她，凡事听从她的命令。那幼童品行端正，家中生养的仆役也极优秀，甚至不惜与主人一同冒险；他们不迟延，也不问缘由。不仅如此，其中为首的那位，竟是如此可敬，以至于连独生子的婚事也被托付给他，并派他远赴他国。\BibleRef{创 24:1}因为正如一位将军，当他的士兵组织良好时，敌人便无隙可乘。同样，我说，这里也是如此：当丈夫、妻子、儿女和仆役都关注同一件事时，家中就有极大的和谐。若情况并非如此，整个家庭往往因一个坏仆役倾覆破裂，而那一人就常常毁坏并彻底摧毁全家。\par

道德教训。那么，让我们对我们的妻子、儿女和仆役倍加谨慎；因为我们知道，这样我们就会为自己建立一个轻松的治理，并将在与他们的账目上温和宽容，并且说：\Quote{看，我和天主所赐给我的孩子们。}\BibleRef{依 8:18}如果丈夫令人起敬，头受尊崇，那么其余肢体就不会受到伤害。现在，妻子的适当态度和丈夫的适当态度如何，他准确地陈述了：命令她敬重他如同头，而他要爱她如同妻子；但人们可能会说，这些事如何可能？他证明了它们确实应当如此。但它们如何可能，我将告诉你们。它们将如此，如果我们轻视钱财，如果我们只看一件事，即灵魂的卓越，如果我们把对天主的敬畏放在眼前。因为他在对仆役的讲话中所说的：\Quote{无论人做什么，或善或恶，都必从主那里领受。}\BibleRef{弗 6:8}这里也是如此。因此，爱她不是为了她的缘故，而是为了基督的缘故。至少，他在说\Quote{如同在主内}时已暗示了这一点。所以，凡事都要如同顺服主，并为他的缘故而做。这就足以激励并劝服我们，不容许有任何嘲弄和纷争。不要相信任何人向妻子毁谤丈夫；同样，也不要让丈夫随意相信任何不利于妻子的事，也不要让妻子无端探询他的外出和归来。不仅如此，在任何情况下，丈夫都不要使自己值得任何猜疑。因为，告诉我，如果你整天与朋友厮守，只把晚间留给妻子，甚至这样还不能使她满足并使她免于猜疑，那又如何？即使你的妻子抱怨，也不要不悦——这是她的爱，不是她的愚昧——这些是热切依恋、燃烧的情感和恐惧的抱怨。是的，她害怕有人夺了她的婚床，害怕有人损害了她福分的高峰，害怕有人从她那里夺走了她的头，害怕有人闯入了她的婚房。\par

还有另一种微小嫉妒的缘由。不要让任何一方过分要求仆役的服务，既不要让丈夫要求女仆，也不要让妻子要求男仆。因为这些事也足以产生猜疑。因为请想一想，我说，我所说的那个义人的家庭。撒辣自己命令先祖娶哈加尔。她亲自指示，没有人强迫她，丈夫也没有试图这样做；不，尽管他长期没有子女，却宁愿永远不做父亲，也不愿使妻子伤心。然而即使如此，撒辣说了什么？\Quote{愿主在你我之间判断。}\BibleRef{创 16:5}现在，我说，如果他换成别人，岂不就会动怒吗？他岂不会也伸出手说：\Quote{你是什么意思？我根本不想与那女人有任何关系；这全都是你做的；而你倒反过来指责我？}但不会，他没有说任何类似的话——而是什么？\Quote{看，你的婢女在你手中，你看怎样好，就怎样待她吧。}\BibleRef{创 16:6}他交出了与他同床的伴侣，为不使撒辣伤心。然而，确实没有什么比这更能产生爱情。因为同席共食甚至能使强盗与仇敌同心（圣咏作者说：\Quote{与我同席吃美味的人}）；更何况成为一体——因为这就是同床共枕的关系——更能有效地吸引我们结合。然而这些事没有一样能克服他；但他将哈加尔交还给妻子，以表明他自己毫无过错。不仅如此，他还把她打发出去，那时她正怀孕。谁会不可怜一个怀了自己孩子的人呢？然而这位义人毫不动摇，因为他将应给妻子的爱置于一切之上。\par

那么，让我们自己效法他。不要让任何人因邻人的贫穷而责备他；不要让任何人贪爱钱财；这样一切困难都将结束。\par

妻子也不可对丈夫说：\Quote{你真是个懦夫，满心的懒惰和迟钝，沉睡不醒！你看某人，出身低贱，父母卑微，他冒险航海，赚了大钱；他妻子佩戴珠宝，驾着一对白骡外出，四处游逛，随从成队，太监成群；而你却畏缩不前，活着毫无用处。}妻子不可说这些话，也不可说类似的话。因为她是身体，不应指挥头，而应顺服和听从。\Quote{可是，}有人会说，\Quote{她如何忍受贫困？到哪里寻求安慰？}让她挑选那些比她更贫穷的人来比较。让她再想想，多少出身高贵的小姐不仅没有从丈夫那里得到什么，反而给了他们嫁妆，并为他们耗尽一切。让她反思从这般财富中产生的危险，她自然会珍惜这平静的生活。简言之，如果她对丈夫有深情，就绝不会说出这类话。不，她宁可选择让他留在身边，即使一无所获，也不愿他赚得万两黄金，却带来那些远航给妻子们带来的种种忧虑和焦虑。\par

然而，丈夫听到这些，也不可因自己拥有最高权威而诉诸辱骂和殴打；而要激励她，劝诫她，因为她不够完美，要用论点说服她。他绝不可举手——这远非高尚的精神——也不可说出侮辱、嘲讽或谩骂；而要管束引导她，因她缺乏智慧。但这如何做到？若她受了真财富、天上学问的教导，就不会发出这样的埋怨。那么，让他教导她贫穷并非邪恶。让他教导她，不仅通过言语，也通过行为。让他教导她轻视荣耀；这样他的妻子就不会谈论任何这类事，也无任何此类渴望。让他如同手中得到一尊可塑造的像，从最初迎她进洞房的那个晚上起，就教导她节制、温柔，以及如何生活，从一开始、从门槛起就铲除贪财之心。让他用智慧训导她，劝她绝不要在耳朵上、脸颊边、颈项上挂金片，也不要在房里收藏，也不要有金贵衣物积存。但让她的房间雅致，却不可让雅致沦为奢靡。不，把这些留给舞台演员。你本人要用一切可能的整洁装饰你的家，使之散发着节制的气息，而非浓香。因为由此将产生两三个好处。首先，当房间打开，薄纱、金饰、银器都被送回各自的主人时，新娘不会悲伤。其次，新郎不会为丢失或积累的财宝的安全而焦虑。第三，除此之外，这是所有这些益处的冠冕，通过这些，他将表明自己的判断：他根本不以这些为乐，而且还将结束与它们一致的其他一切，永不允许存在舞蹈或不雅的歌曲。我知道，给出这些告诫，在许多人看来或许显得可笑。但尽管如此，只要你们听从我，随着时间推移，这些做法的益处临到你们，你们就会明白它的好处。笑声会过去，你们会嘲笑当今的风气，会看到现今的实践其实是愚蠢的孩子和醉汉的作为。而我推荐的则是节制、智慧和最崇高生活方式的部分。那么，我说我们的责任是什么？从婚姻中除去所有那些可耻的、撒殚般的、不雅的歌曲，那些放荡青年的队伍，这将有助于克制你新娘的精神。因为她会立刻这样对自己思量：\Quote{奇妙！这人是个何等优秀的哲人！他把今生视为无物，他带我来到他的家，是要我成为母亲，抚养他的孩子，管理他的家务。} \Quote{是的，但这些东西让新娘反感？} 只是在头一两天，但之后就不会了；不仅如此，她还会从中得到最大的喜乐，并消除所有猜疑。因为一个连笛手、舞者和杂曲都不能忍受的人，尤其是在婚礼大喜之时，那人几乎永不会做或说出任何可耻的事。然后，当你从婚姻中剥离了这一切之后，就带她来，小心塑造和模塑她，长时间鼓励她的羞怯，不要突然破坏它。因为即使少女非常大胆，但在一段时间内，出于对丈夫的尊敬，她会保持沉默，觉得自己是环境中的新手。那么你不可像不贞洁的丈夫那样过于急躁地打破她的含蓄，而要长期鼓励它。因为这对你将是一大益处。在此期间，她不会抱怨，不会对你为她制定的任何规矩吹毛求疵。因此，在那段时期里，羞耻如同勒在灵魂上的缰绳，不让她发出任何怨言，也不抱怨所做的事，你要定下你所有的律法。因为她一旦获得胆量，就会毫无畏惧地推翻和搅乱一切。那么，还有什么比她在尊敬丈夫、仍然胆怯害羞的时候更适合塑造妻子的时机呢？那时你为她定下所有律法，她无论愿意与否，一定会服从。但你怎么能不破坏她的谦逊呢？通过向她表明你自己和她一样谦逊，对她说话不多，且语气庄重沉稳。然后将智慧的言论托付给她，因为她的灵魂会接受。并使她确立在那最可爱的习惯上，即谦逊。如果你愿意，我也将举例告诉你，应该对她说怎样的话。因为如果保禄不羞于说\BibleQuote{你们互相不可欺骗} \BibleRef{格后 7:5}，他说的是伴娘的话，不如说不是伴娘，而是属灵灵魂的话，那么我们更不会羞于说话。那么我们应该对她说怎样的话呢？我们可用极柔和的方式对她说：\Quote{我的孩子，我娶你作为我生命的伴侣，带你进入最亲密最荣耀的纽带中，分享我的儿女和家务的管理。那么我现在该给你什么忠告呢？} 但不如先和她谈谈你对她的爱；因为没有比确信说话者出自挚爱更能让听者真诚接纳所说的话了。那么你如何显示这份爱呢？通过说：\Quote{当我有能力娶许多女子为妻，既有更丰厚的家产，也有高贵的门第，我并没有那样选择，而是恋慕了你，和你美好的生活、你的谦逊、你的温柔、以及你的节制。} 然后立即从这些开头引入关于真正智慧的谈论，绕着弯子对财富表示反对。因为如果你直接针对财富争论，你会打击她太重；但如果你抓住一个时机，你就能完全成功。因为你将显得是在做一种辩解，而非一个乖戾、不雅、对小事过于挑剔的人。但当你从与她自己有关的时机出发时，她甚至会很高兴。那么你就要说（因为现在我必须继续这个话题）：\Quote{我本可以娶一个富有的女人，并有好的命运，但我不能忍受。为什么？不是任性妄为、没有理由；而是我真实地学过：金钱不是真正的财产，而是最可鄙的东西，而且它还属于贼、妓女和盗墓者。所以我放弃了这些东西，一直走到我遇见了你灵魂的卓越，我视之高于一切黄金。因为一个谨慎、诚实、心向虔诚的少女，抵得上整个世界。因此我追求了你，我爱你，将你置于我自己的灵魂之上。因为今生算不得什么。我祈祷、恳求并尽我所能，使我们能堪当度过此生，好在来世也能彼此完全安全地结合。因为我们在世的时间短暂易逝。但如果我们蒙恩悦乐天主，以此生换取那生，那么我们就会永远与基督同在、彼此同在，享受更丰盛的喜乐。我看重你的感情胜过一切，没有什么比与你失和更使我痛苦和悲伤。是的，即使我命中失去一切，变得比伊鲁斯还穷，经历极度的危险，忍受任何痛苦，只要你对我的真情不变，一切都是可容忍、可忍受的。而当你对我充满感情时，我的孩子也将是我最亲爱的。但你也必须尽这些本分。} 然后，也要在你的话语中混合宗徒的话，即\Quote{这样天主愿我们的感情融合在一起；因为请听圣经所说：\BibleQuote{为此，人要离开父亲和母亲，依附自己的妻子。}不要给我们心胸狭窄的嫉妒留下藉口。愿财富、成群仆役和一切外在的虚华都消失。对我来说，这比一切更宝贵。} 什么黄金的重量，多少财宝，能像这些话一样对妻子如此珍贵？不要担心因为她被爱就会对你狂怒，而要承认你爱她。因为那些娼妓，时而依附一人，时而依附另一人，若是听到这样的话，自然会鄙视她们的情人；但一个自由出生的妻子或高贵的少女绝不会被这样的话影响；不，她会更加顺服。也要向她显示，你高度重视她的陪伴，你为了她更愿意在家，而非在市场。并要尊重她胜过你所有的朋友，胜过她所生的子女，并且让这些子女因她而为你所爱。若她做了任何好事，就赞美欣赏；若做了任何蠢事，就像女孩子常会做的，就劝告提醒她。彻底谴责一切财富和奢侈，温和地指出整洁和谦逊中的装饰；并不断教导她有益的事。\par

要让你们的祈祷成为共同的。让每个人都去教会；让丈夫在家询问妻子，妻子也询问丈夫，关于那里所说和所读之事的记述。如果任何贫穷临到你们，就举出那些圣洁的人、\Person{保禄}和\Person{伯多禄}的例子，他们比任何国王或富人都更受尊敬；然而他们是如何度过一生的，在饥饿和干渴中。教导她，生活中没有什么可畏惧的，除了得罪天主。如果有人这样结婚，带着这些看法，他将仅次于隐修士；已婚者也仅次于未婚者。\par

如果你有意设宴款待，就不要有任何不庄重、任何混乱的事。如果你能找到任何一位能祝福你家的贫穷圣人，只要他踏进你家的门，就能带来天主的全部祝福，就邀请他。我还要说另一件事吗？你们中别有人努力娶一个富有的女子，而宁可娶一个贫穷的。当她进门时，她不会从她的财富中带来如此大的快乐，反而会从她的嘲骂、从她要得更多、从她的傲慢、挥霍和烦人的言语中带来烦恼。因为也许她会说：\Quote{我还没有花过你任何东西，我仍然穿着我自己的衣服，是用我父母给我的嫁妆买的。}你说什么，妇人？还穿着你自己的？还有什么比这话更可怜的呢？怎么，你不再有自己的身体，却有自己的钱吗？结婚后你们不再是两个人，而是成为一体，那时你们的财产却是两样，而不是一样？哦！这爱钱之心！你们两人已成为一个人，一个活物，而你仍然说\Quote{我的}？这个可咒骂可恶的词，是魔鬼带来的。天主把比这些更亲近、更亲爱的东西都使我们共享，难道这些就不该共享吗？我们不能说\Quote{我的光，我的太阳，我的水}：我们所有更大的祝福都是共享的，而财富就不共享吗？愿财富消亡一万次！或者不如说，不是财富，而是那些不知道如何使用财富，却把它们看得高于一切的性情。\par

把这些教训连同其他的也教导她，但要十分亲切。因为既然对美德的推荐本身有很多严厉之处，尤其是对于年轻柔弱的少女，每当要讲论真智慧时，你要设法使自己举止充满恩宠和善意。最重要的是，从她灵魂中消除\Quote{我的和你的}这种观念。如果她说\Quote{我的}这个词，就对她说：\Quote{你称什么是你的？因为说实话，我不知道；我本人没有一样是我自己的。那么你怎么说\InnerQuote{我的}呢，既然一切都是你的？}要慷慨地让她说这个词。你没有觉察到我们对孩子就是这样的做法吗？当我们拿着某样东西，孩子抢夺它，并且还想拿到别的东西时，我们允许他，并说：\Quote{是的，这是你的，那也是你的。}同样，让我们也对妻子这样做；因为她的性情或多或少像孩子；如果她说\Quote{我的}，就说：\Quote{为什么，一切都是你的，连我也是你的。}这说法不是奉承，而是极大的智慧。这样你将能平息她的怒气，消除她的失望。因为如果人出于邪恶的目的做出不配的行为，那才是奉承；而这却是至高的哲学。那么就说：\Quote{连我自己也是你的，我的孩子；这个劝告是\Person{保禄}给我的，他说：\BibleQuote{丈夫没有权柄支配自己的身体，而是妻子。}\BibleRef{格前 7:4}如果我对自己的身体没有权柄，而你却有，那么你对我的财产就更有权柄了。}说了这些话，你就能使她安静下来，扑灭火焰，羞辱魔鬼，使她比用钱买来的奴隶更成为你的奴隶，用这些话你就能牢牢绑住她。这样，通过你自己的言语，教导她永远不要说\Quote{我的和你的}。还有，不要只叫她的名字，而要用亲昵的称呼，带着尊重，带着许多爱。尊重她，她就不需要别人的尊重；如果她享有你给予的荣耀，她就不会想要来自别人的荣耀。在一切事上，为她所有的优点，无论美貌还是智慧，都把她放在首位，并赞美她。这样你会说服她不理会外面的人，而蔑视全世界只在乎你。教导她敬畏天主，那么所有美善的事物就会像从泉源一样涌流出来，房屋将充满无数的祝福。如果我们寻求不朽的事物，这些可朽坏的事物也会随之而来。因为祂说：\Quote{因为，}祂说，\Quote{你们先寻求祂的国，这一切都将加给你们。}\BibleRef{玛 6:33}你想，这样的父母的孩子会是什么样的人？这样的主人的仆人是什么样？所有亲近他们的人呢？他们最终不也会获得数不清的祝福吗？因为一般来说，仆人也按主人的品行被塑造，被他们的性情所影响，爱同样的对象，学他们说同样的话，与他们从事同样的追求。如果我们这样规范自己，并专心研读圣经，在大部分事情上我们都会从中得到教导。这样我们将能中悦天主，在今生整个生活中行善，并达到那些为爱祂之人所应许的福分。愿天主恩赐我们众人堪当得到这一切，藉着我们的主耶稣基督的恩宠和慈爱，祂与圣神、圣父一起，现在、永远、直到万世，享有荣耀、权能和尊崇。阿们。\par

\SourceRecord{Translated by Gross Alexander.；Nicene and Post-Nicene Fathers, First Series；Vol. 13.；Edited by Philip Schaff.；Buffalo, NY: Christian Literature Publishing Co.,；1889.；<http://www.newadvent.org/fathers/230120.htm>.}

% structure: p0001 source=h1 target=section reason=page-title

\section{厄弗所书讲道集 第二十一篇}

\BibleRef{弗 6:1-3}\par

\BibleQuote{子女们，你们要在主内服从你们的父母，因为这是理所当然的。要孝敬父亲和母亲（这是第一条带应许的诫命），为使你蒙福，并在地上长寿。}\par

正如人在塑造身体时，先安放头，然后是颈，再后是脚，蒙福的保禄在他的讲论中也是如此进行。他已论及丈夫，也论及妻子（她是次位掌权者），现在他逐步进展到第三等级——即子女的等级。因为丈夫对妻子有权威，而丈夫和妻子对子女有权威。现在请看他在说什么。\par

\BibleQuote{子女们，你们要在主内服从父母，因为这是第一条带应许的诫命。}\par

这里他没有一句关于基督的讲论，也没有一句高深的主题，因为他仍是向柔弱的理解力发言。而且正因如此，他使自己的劝诫简短，因为孩子们不能跟从长篇的论述。为此，他也完全没有论及天国（因为这并不属于幼年的理解力所能明白的），而是说出了孩童的灵魂最特别渴望听到的话，即它将\BibleQuote{长寿}。若有人要问，为什么他省略了关于天国的论述，而向他们提出律法中所规定的诫命，他这样做是因为他对孩童说话，并且他深知如果丈夫和妻子都按他所制定的律法来安排，那么确保子女的顺从就不会有太多麻烦。因为每当一件事有良好、健全、有序的原则和基础时，此后的一切便会顺理成章地、有条不紊地、十分容易地进行：更困难的是奠定基础，建立稳固的根基。他说：\BibleQuote{子女们，你们要在主内服从父母}，也就是，按照主的意思。他的意思是，这是天主命令你们的。但倘若他们命令愚蠢的事呢？一般来说，一个父亲，即使他自己多么愚蠢，也不会命令愚蠢的事。不过，即使那样，宗徒也已经防备了，他说了\BibleQuote{在主内}；也就是说，只要你们不因此而得罪天主。所以，如果父亲是外邦人或异端者，我们就不应再服从，因为那时命令不是\BibleQuote{在主内}。但他说\Quote{哪一条是第一条诫命}？这是怎么回事？因为第一条是：\Quote{不可奸淫——不可杀人}。那么，他说这诫命是第一条，并不是在顺序上，而是在应许上。因为对于其他诫命，并没有附加上赏报，因为它们是被制定来对付恶事以及远离恶事的。而在这些诫命中，既有善行的实践，还进一步附加了应许。请注意，他为德行的道路奠定了多么卓越的基础，即对父母的尊荣和敬意。当他引导我们远离恶行，正要进入德行时，这是他所命令的第一件事：对父母的敬意；因为父母在一切人中，仅次于天主，是我们存在的根源，所以他们理应首先收获我们正确行为的果实；然后才是其余所有的人。因为如果一个人不尊敬父母，他永远不会对与自己无关的人温和。\par

然而，在给了子女必须的训诫后，他转向父亲们，说：\par

第4节。\BibleQuote{你们作父亲的，不要惹你们的子女发怒；却要在主的管教和训诫中养育他们。} \BibleRef{弗 6:4}\par

祂不是说：\Quote{爱他们}，因为本性甚至违背自己的意志也会促使他们这样做，就此立定法律是多余的。但祂说什么？\Quote{不要激怒你们的孩子}，正如许多人通过剥夺继承权、不认他们、以及专横地对待他们，不把他们当作自由人而当作奴隶所做的那样。所以祂说：\Quote{不要激怒你们的孩子。}然后，祂指出了最重要的一点，说明他们将如何被引导到顺从，将一切源头归之于首领和主要权威。正如祂已表明丈夫是妻子服从的原因（这也是为什么祂将大部分论证针对丈夫，劝他通过爱的力量将她与自己结合），同样，我指出，在这里祂也将效力归于丈夫，说：\Quote{但要照主的管教和劝诫养育他们。}你看，哪里有了属灵的纽带，本性的纽带就会随之而来。你希望你的儿子顺从吗？从一开始就\Quote{照主的管教和劝诫养育他}。绝不要认为他勤奋聆听圣经是不必要的事情。因为在那里他首先听到的就是：\Quote{孝敬你的父亲和你的母亲}；所以这对你是有益的。绝不要说：这是隐修士的事。我要使他成为隐修士吗？不。他不需要成为隐修士。为什么如此害怕一件充满如此多益处的事？使他成为一个基督徒。因为在一切事上，对平信徒来说，尤其是对孩子们来说，熟悉来自这一源头的教训是必要的。因为他们的年龄充满愚昧；而在这愚昧之上，又加上了来自异教故事的坏榜样，在那里他们认识了那些他们中间如此钦佩的英雄，那些受制于情欲、对死亡怯懦的人；例如，阿基里斯，当他心软时，当他为他的妾而死时，当另一个人酗酒时，以及许多诸如此类的事。因此，他需要对付这些事的补救方法。把孩子送出去学手艺、上学、为这些目标尽一切努力，却不\Quote{照主的管教和劝诫养育他们}，这难道不荒谬吗？为此，我们确实是首先收获果实的人，因为我们把孩子们培养成傲慢、放荡、不顺从、纯粹的粗俗家伙。所以我们不要这样做；相反，让我们听从这位有福宗徒的劝告。\Quote{让我们照主的管教和劝诫养育他们。}让我们给他们一个榜样。让我们使他们从最年幼时就专心于阅读圣经。唉，我如此不断地重复这一点，却被人看作是在说废话！但我仍不会停止尽我的本分。告诉我，你们为何不效法古时的人？你们妇女，尤其是要效法那些可钦佩的妇女。有人生了孩子吗？效法亚纳的榜样\BibleRef{撒上 1:24}；看看她做了什么。她立刻把他带到圣殿。你们中间谁不愿意他的儿子成为撒慕尔，而不愿意他成为整个世界的国王一万倍呢？\Quote{怎么，}你会说，\Quote{他可能成为这样的人呢？}为什么不可能？因为你自己不选择这样，也不把他托付给那些有能力使他成为这样的人的人。\Quote{那谁，}将会说，\Quote{是这样的人呢？}天主。因为她把他交在天主的手中。因为连厄里自己也不是在很大程度上有资格塑造他的人；（他怎能是呢，他连自己的儿女都不能塑造？）不，是母亲的信德和她热切的热情成就了一切。他是她的第一个孩子，也是唯一的孩子，她不知道是否还会有别的孩子。然而她没有说：\Quote{我要等到孩子长大，让他尝一尝今生的事物，我容许他在童年时代稍微在其中嬉戏一番。}不，这女人摒弃了所有这些想法，她专注于一件事：如何从起初就把属灵的肖像奉献给天主。我们男人在这女人的智慧面前真该脸红。她把他献给了天主，并把他留在那里。因此，她的婚姻状态更为荣耀，因为她把属灵的对象作为首要的关怀，并把初果奉献给天主。因此，她的子宫结出果实，她又得了别的孩子。因此，她看见他在世间也受人尊敬。因为如果人受到尊敬，他们也会回报尊敬，难道天主不更会如此吗？祂即使没有受到尊敬也会这样做。我们要长久地做血肉之块吗？我们要长久地俯伏在地上吗？愿一切事物对我们来说都是次要的，重要的是我们要照顾孩子，以及\Quote{照主的管教和劝诫养育他们}。如果从一开始他就被教导成为真正智慧的热爱者，那么他所获得的财富就大于一切财富，并有一个更威严的名声。你通过教他一种技艺、给他那种外在的学问，使他能获得财富，远不如你教他轻视财富的艺术那样有成就。如果你想使他富有，就这样做。因为富足的人不是那渴望巨大财富、被巨大财富环绕的人，而是那什么都不缺的人。这样训练你的儿子，这样教导他。这是最大的财富。不要寻求如何给他外在学问的名声和高地位，而要深思如何教导他轻视属于今生的荣耀。通过这种方式，他会变得更加杰出、更加真正光荣。穷人和富人都能完成这一点。这些教训不是人从老师那里学到的，也不是通过艺术学到的，而是通过神圣的谕旨。不要寻求他如何在此世享长寿，而要寻求他如何在以后享受无边无际的永生。给他伟大的事物，而不是微小的事物。听保禄说什么：\Quote{照主的管教和劝诫养育他们}；不要致力于使他成为演说家，而要训练他成为哲学家。缺乏前者不会有任何害处；缺乏后者，世间所有的修辞学都无济于事。需要的是性情，而不是口才；是品格，而不是机敏；是行为，而不是言语。这些使人获得天国。这些给予真正的益处。不要磨利他的舌头，而要洁净他的灵魂。我这样说并不是要阻止你教他这些事，而是要阻止你只关注它们。不要以为只有隐修士才需要这些来自圣经的教训。恰恰相反，即将进入世界的孩子们特别需要它们。因为正如那总是停泊在港湾里的人并不需要装备他的船、不需要领航员和水手，而那总是在海上的人则需要；同样，世俗之人和隐修士也是如此。前者仿佛进入了无浪的港湾，过着无忧无虑的生活，远离一切风暴；而后者却总是在海洋上，生活在海洋中央，与无数的风浪搏斗。\newline{}\par

假使他本人或许不需要，但他仍应如此准备，以便堵住他人的口。因此，他在现世生活中越是卓越，就越需要这种教育。如果他在宫廷中生活，便有许多外教人、哲学家以及为现世荣耀所膨胀的人。那里如同一个充满水肿病人的地方。宫廷在某种程度上就是这样。所有人仿佛都膨胀着，处于炎症状态。而那些并非如此的人，也正努力变成那样。那么请想一想，这会是多大的益处：你的儿子进入那里时，应当像一位出色的医生，带着能缓解每个人特有炎症的工具，走到每个人面前，与他交谈，运用从圣经中汲取的良药，倾吐真正哲学的论述，使患病的身体恢复健康。因为隐修士与谁交谈？与他的墙壁和天花板吗？或者与旷野和树林？与飞鸟和树木？因此，他并不那么需要这种训练。然而，他仍致力于完善这项工作，与其说是为了训练他人，不如说是为了训练自己。那么，那些要混迹于现世的人就极需这种训练，因为这样的人比另一个人有更强的犯罪诱惑。如果你愿意理解，他在世本身也将是一个更有用的人。因为所有人都会因这些话而尊敬他，当他们看见他在火中却不被灼烧，且不贪求权力。但当他最不渴望权力时，他反而会获得权力，并成为国王更尊敬的对象；因为这样的品德是不可能隐藏的。在一群健康人中，一个健康的人确实不会被注意；但当一群病人中有一个健康人时，消息很快就会传开，传到国王耳中，国王就会使他统治许多民族。因此，知道了这些事，\BibleQuote{以主的管教和劝诫养育你们的儿女。}\par

\Quote{但是假设一个人是穷人。}他绝不会比那些住在王宫里的人更卑微，因为他不在王宫里；不，他将受人钦佩，并很快获得那种自愿而非强迫的权威。因为如果一群像狗一样毫无价值的希腊人，拿起他们那毫无价值的哲学（因为希腊哲学就是如此），或者说不是哲学本身，而仅仅是它的名号，穿着破旧的外衣，留着长发，就使许多人印象深刻；那么一个真正的哲学家会给人留下多么深的印象呢？如果一种虚假的外表，一种哲学的影子，第一眼就如此吸引我们，如果我们爱真正纯正的哲学，又该如何呢？难道不会所有人都追求它，并完全信任地将房屋、妻子、儿女托付给这样的人吗？但现在没有，不，根本没有这样的哲学家存在。因此，不可能找到这类例子。在隐修士中，确实有这样的人，但在世俗人中不再有了。至于隐修士中有这样的人，可以举出许多实例。然而，我要从许多中提及一个。你们无疑知道、听说过，或许有些人还见过我现在将要提到的那个人。我指的是可敬的\Person{尤利安}。这个人是个乡下人，生活卑微、家世微贱，对于一切外在的才艺完全未受教育，却充满朴实无华的智慧。当他进城时（这很少发生），从未有过这样的聚集，即使演说家、诡辩家或任何其他人进城时也不曾有过。但我在说什么呢？难道他的名字不比任何国王的名更荣耀，甚至直到今日仍被传颂吗？如果这些事发生在这个世界，在这个主没有应许我们任何好处、祂告诉我们我们是客旅的世界，让我们想一想，为我们存留在天上的福分是多么伟大。如果他们在作客旅的地方享有了如此大的荣耀，那么当他们在自己的家乡时，将享有多么大的光荣！如果在祂应许患难的地方，他们得到了如此细心的照顾，那么在那祂应许真正尊荣的地方，他们的安息将是何等伟大！\par

现在，你愿意我向你展示世俗之人的榜样吗？目前，我们确实没有这样的人；但或许也有优秀的世俗之人，尽管他们并未达到最高的哲学境界。因此，我将从古代的圣人中为你引述例子。有多少人有妻子要照顾、有儿女要抚养，却丝毫不逊于那些被提及的人？然而，如今已不再如此，\BibleQuote{为了目前所处的困境} \BibleRef{格前 7:26} ，正如这位有福的宗徒所说。那么，你愿意我提及谁呢？是\Person{诺厄}，还是\Person{亚巴郎}？是\Person{亚巴郎}的儿子，还是\Person{依撒格}？或者，\Person{若瑟}？又或者，你愿意我转向先知们？我指的是\Person{梅瑟}，或是\Person{依撒意亚}？不过，如果你同意，让我们将话题转向\Person{亚巴郎}，所有人都不断地将他置于众人之上。他没有妻子吗？他没有儿女吗？有的，因为我也对你使用同样的言语，正如你对我所说。他有妻子，但他并非因有妻子而如此卓越。他有财富，但他并非因有财富而中悦天主。他生了儿女，但他并非因生儿女而被称有福。他家里有三百一十八个家生的仆人，但他并非因此被视为可敬。但你想知道原因吗？那是因为他的款待，因为他轻视财富，因为他克己的行为。因为，请告诉我，哲学家的职责是什么？难道不是轻视财富与荣耀吗？难道不是超越嫉妒和一切激情吗？那么，现在，让我们把他带出来，剥去他的外衣，向你展示他是怎样的哲学家。首先，他视自己的故乡为无物。天主说：\BibleQuote{你离开你的故乡、你的家族} \BibleRef{创 12:1} ，他立刻出发了。他并未被家庭所束缚（否则他绝不会离开），也未受对亲友之爱的羁绊，更未受任何其他事物的束缚。但还有什么呢？他比任何人都轻视荣耀与金钱。因为当他击退敌人结束战争后，有人请他取走战利品，他却拒绝了。\BibleRef{创 14:21}\par

同样，这位伟人的儿子受到尊敬，并非因他的财富，而是因他的款待；并非因他的儿女，而是因他的服从；并非因他的妻子，而是因他妻子所受的不育之苦。\BibleRef{创 25:21}\par

他们认为今世生命算不得什么，不追求利益，轻视万有。请问，哪一种植物是最好的？难道不是那些自身有力、不因雨水、冰雹、狂风或任何此类变故所伤，反而赤身挺立，无需墙壁或篱笆保护的植物吗？真正的哲学家就是这样，我们所说的财富也是这样。他一无所有，却拥有一切；他拥有一切，却一无所有。因为篱笆不在内，而在外；墙壁不是本性的，只是从外建造的。再者，请问，什么样的身体是强壮的？难道不是健康的、不因饥饱、寒冷或酷热而被打倒的身体吗？还是那种为了健康需要厨师、织工、猎人和医生的身体？他是富人，是真正的哲学家，不需要这些东西。为此，这位蒙福的宗徒说：\Quote{将他们在主的惩戒和劝诫中养育起来。}不要用外在的防御围护他们。因为财富如此，荣耀也是如此；当这些跌倒时——它们确实跌倒——那植物便赤身无蔽，不仅在过去的日子里毫无得益，反而受害。因为那些曾使它不惯于抵御风雨的遮蔽，如今却使它一下子灭亡。因此，财富反而有害，因为它使我们对人生的变故毫无训练。所以让我们训练我们的孩子，使他们能承受一切考验，不为将来可能发生的事情所惊；\Quote{让我们在主的惩戒和劝诫中养育他们。}为我们积攒的赏报将是巨大的。如果人们因为制作帝王的雕像和画像而享有如此的尊荣，那么那些装饰万王之王的肖像的人——因为人是天主的肖像——若能达成真实的相似，岂不领受千万的祝福？因为相似在于灵魂的德行，当我们训练孩子成为善良、温良、宽恕的人（因为这些都是天主的属性），成为仁慈、有人性的人，训练他们视今世如无物时。让这成为我们的任务，塑造和引导我们自己和他人，使之合乎正义。否则，我们有何胆量站在基督的审判台前？如果一个人有不服管教的子女便不配做主教\BibleRef{铎 1:6}，更何况是天国呢？你说什么？如果我们有不服管教的妻子或子女，我们也要交账吗？是的，如果我们不准确尽责地完成自己的本分——因为单靠我们个人的德性不足以得救。如果那个埋没一塔冷通的人一无所获，反而受到如此惩罚，显然单靠个人的德性不足以得救，还需要他人的德性。因此，让我们对妻子怀有极大的关怀，细心照料我们的孩子、仆人和我们自己。在我们管理自己和他们的过程中，让我们祈求天主帮助我们。如果祂看到我们对此事关心和忧虑，祂必帮助我们；但如果祂看到我们漠不关心，祂就不会伸出援手。因为祂不会在我们睡觉时赐予我们援助，只有当我们自己也努力工作时，祂才帮助我们。因为帮助者（顾名思义）不是帮助懒惰的人，而是帮助自己也工作的人。但良善的天主能亲自完成这项工作，使我们众人堪当获得所应许的祝福，藉着祂唯一圣子的恩宠和慈悲，与圣神一起归于圣父的荣耀、权能和尊威，从今世直到永远，万世无疆。阿们。\par

\SourceRecord{Translated by Gross Alexander.；Nicene and Post-Nicene Fathers, First Series；Vol. 13.；Edited by Philip Schaff.；Buffalo, NY: Christian Literature Publishing Co.,；1889.；<http://www.newadvent.org/fathers/230121.htm>.}

% structure: p0001 source=h1 target=section reason=page-title

\section{论\WorkTitle{厄弗所书}第二十二讲}

\BibleRef{弗 6:5-8}\par

\Quote{奴仆们，要战战兢兢，以诚实的心，听从你们肉身的主人，如同听从基督一样；不可只作眼前事奉，像是取悦于人，而要作基督的仆人，从心里遵行天主的旨意；甘心事奉，如同事奉主，而非事奉人，因为你们知道，各人无论作了什么善事，都会从主那里得到赏报，不论是奴仆，或是自主的。}\par

这样看来，不仅是丈夫、妻子、儿女，连善德的奴仆也有助于一个家庭的建立和保护。因此，有福的保禄没有忽略这一部分。不过他是在最后才谈到它，因为它在尊贵和等级中居末位。但他仍向他们讲了许多道理，不再像对儿女那样采用相同的语调，而是更加深奥，因为他对这些人的应许不是在这个世界，而是在将来的世界。\Quote{知道}，他说，\Quote{各人无论作了什么善或恶，都会从主那里得到报应}，这样立刻教导他们热爱智慧。因为他们虽然在尊贵上次于儿女，但在心智上却高于他们。\par

\Quote{奴仆们}，他说，\Quote{要听从你们肉身的主人。}\par

因此他立刻鼓舞、安抚受伤的心灵。他似乎说：不要因你次于妻子和儿女而忧伤。奴役不过是个名号。主从关系是\Quote{按照肉身}，短暂且暂时的；因为凡属于肉体的，都是易逝的。\par

\Quote{恐惧}，他又补充说，\Quote{和战栗。}\par

你看到他没有要求奴仆对主人有如同妻子那样的恐惧吗？因为在那情况下他只是简单地说：\Quote{妻子也当敬畏她的丈夫}；而在这里他加重了语气，说：\Quote{恐惧}，他说，\Quote{和战栗，以诚实的心，如同对基督。}这是他一贯说的。有福的保禄，你这是什么意思？他是兄弟，或者不如说他已经成为兄弟，享有同样的特权，属于同一身体。不仅如此，他不只是他主人的兄弟，也是天主圣子的兄弟，他分享所有同样的特权；而你却说：\Quote{按照肉身，以恐惧和战栗服从你们的主人}？是的，正因为这个原因，他说，我才这样说。因为我若命令自由人因敬畏天主而彼此服从——正如他上面说的：\Quote{要因敬畏基督而彼此服从}；——我若还命令妻子敬畏和尊敬她的丈夫，尽管她与他是平等的；那么，我岂不更应对奴仆这样说？这并非出身低微的记号，反而是最真实的尊贵，懂得如何谦卑自己，谦虚谨慎，向邻人让步。而且自由人曾怀着极大的恐惧和战栗服侍自由人。\par

\Quote{以诚实的心}，他说。\par

他说得好，因为人可能怀着恐惧和战栗服侍，却并非出于善意，而是以任何可能的方式。许多奴仆常常暗地里欺诈他们的主人。因此他消除这种欺诈，说：\Quote{以诚实的心如同对基督，不可只作眼前事奉，像是取悦于人，而要作基督的仆人，从心里遵行天主的旨意；甘心事奉，如同事奉主，而非事奉人。}你看到他为了植入这个好原则，说了多少话？我指的是\Quote{善意}和\Quote{从心里}。至于另一种服侍——我指的是\Quote{恐惧和战栗}——我们看到许多人向他们的主人提供，主人的威胁大大促进了这一点。但是，他说，要表明你服侍是\Quote{基督的仆人}，而不是人的仆人。使这正当行为成为你自己的，而非出于强迫。正如在接下来的话中，他劝说并教导那个被他人恶待的人，使正当行为成为他自己的，是他自由选择的行动。因为既然那个打脸颊的人，并不是因为被打者的任何意图而做出那行动，只是出于他自己的恶意，主说什么？\Quote{把另一边也转给他}\BibleRef{玛 5:39}；向他表明，你在承受第一下时并不是不情愿的。因为那慷慨忍受冤屈的人，通过让自己也被打另一边，而不是仅仅忍受第一下，使得那不属于他的行动成为他自己的。因为后者也许甚至显得怯懦；而前者才是高尚的哲学。——这样你就将表明，你承受第一下也是出于智慧。所以在当前情况下，也要表明，你甘愿承受这奴役。那么，取悦于人者就不是基督的仆人。基督的仆人不取悦于人。\BibleRef{迦 1:10}因为谁是事奉天主的，以取悦人为目的？谁取悦于人，又能事奉天主？\par

\Quote{从心里}，他说，\Quote{怀着善意服务}。因为既可能以纯一的心、并非不正当地服务，却又不尽全力，只尽本分，所以他吩咐要欣然去做，不是出于勉强，而是出于原则，不是出于强制。你若这样服务，就不是奴隶；若是出于原则，若是怀着善意，若是从心里，且是为基督的缘故。因为就是那自由人保禄也这样服务，并高呼说：\Quote{我们并不是宣传我们自己，而是宣传耶稣基督为主，并因耶稣的缘故作你们的奴仆。}\BibleRef{格后 4:5}你看他是如何抹去你奴役中的卑贱。正如一个被抢的人，若他反而更多地给与那抢夺者，他就不被视为被抢者，而被视为慷慨的施与者；不算是受害者，而是行善者；并且反而以自己的慷慨使对方蒙羞，而非因被抢而蒙羞——同样，我说，在这种情况下，他的慷慨会立刻显出他更高尚，并且因表明自己不受伤害而令对方羞愧。\par

那么，我们要为基督的缘故服事我们的主人，\Quote{知道}，他接着说，\Quote{各人无论行了什么善事，都必从主那里得到报应，无论是为奴的或是自主的。}因为很可能许多主人作为不信者，不知羞耻，也不回报奴隶的服从，请注意他是如何鼓励他们，使他们对于报酬不存疑虑，而充分相信那赏报。因为正如那些受惠者若不回报，便使天主成为施惠者的债务人一样，我说，主人也是这样，若他们受了好待却不回报，反因使天主成为你的债务人而更报答你。\par

第9节。\BibleRef{弗 6:9}\Quote{你们作主人的，}他接着说，\Quote{也要照样对待他们。}\par

照样。这是什么？\Quote{怀着善意服务}。不过他并没有实际说\Quote{服务}，但藉着说\Quote{照样}，他清楚表明这是他的意思。因为主人自己也是仆人。\Quote{不是讨人喜欢的，}他意指，\Quote{且存着敬畏和战栗}：即对天主，恐怕祂有一天因你们对奴隶的疏忽而指责你们。\par

\Quote{且不可威吓；}不要激怒他们，他意指，也不要压迫他们。\par

\Quote{知道他们和你一样的主人在天上。}\par

啊！他在这里暗示了何等大能的主！这提示何等惊心！是这样的：\Quote{你们用什么量器量给人，也必用什么量器量给你们。}\BibleRef{玛 7:2}免得你听到这句话：\Quote{你这恶仆！我免了你那一切债。}\BibleRef{玛 18:32}\par

\Quote{祂并不偏待人。}他说。\par

不要以为，他要说，向仆人做的事，祂会因此宽恕，因为是对仆人做的。外邦人的法律，作为人的法律，的确承认这类罪行之间有差别。但众主的共同主的法律，因为同等地善待众人，赐予众人同样的权利，并不认这样的差别。\par

但若有人问，奴役从何而来，它为何进入人类生活（我知道许多人既喜欢问这样的问题，又渴望得到解答），我将告诉你们。奴役是贪婪、堕落、野蛮的果实；因为我们知道，诺厄没有仆人，亚伯尔也没有，舍特也没有，他们之后的人也没有。这东西是罪的果实，是违背父母的果实。愿儿女们听这话：每当我们不孝敬父母，他们就该做仆人。这样的孩子剥夺了自己的尊贵出身；因为那悖逆他父亲的，就不再是儿子；那悖逆自己父亲的若不是儿子，那么悖逆我们真父的怎能算是儿子呢？他已离弃了自己的尊贵出身，侮辱了本性。接着就来了战争、战斗，并俘虏了他们的战俘。可是，你会说，亚巴郎有仆人。是的，但他并不像使用仆人那样使用他们。\par

请注意万事如何取决于头：妻子，藉着告诉他\Quote{要爱她}；儿女，藉着告诉他\Quote{要在主的管教和劝诫中养育他们}；仆人，藉着这话：\Quote{知道他们和你一样的主人在天上。}所以他说道，你们也要照样，因为是仆人，要仁慈宽厚。\Quote{最后，你们要在主内，并依靠祂强有力的德能，坚强起来。}\par

但若在考虑这下一件事之前，你们愿意听，我将就仆人发表同样的意见，如同我先前就儿女所说的一样。教导他们虔诚，其余一切自会随之而来。但现在，当有人要去剧场或去浴场时，他拖着他所有的仆人跟在后面；但当他去教堂时，片刻也不曾；他也不强迫他们参加并聆听。现在，当你的主人正忙于别的事情时，你的仆人如何聆听呢？你买了、你赎了你的奴隶吗？在万事先吩咐他天主所愿他做的事：要温良对待同伴仆人，并重视德行。\par

每个人的家都是一座城；每个人在自己的家中都是君王。富人的家具有这种性质，这是显而易见的，那里有田产、管家和管家的管家。但我要说，穷人的家也是一座城。因为这里也有权力的职分：例如，丈夫对妻子有权威，妻子对仆人有权威，仆人又对自己的妻子有权威；再者，妻子和丈夫对儿女有权威。岂不是在你看来，他好像一位君王，拥有如此众多的权威在自己的权下吗？难道他不应当比其余所有人更善于管理家事和国事吗？因为懂得如何按各自的职务管理这些人的人，也必知道如何选拔最合适的人担任职分，并且会选出杰出的人。如此，妻子在家中是第二位君王，只缺少冠冕；而懂得如何选择这位君王的人，必能卓越地管理其余一切。\par

第十节。他说：\Quote{末了，你们要因主刚强}。\BibleRef{弗 6:10}\par

每当讲论即将结束时，他总是使用这样的转折。我从一开始岂不是说得很好吗？每一个人的家本身就是一座军营。因为看，他处理了各种职务之后，接着就武装他们，领他们上战场。若是没有人侵占别人的职务，而是各守其位，一切就都会井然有序。\par

他说：\Quote{你们要因主刚强，并因祂大能的力。}\par

就是说，靠着我们对祂所怀的希望，借祂的助佑。因为他已命令了许多必须履行的职责，不要害怕，他似乎在说，把你的希望寄托于主，祂必使一切容易。\par

第十一节。他说：\Quote{穿上天主的全副\OriginalTerm{武装}{armor}，使你们能抵挡魔鬼的\OriginalTerm{诡计}{wiles}。}\BibleRef{弗 6:11}\par

他没有说\Quote{抗拒战斗}或\Quote{抗拒敌意}，而是\Quote{抗拒\InnerQuote{诡计}}。因为这位仇敌与我们作战，不是简单的，也不是公开的，而是通过\Quote{诡计}。诡计是什么意思？使用\Quote{诡计}，就是欺骗，用巧计或计谋来擒拿；这既发生在技艺上，也发生在言语、行动和策略上，在那些诱惑我们的人中。我的意思是这样：魔鬼从不以本来的颜色向我们提出罪过；他不提偶像崇拜，而是用另一种装扮来提出，使用\Quote{诡计}，即他的言辞貌似有理，运用伪装。因此，宗徒现在既以此激励士兵，又使他们警醒，说服并教导他们：我们的战斗是与一位精通战争技艺的敌人作战，他作战不是简单、直接的方式，而是极其阴险。首先，他借魔鬼的技巧来唤醒门徒；其次，借其本性及其势力的多寡。他并不是要使他麾下的士兵灰心，而是要激发并唤醒他们，才提到这些策略，并预备他们保持警醒；因为如果他仅仅详述他们的力量，然后就止住话语，必定会使他们灰心。然而现在，他在此之前和之后都表明，战胜这样的仇敌是可能的，这反而提升了他们的勇气；因为越是清楚地向我方人民陈述我们敌人的力量，就越会使我们的士兵变得认真。\par

第十二节。他说：\Quote{因为我们战斗不是对抗血和肉，而是对抗\OriginalTerm{率领者}{principalities}，对抗\OriginalTerm{掌权者}{powers}，对抗这\OriginalTerm{黑暗世界的霸主}{world rulers of this darkness}，对抗\OriginalTerm{天界里邪恶的属灵军旅}{spiritual hosts of wickedness in the heavenly places}。}\BibleRef{弗 6:12}\par

他用战斗的特性激励了他们之后，接着又用摆在他们面前的奖赏来唤醒他们。他的论点是什么？他说了敌人凶猛之后，又补充说，他们剥夺我们巨大的福分。这些是什么？战斗是在\Quote{在天上的事物}之中；斗争不是关于财富，不是关于荣耀，而是关于我们是否被奴役。因此敌意是不可调和的。当涉及重大利益时，争斗和冲突更为激烈；因为\Quote{在天上的事物}这个说法等于\Quote{为了天上的事物}。敌人并不是为了通过征服获得什么，而是为了剥夺我们。就好像有人说：\Quote{合同在哪里？}在金子里。\Quote{在}这个词的意思是\Quote{为了}；\Quote{在}这个词也有\Quote{因}的意思。看，敌人的力量如何使我们震惊；它使我们事事谨慎，知道冒险是为了重大利益，胜利是为了巨大回报。因为他正在竭力将我们从天上赶出去。\par

他提到某些\Quote{率领者、掌权者和这黑暗世界的霸主}。什么黑暗？是黑夜的吗？不，是邪恶的。因为他说：\Quote{你们从前曾是黑暗}，\BibleRef{弗 5:8}；如此称那属于现世生活的邪恶；因为在此世之外，它没有地方，不在天上，也不在来世。\par

他称他们为\Quote{世界统治者}，不是因为他们主宰世界，而是因为圣经惯常称恶行作\Quote{世界}，例如，基督说：\Quote{他们不属于世界，如同我不属于世界一样。}\BibleRef{若 17:16} 那么，难道他们不属于世界吗？难道他们没有穿上肉体？难道他们不是那些在世界里的人吗？又说：\Quote{世界恨我，却不能恨你们。}\BibleRef{若 7:7} 这里他又用这名称称恶行。因此，宗徒在这里用世界指恶人，而恶灵对他们有特别大的权能。他说：\Quote{对抗天界里邪恶的属灵军旅。}他提到\Quote{率领者和掌权者}；正如在天上有\Quote{宝座和宰制、率领者和掌权者}。\BibleRef{哥 1:16}\par

第十三节。他说：\Quote{因此，拿起天主的\OriginalTerm{全副武装}{whole armor of God}，使你们能在邪恶的日子站立得住，成就了一切，还能站立。}\BibleRef{弗 6:13}\par

所谓\Quote{邪恶的日子}，他是指现世的生命，又称之为\Quote{这邪恶的现世}\BibleRef{迦 1:4}，源于其中所行的邪恶。这无异于说：要常备不懈。他又说：\Quote{做成了一切}；也就是说，既克服了情欲、卑鄙的贪恋，也克服了所有困扰我们的事。他所说的不仅是做出行动，而是彻底完成，不仅要杀戮，还要在杀戮之后站立。因为许多人虽已赢得这场胜利，却又跌倒了。他说：\Quote{做成了这一切}，不是做了一件而没做另一件。因为即使在胜利之后，我们还必须站立。敌人可能被击中，但被击中的东西如果我们不站立，它们会重新复活。但如果它们跌倒后又再起来，只要我们站立，它们就仍倒在地上。只要我们不摇摆，仇敌就不会再起。\par

\Quote{让我们穿上天主的全副武装。}你看见他如何驱除一切恐惧吗？因为如果可能\Quote{做成一切并站立}，他详细描述敌人的力量并不会造成怯懦和恐惧，反而会驱除懈怠。他说：\Quote{使你们能在邪恶的日子抵挡得住}。他还从时间上给予他们鼓励；他似乎说：时间短暂；所以你们必须站立；当杀戮完成时，不要灰心。\par

论道德。如果这是一场战争，如果摆在我们面前的敌军如此强大，如果\Quote{\OriginalTerm{率领者}{principalities}}是无形的，如果他们是\Quote{\OriginalTerm{世界的统治者}{rulers of the world}}，如果他们是\Quote{\OriginalTerm{邪恶的鬼神}{spiritual hosts of wickedness}}，那么，请告诉我，你怎么能过着自我放纵的生活？你怎么能放荡不羁？如果我们赤手空拳，怎能战胜？愿每人每天对自己重复这些话，每当他在愤怒或情欲的影响下，每当他追求这种萎靡的生活而徒劳无功时。愿他聆听蒙福的保禄对他说的话：\BibleQuote{我们的战斗不是对抗血肉，而是对抗率领者，对抗掌权者。}这是一场比感官之战更艰难的战争，更激烈的冲突。想想这个敌人角力了多久，它为何而战，并比以前更加警惕。有人会说：\Quote{可是，既然他是魔鬼，就该被清除掉，那样所有人都得救了。}这是你们中一些懒惰者出于自卫而发出的借口。人啊，如果你愿意，你本应感谢天主能战胜这样的敌人，但你反而心怀不满，说出一些懒散困倦的士兵的话。如果你愿意，你知道攻击点。从四面八方侦察，巩固你自己。这场冲突不仅是对抗魔鬼，也是对抗他的权势。那么，你可能会说，我们如何与黑暗角力？通过成为光。如何与\Quote{邪恶的鬼神}角力？通过成为善。因为恶与善对立，光驱散黑暗。但如果我们自己也是黑暗，我们必然会被俘。那么我们如何战胜它们？如果它们本性如此，我们通过选择成为那样，摆脱血肉，这样我们就能战胜它们。因为从前门徒很可能会受到许多迫害者，他（保禄）会说：\Quote{不要以为是与你们作战的他们。真正与你们作战的是在他们内工作的鬼神。我们的战斗就是对抗它们。}通过这些考虑，他提供了两点；他使他们自己更加勇敢，并让他们把怒气转向那些与他们作战者。为什么我们的战斗是对抗这些？因为我们也有一个无敌的盟友——圣神的恩宠。我们已被教导一种技艺，使我们可以对抗的不仅是人，而是鬼神。不，如果我们愿意，我们甚至根本不需要角力；因为我们选择，才有了争战，因为住在我们内的那一位的力量如此之大，以至于祂说：\BibleQuote{看，我已赐给你们权柄践踏蛇和蝎子，并胜过仇敌的一切能力。}\BibleRef{路 10:19}祂将一切权柄都赐给我们，既有角力的权柄，也有不角力的权柄。是因为我们懒惰，才不得不与他们角力；因为保禄并没有角力，请听他自己说的话：\BibleQuote{谁能使我们与基督的爱隔绝？是困苦吗？是窘迫吗？是迫害吗？是饥饿吗？是赤身露体吗？是危险吗？是刀剑吗？}\BibleRef{罗 8:35}再听他的话：\BibleQuote{天主快要将撒殚践踏在你们脚下。}\BibleRef{罗 16:20}因为他已将他制服；因此他也说：\BibleQuote{我因耶稣基督的名命你从她身上出来。}\BibleRef{宗 16:18}这不是角力者的话；因为角力者尚未得胜，而得胜者不再角力；他已经制服、俘虏了对方。同样，伯多禄也没有与魔鬼角力，而是做了比角力更好的事。在信友、服从者、慕道者身上，他们大大胜过魔鬼和它的权势。因此，蒙福的保禄也说：\BibleQuote{因为我们不是不知道他的诡计}\BibleRef{格后 2:11}，这尤其是他胜过他的方式；再听他的话：\BibleQuote{这也不足为怪——连他的差役也冒充正义的差役。}\BibleRef{格后 11:14}他是如此熟悉争战的每一个部分，没有什么能逃脱他。又说：\BibleQuote{因为不法之奥秘已经在运作。}\BibleRef{得后 2:7}\par

但我们的战斗是与之对抗；因为请再听他所说：\Quote{我深信，无论是天使、掌权者、现今的事、将来的事、异能、或其他受造之物，都不能使我们与基督的爱分离。}\BibleRef{罗 8:38} 他并非简单地说\Quote{与基督}，而是\Quote{与基督的爱}。因为有许多人虽然与基督结合，却并不爱祂。他说，你们不但不能说服我否认祂，甚至不能使我减少对祂的爱。如果天上的掌权者都不能做到这一点，还有谁能动摇他呢？然而，他这样说并非因为他们实际在尝试，而是出于假设；因此他也说：\Quote{我深信。} 所以他没有挣扎，但仍惧怕他的诡计；因为请听他所说：\Quote{我怕你们的心意被败坏，失去那向基督的纯一，就像蛇用它的狡猾欺骗了厄娃一样。}\BibleRef{格后 11:3} 诚然，你会说，但他在论及自己时也用了这话，如他所说：\Quote{我恐怕我传给别人，自己反而被弃绝。} 那么，你如何\Quote{深信没有人能分离你}呢？你岂不看出这话是出自谦卑和自抑吗？因为他已住在天上。由此他也说：\Quote{我虽不觉得自己有错}\BibleRef{格前 4:4}；又说：\Quote{我跑完了全程。}\BibleRef{弟后 4:7} 所以，魔鬼所设的障碍并非针对这些事，而是针对门徒们的利益。为什么呢？因为在这些事上，他并非独自作主，他们的意志也起作用。魔鬼在某些情况下占了上风；但即使在那里，也不是胜过了他，而是胜过了那些漫不经心之人的怠惰。确实，如果因懒惰或其他类似原因，他没有尽自己的本分，那么魔鬼就胜过了他；但如果他自己尽了一切努力，而他们却不服从，那么魔鬼胜过的不是他，而是他们的不顺服；疾病胜过的不是医生，而是病人的不驯服；因为当医生采取一切预防措施，而病人却破坏一切时，失败的是病人，不是医生。这样，魔鬼在任何情况下都没有胜过保禄。但在我们自己的情况下，我们能与之搏斗就是满足。因为对罗马人，他没有要求这件事，而是什么？\Quote{天主快要将撒殚踏碎在你们脚下。}\BibleRef{罗 16:20} 他为这些厄弗所人祈求：\Quote{那位能按着运行在我们里面的大能，充充足足地成就一切，超过我们所求所想的。}\BibleRef{弗 3:20} 搏斗的人仍然被抓住，但对他来说，没有跌倒就够了。当我们离开这个世界时，而不是之前，荣耀的胜利才会实现。例如，以某种邪恶的私欲为例。非凡的事是不起私欲，甚至扼杀它。然而，如果这不可能，那么即使我们不得不与之搏斗，并坚持到底，但如果我们离开时仍在搏斗，我们就是得胜者。因为这里的情况与摔跤手不同；在那里，如果你没有摔倒对手，你就没有得胜；但在这里，如果你没有被摔倒，你就得胜了；如果你没有被摔倒，你就摔倒了它；这有道理，因为在那里双方都争取胜利，当一方被摔倒时，另一方得冠冕；但在这里，并不如此，魔鬼是谋求我们的失败；那么当我不让他达成他的意图时，我就是得胜者。因为他不是要推翻我们，而是要使我们与他一同毁灭而急切。我已经是得胜者了，因为他已经被摔下，处于毁灭状态；他的胜利不在于自己被加冕，而在于造成我的毁灭；所以，虽然我没有推翻他，但如果我没有被推翻，我就得胜了。那么什么是光荣的胜利呢？那就是更进一步，将他踏在脚下，如同保禄那样，视今世之事为无物。我们也要效法他，努力超脱它们，决不让他抓住我们。财富、财产、虚荣，都给他可乘之机。这常常激起他，也常常激怒他。但何需搏斗？何需与他较量？参与搏斗的人结果不确定，他可能自己失败被擒。而将他踏在脚下的人，胜利是确定的。\par

那么，让我们将魔鬼的力量踏在脚下吧；让我们将我们的罪，即今世一切相关的事——愤怒、私欲、虚荣、一切情欲——踏在脚下；这样当我们离开到那世界时，我们就不会因背叛天主所赐给我们的力量而被定罪；因为这样我们也将获得未来的福分。但如果我们在这些事上不忠信，谁会把更大的事托付给我们呢？如果我们不能将一个已跌倒、已蒙羞、已受轻视、在我们脚下的人踏碎，父怎能赐给我们父的赏报呢？如果我们没有制伏一个如此被置于我们权下的人，我们有什么信心进入父的家呢？因为告诉我，假如你有一个儿子，他不顾你家中那些善意的人，却与那些使你烦恼的人、那些被逐出父家的人、那些在赌场消磨时光的人交往，并且他这样一直到底；他难道不会被剥夺继承权吗？很明显，他会。我们也一样；如果我们不顾那些取悦我们父、并被他派来管理我们的天使，却与魔鬼交往，我们必然会被剥夺继承权，愿天主禁止；但让我们与他进行我们必须进行的战争吧。\par

若有人有仇敌，若有人曾受他亏待，若有人被激怒，让他把一切愤怒、一切猛烈都聚集起来，倾倒在魔鬼的头上。在此，愤怒是好事，在此，怒气有益，在此，报复值得称赞，因为正如在外邦人中，报复是恶习，同样，在此，报复确实是美德。所以，若你有什么过失，在此除去它们。若你自己不能除去它们，就去做，即便用你的肢体也当这样。有人打了你吗？向魔鬼怀恨，永不放弃对他的仇恨。或者，没人打你？但还是要对他怀恨，因为他侮辱了、冒犯了你的主和师傅，因为他伤害并攻击你的弟兄。对他要永远仇恨，永远不和解，永远不仁慈。如此，他将被谦卑，如此，他将变得可鄙，如此，他将成为易得的猎物。若我们对他凶猛，他绝不会对我们凶猛。若我们顺从，他就会凶猛；他与我们的弟兄不同。他是生命和救恩的仇敌和敌人，既是我们也是他自己的。若他不爱自己，怎能爱我们？让我们排兵布阵，击伤他，因我们有强大的同盟主耶稣基督，祂能既使我们对他的陷阱坚不可摧，也使我们堪当将来的美善；愿天主恩赐我们众人皆能达到，借着我们主耶稣基督的恩宠和慈爱，与祂一起，偕同圣神，归于圣父，荣耀、权能和尊威，自今而始，以至永远，及世之世。阿们。\par

\SourceRecord{Translated by Gross Alexander.；Nicene and Post-Nicene Fathers, First Series；Vol. 13.；Edited by Philip Schaff.；Buffalo, NY: Christian Literature Publishing Co.,；1889.；<http://www.newadvent.org/fathers/230122.htm>.}

% structure: p0001 source=h1 target=section reason=page-title

\section{厄弗所书讲道集 第二十三篇}

\BibleRef{弗 6:14}\par

\BibleQuote{所以要站稳，用真理束上你们的腰。}\par

在这样部署了这支军队，并激发了他们的热忱之后——因为这两件事都是必需的，即他们应当被排列有序并彼此服从，以及他们的精神应当被激发——并给予他们勇气，因为这也是必需的，他接着又给他们装备武器。因为武器是没有用的，除非他们先各就各位，除非士兵灵魂的精神被激发；因为我们必须先在内在武装他，然后在外在武装他。\par

\BibleQuote{所以要站稳，}他说。最首要的战术要点就是懂得如何站稳，许多事情都依赖于此。因此他多次论及站稳，在别处也说：\BibleQuote{你们要警醒，站稳。}\BibleRef{格前 16:13}并且又说：\BibleQuote{你们要这样在主内站稳。}\BibleRef{斐 4:1}并且又说：\BibleQuote{自以为站得稳的，要小心，免得跌倒。}\BibleRef{格前 10:12}并且又说：\BibleQuote{为使你们能完成一切之后，得以站稳。}\BibleRef{弗 6:13}无疑，他指的不是随便什么姿势，而是一种正确的姿势，许多有战争经验的人都知道懂得如何站稳是多么重要。因为如果在拳击手和摔跤手的情况下，教练首先推荐的就是站稳，那么在战争和军事事务中，这更是首要之事。\par

真正站稳的人，是正直的；他不是懒散地站着，也不倚靠任何东西。完全的直立通过站立的方式显示出来，因此完全正直的人，他们就站立。而不站立的人，就不能正直，而是松弛和脱节。奢侈的人不直立，而是弯曲；放荡的人也是如此，贪财的人也是如此。懂得如何站立的人，从站立本身，如同从一个地基，会发现冲突的每个部分都容易应对。\par

\BibleQuote{所以要站稳，}他说，\BibleQuote{用真理束上你们的腰。}\par

他不是在谈论字面意义上的、物质的束腰带，因为他在这一段中所有语言都采用属灵的意义。注意他如何有条不紊地进行。首先他束紧他的士兵。那么这是什么意思呢？一个生活松散、沉溺于情欲、思想拖在地上的人，他借着这腰带将他束紧，不让衣服绊住他的腿，而是让他的脚自由奔跑。\BibleQuote{所以要站稳，束上你们的腰，}他说。这里他所说的\BibleQuote{腰}，指的是：如同船上的龙骨一样，腰对我们来说也是如此，是整个身体的基础或根基：因为腰，正如医生学校告诉你的，整个身体结构建立在腰上。因此，\BibleQuote{束紧腰}就是巩固我们灵魂的基础；因为他当然不是谈论我们身体上的腰，而是进行属灵的论述：正如腰是上下部分的基础一样，在这些属灵的腰中也是如此。我们知道，当人疲劳时，常常把手放在腰上，如同放在一个基础上，以此来支撑自己；正因如此，在战争中使用腰带，以便它能够束缚并固定我们身体中的这个基础；也正因如此，当我们奔跑时我们会束紧自己。它守护着我们的力量。那么，他说，也要对灵魂这样做，这样无论做什么，我们都会强壮；这对于士兵来说尤其合适。\par

诚然，你或许会说，我们天然的腰用皮带束腰；但我们这些属神的士兵，用什么束腰呢？我回答说：用那作为一切思想之首与冠冕的，即\Quote{以真理}。他说：\Quote{要束上腰，}\Quote{以真理。}那么，\Quote{以真理}是什么意思呢？让我们不要喜爱任何虚假的事，让我们\Quote{以真理}从事一切本分，不要彼此说谎。无论是意见，让我们寻求真理；无论是生活方式，让我们寻求真道。如果我们以这真理武装自己，如果我们\Quote{以真理束腰}，那么就没有人能胜过我们。寻求真理教义的人决不会跌倒在地；因为非真理的事属于地，这从以下事实显而易见：所有那些在外面的人都受制于情欲，跟随他们自己的推理；因此，如果我们清醒，就不需要从希腊人的传说中接受教训。你看见他们是多么软弱和轻浮吗？对于天主，他们无法有一个严肃的思想或任何超越人类推理的观念。为什么呢？因为他们没有\Quote{以真理束腰}；因为他们的腰——生命的种子所藏之处和推理的主要力量——没有束上；所以没有比他们更软弱的了。再看那些\Person{摩尼教徒}，你看见他们大胆说出的一切，不都是出于他们自己的推理吗？他们说：\Quote{没有物质，天主不可能创造世界。}这一点从哪里如此明显呢？他们说这些话，是匍匐在地、从地上、从我们中间发生的事情而来；因为人，他们说，不能以别的方式创造。再看\Person{玛尔基翁}，看他说什么。\Quote{如果天主取了血肉，祂就不能保持纯洁。}这一点从哪里明显呢？他说：\Quote{因为人也不能。}但人能够做到这一点。再看\Person{瓦伦提诺}，他的推理全都拖曳在地，说的是地上的事；同样还有\Person{萨莫萨塔的保禄}。\Person{亚略}，他说什么？\Quote{天主在生的时候，不可能无受生地生。}亚略，你从哪里有胆量断言这一点？仅仅从我们中间发生的事情吗？你看见所有这些人的推理是如何在地上拖曳吗？全都，仿佛松开了、不受约束、带着地上的气味吗？关于教义方面就说到这里。关于生活和行为，同样，淫乱者、贪财者、贪图虚荣者以及一切其他的，都拖曳在地。他们没有自己的腰站立稳固，好叫他们在疲倦时能倚靠其上；但他们在疲倦时，并不把手放在上面站直，而是下垂。然而，那\Quote{以真理束腰}的人，第一，永不疲倦；第二，如果他疲倦了，他会在真理本身之上安息。什么？告诉我，贫穷会让他疲倦吗？不，决不会；因为他会安息于真正的财富，并借着这贫穷明白什么是真正的贫穷。或者，奴役会让他疲倦吗？不，决不会，因为他会知道什么是真正的奴役。或者疾病？不，连疾病也不能。基督说：\Quote{你们的腰要束上，灯也要点燃，}\BibleRef{路 12:35}，用那永不熄灭的光。这也是以色列人出埃及时\BibleRef{出 12:11}所受的命令。因为他们为什么束着腰吃逾越节羊？你愿意听其中的理由吗？按历史事实，还是按神秘意义，我该陈述呢？但我会两者都陈述，你们要记在心里，因为我这样做不是没有目的的，仅仅为了告诉你们解答，而是为了让我的话在你们中成为现实。经上记载，他们束着腰，手里拿着杖，脚上穿着鞋，这样吃逾越节羊。可敬畏和可怖的奥秘，极其深奥；如果预象如此可畏，现实岂不更可畏？他们从埃及出来，吃逾越节。请注意。经上说：\Quote{我们的逾越节羔羊基督已经被祭杀。}为什么他们束着腰？他们的装束是行路人的装束；因为他们穿着鞋，手里拿着杖，站着吃，正是说明了这一点。你们愿意先听历史，还是先听奥秘？最好是先听历史。那么历史的意义是什么呢？犹太人不断忘记天主对他们的恩惠。所以天主将这些恩惠的意义不仅系于时间，也系于那些吃的人本身的装束。这就是为什么他们要束着腰、穿着鞋吃，这样当他们被问及原因时，他们可以回答说：\Quote{我们准备好了旅行，我们即将出埃及到应许之地，我们准备好了出埃及。}这就是历史的预象。但现实是这样的：我们也吃一逾越节，就是基督；因为他说：\Quote{我们的逾越节羔羊基督已经被祭杀。}\BibleRef{格前 5:7}那么怎样？我们也应当穿着鞋、束着腰去吃。为什么？为叫我们也准备好我们的出埃及，即从此世离去。\par

Moral. 凡吃这逾越节晚餐的人，不要观望埃及，而要观望天上，就是\BibleQuote{那在上耶路撒冷}。\BibleRef{迦 4:26} 因此你们束着腰吃，因此你们穿着鞋吃，好叫你们知道，从你们开始吃逾越节那一刻起，就该出发，踏上旅程。这也意味着两件事：我们必须离开埃及，并且当我们停留时，从此要如同在异乡一般暂时居住；\BibleQuote{我们的户籍}，他说，\BibleQuote{是在天上}，\BibleRef{斐 3:20} 并且我们一生都该常常准备，好叫我们被召时不会拖延，而是说：\BibleQuote{我的心已坚定。}\BibleRef{咏 107:1} \Quote{是的，但保禄确实能说这话，他对自己一无所知；但我，需要长时间悔改，我不能说这话。} 然而束腰是警觉灵魂的特征，请听天主对那义人所说的话：\BibleQuote{你要像人一样束上你的腰，因为我要问你，你要向我说明。}\BibleRef{约 38:3} 祂也对众先知说这话，又对\Person{梅瑟}说要束腰。祂自己也向\Person{厄则克尔}显现时是束着腰的。\BibleRef{则 9:11} 不但如此，天使们也束着腰向我们显现，\BibleRef{默 15:6} 如同战士。因我们束腰，我们也勇敢站立；正如因我们站立，束腰而来。\par

因为我们也要出发，途中有许多困难。当我们渡过这个平原，魔鬼立刻攻击我们，用尽一切方法、施尽一切诡计，为要使那些从埃及得救、渡过红海、从恶鬼和无数灾祸中解脱出来的人，被他擒获毁灭。但是，如果我们警醒，我们也有火柱，就是圣神的恩宠。这恩宠既光照我们，也荫庇我们。我们有玛纳——不，不是玛纳，而是远胜玛纳。我们有属神的饮料，不是从磐石涌出的水。我们也有我们的营地，\BibleRef{默 20:9} 并且现在也住在旷野；因为现今的大地确实是缺德的旷野，甚至比那旷野更荒凉。那旷野为何如此可怕？岂不是因为其中有蝎子和毒蛇吗？\BibleRef{申 8:15} \BibleQuote{一片土地}，据说，\BibleQuote{无人经过。}\BibleRef{耶 2:6} 然而那旷野并不像这人性这样不结果实。此时此刻，这旷野中有多少蝎子、多少蝰蛇、多少蛇、多少\BibleQuote{毒蛇的种类}\BibleRef{玛 3:7} 是我们经过时要面对的！但让我们不要害怕，因为这次出谷的领队不是\Person{梅瑟}，而是\Person{耶稣}。\par

那么我们怎能不遭受同样的事呢？让我们不要犯同样的行为，这样我们就不会受同样的惩罚。他们抱怨过，他们忘恩负义；所以让我们不要怀有这些激情。他们如何全都跌倒了呢？\BibleQuote{他们蔑视了那佳美之地。}\BibleRef{咏 105:24} \Quote{怎么\InnerQuote{蔑视}呢？他们当然非常珍视它。} 由于懒惰和胆怯，不愿付出任何劳苦去获得它。那么让我们不要\Quote{蔑视}天堂！这就是\Quote{蔑视}的意思。同样，在我们中间也有果实被带来，是从天上来的果实，不是托在杠上的葡萄串，\BibleRef{户 13:23} 而是\BibleQuote{圣神的凭据}\BibleRef{格后 1:22} 、\BibleQuote{那在天上的户籍}\BibleRef{斐 3:20} ，这是\Person{保禄}和全体宗徒，那些奇妙的农夫，所教导我们的。带来这些果实的不是耶孚乃的儿子\Person{加肋布}，也不是农的儿子\Person{若苏厄}，而是\Person{耶稣}——\BibleQuote{仁慈的圣父}\BibleRef{格后 1:3} 之子，真天主的儿子，带来了每一种德行，从天上带来了所有来自那里的果实，带来了天上的赞歌。因为基路伯在上所说的言语，祂也命令我们这样说：\Quote{圣、圣、圣。} 祂给我们带来了天使的德行。\BibleQuote{天使不婚不嫁，}\BibleRef{玛 22:30} 祂也将这美好的植物栽种在这里。他们不爱钱财，也不爱任何类似的东西；这也播在我们中间。他们永远不死；同样，祂也白白赐给了我们，因为死亡不再是死亡，而是睡眠。因为请听祂所说的话：\BibleQuote{我们的朋友拉匝禄睡着了。}\BibleRef{若 11:11}\par

你们看见那\BibleQuote{在上面的耶路撒冷}的果实吗？\BibleRef{迦 4:26} 而最令人惊叹的是，我们的争战尚未结束，这一切却已在应许实现之前赐给了我们！因为他们在进入应许之地后确实劳苦——其实他们并未劳苦，因为他们若选择服从天主，就能毫无武器、不列阵势地夺取所有城邑。我们知道，他们推翻了\Place{耶里哥}，与其说是战士，不如说是舞者。然而，我们进入应许之地，即天堂之后，不再有争战，只在旷野中，即现世生活中才有。因为\BibleQuote{那进入安息的人，自己也歇了自己的工作，正如天主歇了他的工作一样。}\BibleRef{希 4:10} \BibleQuote{所以我们不要灰心行善，因为若不疲倦，到了时候我们必要收获。}\BibleRef{迦 6:9} 你们看见祂如何带领他们，也照样带领我们吗？在他们身上，关于玛纳和旷野，经上说：\BibleQuote{多收的没有剩余，少收的也没有缺乏。}\BibleRef{出 16:18} 我们也受命：\BibleQuote{不要在地上积蓄财宝。}\BibleRef{玛 6:19} 但如果我们积蓄财宝，腐蚀它的就不再是地上的虫子，如同玛纳那样，而是那永远与火同在的。那么，让我们\Quote{征服一切}，免得给这虫子提供食物。因为\Quote{他,}，经上说，\BibleQuote{多收的没有剩余。} 因为这也天天发生在我们身上。我们所有人都只有同样的饥饿容量需要满足。超过这个的，不过是增加忧虑。因为祂后来要传讲的话，即\BibleQuote{一天的难处一天当就够了}\BibleRef{玛 6:34}，这从起初祂就这样教导，但即使这样，他们还是没领受。至于我们，让我们不要贪得无厌，不要不知足，不要寻求华丽的房屋；因为我们是在朝圣途中，并非在家；所以，如果有人知道现世生活是一种旅程、远征，或者可以说，是所谓的野营，他就不会寻求华丽的建筑。因为，请问，即使再富有，谁会在营地里建造华丽房屋呢？没有人；他会成为笑柄，他会为敌人建造，更有效地引他们前来；所以，如果我们理智，我们也不会。现世生活不过是一次行军和一个营地。\par

因此，我恳求你们，尽一切努力，不要在这里积蓄财宝；因为如果盗贼来了，我们必须立刻起身离开。祂说：\BibleQuote{你们要醒寤，因为你们不知道盗贼什么时候来到。}\BibleRef{玛 24:42} 祂这样称呼死亡。那么，在他来到之前，让我们把一切先遣送到我们的祖国；但在这里，让我们\Quote{束好腰}，好能战胜我们的敌人，愿天主使我们能战胜他们，藉着我们的主耶稣基督的恩宠和慈爱，同祂一起，与圣神，愿光荣、德能、尊威归于圣父，直到永远。阿们。\par

\SourceRecord{Translated by Gross Alexander.；Nicene and Post-Nicene Fathers, First Series；Vol. 13.；Edited by Philip Schaff.；Buffalo, NY: Christian Literature Publishing Co.,；1889.；<http://www.newadvent.org/fathers/230123.htm>.}

% structure: p0001 source=h1 target=section reason=page-title

\section{论厄弗所书讲道词24}

\BibleRef{弗6:14-17}\par

\BibleQuote{所以，你们要站稳！用真理作带，束起你们的腰，穿上正义作甲，以和平的福音作预备走路的鞋，穿在脚上；此外，还要拿起信德作盾牌，使你们能以此扑灭恶者的一切火箭；并戴上救恩的头盔，拿着圣神的利剑，即天主的话。}\par

他说的\Quote{束起你们的腰}，然后\Quote{用真理}。这是什么意思？我在前一篇讲道中说过，他应当轻装上阵，以免跑步时受到任何阻碍。\par

他继续说道：\Quote{穿上}，\Quote{正义的甲}。如甲胄一样坚不可摧，正义也是如此，此处正义指一切美德的生活。这样的人永远不会被推翻，因为许多人会伤害他，但无人能击溃他，连魔鬼也不能。正如有人说：\Quote{正义的行为铭刻在胸中}。关于这些，基督说：\BibleQuote{饥渴慕义的人是有福的，因为他们要得饱饫。}\BibleRef{玛5:6}这样他就坚固如甲胄。这样的人永远不会发脾气。\par

\Quote{以和平的福音作预备走路的鞋，穿在脚上。}这话在何种意义上说，更不确定。那么它的意思是什么？无疑，这是高贵的胫甲，他用它来装备我们。要么他意指我们应为福音做好准备，并为此使用我们的脚，为它预备和铺平道路；否则，至少我们自己应为离世做好准备。那么，\Quote{和平福音的准备}不是别的，而是一种最道德的生活；依照先知所说的：\Quote{祢将预备他们的心，祢将使祢的耳垂听。}\BibleRef{咏9:17}他说：\Quote{和平的福音}，这是有理由的；因为他既提到了战争和战斗，就向我们表明，与恶灵的这场冲突我们必须进行：因为福音是\Quote{和平的福音}；我们与他们之间的这场战争，结束了另一场战争，即我们与天主之间的战争；如果我们与魔鬼作战，我们就与天主和好。所以，亲爱的，不要害怕；这是\Quote{福音}，即好消息；胜利已经赢得了。\par

\Quote{此外，还要拿起信德作盾牌。}\par

此处\Quote{信德}非指知识（因为那样他绝不会将其排在最后），而是指行奇迹的恩赐。他有理由称此为\InnerQuote{信德的盾牌}；因为正如盾牌遮蔽全身，如同壁垒，这种信德也是如此；因为万事都屈服于它。\par

他说：\Quote{你们就能以此}，\Quote{扑灭恶者的一切火箭。}\par

因为没有什么能抵抗这盾牌；请听基督对祂门徒说的话：\BibleQuote{如果你们有信德像芥子那样大，你们向这座山说：从这里移到那里去！它必会移过去。}\BibleRef{玛17:20}但我们是怎样得到这种信德呢？当我们妥善完成所有那些本分时。\par

\Quote{恶者的火箭}他指诱惑和邪恶的欲望；他说\Quote{火一般的}，因为这类欲望具有这样的特质。然而，如果信德能命令恶灵，那么它更能命令灵魂的情欲。\par

他继续说：\Quote{戴上救恩的头盔}，即你们的救恩。因为他正在为他们装备盔甲。\par

\Quote{拿着圣神的剑，即天主的话。}他或指圣神，或指\Quote{属神的剑}；因为藉此万事被分割，藉此万事被劈开，我们甚至能砍断蛇的头。\par

第18、19、20节。他说：\BibleQuote{时时靠着圣神，以各种祈求和哀祷祈祷；且要醒寤不倦，为众圣徒祈求；也为我祈求，使我在开口的时候，能放胆传扬福音的奥秘——为这福音我作了带锁链的使者——使我能放胆讲论，照我所该讲的。}\BibleRef{弗6:18}\par

既然天主的圣言有能力成就一切，同样，拥有神恩的人也是如此。\BibleQuote{因为天主的圣言是生活的}，他说，\BibleQuote{是有效的，比任何双刃剑更锋利。}\BibleRef{希 4:12} 现在请注意这位蒙福宗徒的智慧。他装备了他们，使他们万无一失。此后还需要什么呢？呼求君王，好使他伸出援手。\Quote{以各种祈祷和恳求，时时在圣神内祈祷}；因为有可能\Quote{祈祷}不是\Quote{在圣神内}，而当人\BibleQuote{用重复的空言}\BibleRef{玛 6:7}；\Quote{且为此醒寤}，他又加上，即保持清醒；因为武装的战士，那站在君王身旁的人，应当如此：警醒而节制——\Quote{在一切坚忍和恳求中，为众圣徒；并也为我，使我在开口时，能获得言语。}你说什么，蒙福的保禄？那么，你需要你的门徒吗？他说\Quote{在开口时}这话说得对。他当时并不是研究自己要说什么，而是照基督所说的：\BibleQuote{当他们把你们交出时，不要担心怎样或说什么：因为在那时刻，必会赐给你们当说的话。}\BibleRef{玛 10:19} 他确实凡事凭信德，凡事凭恩宠行事。\Quote{以大胆}，他继续说，\Quote{讲明福音的奥秘}；也就是说，使我能为它辩护，如同我应当做的。而你身缠锁链，还需要他人的帮助吗？他说，是的，因为伯多禄也曾被锁链束缚，但尽管如此，\BibleQuote{为他恳切祈祷}\BibleRef{宗 12:5}。\Quote{我为了这福音，成了带锁链的使者，好使我放胆讲论，如同我应当讲的。}也就是说，使我以信心、勇气和极大的智慧来回答。\par

第21节。\BibleQuote{但是，为叫你们也知道我的事，我怎样行事，有一切事，主内亲爱的弟兄、忠信的仆役提希苛，必会叫你们知道。}\BibleRef{弗 6:21}\par

他一提到自己的锁链，也就为提希苛留下些事，让他主动向他们讲述。因为凡是教义和劝勉的话题，他都藉信件解释；但那些仅是叙述的事，他就委托送信的人。\Quote{为叫你们知道我的事}，即得以了解。这显明了他对他们的爱，以及他们对他的爱。\par

第22节。\BibleQuote{我打发他到你们那里去，正是为此，为叫你们知道我们的情况，并安慰你们的心。}\BibleRef{弗 6:22}\par

他说这话并非无的放矢，而是因先前的教训：\Quote{用真理束腰，穿上正义的甲}等等，这些是恒常和不停前进的记号；因为听先知说什么：\BibleQuote{愿这事如同他披上衣服，如同他常束的腰带}\BibleRef{咏 108:19}；依撒意亚先知又说，天主\BibleQuote{穿上正义作甲}\BibleRef{依 59:17}；这些表达教导我们，这些事我们必须常有，不仅短时间，而是持续不断，因为常有战争的需要。\BibleQuote{因为经上说：义人胆大如狮子。}\BibleRef{箴 28:1} 因为那披戴这样甲的人，不可能害怕那对抗的阵营，反而要跳入敌阵中。依撒意亚又说：\BibleQuote{那传报喜讯者的脚步是多么美丽啊！}\BibleRef{依 52:7} 谁不愿奔跑，谁不在这样的使命中服务？传报平安的喜讯，天主与人和好，平安，人没有劳苦，而天主完成了这一切。\par

但是，什么是\Quote{福音的预备}？让我们听从若翰的话：\BibleQuote{预备主的道路，修直他的途径。}\BibleRef{玛 3:3} 但在圣洗后还需要另一种\Quote{预备}，好使我们不做任何不配\Quote{平安}的事。再者，既然脚通常象征生活方式，因此他常用这话劝勉：\BibleQuote{所以你们应谨慎行事，不要像愚昧人，却要像明智人。}\BibleRef{弗 5:15} 因此，他要我们展现一种配得福音的行径和榜样；即使我们的生活和行为纯洁。平安的喜讯已传报给你们，要给这喜讯一条预备好的道路；因为若你们再次成为仇敌，就不再是\Quote{平安的预备}。要准备好，不要退缩接受这平安。正如你们已准备好并倾向平安和信德，也要如此持续。盾牌是首先承受敌人攻击的，并保护盔甲不受损伤。所以，只要信德正当，生活正当，盔甲就保持无损。\par

然而，他谈论信德很多，尤其在写给希伯来人的书中，如同他也谈论望德。相信未来的美物，他说，这样这些盔甲没有一件会受到损伤。在危险中，在劳苦中，用你们的望德和信德来保护自己，你们就能保持盔甲无损。\BibleQuote{凡来到天主面前的人，应该信他存在，且信他成为寻求他者的赏报者。}\BibleRef{希 11:6} 信德是盾牌；但哪里有诡辩、推理和仔细审查，那里就不再是盾牌，反而妨碍我们。让我们的信德如此，好能覆盖并庇护整个身体。不要让它过于狭小，以致暴露脚或任何部分，而要让盾牌与全身相称。\par

\Quote{火焰的箭矢。}因为有许多令人怀疑的推理，它们仿佛点燃了灵魂，有许多困难，有许多困惑，但信德使这一切完全平息；魔鬼投掷许多事物，为要焚烧我们的灵魂，使我们陷入不确定之中；例如，当有些人说：\Quote{那么，有复活吗？}\Quote{有审判吗？}\Quote{有报应吗？}\Quote{但是，有信德吗？}宗徒会说：\Quote{你要以信德熄灭魔鬼的箭矢。是否有卑劣的情欲袭击了你？在你面前持守对未来美善事物的信德，它甚至不会显现，是的，它将消亡。}\Quote{所有的箭矢}；不是一些被熄灭，而另一些没有。请听保禄所说的：\BibleQuote{因为我算定，现时的苦楚，与将来要在我们身上显示的光荣，是不能相比的。}\BibleRef{罗 8:18}你看见那些日子里义人熄灭了多少箭矢吗？在你看来，当圣祖献上他的儿子时，他内心燃烧着烈火，这岂不正是\Quote{火焰的箭矢}吗？是的，其他义人也熄灭了\Quote{他所有的箭矢}。那么，无论是攻击我们的推理，还是卑劣的欲望，还是劳苦和困苦，我们都应持守这个。在所有其他武器中，这是保障；如果我们没有这个，它们将迅速被刺穿。\Quote{此外，}他说，\Quote{拿起信德的盾牌。}\Quote{此外}是什么意思？它既指\Quote{在真理中}，也指\Quote{在义德中}，也指\Quote{在福音的准备中}；这就是说，所有这些都需要信德的帮助。\par

因此，他又补充说：\BibleQuote{并戴上救恩的头盔}；这就是说，最终藉此你便能安全。接受救恩的头盔就是逃避危险。因为正如头盔完美地覆盖头部的每一部分，不让它受到任何伤害，反而保护它，同样，信德也同时起到盾牌和头盔的作用，保护我们。因为如果我们熄灭了他的箭矢，我们也会迅速接受那些拯救性的思想，这些思想不让我们的主导原则受到任何伤害；因为如果这些敌对於我们得救的思想被熄灭了，那些并非如此、反而有助於我们的得救、并给我们带来美好希望的思想，就会在我们内产生，并安息在我们的主导原则上，如同头盔安息在头上。\par

不仅如此，我们还要拿起\Quote{圣神的利剑}，这样不仅能抵挡他的投射物，还能击伤魔鬼自身。因为一个不自暴自弃、能抵挡那些火焰箭矢的灵魂，会毫无畏惧地面对敌人，并用这把剑刺穿他的胸甲——保禄正是用这把剑刺穿了他的胸甲，\BibleQuote{把他的计谋俘获了}\BibleRef{格后 10:5}；他将砍断并斩下那蛇的头。\par

\Quote{就是天主的圣言。}\par

在此处，\Quote{天主的圣言}一方面指天主的诫命，或命令之言；另一方面指那是因基督之名。因为如果我们遵守他的诫命，藉此我们就能杀死并诛灭那龙本身，\BibleQuote{那弯曲的蛇}\BibleRef{依 27:1}。正如他所说：\Quote{你们将能熄灭那恶者一切的火焰箭矢}；为使他不致使他们骄傲，他表明他们最需要的是天主；因为他说了什么？\par

\Quote{藉着各种祈祷和恳求，}他说，这些事将要成就，你们藉着祈祷将完成一切。但当你走近时，绝不要只为自己祈求：这样你才能得到天主的眷顾。\par

\Quote{藉着各种祈祷和恳求，时时在圣神内祈祷，并为此警醒，恒心为众圣人祈求。}我说，不要限定在一天的某个时辰，请听他所说的：要时时亲近；\BibleQuote{不断祈祷}\BibleRef{得前 5:17}。你岂没有听说过那个寡妇，她如何因她的恒心而成功？\BibleRef{路 18:1}你岂没有听说过那个朋友，他在半夜因坚持而羞愧地使他的朋友让步？\BibleRef{路 11:5}你岂没有听说过那个客纳罕妇人\BibleRef{谷 7:25}，她如何因恳求的恒心而唤起了主的怜悯？这些人都是因他们的坚持而达到了目的。\par

\Quote{时时在圣神内祈祷}，他说。\par

这就是说，让我们寻求那合乎天主的事物，不要寻求这世界的事物，不要寻求与此生命相关的事物。\par

因此，我们不仅需要\Quote{不断祈祷}，而且还要\Quote{警醒——并为此警醒}，他说。无论他在这里是指守夜，还是指灵魂的警觉，我都接受这两种含义。你看见那客纳罕妇人如何警醒祈祷了吗？虽然主没有回答她，反而拒绝她，称她为狗，她却说：\BibleQuote{是啊，主！可是小狗也吃主人桌子上掉下来的碎屑}\BibleRef{玛 15:27}，并且没有停止，直到她得到了所求的。同样，那寡妇如何呼喊并坚持如此之久，直到她羞愧地使那个不怕天主也不尊重人的判官让步？\BibleRef{路 18:1}再者，那朋友如何坚持，在深夜一直守在门口，直到他因坚持而使对方羞愧地让步，叫他起来。\BibleRef{路 11:5}这就是警醒。\par

你想知道祈祷中的警醒是什么吗？到\Person{哈纳}那里去，听她的话语：\Quote{\OriginalTerm{阿多乃·厄罗依·撒巴奥特}{Adonai Eloi Sabaoth}}。\BibleRef{撒上 1:11}不，不如先听她的话之前的事；历史记载说：\BibleQuote{他们都站起身来}，\BibleQuote{离开了饭桌}\BibleRef{撒上 1:9}，而她立即没有去睡觉或休息。在我看来，她甚至在坐席时也吃得很少，没有被食物所拖累。否则她绝不会流那么多眼泪；因为如果我们，在禁食和无食物时，很难如此祈祷，或者更确切地说，从不如此祈祷，那么她更不可能在饭后如此祈祷，除非她在吃饭时也如不吃的人一样。让我们这些做男人的，以这个女人的榜样为羞耻；让我们这些追求和渴慕天国的人，因她为一个小孩祈祷哭泣而羞耻。经上说：\BibleQuote{她站在}，\BibleQuote{上主面前}\BibleRef{撒上 1:10}；她的话是什么？\Quote{\OriginalTerm{阿多乃，主，厄罗依·撒巴奥特}{Adonai, Lord, Eloi Sabaoth}}！这就是翻译过来的：\Quote{上主，万军的天主。}她的眼泪比舌头先到；她以此希望说服天主垂听她的祈求。哪里，眼泪总是伴随着痛苦；哪里，痛苦就伴随着大智慧和警觉。她继续说：\BibleQuote{你若垂顾婢女的苦情，赐给婢女一个男孩，我必将他终身献给上主。}\BibleRef{撒上 1:11}她没有说\Quote{一年}或\Quote{两年}，像我们说的——也没有说\Quote{你若给我一个孩子，我就给你钱}；而是说\Quote{我把这礼物本身完整地归还给你，我的长子，我祈祷的儿子。}这真是\Person{亚巴郎}的女儿。他是在被要求时给出；她甚至在未被要求之前就献上。\par

但请注意她在此之后深切的敬畏。经上说：\BibleQuote{只见她的嘴唇在动，却听不到她的声音。}\BibleRef{撒上 1:13}凡祈求的人就这样亲近天主；不是顾及自己安逸，不是张口、懒散、挠头，也不是完全无精打采。怎么，难道天主不能无需任何祈祷就赐予吗？怎么，难道祂在女子祈求之前不知道她的愿望吗？然而，如果他在她祈求之前就赐予，那么女子的热忱就不会显现，她的德行就不会彰显，她也得不到如此大的赏报。所以，延迟不是出于嫉妒或巫术，而是出于上智的安排。因此，当你们听到圣经说\Quote{上主封闭了她的子宫}（第5、6节），以及\Quote{她的情敌深深激怒她}时，当知道这是天主有意证明女子的认真。因为，请注意，她有一个深爱她的丈夫，因为他说：\BibleQuote{有我在你身边，岂不比十个儿子还好？}\BibleRef{撒上 1:8}经上说：\Quote{她的情敌深深激怒她}，即责骂她，凌辱她。然而她从未报复，也没有咒骂她，也没有说\Quote{为我报仇，因为我的情敌辱骂我}。另一个（情敌）有孩子，但这个女子有丈夫的爱作为补偿。至少，她的丈夫用同样的话安慰她。\par

然而，让我们再看看这个妇女的大智慧。经上说：\BibleQuote{厄里以为她喝醉了。}\BibleRef{撒上 1:13}但她对\Person{厄里}说的话：\BibleQuote{我主，不要将你的婢女看作一个\OriginalTerm{贝里雅耳}{Belial}的女子，因为我因满怀愁苦和激愤，才一直说到现在。}\BibleRef{撒上 1:16}这真是痛悔之心的明证：当我们不被辱骂我们的人激怒，不愤恨他们，只为自己辩护。没有什么比受苦更能使心灵有智慧；没有什么比\BibleQuote{依天主的心而忧愁}\BibleRef{格后 7:10}更甜美的。她说：\Quote{我因满怀愁苦和激愤，才一直说到现在。}让我们众人效仿她。你们这些不育的，听啊！你们这些想要孩子的，听啊！丈夫和妻子，都听啊！是的，因为丈夫也常常尽他们的本分；因为听圣经说：\BibleQuote{依撒格为自己的妻子黎贝加祈求上主，因为她不能生育。}\BibleRef{创 25:21}因为祈祷能成就大事。\par

他说：\Quote{以各种祈祷和恳求，为众圣人，也为我，}把自己放在最后。哦，蒙福的\Person{保禄}，你为何如此把自己放在最后？是的，他说：\Quote{使我能得着宣讲的恩赐，大胆开口，阐明福音的奥秘，我为此作了带锁链的使者。}你在哪里作使者？他说：\Quote{向人类。}哦！天主奇妙的慈爱！祂从天上以自己的名派和平的使者，然而人们抓住他们，捆绑他们，甚至不尊重国际法（使者不应受任何伤害）。他说：\Quote{但无论如何，我是带锁链的使者。锁链如缰绳般束缚着我，限制我的胆量，但你们的祈祷会开口}，好使我讲出所被派遣的全部话语。\par

\Quote{但为叫你们也知道我的事，我现在的景况，\Person{提希苛}，我所亲爱、在主里忠信的执事，将会告诉你们一切。}若他\Quote{忠信}，他不会说谎，他会在一切事上说实话——\Quote{我派他到你们那里去，正是为此，为叫你们知道我们的景况，并使你们的心得安慰。}奇妙超绝的爱！他意思是：\Quote{这样，那些想恐吓你们的人就不能得逞。}因为很可能他们正处患难中；因为\Quote{使你们的心得安慰}这句话暗示了这一点；即\Quote{不叫你们沉沦。}\par

第23节。\BibleQuote{愿平安与爱德偕信德，从天主父及主耶稣基督，赐与弟兄们。}\BibleRef{弗 6:23}\par

他以\Quote{平安与爱德偕同信德}呼求他们。他说得好：因为他不愿他们只顾爱德，而与异信者混杂。他或是指此意，或是指上述的，即他们也当有信德，好对将来的福乐怀有愉快的信赖。这\Quote{平安}是对天主的，以及这\Quote{爱德}。若有平安，也必有爱德；若有爱德，也必有平安。\Quote{偕同信德}，因为没有信德，爱德便毫无价值；或者说，没有信德，爱德根本不可能存在。\par

第24节。\BibleQuote{愿恩宠与所有在不变中爱我们的主耶稣基督的人同在。}（\BibleRef{弗 6:24}）\par

他为何在此将二者分开，把\Quote{平安}单独放置，而把\Quote{恩宠}也单独放置？\par

\Quote{在不变中}，他这样结束。\par

\Quote{在不变中}是什么意思？它或指\Quote{在纯洁中}；或指\Quote{为了那些不朽坏的事物}，例如，不是在财富中，也不是在荣耀中，而是在那些不朽坏的宝藏中。\Quote{在}字的意思是\Quote{藉着}。\Quote{藉着不变}，即\Quote{藉着德行}。因为一切罪都是败坏。正如我们说一个童贞女被败坏，我们也这样论及灵魂。因此保禄说：\Quote{免得你们的心意被败坏。}（\BibleRef{格后 11:3}）他在别处又说：\Quote{在教训上，显出不变。}因为，请告诉我，身体的败坏是什么？岂非整个架构及其结合的解体？当罪进入时，灵魂中也发生同样的事。灵魂的美是节制与正义；灵魂的健康是勇敢与明智；因为卑鄙的人在我們眼中是丑陋的，贪婪的人也是如此，沉溺于恶行的人也是如此，懦弱无丈夫气的人是有病的，愚昧的人是不健康的。罪造成败坏，从这一点可以显明：它们使人卑鄙、软弱，使人患病、不健康。不仅如此，当我们说一个童贞女被败坏时，严格地说，也是为此缘故，不仅因为身体被玷污，而且因为犯过。因为单纯的行为是自然的；如果败坏在于那行为，那么婚姻便是败坏了。因此，败坏的不是行为，而是罪，因为它羞辱她，使她蒙羞。再者，在一所房屋的情况下，什么是败坏？它的解体。所以，普遍来说，败坏是一种向恶的变化，一种进入另一种状态的变化，以致前一种状态完全消失。因为请听圣经所说：\BibleQuote{凡有血肉的人，品行败坏了}（\BibleRef{创 6:12}）；又说：\BibleQuote{在不能忍受的败坏中}（\BibleRef{出 18:18}）；又说：\BibleQuote{心术败坏了的人}（\BibleRef{弟后 3:8}）。我们的身体是可朽坏的，但我们的灵魂是不朽坏的。那么，我们不要使灵魂也成为可朽坏的。身体的败坏是先前罪的后果；但圣洗之后的罪也有能力使灵魂成为可朽坏的，并使它容易成为\Quote{不死的虫}的猎物。因为那虫若没有发现灵魂是可朽坏的，就绝不会碰它。虫碰不着金刚石，即使碰了，也不能伤害它。所以，不要败坏灵魂；因为败坏之物充满恶臭；请听先知说：\BibleQuote{因我的愚妄，我的疮口流脓溃烂。}（\BibleRef{咏 37:5}）\par

然而，身体的\BibleQuote{这败坏}将\BibleQuote{穿上不朽}（\BibleRef{格前 15:53}），但灵魂的败坏则永不如此；因为哪里有不朽，哪里就没有败坏。如此，它是一种不朽的败坏，没有尽头，一种不死的死亡；倘若身体不死，那就会是这样。如果我们离世进入来世时没有败坏，我们便有那种不朽且无尽的败坏；因为永远燃烧而不被烧毁，永远被虫蛀蚀，便是不朽的败坏；正如蒙福的约伯的情形。他败坏了，却没有死，而且持续了很长一段时间，\Quote{不断消耗，从伤口刮去泥块。}当虫子环绕并吞噬它时，它将遭受类似的折磨，不是两年、三年、十年、一万年，而是无穷尽的年岁；因为祂说：\Quote{他们的虫总不死。}\par

劝言：让我们警醒，我恳求你们，让我们畏惧这些话，以免遭遇现实。贪婪是一种败坏，比其他败坏更危险，并导致偶像崇拜。让我们逃避这败坏，选择不朽。你是否在贪婪中欺骗并欺诈了某人？你贪婪的果实腐朽，但贪婪本身却存留；这是一种败坏，是不朽之败坏的基础。享乐确实消逝，但罪却永不腐朽。我们在此世不抛弃一切，乃是可怕的灾祸；带着罪的沉重负累离开尘世，乃是巨大的灾难。经上说：\Quote{因为在阴府，} 经上说 \Quote{谁会称谢祢？} \BibleRef{咏 6:5} 那是审判的地方；那时再无悔改的时期。那时那富人哀叹了多少事？\BibleRef{路 16:23} 却无济于事。那些疏忽喂养基督的人说了多少话？\BibleRef{玛 25:41} 他们仍被带入永火。那时他们有多少话要说：\Quote{那作恶的人}；\Quote{主啊，我们不是因祢的名字预言，因祢的名字驱魔吗？} 然而他们却不被承认。因此这一切届时将会发生；但若现在不实行，便毫无益处。让我们畏惧，免得那时我们不得不说：\Quote{主啊，我们几时见了祢饥饿，而没有供养祢呢？} \BibleRef{玛 25:44} 让我们现在就供养祂，不是一天、两天、三天。\Quote{不可让慈爱和诚实离开你，} 智者说，\Quote{要系在你颈项上。} \BibleRef{箴 3:3} 祂不是说\Quote{行一次，也不两次。} 我们知道，那几位童女有油，却不够持久。\BibleRef{玛 25:3} 因此我们需要大量的油，并且应当如同天主殿中的青橄榄树。\BibleRef{咏 51:8} 让我们反省自己每人身上有多少罪的重量，并让我们的慈善行为与之平衡；不，更要远远超过它们，好使不仅罪恶被消灭，而且正义行为也被算为我们的义。因为如果善行在数量上不足以抵消指控我们的罪行，并且剩余部分不足以被算为我们的义，那么无人能从那种惩罚中拯救我们。愿天主恩赐我们所有人都能逃脱那惩罚，借着我们的主耶稣基督的恩宠和慈爱，愿光荣归于祂，偕同圣父，等等。\par

\SourceRecord{Translated by Gross Alexander.；Nicene and Post-Nicene Fathers, First Series；Vol. 13.；Edited by Philip Schaff.；Buffalo, NY: Christian Literature Publishing Co.,；1889.；<http://www.newadvent.org/fathers/230124.htm>.}
